\documentclass[11pt,a4paper]{article}

\usepackage[margin=2.2cm]{geometry}
\usepackage{booktabs}
\usepackage{array}
\usepackage{longtable}
\usepackage{tabularx}
\usepackage{graphicx}
\usepackage{xcolor}
\usepackage{hyperref}
\usepackage{siunitx}
\usepackage{enumitem}
\usepackage{amsmath}
\usepackage{amssymb}
% Packages the pasted reports assumed and the diary's own preamble never loaded.
\usepackage{bbm}
\usepackage{algorithm}
\usepackage{algpseudocode}
\usepackage{subcaption}
\usepackage{multirow}
\usepackage{makecell}
\usepackage{threeparttable}
\usepackage{listings}
\usepackage{caption}
\usepackage{microtype}
\usepackage{tabularx}
\usepackage{siunitx}
\usepackage{listings}
\hypersetup{
    colorlinks=true,
    linkcolor=blue!60!black,
    urlcolor=blue!60!black,
    citecolor=blue!60!black
}

\lstdefinestyle{bashstyle}{
    basicstyle=\ttfamily\small,
    breaklines=true,
    frame=single,
    columns=fullflexible,
    backgroundcolor=\color{gray!5}
}

\lstset{style=bashstyle}

\title{
    \textbf{Cache-Only Temporal Refinement of S2M2-S on SCARED}\\
    \large Current Findings, Negative Results, and Next Research Direction
}

\author{ARGOS Internal Research Note}
\date{\today}

% The diary quotes reports written for venues that use natbib, but it carries no
% bibliography of its own, so \citep would be an undefined control sequence. Printing the
% key is the honest rendering: the reference is a note about which paper, not a citation.
% Notation macros carried over from ICRA.tex, so blocks pasted from the paper render
% here without keeping two copies of the definitions in step.
\providecommand{\steer}{\textsc{TETHER}}
\providecommand{\raw}{\mathrm{raw}}
\providecommand{\mem}{\mathrm{mem}}
\providecommand{\fused}{\mathrm{fused}}
\providecommand{\warp}{\mathcal{W}}
\providecommand{\EPE}{\mathrm{EPE}}
\providecommand{\HUR}{\mathrm{HUR}}
\providecommand{\BUR}{\mathrm{BUR}}
\providecommand{\HPlus}{H^{+}}
\providecommand{\BPlus}{B^{+}}
\providecommand{\Bad}{\mathrm{Bad}}
\providecommand{\DTCE}{\mathrm{DTCE}}
\providecommand{\TEPE}{\mathrm{TEPE}}
\providecommand{\OPW}{\mathrm{OPW}}
\providecommand{\AURC}{\mathrm{AURC}}
\providecommand{\R}{\mathbb{R}}
\providecommand{\citep}[1]{[#1]}
\providecommand{\citet}[1]{[#1]}

% Carried up from a report that was pasted in below with its own preamble.
\definecolor{good}{RGB}{32,120,75}
\definecolor{bad}{RGB}{180,45,45}
\definecolor{accent}{RGB}{35,85,150}
\definecolor{lightgray}{RGB}{245,245,245}
\newcommand{\good}[1]{\textcolor{good}{#1}}
\newcommand{\bad}[1]{\textcolor{bad}{#1}}
\newcommand{\method}{EGBM-CARE}
\newcommand{\smodel}{S2M2}
\hypersetup{
    colorlinks=true,
    linkcolor=accent,
    urlcolor=accent,
    citecolor=accent
}

\begin{document}

\maketitle

\begin{abstract}
This document summarizes the current experimental findings on temporal refinement for surgical stereo depth/disparity estimation within the ARGOS project. The current benchmark evaluates cached S2M2-S predictions on an audited SCARED temporal-GT sequence, comparing raw predictions, fixed causal exponential moving average (EMA), no-flow adaptive EMA variants, RAFT-based motion-compensated EMA, and StereoAnyVideo as an external video-model comparison. The main finding is that a very simple fixed EMA with $\alpha=0.35$ is an unexpectedly strong deployment baseline, substantially improving temporal stability and depth error over raw S2M2-S at negligible cost. No-RAFT adaptive heuristics based on raw-vs-memory disagreement and disparity gradients did not outperform the fixed EMA baseline. RAFT-warped EMA remains the best quality method, but its online cost is significantly higher. These results suggest that simple adaptive rules without explicit motion compensation are insufficient, and that the next meaningful research direction should focus either on lightweight flow-based motion compensation, such as RAFT-Small, or on distilling a RAFT-based temporal teacher into a lightweight flow-free recurrent student.
\end{abstract}

\section{Context and Research Question}

The broader goal is to build a practical temporal refinement pipeline for surgical stereo reconstruction. The target setting is not generic video depth estimation, but lightweight and causal temporal refinement for surgical scenes, where deployment constraints matter.

The current research question is:

\begin{quote}
\emph{Can we improve the temporal consistency and geometric accuracy of S2M2-S disparity predictions using lightweight causal temporal fusion, without running a heavy video model or optical flow estimator online?}
\end{quote}

This question emerged after observing that a full RAFT-warped EMA method improves quality, but is too expensive for a lightweight deployment scenario. The key tension is therefore:

\begin{center}
\textbf{motion-aware temporal fusion improves quality} \quad vs. \quad \textbf{online optical flow is expensive}.
\end{center}

\section{Dataset and Cached Inputs}

The current benchmark is cache-only: no model inference is rerun, and no training is performed.

\subsection{Audited SCARED Temporal-GT Sequence}

The audited sequence is:

\begin{lstlisting}
dataset/SCARED/curated/temporal_gt/test_dataset_9_keyframe_3
\end{lstlisting}

The audit found:

\begin{itemize}[leftmargin=*]
    \item 130 left/right RGB frames.
    \item 130 ground-truth disparity maps.
    \item 130 ground-truth depth maps.
    \item 130 valid masks.
    \item 130 calibration files.
    \item 129 adjacent frame pairs.
    \item 103 valid frames retained after validity filtering.
    \item 102 valid temporal pairs used for temporal metrics.
\end{itemize}

\subsection{Prediction Caches}

The benchmark uses cached prediction folders:

\begin{lstlisting}
results/03_temporal_refinement/evaluation/gt_temporal_test_dataset_9_keyframe_3/predictions/S2M2-S_512
\end{lstlisting}

\begin{lstlisting}
results/03_temporal_refinement/evaluation/gt_temporal_test_dataset_9_keyframe_3/predictions/StereoAnyVideo_384x640
\end{lstlisting}

The RAFT flow cache is used only for metric computation and for evaluating already-defined RAFT-based baselines:

\begin{lstlisting}
results/04_dataset_derivatives/SCARED/temporal_gt_flow_cache/test_dataset_9_keyframe_3/raft
\end{lstlisting}

For no-RAFT adaptive methods, the flow cache is explicitly \textbf{not} used in the prediction formula.

\section{Benchmark Design}

\subsection{Cache-Only Evaluation}

The benchmark is deliberately cache-only. It does not run:

\begin{itemize}[leftmargin=*]
    \item S2M2-S inference.
    \item StereoAnyVideo inference.
    \item RAFT inference.
    \item Training.
    \item Dataset modification.
    \item Prediction-cache modification.
\end{itemize}

This design isolates the effect of the temporal fusion/refinement rule.

\subsection{Main Evaluated Methods}

The evaluated families are:

\begin{enumerate}[leftmargin=*]
    \item \textbf{S2M2-S raw}: cached single-frame S2M2-S predictions.
    \item \textbf{Fixed EMA}: causal temporal smoothing with fixed global $\alpha$.
    \item \textbf{RAFT-warped EMA}: motion-compensated EMA using cached RAFT flow.
    \item \textbf{Confidence/occlusion reset}: RAFT-based warped EMA with reset logic.
    \item \textbf{Conservative adaptive EMA with RAFT}: RAFT-based adaptive variant.
    \item \textbf{StereoAnyVideo}: cached video-model comparison, not ground truth.
    \item \textbf{No-RAFT adaptive EMA}: new lightweight adaptive heuristics without flow.
\end{enumerate}

\section{Temporal Fusion Methods}

\subsection{Fixed Causal EMA}

The fixed EMA baseline is:

\begin{equation}
    \hat{d}_t = \alpha d_t + (1-\alpha)\hat{d}_{t-1},
\end{equation}

where $d_t$ is the raw S2M2-S disparity prediction at time $t$, and $\hat{d}_t$ is the refined disparity.

The best fixed EMA found so far is:

\begin{equation}
    \alpha = 0.35.
\end{equation}

This means the method strongly reuses temporal memory while still incorporating the current prediction.

\subsection{RAFT-Warped EMA}

The RAFT-warped EMA uses optical flow to warp the previous refined disparity into the current frame before fusion:

\begin{equation}
    \hat{d}_t = \alpha d_t + (1-\alpha)W(\hat{d}_{t-1}, F_{t-1 \rightarrow t}),
\end{equation}

where $W(\cdot)$ denotes warping and $F_{t-1 \rightarrow t}$ is the forward optical flow.

This method is useful as a motion-compensated reference, but it requires online flow if deployed directly.

\subsection{No-RAFT Adaptive EMA: Raw-vs-Memory Disagreement}

The first no-flow adaptive variant uses pixel-wise disagreement between the current raw prediction and previous temporal memory:

\begin{align}
    r_t &= \operatorname{clip}\left(\frac{|d_t - \hat{d}_{t-1}|}{s_{\mathrm{diff}}}, 0, 1\right),\\
    \alpha_t &= \alpha_{\min} + (\alpha_{\max} - \alpha_{\min})r_t,\\
    \hat{d}_t &= \alpha_t d_t + (1-\alpha_t)\hat{d}_{t-1}.
\end{align}

The intended interpretation is:

\begin{itemize}[leftmargin=*]
    \item If raw and memory agree, trust memory more.
    \item If raw and memory disagree, trust raw more.
\end{itemize}

The tested sweep was:

\begin{itemize}[leftmargin=*]
    \item $\alpha_{\min} \in \{0.25,0.30,0.35,0.40\}$.
    \item $\alpha_{\max} \in \{0.70,0.80,0.90\}$.
    \item $s_{\mathrm{diff}} \in \{2.0,3.0,5.0,8.0\}$.
\end{itemize}

\subsection{No-RAFT Adaptive EMA: Disagreement + Disparity Gradient}

The second no-flow adaptive variant adds disparity-gradient protection:

\begin{align}
    r_{\mathrm{diff}} &= \operatorname{clip}\left(\frac{|d_t - \hat{d}_{t-1}|}{s_{\mathrm{diff}}}, 0, 1\right),\\
    r_{\mathrm{grad}} &= \operatorname{clip}\left(\frac{\|\nabla d_t\|}{s_{\mathrm{grad}}}, 0, 1\right),\\
    r_t &= \frac{w_{\mathrm{diff}}r_{\mathrm{diff}} + w_{\mathrm{grad}}r_{\mathrm{grad}}}{w_{\mathrm{diff}} + w_{\mathrm{grad}}},\\
    \alpha_t &= \alpha_{\min} + (\alpha_{\max}-\alpha_{\min})r_t,\\
    \hat{d}_t &= \alpha_t d_t + (1-\alpha_t)\hat{d}_{t-1}.
\end{align}

The intended interpretation is:

\begin{itemize}[leftmargin=*]
    \item Smooth more on flat/stable tissue regions.
    \item Trust raw predictions more near disparity edges and discontinuities.
\end{itemize}

The tested sweep was:

\begin{itemize}[leftmargin=*]
    \item $\alpha_{\min} \in \{0.25,0.30,0.35\}$.
    \item $\alpha_{\max} \in \{0.75,0.85,0.90\}$.
    \item $s_{\mathrm{diff}} \in \{3.0,5.0\}$.
    \item $s_{\mathrm{grad}} \in \{3.0,5.0,8.0\}$.
    \item $w_{\mathrm{grad}} \in \{0.25,0.50,1.00\}$.
\end{itemize}

\section{Metrics}

The benchmark reports both geometric and temporal metrics.

\subsection{Geometric Metrics}

The main geometric metrics are:

\begin{itemize}[leftmargin=*]
    \item Depth MAE in millimeters.
    \item Disparity MAE in pixels.
    \item Disparity RMSE in pixels.
    \item Bad-1px, Bad-2px, Bad-3px.
    \item Bad-1mm, Bad-2mm, Bad-5mm.
    \item Valid frame count.
    \item Valid pixel coverage.
\end{itemize}

\subsection{Temporal Metrics}

The temporal diagnostics are:

\begin{itemize}[leftmargin=*]
    \item Raw temporal disparity difference:
    \[
        | \hat{d}_t - \hat{d}_{t-1} |.
    \]
    \item Motion-compensated temporal MAE:
    \[
        | \hat{d}_t - W(\hat{d}_{t-1}, F_{t-1\rightarrow t}) |.
    \]
\end{itemize}

The second metric uses RAFT flow as a \textbf{metric-only} tool for no-RAFT methods. It does not imply that no-RAFT methods use flow in their prediction formula.

\subsection{Runtime Metrics}

Runtime is reported fairly according to the expected deployment cost:

\begin{itemize}[leftmargin=*]
    \item No-flow methods:
    \[
        T_{\mathrm{online}} = T_{\mathrm{S2M2}} + T_{\mathrm{postprocess}}.
    \]
    \item Forward-flow methods:
    \[
        T_{\mathrm{online}} = T_{\mathrm{S2M2}} + T_{\mathrm{flow,fwd}} + T_{\mathrm{postprocess}}.
    \]
    \item Forward-backward-flow methods:
    \[
        T_{\mathrm{online}} = T_{\mathrm{S2M2}} + T_{\mathrm{flow,fwd}} + T_{\mathrm{flow,bwd}} + T_{\mathrm{postprocess}}.
    \]
\end{itemize}

The current measured averages are:

\begin{itemize}[leftmargin=*]
    \item S2M2-S inherited runtime: \SI{60.97}{ms}.
    \item RAFT forward runtime: \SI{71.63}{ms}.
    \item RAFT backward runtime: \SI{68.84}{ms}.
    \item S2M2-S peak VRAM: \SI{371.33}{MB}.
    \item RAFT peak VRAM: \SI{3429.09}{MB}.
    \item StereoAnyVideo cached runtime: \SI{146.30}{ms}.
    \item StereoAnyVideo peak VRAM: \SI{10132.18}{MB}.
\end{itemize}

\section{Main Results}

\subsection{Key Method Comparison}

Table~\ref{tab:main-results} summarizes the most important rows from the benchmark.

\begin{table}[htbp]
\centering
\caption{Main comparison of representative methods on the audited SCARED temporal-GT sequence. Lower is better for all error and temporal metrics.}
\label{tab:main-results}
\small
\begin{tabularx}{\textwidth}{lrrrrr}
\toprule
Method & Depth MAE mm & Disp MAE px & MC Temp MAE px & Runtime ms & Peak VRAM MB \\
\midrule
S2M2-S raw & 2.5840 & 8.3835 & 1.7244 & 60.97 & 371 \\
Fixed EMA $\alpha=0.35$ & \textbf{2.4516} & \textbf{8.2498} & \textbf{1.4950} & \textbf{63.16} & \textbf{371} \\
Fixed EMA $\alpha=0.50$ & 2.4902 & 8.2511 & 1.5754 & 63.11 & 371 \\
RAFT-warped EMA $\alpha=0.50$ & 2.4377 & 8.2217 & 1.4826 & 157.03 & 3429 \\
Conf./occ. reset $\alpha=0.50$ & 2.5160 & 8.2827 & 1.6218 & 212.78 & 3429 \\
Conservative adaptive EMA with RAFT & 2.4612 & 8.2077 & 1.5246 & 214.40 & 3429 \\
StereoAnyVideo & 2.5874 & 8.2502 & 1.6618 & 146.30 & 10132 \\
Best no-RAFT adaptive diff+grad & 2.4955 & 8.2813 & 1.5575 & 70.87 & 371 \\
\bottomrule
\end{tabularx}
\end{table}

\subsection{Best No-RAFT Adaptive Variant}

The best no-RAFT adaptive variant observed in the sweep is approximately:

\begin{itemize}[leftmargin=*]
    \item Method: adaptive no-RAFT diff+grad.
    \item $\alpha_{\min}=0.25$.
    \item $\alpha_{\max}=0.75$.
    \item $s_{\mathrm{diff}}=3.0$.
    \item $s_{\mathrm{grad}}=8.0$.
    \item $w_{\mathrm{grad}}=1.0$.
\end{itemize}

Its metrics are:

\begin{itemize}[leftmargin=*]
    \item Depth MAE: \SI{2.4955}{mm}.
    \item Disparity MAE: \SI{8.2813}{px}.
    \item Motion-compensated temporal MAE: \SI{1.5575}{px}.
    \item Estimated online runtime: \SI{70.87}{ms}.
    \item Peak VRAM: \SI{371}{MB}.
\end{itemize}

\subsection{Fixed EMA vs Best No-RAFT Adaptive}

Table~\ref{tab:fixed-vs-adaptive} shows the decisive comparison.

\begin{table}[htbp]
\centering
\caption{Fixed EMA remains better than the best no-RAFT adaptive variant.}
\label{tab:fixed-vs-adaptive}
\small
\begin{tabularx}{\textwidth}{lrrrr}
\toprule
Method & Depth MAE mm & Disp MAE px & MC Temp MAE px & Runtime ms \\
\midrule
Fixed EMA $\alpha=0.35$ & \textbf{2.4516} & \textbf{8.2498} & \textbf{1.4950} & \textbf{63.16} \\
Best no-RAFT adaptive diff+grad & 2.4955 & 8.2813 & 1.5575 & 70.87 \\
\bottomrule
\end{tabularx}
\end{table}

The conclusion is direct:

\begin{quote}
\textbf{No-RAFT adaptive EMA did not outperform fixed EMA.}
\end{quote}

The adaptive method is worse in depth accuracy, disparity accuracy, temporal consistency, and runtime.

\section{Interpretation}

\subsection{Fixed EMA Is a Strong Baseline, Not a Novel Method}

The fixed EMA result is important but not itself a strong methodological novelty. It shows that S2M2-S predictions contain temporally coherent signal that can be improved with simple causal smoothing.

Compared with S2M2-S raw, fixed EMA $\alpha=0.35$ improves:

\begin{itemize}[leftmargin=*]
    \item Depth MAE from \SI{2.5840}{mm} to \SI{2.4516}{mm}.
    \item Disparity MAE from \SI{8.3835}{px} to \SI{8.2498}{px}.
    \item Motion-compensated temporal MAE from \SI{1.7244}{px} to \SI{1.4950}{px}.
\end{itemize}

This is a very strong deployment baseline because it adds only about \SI{2.2}{ms} of post-processing and does not require flow or a video model.

However, fixed EMA is too simple to be claimed as the main methodological contribution.

\subsection{Why No-RAFT Adaptive Failed}

The no-RAFT adaptive rules attempted to use local risk estimates to decide how much memory to keep. The underlying assumption was:

\begin{quote}
If raw and memory disagree, trust raw more; if they agree, trust memory more.
\end{quote}

This assumption is too weak in surgical temporal stereo because disagreement can mean many different things:

\begin{itemize}[leftmargin=*]
    \item True scene motion.
    \item Camera motion.
    \item Tissue deformation.
    \item S2M2-S flicker.
    \item Specular instability.
    \item Disparity-edge motion.
    \item Occlusion/disocclusion.
    \item Noise or invalid predictions.
\end{itemize}

Without explicit motion compensation or learned temporal context, the adaptive rule cannot reliably distinguish these cases. As a result, it often uses too much raw prediction or too much memory in the wrong regions. The fixed EMA, although simpler, is more robust because it does not overreact to noisy risk maps.

\subsection{RAFT-Warped EMA Is the Quality Reference}

RAFT-warped EMA remains the strongest quality reference:

\begin{itemize}[leftmargin=*]
    \item Depth MAE: \SI{2.4377}{mm}.
    \item Disparity MAE: \SI{8.2217}{px}.
    \item Motion-compensated temporal MAE: \SI{1.4826}{px}.
\end{itemize}

Compared with fixed EMA, the gain is real but small:

\begin{itemize}[leftmargin=*]
    \item Depth MAE improves by about \SI{0.014}{mm}.
    \item MC temporal MAE improves by about \SI{0.012}{px}.
\end{itemize}

The cost difference is much larger:

\begin{itemize}[leftmargin=*]
    \item Fixed EMA runtime: about \SI{63}{ms}.
    \item RAFT-warped EMA runtime: about \SI{157}{ms} in the latest fair accounting.
    \item Fixed EMA peak VRAM: about \SI{371}{MB}.
    \item RAFT-warped EMA peak VRAM: about \SI{3429}{MB}.
\end{itemize}

This means RAFT motion compensation is useful, but full RAFT is not ideal for lightweight deployment.

\subsection{StereoAnyVideo Is Not Competitive Here}

StereoAnyVideo is useful as a cached video-model comparison, but on this audited sequence it is not superior to fixed EMA:

\begin{itemize}[leftmargin=*]
    \item Depth MAE is worse than fixed EMA.
    \item MC temporal MAE is worse than fixed EMA.
    \item Runtime is higher.
    \item VRAM is much higher.
\end{itemize}

This is an important empirical observation: a general video stereo/depth model is not automatically better than a strong single-frame stereo model plus simple temporal smoothing in this surgical setting.

\section{Current Scientific Status}

The current findings can be summarized as:

\begin{enumerate}[leftmargin=*]
    \item S2M2-S is a strong single-frame baseline.
    \item Fixed EMA is a surprisingly strong temporal baseline.
    \item No-RAFT adaptive heuristics do not beat fixed EMA.
    \item RAFT-warped EMA is the best quality reference but too expensive online.
    \item StereoAnyVideo is not superior on this audited sequence.
\end{enumerate}

Therefore:

\begin{quote}
\textbf{The current work has strong empirical insights, but the final methodological novelty is not yet fixed.}
\end{quote}

\section{Implications for Novelty}

\subsection{What Is Not Enough}

The following are not sufficient as the main novelty:

\begin{itemize}[leftmargin=*]
    \item Fixed EMA alone.
    \item Raw-vs-memory adaptive EMA alone.
    \item Disparity-gradient-protected adaptive EMA alone.
\end{itemize}

These are valuable baselines and diagnostics, but not enough for a strong method paper.

\subsection{Potential Novel Direction 1: Fast Flow-Based Temporal Fusion}

The first next direction is to replace full RAFT with a faster optical-flow backend and evaluate whether motion compensation can become practical.

A minimal next experiment is:

\begin{itemize}[leftmargin=*]
    \item RAFT-Small with 6 iterations.
    \item RAFT-Small with 12 iterations.
    \item Warped EMA with $\alpha=0.50$.
\end{itemize}

This would answer:

\begin{quote}
\emph{Is full RAFT necessary, or can RAFT-Small recover most of the temporal benefit at much lower cost?}
\end{quote}

The comparison table should include:

\begin{itemize}[leftmargin=*]
    \item S2M2-S raw.
    \item Fixed EMA $\alpha=0.35$.
    \item Best no-RAFT adaptive.
    \item RAFT-Small 6-iters warped EMA.
    \item RAFT-Small 12-iters warped EMA.
    \item Full RAFT warped EMA.
    \item StereoAnyVideo.
\end{itemize}

This is a clean engineering/scientific baseline, but it is probably not the strongest novelty because it still uses optical flow online.

\subsection{Potential Novel Direction 2: RAFT Teacher to Lightweight Student}

The more novel direction is to use the expensive RAFT-based method as an offline teacher and train a lightweight recurrent student for online inference.

The core idea is:

\begin{center}
\textbf{offline RAFT-based temporal teacher} $\rightarrow$ \textbf{online flow-free lightweight student}.
\end{center}

At training time:

\begin{itemize}[leftmargin=*]
    \item Use raw S2M2-S predictions.
    \item Use RAFT-warped EMA outputs as temporal teacher targets or regularizers.
    \item Use ground-truth disparity where available.
\end{itemize}

At inference time:

\begin{itemize}[leftmargin=*]
    \item Use S2M2-S raw prediction.
    \item Use a small ConvGRU or alpha/reset fusion module.
    \item Do not run RAFT.
\end{itemize}

A safe student design would predict pixel-wise fusion weights rather than an unconstrained residual:

\begin{align}
    \alpha_t &= f_\theta(d_t, \hat{d}_{t-1}, I_t, h_{t-1}),\\
    \hat{d}_t &= \alpha_t d_t + (1-\alpha_t)\hat{d}_{t-1}.
\end{align}

The student could be supervised with a mixed loss:

\begin{equation}
    \mathcal{L}
    =
    \lambda_{\mathrm{gt}}\mathcal{L}_{\mathrm{gt}}
    +
    \lambda_{\mathrm{teacher}}\mathcal{L}_{\mathrm{teacher}}
    +
    \lambda_{\mathrm{raw}}\mathcal{L}_{\mathrm{raw\ guard}}
    +
    \lambda_{\mathrm{tv}}\mathcal{L}_{\mathrm{alpha\ smooth}}.
\end{equation}

A reasonable initial weighting would be:

\begin{itemize}[leftmargin=*]
    \item $\lambda_{\mathrm{gt}}=1.0$.
    \item $\lambda_{\mathrm{teacher}}=0.2$.
    \item $\lambda_{\mathrm{raw}}=0.05$.
    \item $\lambda_{\mathrm{tv}}=0.01$.
\end{itemize}

The key claim would be:

\begin{quote}
\emph{A lightweight recurrent student can approximate motion-compensated temporal behavior learned from an expensive RAFT teacher without requiring optical flow at inference.}
\end{quote}

This is stronger as a methodological novelty, but it requires more data and a proper train/validation protocol.

\section{Recommended Next Steps}

\subsection{Immediate Decision}

The no-RAFT adaptive branch should be closed as a negative result:

\begin{quote}
\emph{Simple no-flow adaptive fusion based on raw-memory disagreement and disparity gradients does not improve over fixed EMA.}
\end{quote}

The fixed EMA $\alpha=0.35$ should remain the deployment baseline.

\subsection{Next Controlled Experiment}

The next low-effort controlled experiment should be:

\begin{quote}
\textbf{Evaluate RAFT-Small 6/12 iterations as a lightweight motion-compensation backend for warped EMA.}
\end{quote}

This is the cleanest way to determine whether the current cost barrier comes from full RAFT specifically, or from online optical flow in general.

\subsection{Next Novel Method Direction}

If RAFT-Small still does not provide a satisfactory quality/runtime trade-off, the most meaningful method direction is:

\begin{quote}
\textbf{RAFT-warped EMA teacher $\rightarrow$ lightweight causal recurrent student.}
\end{quote}

This would turn the current empirical observations into a coherent method paper.

\section{Current Takeaway}

The current takeaway is:

\begin{quote}
\emph{Fixed causal EMA is a very strong deployment baseline for temporal refinement of S2M2-S on the audited SCARED sequence. Simple no-flow adaptive heuristics do not outperform it. RAFT-based motion compensation provides the best quality, but at significantly higher runtime and memory cost. The next meaningful step is therefore either to test lightweight flow backends such as RAFT-Small, or to distill the RAFT-based temporal behavior into a lightweight flow-free recurrent student.}
\end{quote}



\appendix

\section{Relevant Commands}

\subsection{Inspect Output Folder}

\begin{lstlisting}
ls -lh results/03_temporal_refinement/scared_s2m2_temporal_baselines_v2_no_raft_adaptive
\end{lstlisting}

\subsection{Inspect Summary CSV}

\begin{lstlisting}
cat results/03_temporal_refinement/scared_s2m2_temporal_baselines_v2_no_raft_adaptive/summary.csv
\end{lstlisting}

\subsection{Inspect Adaptive Sweep}

\begin{lstlisting}
cat results/03_temporal_refinement/scared_s2m2_temporal_baselines_v2_no_raft_adaptive/adaptive_sweep_summary.csv
\end{lstlisting}

\subsection{Inspect Videos}

\begin{lstlisting}
ls -lh results/03_temporal_refinement/scared_s2m2_temporal_baselines_v2_no_raft_adaptive/videos/
\end{lstlisting}

\section{Canonical Benchmark Folder}

\begin{lstlisting}
results/03_temporal_refinement/scared_s2m2_temporal_baselines_v2_no_raft_adaptive
\end{lstlisting}

\section{Files Involved}

\begin{lstlisting}
scripts/temporal_refinement/lib/temporal_baselines.py
scripts/temporal_refinement/eval_scripts/benchmark_scared_s2m2_temporal_baselines.py
scripts/temporal_refinement/lib/video_qualitative.py
\end{lstlisting}

\subsection{RAFT-Small Lightweight Flow Results}

After the failure of the no-RAFT adaptive EMA variants to outperform the fixed EMA baseline, we evaluated RAFT-Small as a lightweight optical-flow backend for motion-compensated temporal fusion. The goal was to determine whether the quality improvement of full RAFT-warped EMA depends on the full RAFT model, or whether a smaller and faster RAFT variant is sufficient.

The evaluated RAFT-Small variants were:

\begin{itemize}[leftmargin=*]
\item RAFT-Small with 6 update iterations.
\item RAFT-Small with 12 update iterations.
\end{itemize}

Both were used inside the same warped EMA formulation:

\begin{equation}
\hat{d}t = \alpha d_t + (1-\alpha)W(\hat{d}{t-1}, F_{t-1 \rightarrow t}),
\end{equation}

with $\alpha=0.50$, where $F_{t-1 \rightarrow t}$ is estimated by RAFT-Small rather than full RAFT.

The key result is that RAFT-Small essentially matches full RAFT quality while substantially reducing optical-flow runtime. In fact, the 6-iteration RAFT-Small variant slightly outperformed full RAFT in the current benchmark, although the difference is small enough to be treated as practically equivalent.

\begin{table*}[t]
\centering
\caption{Comparison between fixed EMA, no-RAFT adaptive EMA, RAFT-Small, full RAFT, and StereoAnyVideo. Lower is better for all error and temporal metrics.}
\label{tab:raft-small-results}
\scriptsize
\begin{tabular}{lrrrrr}
\toprule
\textbf{Method} & \textbf{Depth MAE} & \textbf{Disp MAE} & \textbf{MC Temp MAE} & \textbf{Runtime} & \textbf{Peak VRAM} \\
 & \textbf{mm} & \textbf{px} & \textbf{px} & \textbf{ms} & \textbf{MB} \\
\midrule
S2M2-S raw & 2.5840 & 8.3835 & 1.7244 & 60.97 & 371 \\
Fixed EMA $\alpha=0.35$ & 2.4516 & 8.2498 & 1.4950 & 63.13 & 371 \\
Best no-RAFT adaptive & 2.4955 & 8.2813 & 1.5575 & 74.25 & 371 \\
RAFT-Small 6-iters warped EMA & \textbf{2.4360} & \textbf{8.2189} & \textbf{1.4819} & 117.72 & 3393 \\
RAFT-Small 12-iters warped EMA & 2.4371 & 8.2206 & 1.4822 & 110.42 & 3394 \\
Full RAFT warped EMA & 2.4377 & 8.2217 & 1.4826 & 139.28 & 3429 \\
StereoAnyVideo & 2.5874 & 8.2502 & 1.6618 & 146.30 & 10132 \\
\bottomrule
\end{tabular}
\end{table*}

The most important observation is that full RAFT is not necessary to obtain the motion-compensated temporal gain. RAFT-Small with 6 iterations achieved:

\begin{itemize}[leftmargin=*]
\item Depth MAE: \SI{2.4360}{mm}.
\item Disparity MAE: \SI{8.2189}{px}.
\item Motion-compensated temporal MAE: \SI{1.4819}{px}.
\item Estimated online runtime: \SI{117.72}{ms}.
\item Peak VRAM: approximately \SI{3393}{MB}.
\end{itemize}

Compared with full RAFT-warped EMA, RAFT-Small 6-iters is practically equivalent in quality and faster in runtime:

\begin{itemize}[leftmargin=*]
\item Full RAFT forward-flow runtime: \SI{71.63}{ms}.
\item RAFT-Small 6-iters forward-flow runtime: \SI{33.09}{ms}.
\item RAFT-Small 12-iters forward-flow runtime: \SI{41.70}{ms}.
\end{itemize}

However, the memory reduction is limited:

\begin{itemize}[leftmargin=*]
\item Full RAFT peak VRAM: approximately \SI{3429}{MB}.
\item RAFT-Small peak VRAM: approximately \SI{3393}{MB}.
\end{itemize}

Therefore, the main advantage of RAFT-Small is reduced runtime, not reduced VRAM.

\paragraph{Interpretation.}
RAFT-Small 6-iters is currently the best quality method in the benchmark. It slightly improves over fixed EMA and full RAFT, while being faster than full RAFT. However, it is still much more expensive than fixed EMA in both runtime and memory. Compared with fixed EMA $\alpha=0.35$, RAFT-Small 6-iters improves Depth MAE by approximately \SI{0.0156}{mm} and motion-compensated temporal MAE by approximately \SI{0.0131}{px}, but increases runtime from approximately \SI{63.13}{ms} to \SI{117.72}{ms} and peak VRAM from approximately \SI{371}{MB} to approximately \SI{3393}{MB}.

This creates a clear quality–deployment trade-off:

\begin{quote}
\emph{Fixed EMA remains the best cheap deployment baseline, while RAFT-Small 6-iters is the best motion-compensated quality baseline.}
\end{quote}

The remaining research opportunity is to close the gap between these two extremes: obtain quality close to RAFT-Small while keeping runtime and memory closer to fixed EMA.

\section{Experiment Paths and Artifacts}

This section records the canonical paths of the current ARGOS temporal-refinement experiments.

\subsection{Audited Input Sequence}

\begin{lstlisting}
dataset/SCARED/curated/temporal_gt/test_dataset_9_keyframe_3
\end{lstlisting}

This folder contains the audited SCARED temporal-GT sequence used for all current cache-only temporal-refinement experiments.

\subsection{Cached S2M2-S Predictions}

\begin{lstlisting}
results/03_temporal_refinement/evaluation/gt_temporal_test_dataset_9_keyframe_3/predictions/S2M2-S_512
\end{lstlisting}

This folder contains the cached S2M2-S predictions used as the raw disparity input for the temporal-refinement benchmark.

\subsection{Cached StereoAnyVideo Predictions}

\begin{lstlisting}
results/03_temporal_refinement/evaluation/gt_temporal_test_dataset_9_keyframe_3/predictions/StereoAnyVideo_384x640
\end{lstlisting}

This folder contains the cached StereoAnyVideo predictions used as a video-model comparison. StereoAnyVideo is treated as a comparison/teacher candidate, not as ground truth.

\subsection{Full RAFT Flow Cache}

\begin{lstlisting}
results/04_dataset_derivatives/SCARED/temporal_gt_flow_cache/test_dataset_9_keyframe_3/raft
\end{lstlisting}

This folder contains the full RAFT optical-flow cache. For no-RAFT adaptive methods, this cache is used only for motion-compensated metric computation, not for prediction.

\subsection{No-RAFT Adaptive Benchmark Output}

\begin{lstlisting}
results/03_temporal_refinement/scared_s2m2_temporal_baselines_v2_no_raft_adaptive
\end{lstlisting}

This folder contains the no-RAFT adaptive EMA benchmark. It includes the fixed EMA baselines, no-flow adaptive sweeps, summary CSVs, and qualitative videos.

Important files:

\begin{lstlisting}
summary.csv
adaptive_sweep_summary.csv
per_frame_metrics.csv
per_pair_temporal_metrics.csv
method_config.json
cache_audit.json
README.md
run.log
videos/
\end{lstlisting}

The main conclusion from this experiment is that no-RAFT adaptive EMA variants did not outperform fixed EMA $\alpha=0.35$.

\subsection{RAFT-Small Benchmark Output}

\begin{lstlisting}
results/03_temporal_refinement/scared_s2m2_temporal_baselines_v3_raft_small
\end{lstlisting}

This folder contains the RAFT-Small benchmark comparing RAFT-Small 6-iters, RAFT-Small 12-iters, full RAFT, fixed EMA, best no-RAFT adaptive, S2M2-S raw, and StereoAnyVideo.

Important files:

\begin{lstlisting}
summary.csv
per_frame_metrics.csv
per_pair_temporal_metrics.csv
flow_cache_audit.json
method_config.json
README.md
run.log
videos/
\end{lstlisting}

The main conclusion from this experiment is that RAFT-Small 6-iters matches or slightly improves over full RAFT-warped EMA while reducing forward-flow runtime substantially. However, its VRAM cost remains close to full RAFT.

\subsection{Code Files}

The current benchmark code is organized around the following files:

\begin{lstlisting}
scripts/temporal_refinement/lib/temporal_baselines.py
scripts/temporal_refinement/eval_scripts/benchmark_scared_s2m2_temporal_baselines.py
scripts/temporal_refinement/lib/video_qualitative.py
scripts/temporal_refinement/lib/flow_cache.py
scripts/temporal_refinement/eval_scripts/build_scared_temporal_gt_flow_cache.py
\end{lstlisting}

The temporal baselines are implemented in:

\begin{lstlisting}
scripts/temporal_refinement/lib/temporal_baselines.py
\end{lstlisting}

The main cache-only benchmark entry point is:

\begin{lstlisting}
scripts/temporal_refinement/eval_scripts/benchmark_scared_s2m2_temporal_baselines.py
\end{lstlisting}

The flow-cache builder is:

\begin{lstlisting}
scripts/temporal_refinement/eval_scripts/build_scared_temporal_gt_flow_cache.py
\end{lstlisting}

\subsection{Artifact-Aware Evaluation of RAFT-Small}

The RAFT-Small experiment initially appeared highly favorable under standard global metrics. RAFT-Small 6-iters warped EMA achieved the best numerical performance among the evaluated methods:

\begin{itemize}[leftmargin=*]
\item Depth MAE: \SI{2.4360}{mm}.
\item Disparity MAE: \SI{8.2189}{px}.
\item Motion-compensated temporal MAE: \SI{1.4819}{px}.
\end{itemize}

However, qualitative video inspection revealed strong ghosting, blur, and loss of sharpness. This motivated a new artifact-aware evaluation pass. The goal was to test whether the numerically best method was also visually reliable enough for deployment.

The artifact-aware benchmark added the following metrics:

\begin{itemize}[leftmargin=*]
\item \textbf{Edge sharpness ratio}, measuring whether disparity edges are preserved or blurred.
\item \textbf{Boundary disparity MAE}, measuring error specifically near strong GT disparity boundaries.
\item \textbf{Motion-region ghosting score}, measuring disagreement with current-frame evidence in dynamic regions.
\item \textbf{GT-based ghosting error}, measuring whether ghosting regions are also wrong relative to GT.
\item \textbf{Occlusion/disocclusion MAE}, measuring error in regions where temporal warping is unreliable.
\item \textbf{Temporal lag rate}, measuring whether the prediction resembles the previous GT frame more than the current GT frame.
\item \textbf{RGB–disparity edge alignment}, measuring whether disparity edges remain aligned with image edges.
\end{itemize}

Table~\ref{tab:artifact-aware-results} summarizes the main artifact-aware findings.

\begin{table*}[t]
\centering
\tiny
\setlength{\tabcolsep}{3pt}
\caption{Artifact-aware comparison of temporal-refinement methods. Lower is better for error, ghosting, lag, and runtime. Edge sharpness values closer to 1 indicate stronger preservation of raw disparity edges.}
\label{tab:artifact-aware-results}
\resizebox{\textwidth}{!}{
\begin{tabular}{lrrrrrrrr}
\toprule
Method & Depth & Disp & MC Temp & EdgeSharp & Boundary & GhostGT & Lag & Runtime \\
 & MAE mm & MAE px & MAE px & RawEdges & MAE px & $\tau=2$ px & Rate & ms \\
\midrule
S2M2 raw & 2.5840 & 8.3835 & 1.7244 & 1.0000 & 13.4354 & 13.6517 & 0.5196 & 60.97 \\
Fixed EMA & 2.4516 & 8.2498 & 1.4950 & 0.6523 & 13.2133 & \textbf{10.6972} & 0.6667 & 62.59 \\
Best adaptive & 2.4955 & 8.2813 & 1.5575 & 0.8589 & 13.3289 & 11.0833 & 0.6373 & 73.19 \\
RAFT-Small 6 & \textbf{2.4360} & \textbf{8.2189} & \textbf{1.4819} & 0.7022 & 13.2474 & 11.2446 & 0.6471 & 120.60 \\
RAFT-Small 12 & 2.4371 & 8.2206 & 1.4822 & 0.7025 & 13.2494 & 11.2468 & 0.6471 & 110.24 \\
Full RAFT & 2.4377 & 8.2217 & 1.4826 & 0.7024 & 13.2501 & 11.2488 & 0.6471 & 140.00 \\
SAV & 2.5874 & 8.2502 & 1.6618 & 0.9020 & \textbf{12.7958} & 12.6591 & 0.5098 & 146.30 \\
\bottomrule
\end{tabular}
}
\end{table*}

The results show that RAFT-Small remains the best method under global numerical metrics, but it is not artifact-free. In particular, RAFT-Small 6 produces higher GT-based ghosting error than fixed EMA in motion regions:

\begin{itemize}[leftmargin=*]
\item Fixed EMA ghosting GT error at $\tau=2$: \SI{10.6972}{px}.
\item RAFT-Small 6 ghosting GT error at $\tau=2$: \SI{11.2446}{px}.
\end{itemize}

This confirms the qualitative observation: motion-compensated smoothing can reduce global error while still producing stale or ghosted geometry in dynamic regions.

The edge-sharpness results also show that all temporal smoothing methods reduce high-frequency disparity structure. Raw S2M2-S has an edge-sharpness ratio of 1.0 by definition. Fixed EMA drops to 0.6523 on strong raw edges, while RAFT-Small variants remain around 0.702. This suggests that RAFT-based warping preserves more edge structure than fixed EMA, but still introduces substantial smoothing compared with the raw prediction.

StereoAnyVideo has the best boundary MAE and relatively high edge preservation, but its global depth and temporal metrics are worse. This indicates that it may preserve some boundary structure, but does not provide the best metric geometry on this sequence.

Overall, the artifact-aware benchmark changes the interpretation of RAFT-Small:

\begin{quote}
\emph{RAFT-Small is the best numerical motion-compensated baseline, but should be treated as an artifact-prone teacher or upper bound rather than a direct deployment method.}
\end{quote}

This finding motivates selective distillation rather than naive teacher copying.

\subsection{From Artifact-Aware Benchmarking to Selective Distillation}

The artifact-aware results suggest that naive distillation from RAFT-Small would be risky. Although RAFT-Small improves global Depth MAE and motion-compensated temporal MAE, it can introduce ghosting and blur in visually critical regions. Therefore, the future student should not copy RAFT-Small everywhere.

The next step is to build an oracle teacher-selection benchmark. The oracle will choose, for each pixel or region, the best source among:

\begin{itemize}[leftmargin=*]
\item raw S2M2-S,
\item fixed EMA,
\item best no-RAFT adaptive EMA,
\item RAFT-Small warped EMA,
\item full RAFT warped EMA,
\item StereoAnyVideo.
\end{itemize}

This oracle is not a deployable method because it uses ground truth. Its purpose is to estimate the headroom available for a learned student. If the oracle significantly improves over RAFT-Small while avoiding artifact-prone regions, this would justify training an artifact-aware student.

The intended student direction is:

\begin{quote}
\emph{Use RAFT-Small as an offline motion-compensated teacher only where it is artifact-safe, while preserving raw or fixed-EMA behavior near boundaries, occlusions, and dynamic regions.}
\end{quote}

This leads to the following research hypothesis:

\begin{quote}
\emph{A lightweight causal student can approach the numerical performance of RAFT-Small while avoiding the ghosting, blur, and lag revealed by artifact-aware diagnostics.}
\end{quote}

The resulting method would not simply distill RAFT. It would perform artifact-aware selective temporal distillation.

\subsection{Oracle Teacher Selection: Meaning and Interpretation}

After the artifact-aware evaluation, we introduced an oracle teacher-selection benchmark. The goal of this experiment is not to define a deployable method, but to estimate the theoretical headroom available if a future model could learn when to trust each prediction source.

An oracle is an upper-bound baseline that uses privileged information, in this case the ground-truth disparity, to select the best prediction source. Therefore, it cannot be used at inference time. Its purpose is diagnostic: it answers the question of whether different methods contain complementary local information.

Given a set of candidate predictions
\[
\mathcal{M}
=
\{
d^{raw},
d^{fixed},
d^{adaptive},
d^{raft\text{-}small},
d^{raft},
d^{sav}
\},
\]
the pixel-wise oracle selects, for each valid pixel \(p\) at time \(t\), the method with the lowest absolute disparity error with respect to ground truth:
\[
m^*(t,p)
=
\arg\min_{m \in \mathcal{M}}
\left|
d_m(t,p) - d_{gt}(t,p)
\right|.
\]
The oracle prediction is then constructed as
\[
d_{oracle}(t,p)
=
d_{m^*(t,p)}(t,p).
\]

This means that the oracle may choose S2M2-S raw in one pixel, fixed EMA in another pixel, RAFT-Small in another region, and StereoAnyVideo elsewhere. Such a method is not deployable because it uses the ground truth to make the selection. However, it is extremely useful because it tells us whether a learned selector or student model has meaningful potential.

\paragraph{Key finding.}
The best single temporal baseline before oracle selection was RAFT-Small 6-iters warped EMA:
\[
\text{Depth MAE}_{raft\text{-}small}
=
\SI{2.4360}{mm}.
\]
The all-source pixel-wise oracle achieved:
\[
\text{Depth MAE}_{oracle}
=
\SI{2.1395}{mm}.
\]
This corresponds to a headroom of:
\[
\Delta_{\text{oracle}}
=
2.4360 - 2.1395
=
\SI{0.2965}{mm}.
\]

This is a substantial gap. It shows that no single prediction source is locally optimal across the whole surgical scene. Even though RAFT-Small is the best single method globally, other sources are better in many local regions.

\paragraph{Oracle selection distribution.}
The all-source oracle selected the candidate methods with the following average proportions:
\[
\begin{aligned}
\text{StereoAnyVideo} &: 29.59\%,\\
\text{Fixed EMA} &: 22.44\%,\\
\text{Raw S2M2-S} &: 17.63\%,\\
\text{RAFT-Small} &: 9.00\%,\\
\text{Full RAFT} &: 7.99\%.
\end{aligned}
\]

This result is important because RAFT-Small, despite being the best global baseline, is selected in only a small fraction of pixels. Conversely, StereoAnyVideo is selected very frequently despite having weak global depth performance. This suggests that StereoAnyVideo may preserve useful local structure or boundary information, while fixed EMA and raw S2M2-S retain useful information in other regions.

Therefore, the problem should not be framed as choosing a single best teacher. Instead, the results motivate a multi-teacher selection strategy.

\begin{quote}
\emph{The best temporal refinement strategy is likely not a single global method, but a learned selective fusion mechanism that decides when to trust raw current-frame evidence, temporal smoothing, motion-compensated refinement, or video-model priors.}
\end{quote}
\begin{table*}[t]
\centering
\scriptsize
\caption{Oracle teacher-selection results. Oracle methods use ground truth for source selection and are therefore upper bounds, not deployable methods. Lower is better for all error, ghosting, and lag metrics.}
\label{tab:oracle-selection-results}
\begin{tabular}{lrrrrrr}
\toprule
\textbf{Method} &
\textbf{Depth MAE} &
\textbf{Disp MAE} &
\textbf{MC Temp} &
\textbf{Boundary MAE} &
\textbf{GhostGT} &
\textbf{Lag Rate} \\
&
\textbf{mm} &
\textbf{px} &
\textbf{px} &
\textbf{px} &
\textbf{$\tau=2$ px} &
\\
\midrule
RAFT-Small 6 warped EMA &
2.4360 & 8.2189 & 1.4819 & 13.2474 & 11.2446 & 0.6471 \\
Oracle raw/fixed/RAFT-Small &
2.2079 & 7.3824 & \textbf{1.3547} & 11.8713 & 10.6597 & 0.1176 \\
Oracle all sources &
\textbf{2.1395} & \textbf{7.0901} & 1.4585 & \textbf{11.0762} & \textbf{10.2597} & \textbf{0.0000} \\
Region edge-aware oracle &
2.5046 & 8.2610 & 2.0156 & 13.4354 & 10.8990 & 0.5490 \\
Artifact-safe RAFT selector &
2.2365 & 7.4697 & 1.3887 & 12.0842 & 10.7794 & 0.1275 \\
\bottomrule
\end{tabular}
\end{table*}
\paragraph{Interpretation.}
The oracle results reveal substantial headroom above the best single temporal baseline. The all-source oracle improves Depth MAE from \SI{2.4360}{mm} to \SI{2.1395}{mm}, and Disparity MAE from \SI{8.2189}{px} to \SI{7.0901}{px}. This confirms that the evaluated prediction sources are complementary.

However, the pixel-wise oracle should not be interpreted as a deployable visual reconstruction method. Since it switches prediction source independently at each pixel, it can introduce artificial spatial discontinuities. This is reflected by the very high edge-sharpness ratios observed for the oracle outputs. These high values do not necessarily indicate better visual quality; instead, they may indicate source-switching artifacts.

Therefore, the oracle should be interpreted only as a theoretical upper bound. Its role is to quantify whether a learned, regularized, artifact-aware selector could improve over any individual teacher.

\paragraph{Artifact-safe selector.}
The artifact-safe RAFT selector provides a more constrained upper bound. It uses RAFT-Small only where it is more accurate than raw and fixed EMA, and where artifact diagnostics indicate that RAFT-Small is safe. This selector achieves:
\[
\text{Depth MAE}
=
\SI{2.2365}{mm},
\]
while selecting RAFT-Small in only:
\[
15.88\%
\]
of pixels.

This is highly informative. It suggests that motion-compensated RAFT-Small behavior is useful, but only in selected regions. Copying RAFT-Small everywhere would be suboptimal and may reproduce ghosting or blur artifacts.

The correct direction is therefore:
\begin{quote}
\emph{artifact-aware multi-teacher selective distillation, not naive RAFT-Small distillation.}
\end{quote}
\subsection{Current Evidence and Dataset Limitations}

The current results should be interpreted as a controlled pilot benchmark, not as a final general claim. All temporal-refinement, artifact-aware, and oracle-selection results reported so far are based on a single audited SCARED temporal-GT sequence:
\begin{lstlisting}
dataset/SCARED/curated/temporal_gt/test_dataset_9_keyframe_3
\end{lstlisting}

After validity filtering, the benchmark contains:
\[
103
\]
valid frames and:
\[
102
\]
valid temporal pairs.

This is sufficient to reveal important failure modes and to motivate the next methodological direction, but it is not sufficient to claim generality across surgical scenes, anatomies, lighting conditions, camera motions, datasets, or sensing setups.

The correct interpretation is therefore:
\begin{quote}
\emph{On this audited sequence, global metrics favor RAFT-Small, artifact-aware metrics reveal ghosting and blur, and oracle teacher selection shows substantial headroom for selective multi-teacher fusion.}
\end{quote}

The current results do not yet prove that the same oracle selection distribution, artifact behavior, or teacher complementarity will hold across all surgical datasets. In particular, the high selection rate of StereoAnyVideo, the artifact-prone behavior of RAFT-Small, and the magnitude of the oracle headroom must be validated on additional sequences.

A stronger benchmark protocol should include:
\begin{itemize}[leftmargin=*]
    \item multiple SCARED temporal-GT sequences;
    \item sequence-level splits;
    \item evaluation across different motion, occlusion, and boundary regimes;
    \item additional surgical datasets where possible;
    \item global geometric metrics;
    \item artifact-aware diagnostics;
    \item runtime and memory evaluation under a fair deployment model.
\end{itemize}

Thus, the current benchmark is best viewed as a diagnostic proof of concept. It identifies a clear research direction, but multi-sequence validation is required before making broad claims.

\paragraph{Next methodological implication.}
The oracle benchmark motivates the construction of a lightweight student model trained with selective multi-teacher supervision. The future student should learn to approximate the oracle selection behavior without access to ground truth at inference time. It should learn when to preserve raw current-frame evidence, when to use fixed temporal smoothing, and when to imitate motion-compensated teacher behavior.

This leads to the following hypothesis:
\begin{quote}
\emph{A lightweight causal student can approach the headroom exposed by oracle teacher selection while avoiding the ghosting, blur, and discontinuities of direct temporal smoothing or naive teacher copying.}
\end{quote}

\subsection{Rectified Temporal-GT Conversion from Raw SCARED Keyframes}
\label{sec:scared_rectified_temporal_gt}

To extend the temporal evaluation beyond the initial audited sequence, we investigated whether additional raw SCARED keyframes could be converted into temporally consistent stereo sequences with dense geometric supervision. The raw SCARED keyframes contain a vertically stacked stereo video, per-frame calibration metadata, and per-frame 3D scene point maps.

A first conversion extracted the top and bottom halves of the stereo video as left and right views, and derived depth supervision directly from the third channel of the scene point maps. For each valid point, disparity was initially computed as
\begin{equation}
    d = \frac{f_x B}{Z},
\end{equation}
where \(f_x\) is the focal length, \(B\) is the stereo baseline, and \(Z\) is the depth value from the scene point map.

This initial conversion produced visually plausible depth and validity overlays. For \texttt{dataset\_1\_keyframe\_1}, the converted sequence contained 197 frames with complete left/right images, depth maps, disparity maps, validity masks, and calibration files. The median depth was approximately \(70.3\,\mathrm{mm}\), the median disparity was approximately \(61.3\,\mathrm{px}\), and the mean valid-pixel ratio was \(46.3\%\). However, when S2M2-S at 512 px input width was evaluated on this unrectified conversion, the results were catastrophic:
\begin{equation}
    \mathrm{Disp.\ MAE} = 58.74\,\mathrm{px},
    \qquad
    \mathrm{Bad\text{-}1px} \approx 99.99\%,
    \qquad
    \mathrm{Depth\ MAE} = 2798.65\,\mathrm{mm}.
\end{equation}
Since the image shapes, prediction counts, masks, and calibration files were all valid, this failure was not caused by an I/O mismatch.

We traced the issue to a stereo geometry mismatch. The previous working temporal-GT sequence, \texttt{test\_dataset\_9\_keyframe\_3}, had been generated in rectified stereo coordinates and stored the rectified calibration matrices \(\mathbf{P}_1\), \(\mathbf{P}_2\), \(\mathbf{R}_1\), \(\mathbf{R}_2\), and \(\mathbf{Q}\). In contrast, the newly extracted \texttt{dataset\_1\_keyframe\_1} sequence was still in the raw top/bottom camera geometry. This is incompatible with standard stereo networks such as S2M2, RAFT-Stereo, and CREStereo, which assume rectified stereo pairs where disparity corresponds to horizontal displacement.

We therefore implemented a rectified temporal-GT converter. The converter first splits the vertically stacked raw stereo video into left and right views, then uses the per-frame calibration parameters \(\mathbf{K}_L\), \(\mathbf{K}_R\), \(\mathbf{D}_L\), \(\mathbf{D}_R\), \(\mathbf{R}\), and \(\mathbf{T}\) to compute stereo rectification using OpenCV. Rectification maps are produced with \texttt{cv2.stereoRectify} and \texttt{cv2.initUndistortRectifyMap}. The left and right images are then remapped into a common rectified coordinate system. The 3D scene points are projected into the rectified left camera, using a minimum-depth scatter operation to resolve many-to-one projections. The final rectified depth map, validity mask, and disparity map are saved in the same coordinate system as the rectified left image. Rectified disparity is then computed as
\begin{equation}
    d_{\mathrm{rect}} =
    \frac{f_{x,\mathrm{rect}} B_{\mathrm{rect}}}{Z_{\mathrm{rect}}}.
\end{equation}

The resulting rectified sequence is saved with the same layout convention as the previously validated temporal-GT sequence:
\begin{verbatim}
gt/DepthL_float32
gt/Disparity_float32
gt/ValidMask
\end{verbatim}
Each frame also stores a rectified calibration file containing \(\mathbf{P}_1\), \(\mathbf{P}_2\), \(\mathbf{R}_1\), \(\mathbf{R}_2\), \(\mathbf{Q}\), \(f_x\), and the stereo baseline, together with the original calibration parameters. The unrectified conversion was kept as an audit artifact and was not overwritten.

To validate the rectified conversion, we reran S2M2-S at 512 px input width on the rectified version of \texttt{dataset\_1\_keyframe\_1}. The same model, evaluated on the rectified sequence, improved dramatically:
\begin{equation}
    \mathrm{Disp.\ MAE} = 0.891\,\mathrm{px},
    \qquad
    \mathrm{Disp.\ RMSE} = 1.687\,\mathrm{px},
\end{equation}
\begin{equation}
    \mathrm{Bad\text{-}1px} = 33.55\%,
    \qquad
    \mathrm{Bad\text{-}2px} = 6.56\%,
    \qquad
    \mathrm{Bad\text{-}3px} = 1.27\%,
\end{equation}
\begin{equation}
    \mathrm{Depth\ MAE} = 0.942\,\mathrm{mm}.
\end{equation}
The prediction count was 197/197, the prediction, ground-truth disparity, and valid-mask shapes all matched at \(1024 \times 1280\), no resizing was required for metric computation, and no warnings were produced. This confirmed that the previous catastrophic error was caused by missing rectification rather than by invalid raw data.

Finally, we implemented a batch wrapper for rectified temporal-GT conversion. The wrapper scans raw SCARED folders using the strict pattern
\begin{verbatim}
dataset/SCARED/raw/extracted/<dataset>/<dataset>/keyframe_*
\end{verbatim}
selects keyframes containing \texttt{rgb.mp4}, \texttt{frame\_data.tar.gz}, \texttt{scene\_points.tar.gz}, and stereo reference images, and calls the validated rectified converter for each candidate. Test datasets are skipped by default because the available raw test folders do not contain the required temporal scene-point archives. The wrapper produces per-sequence summaries, a conversion CSV, a skipped-sequence CSV, a batch summary JSON, and a README. This keeps dataset preparation, model inference, and temporal benchmarking as separate stages of the pipeline.

This validation establishes the correct preprocessing path for multi-sequence SCARED temporal evaluation:
\begin{equation}
\begin{split}
    \text{raw vertical stereo video}
    &\rightarrow
    \text{split left/right}
    \rightarrow
    \text{stereo rectification} \\
    &\rightarrow
    \text{GT reprojection}
    \rightarrow
    \text{rectified depth/disparity/mask}
    \rightarrow
    \text{stereo model evaluation}.
\end{split}
\end{equation}

\begin{table}[t]
\centering
\caption{Effect of stereo rectification on S2M2-S@512 evaluation for \texttt{dataset\_1\_keyframe\_1}. The unrectified conversion is geometrically incompatible with standard disparity-based stereo models, while the rectified conversion produces meaningful errors.}
\label{tab:rectification_sanity_check}
\begin{tabular}{lccccc}
\toprule
Conversion &
Disp. MAE $\downarrow$ &
Disp. RMSE $\downarrow$ &
Bad-1px $\downarrow$ &
Bad-3px $\downarrow$ &
Depth MAE $\downarrow$ \\
\midrule
Unrectified raw split &
58.74 px &
59.70 px &
99.99\% &
99.99\% &
2798.65 mm \\
Rectified + GT reprojection &
0.891 px &
1.687 px &
33.55\% &
1.27\% &
0.942 mm \\
\bottomrule
\end{tabular}
\end{table}
\subsection{Towards Teacher-Oriented Lightweight Temporal Stereo Refinement}

The rectified multi-sequence evaluation shows that a strong single-frame stereo
foundation model, such as S2M2, can produce geometrically plausible disparities on
individual frames, but still suffers from temporal instability and artifact-sensitive
failure modes in surgical video. In particular, the model may preserve local
stereo geometry while exhibiting frame-to-frame flickering, boundary degradation,
and inconsistent behavior around tissue edges, specularities, occlusion-like
regions, and rapidly changing surgical structures.

This motivates moving beyond a purely frame-wise stereo benchmark. Instead of
asking whether a single model is globally better, we consider the problem of
\emph{local teacher selection}: different temporal or refinement strategies may be
preferable in different spatial and temporal regimes. A raw S2M2 prediction may
preserve sharp boundaries, an exponential moving average may reduce flicker in
stable regions, a motion-aware refinement may better handle temporal consistency,
and a video-based teacher may provide smoother behavior across difficult clips.
However, no single candidate is expected to dominate everywhere.

\paragraph{Oracle teacher selection as an upper bound.}
To quantify the potential benefit of local teacher selection, we use an oracle
analysis. Given a set of candidate disparity predictions
\[
\mathcal{D}_t = \{D_t^{(1)}, D_t^{(2)}, \ldots, D_t^{(K)}\},
\]
and a ground-truth disparity map $D_t^{gt}$, the pixel-wise oracle selects, for
each pixel $p$, the candidate with the lowest absolute disparity error:
\[
k^*(p,t) =
\arg\min_{k \in \{1,\ldots,K\}}
\left|D_t^{(k)}(p) - D_t^{gt}(p)\right|.
\]
The oracle disparity is therefore
\[
D_t^{oracle}(p) = D_t^{(k^*(p,t))}(p).
\]

This oracle is not deployable, since it requires ground truth at inference time.
Its purpose is diagnostic: it measures the maximum achievable gain if a perfect
selector could choose the best local teacher. In previous small-scale oracle
experiments, the oracle improved over the raw S2M2 baseline, showing that there
is local headroom. Importantly, the oracle analysis also showed that the best
candidate depends on the region type: raw predictions may be preferable near
edges, smoothing may help in stable tissue regions, and motion-aware or video
teachers may help in temporally unstable regions. This indicates that the problem
is not simply solved by selecting one globally superior model.

\paragraph{Artifact-aware teacher analysis.}
To understand where each teacher is useful, we evaluate not only standard
geometric metrics, but also temporal and artifact-aware metrics. These include
standard disparity and depth errors:
\[
\mathrm{MAE}_{disp}
=
\frac{1}{|\Omega|}
\sum_{p \in \Omega}
\left|D_t(p)-D_t^{gt}(p)\right|,
\]
as well as bad-pixel rates, depth MAE, and RMSE. In addition, we consider
artifact-sensitive diagnostics such as boundary disparity error, edge sharpness
ratio, temporal disparity variation, motion-region error proxies, lag proxies,
RGB-disparity edge correlation, and ghosting/occlusion indicators when the
required motion information is available.

For the new streaming rectified benchmark, the evaluation is deliberately
performed without dense prediction caches. Predictions are computed on the fly
and immediately consumed for metric computation. Temporal metrics are computed
using only the previous valid frame in memory. Flow-dependent metrics, such as
motion-compensated temporal error, are not reported unless optical flow is
explicitly computed; otherwise, same-pixel no-flow proxies are used and clearly
reported as such.

\paragraph{From oracle to lightweight distillation.}
The oracle itself is not intended as the final method. Instead, it provides
supervision for a deployable lightweight refinement module. The central question
becomes:
\[
\textit{Can a small online module recover part of the oracle/teacher headroom
without running heavy teachers at inference time?}
\]

Heavy teachers such as RAFT-based motion compensation or video stereo models can
be useful offline, but they are too expensive for a real-time surgical pipeline.
Therefore, we use them only during training or analysis to generate targets. At
inference time, the goal is to keep the system close to:
\[
\text{S2M2 raw prediction}
\quad + \quad
\text{lightweight temporal refinement head}.
\]

A practical formulation is to train a small refinement head that predicts either
a residual correction, a confidence/artifact mask, or blending weights among
cheap online candidates. For residual distillation, the module predicts
\[
\Delta D_t = D_t^{target} - D_t^{S2M2},
\]
where $D_t^{target}$ may be an oracle prediction, a selected teacher, or a
teacher-refined disparity. The final refined disparity is then
\[
D_t^{refined} = D_t^{S2M2} + \Delta \hat{D}_t.
\]
Alternatively, the module may predict a confidence map $C_t$ and apply the
correction only where the raw stereo prediction is unreliable:
\[
D_t^{refined}
=
D_t^{S2M2}
+
C_t \cdot \Delta \hat{D}_t.
\]

Another lightweight option is to predict blending weights over cheap online
candidates:
\[
D_t^{refined}
=
w_{raw}D_t^{raw}
+
w_{ema}D_t^{ema}
+
w_{adaptive}D_t^{adaptive},
\]
where
\[
w_{raw} + w_{ema} + w_{adaptive} = 1.
\]
In this setting, expensive teachers such as RAFT or video stereo are used only
offline to supervise the desired behavior, while the online system uses only
fast S2M2-derived candidates.

\paragraph{DINOv3 as a tissue- and artifact-aware prior.}
A promising extension is to condition the lightweight refinement module on dense
foundation features extracted from a vision backbone such as DINOv3. The goal is
not to replace stereo geometry with monocular depth. Instead, DINOv3 features can
provide semantic and structural priors that are difficult to obtain from disparity
alone. In surgical scenes, these features may encode tissue appearance,
instrument/tissue separation, specular regions, boundaries, texture changes, and
other cues correlated with stereo failure modes.

The proposed architecture is therefore:
\[
\begin{aligned}
(I_t^L, I_t^R) &\rightarrow \text{S2M2} \rightarrow D_t^{S2M2}, \\
I_t^L &\rightarrow \text{DINOv3 encoder} \rightarrow F_t^{DINO}, \\
(D_t^{S2M2}, D_{t-1}^{S2M2}, F_t^{DINO}, \phi_t)
&\rightarrow \text{lightweight refinement head}
\rightarrow
(D_t^{refined}, C_t),
\end{aligned}
\]
where $\phi_t$ denotes additional cheap features such as disparity gradients,
RGB gradients, temporal disparity differences, edge maps, and artifact proxies.

This design separates geometry from semantic/artifact reasoning. S2M2 remains
responsible for stereo matching, while DINOv3 provides dense tissue-aware features
that help the refinement head decide where to preserve sharp stereo boundaries,
where to smooth flicker, and where to reduce confidence. The DINOv3 encoder can
be frozen and evaluated at reduced resolution, while only a small decoder or
refinement head is trained. This keeps the online cost significantly lower than
running heavy temporal teachers.

\paragraph{Minimal experimental plan.}
Before training a full refinement model, we first test whether DINOv3 features
are predictive of S2M2 failure modes. A small diagnostic head can be trained to
predict high-error or high-flicker regions from DINO features and simple stereo
features:
\[
(F_t^{DINO}, D_t^{S2M2}, \nabla D_t^{S2M2}, |D_t^{S2M2}-D_{t-1}^{S2M2}|)
\rightarrow
\hat{E}_t,
\]
where $\hat{E}_t$ is a predicted error, artifact, or low-confidence map. If these
features correlate with disparity error, boundary error, or temporal instability,
then DINO-guided refinement is justified.

The final method can then be trained using oracle or teacher-derived targets:
\[
\mathcal{L}
=
\lambda_{disp}\mathcal{L}_{disp}
+
\lambda_{temp}\mathcal{L}_{temp}
+
\lambda_{edge}\mathcal{L}_{edge}
+
\lambda_{conf}\mathcal{L}_{conf}
+
\lambda_{distill}\mathcal{L}_{distill}.
\]
Here, $\mathcal{L}_{distill}$ encourages the lightweight model to recover part of
the oracle or teacher improvement, while $\mathcal{L}_{edge}$ discourages
oversmoothing across tissue boundaries and $\mathcal{L}_{temp}$ penalizes
unwanted flickering.

\paragraph{Scientific interpretation.}
This reframes the contribution from simply benchmarking stereo models to
understanding and exploiting local teacher complementarity. The key hypothesis is
that surgical stereo errors are structured: they depend on tissue appearance,
boundaries, motion, specularities, occlusions, and temporal instability. Heavy
teachers and oracle selection reveal where improvements are possible, while a
DINO-guided lightweight module aims to distill this behavior into a deployable
real-time refinement stage.

In summary, the proposed direction is:
\[
\text{fast stereo foundation model}
+
\text{offline teacher/oracle analysis}
+
\text{DINO-guided lightweight distillation}
\rightarrow
\text{temporally stable surgical depth}.
\]

\subsection{Rectified Multi-Sequence Evaluation and Teacher-Oriented Distillation}

After validating the rectified SCARED temporal-GT conversion on a single clean
sequence, we extended the evaluation to the full rectified multi-sequence subset.
The final rectified temporal-GT dataset contains 27 complete sequences and
22,950 frames. An integrity audit identified 20,621 frames suitable for evaluation
after filtering low-valid or suspicious frames, while 2,329 frames were skipped.
All evaluated sequences share a consistent rectified stereo representation with
calibrated projection matrices and rectified ground-truth disparity/depth maps.

\paragraph{Streaming S2M2-S evaluation.}
We evaluated S2M2-S in a streaming, no-cache setting. Predictions were computed
on the fly and immediately consumed for metric computation, avoiding dense
prediction storage. The evaluation was performed over all 27 sequences and
20,621 valid frames, without saving per-frame prediction arrays. This avoided an
estimated 100.69 GiB of prediction cache storage.

The full S2M2-S artifact-temporal evaluation produced the following aggregate
results:
\[
\begin{aligned}
\mathrm{MAE}_{disp} &= 6.829~\mathrm{px}, \\
\mathrm{Bad}\text{-}3\mathrm{px} &= 42.72\%, \\
\mathrm{MAE}_{depth} &= 3.544~\mathrm{mm}.
\end{aligned}
\]
The runtime was approximately 6,521.3 seconds for the full evaluation, with a
median runtime of 62.45 ms per frame and a peak GPU memory usage of 395 MB.

Compared with the initial clean sequence, where S2M2-S achieved a disparity MAE
of 0.891 px and a Bad-3px rate of 1.27\%, the full multi-sequence benchmark
reveals substantially broader degradation. This indicates that the original
single-sequence result was not representative of the full rectified SCARED
temporal setting.

\paragraph{Artifact-aware and temporal diagnostics.}
In addition to standard geometry metrics, we computed temporal and artifact-aware
diagnostics. Since the full streaming run avoids optical flow, all temporal
metrics are reported as same-pixel, no-flow proxies. Motion-compensated metrics,
flow-based ghosting, and flow-based occlusion metrics are intentionally not
computed in this pass.

The main temporal and artifact-aware aggregate results were:
\[
\begin{aligned}
\mathrm{TemporalPairs} &= 20{,}594, \\
\mathrm{TemporalCoverage} &= 1.0, \\
\Delta D^{samepixel}_{temporal} &= 1.648~\mathrm{px}, \\
\mathrm{MotionMismatch}^{samepixel} &= 2.026~\mathrm{px}, \\
\mathrm{BoundaryMAE}_{disp} &= 8.982~\mathrm{px}, \\
\mathrm{LagRate} &= 0.656.
\end{aligned}
\]

These diagnostics show that S2M2-S is not only affected by frame-wise geometric
errors, but also by temporal instability, boundary-sensitive errors, and
lag/flicker-like behavior. In particular, the boundary disparity MAE is higher
than the global disparity MAE, suggesting that naive temporal smoothing may
reduce flicker but risks degrading sharp tissue boundaries.

\paragraph{Failure taxonomy and clip selection.}
Using the full S2M2-S artifact-temporal run, we generated a distillation planning
set. Sequences were grouped using quantile-based thresholds into:
\[
\{\text{clean core},\ \text{high boundary error},\ \text{high temporal flicker},
\ \text{high motion mismatch},\ \text{low-valid stress},\
\text{catastrophic geometry}\}.
\]
From this taxonomy, six compact contiguous clips were selected for future
teacher/oracle distillation. The selected clips cover clean and stress cases and
contain 502 frames in total. This subset is designed to be small enough for
offline teacher generation while still representing the dominant failure modes
observed in the full benchmark.

\paragraph{Teacher/oracle target generation.}
To measure the local headroom available from temporal and video teachers, we
generated selected-clip teacher/oracle distillation targets. Unlike the full
streaming evaluation, this stage is allowed to use heavier teachers offline, but
only on the selected clips. The online method is not expected to run these
teachers at inference time.

For each selected frame, the following candidate predictions were generated:
\[
\mathcal{D}_t =
\{
D_t^{raw},
D_t^{ema},
D_t^{adaptive},
D_t^{raft},
D_t^{sav}
\},
\]
where \(D_t^{raw}\) is the raw S2M2-S prediction, \(D_t^{ema}\) is a fixed EMA
candidate with \(\alpha=0.35\), \(D_t^{adaptive}\) is the best previous no-RAFT
adaptive candidate, \(D_t^{raft}\) is a RAFT-Small warped EMA candidate, and
\(D_t^{sav}\) is a StereoAnyVideo prediction.

StereoAnyVideo was processed in temporal chunks of 32 frames. The last short
chunk was padded by repeating the last valid frame, and padded outputs were
discarded after inference. This allowed SAV to be used consistently on clips
longer than its internal temporal window.

Given the set of candidates, the pixel-wise oracle selects the candidate with
minimum ground-truth disparity error:
\[
k^*(p,t)
=
\arg\min_{k}
\left|D_t^{(k)}(p) - D_t^{gt}(p)\right|,
\]
and the oracle disparity is:
\[
D_t^{oracle}(p) = D_t^{(k^*(p,t))}(p).
\]

We compute three oracle variants:
\[
\begin{aligned}
D_t^{oracle,no\text{-}flow}
&=
\mathrm{Oracle}
\left(
D_t^{raw}, D_t^{ema}, D_t^{adaptive}
\right), \\
D_t^{oracle,raft}
&=
\mathrm{Oracle}
\left(
D_t^{raw}, D_t^{ema}, D_t^{raft}
\right), \\
D_t^{oracle,all}
&=
\mathrm{Oracle}
\left(
D_t^{raw}, D_t^{ema}, D_t^{adaptive}, D_t^{raft}, D_t^{sav}
\right).
\end{aligned}
\]

The final selected-clip teacher/oracle run processed all 6 clips and 502 frames.
RAFT-Small and SAV were available for all clips. The target outputs were saved as
compressed downsampled arrays at scale 0.25, using float16 for dense maps and
uint8 for masks and labels. No full-resolution teacher cache was created.

The selected-clip aggregate results were:
\[
\begin{aligned}
\mathrm{MAE}_{raw} &= 13.2051~\mathrm{px}, \\
\mathrm{MAE}_{oracle,no\text{-}flow} &= 11.9220~\mathrm{px}, \\
\mathrm{MAE}_{oracle,all} &= 8.9466~\mathrm{px}.
\end{aligned}
\]
Therefore, the no-flow oracle improves over raw S2M2-S by:
\[
13.2051 - 11.9220 = 1.2831~\mathrm{px},
\]
corresponding to approximately:
\[
\frac{1.2831}{13.2051} \times 100 \approx 9.7\%.
\]
The all-source oracle improves over raw S2M2-S by:
\[
13.2051 - 8.9466 = 4.2585~\mathrm{px},
\]
corresponding to approximately:
\[
\frac{4.2585}{13.2051} \times 100 \approx 32.2\%.
\]

This result confirms that local teacher selection exposes substantial headroom.
Cheap no-flow candidates provide a moderate improvement, while the inclusion of
RAFT-Small and StereoAnyVideo as offline teachers provides a much larger upper
bound.

\paragraph{Motivation for lightweight distillation.}
The oracle itself is not deployable, since it uses ground truth to select the
best teacher. Similarly, RAFT-Small and StereoAnyVideo are too expensive to run
as part of the intended online surgical stereo pipeline. Their role is therefore
limited to offline supervision.

The selected-clip teacher/oracle experiment motivates a lightweight distillation
strategy. The goal is to train a small online refinement module that recovers
part of the oracle improvement without running heavy teachers at inference time.
A natural formulation is residual distillation:
\[
\Delta D_t^{target}
=
D_t^{oracle,all} - D_t^{raw},
\]
where the model predicts:
\[
\Delta \hat{D}_t
=
f_{\theta}
\left(
I_t^L,
D_t^{raw},
D_{t-1}^{raw},
\phi_t
\right),
\]
with \(\phi_t\) denoting cheap temporal and artifact features such as disparity
gradients, RGB gradients, temporal disparity differences, valid masks, and
boundary indicators.

The refined disparity can then be written as:
\[
D_t^{refined}
=
D_t^{raw} + C_t \cdot \Delta \hat{D}_t,
\]
where \(C_t\) is a predicted confidence or correction mask. This formulation
allows the network to correct unstable or artifact-prone regions while preserving
raw S2M2 predictions where they are already reliable.

A second formulation is blending-weight distillation:
\[
D_t^{refined}
=
w_{raw}D_t^{raw}
+
w_{ema}D_t^{ema}
+
w_{adaptive}D_t^{adaptive},
\]
with
\[
w_{raw} + w_{ema} + w_{adaptive} = 1.
\]
In this case, the heavy teachers are used only to define the desired behavior
during training, while the online model blends only cheap S2M2-derived
candidates.

\paragraph{Compressed temporal context.}
To improve temporal consistency without online optical flow or video teachers, we
also consider a compressed causal temporal context window. Instead of providing
only the current disparity and the immediately previous frame, the lightweight
refiner may receive a small history:
\[
\mathcal{C}_t =
\left[
D_t^{raw},
D_{t-1}^{raw},
D_{t-2}^{raw},
D_{t-3}^{raw}
\right],
\]
stored or reconstructed at reduced resolution. Additional temporal statistics
can be computed from this window:
\[
\begin{aligned}
\mu_t &= \mathrm{mean}(\mathcal{C}_t), \\
m_t &= \mathrm{median}(\mathcal{C}_t), \\
\sigma_t^2 &= \mathrm{var}(\mathcal{C}_t), \\
\Delta_t^{prev} &= |D_t^{raw} - D_{t-1}^{raw}|, \\
\Delta_t^{median} &= |D_t^{raw} - m_t|.
\end{aligned}
\]
These compressed temporal features can help the refiner distinguish true
geometric change from flicker. The design remains causal and lightweight, making
it more suitable for surgical robotics than online RAFT or full video stereo.

\paragraph{DINO-guided feature probing.}
In parallel, we identified a second direction based on dense self-supervised
visual features. Rather than using DINO as a monocular depth model, we treat DINO
features as tissue-, boundary-, and artifact-aware priors. The hypothesis is that
frozen DINO features extracted from the left image encode appearance cues
correlated with S2M2 failure modes, such as tissue boundaries, specularities,
instrument regions, and texture changes.

The proposed diagnostic probe is:
\[
\left(
F_t^{DINO},
D_t^{raw},
\nabla D_t^{raw},
|D_t^{raw} - D_{t-1}^{raw}|
\right)
\rightarrow
\hat{E}_t,
\]
where \(F_t^{DINO}\) are frozen DINO features and \(\hat{E}_t\) is a predicted
error, artifact, or low-confidence map. The probe should first be evaluated on a
small balanced subset of frames selected from the full S2M2 artifact-temporal
run. The goal is to test whether DINO features can predict high-error,
high-boundary-error, high-flicker, or high-motion-mismatch regions.

If the probe is successful, DINO features can be integrated into the lightweight
refinement architecture:
\[
\begin{aligned}
(I_t^L, I_t^R) &\rightarrow \mathrm{S2M2} \rightarrow D_t^{raw}, \\
I_t^L &\rightarrow \mathrm{DINO} \rightarrow F_t^{DINO}, \\
(D_t^{raw}, \mathcal{C}_t, F_t^{DINO}, \phi_t)
&\rightarrow
\mathrm{TinyRefiner}
\rightarrow
(D_t^{refined}, C_t).
\end{aligned}
\]
This keeps S2M2 responsible for stereo geometry, while DINO provides structural
and tissue-aware context for confidence prediction and residual correction.

\paragraph{Summary.}
The current results support the following research direction:
\[
\text{fast stereo foundation model}
+
\text{offline heavy teachers}
+
\text{oracle target generation}
+
\text{lightweight distillation}
\rightarrow
\text{temporally stable surgical depth}.
\]

The full S2M2-S evaluation shows that a fast stereo foundation model alone is not
sufficient on multi-sequence surgical video. The selected-clip oracle analysis
shows that RAFT-Small and StereoAnyVideo provide substantial offline teacher
headroom, reducing disparity MAE from 13.2051 px to 8.9466 px on the selected
distillation clips. The next step is therefore not to deploy heavy teachers, but
to distill their local behavior into a lightweight online temporal refinement
module, optionally guided by compressed temporal context and DINO-based
tissue/artifact features.

\subsection{From Selected-Clip Oracle Distillation to Full-Scale Refiner Training}

The selected-clip teacher/oracle experiment demonstrated that heavy video
teachers, especially StereoAnyVideo, expose substantial local headroom over raw
S2M2 predictions. However, training a lightweight refiner only on the selected
oracle clips is not sufficient for generalization. The selected distillation set
contains only 6 clips and 502 frames, and early tiny-refiner experiments showed
that a model can learn the oracle residual on validation-like data but fails to
generalize reliably to held-out failure modes.

For this reason, we separate the refinement problem into two stages. The first
stage uses the full rectified temporal-GT dataset to learn a general geometric
correction of S2M2. The second stage uses the selected teacher/oracle clips to
inject failure-mode-specific and temporally informed behavior from heavy offline
teachers such as StereoAnyVideo.

\paragraph{Full-scale S2M2--GT target generation.}
We generated a large-scale supervised training target dataset using all valid
frames from the rectified temporal-GT SCARED subset. Unlike the selected-clip
teacher/oracle generation, this stage does not use RAFT, StereoAnyVideo, DINO, or
any other heavy teacher. Instead, it uses only S2M2-S predictions and the
available ground-truth disparity.

For each valid stereo frame pair, we run S2M2-S once on the rectified left/right
images:
\[
(I_t^L, I_t^R) \xrightarrow{\mathrm{S2M2}} D_t^{raw}.
\]
The raw S2M2 disparity, ground-truth disparity, and valid mask are then
downsampled to the target training resolution using the same valid-mask-aware
downsampling procedure introduced for the corrected oracle targets. The training
target is the residual between the low-resolution ground truth and the
low-resolution raw S2M2 prediction:
\[
\Delta D_t^{GT} = D_t^{GT} - D_t^{raw}.
\]
Thus, the large-scale training set contains compact low-resolution tensors:
\[
\left\{
D_t^{raw},\ D_t^{GT},\ M_t,\ \Delta D_t^{GT}
\right\},
\]
where \(M_t\) is the downsampled valid mask. Full-resolution S2M2 predictions are
not cached.

The generated full-scale target dataset contains:
\[
\begin{aligned}
N_{\mathrm{seq}} &= 27, \\
N_{\mathrm{frames}} &= 20{,}621, \\
N_{\mathrm{shards}} &= 27, \\
\mathrm{storage} &= 2.9~\mathrm{GB}.
\end{aligned}
\]
The raw low-resolution S2M2 metrics over this dataset are:
\[
\begin{aligned}
\mathrm{MAE}_{raw} &= 7.025~\mathrm{px}, \\
\mathrm{Bad}\text{-}1\mathrm{px}_{raw} &= 73.57\%, \\
\mathrm{Bad}\text{-}3\mathrm{px}_{raw} &= 44.15\%, \\
\overline{r}_{valid} &= 45.53\%.
\end{aligned}
\]
The target generation took approximately 3650.5 seconds, corresponding to about
60.8 minutes, with an average runtime of 58.18 ms per frame. This is comparable
to the previous full S2M2 streaming evaluation, but instead of discarding the
prediction after metric computation, the pipeline stores only compact
low-resolution supervision targets.

\paragraph{Geometry-first lightweight refinement.}
The first full-scale refiner is formulated as a geometric residual correction
module. Given the raw S2M2 prediction, the model predicts a residual correction:
\[
\widehat{\Delta D}_t = f_{\theta}(D_t^{raw}, M_t),
\]
and the refined disparity is:
\[
D_t^{refined} = D_t^{raw} + \widehat{\Delta D}_t.
\]
The model is supervised directly with the ground-truth residual:
\[
\mathcal{L}_{GT}
=
\frac{1}{|M_t|}
\sum_{p \in M_t}
\rho
\left(
\widehat{\Delta D}_t(p) -
\Delta D_t^{GT}(p)
\right),
\]
where \(\rho(\cdot)\) is a robust regression loss such as masked \(L_1\) or
Charbonnier loss.

This stage is intentionally teacher-free. Its goal is not yet to imitate
StereoAnyVideo or RAFT, but to answer a simpler and more fundamental question:
can a small residual model systematically improve S2M2 on held-out surgical
sequences when trained on the full 20k-frame temporal-GT target set?

The expected evaluation compares raw and refined predictions against the saved
low-resolution ground truth:
\[
\begin{aligned}
\mathrm{MAE}_{raw}
&=
\left|D_t^{raw} - D_t^{GT}\right|, \\
\mathrm{MAE}_{refined}
&=
\left|D_t^{refined} - D_t^{GT}\right|.
\end{aligned}
\]
A successful first-stage refiner should satisfy:
\[
\mathrm{MAE}_{refined} < \mathrm{MAE}_{raw},
\]
and should also reduce Bad-1px and Bad-3px rates on held-out sequences.

\paragraph{Why this stage does not yet solve temporal consistency.}
The full-GT residual refiner is primarily a frame-wise geometric correction
module. If the input is only the current raw disparity \(D_t^{raw}\), the model
can correct systematic spatial errors, but it has limited ability to reason about
temporal instability, flicker, lag, or frame-to-frame inconsistency. Therefore,
this first stage should be interpreted as a geometry refinement baseline rather
than a complete temporal refinement method.

The temporal problem requires the model to observe a causal history of S2M2
predictions:
\[
\mathcal{C}_t =
\left[
D_{t-3}^{raw},
D_{t-2}^{raw},
D_{t-1}^{raw},
D_t^{raw}
\right].
\]
From this context, additional temporal features can be derived:
\[
\begin{aligned}
\mu_t &= \mathrm{mean}(\mathcal{C}_t), \\
m_t &= \mathrm{median}(\mathcal{C}_t), \\
\sigma_t^2 &= \mathrm{var}(\mathcal{C}_t), \\
\Delta_t^{prev} &= |D_t^{raw} - D_{t-1}^{raw}|, \\
\Delta_t^{median} &= |D_t^{raw} - m_t|.
\end{aligned}
\]
A temporal refiner can then be written as:
\[
\widehat{\Delta D}_t
=
f_{\theta}
\left(
D_t^{raw},
D_{t-1}^{raw},
D_{t-2}^{raw},
D_{t-3}^{raw},
\phi_t
\right),
\]
where \(\phi_t\) may include temporal statistics, valid masks, disparity
gradients, RGB gradients, or artifact indicators.

This formulation allows the model to distinguish true geometric change from
unstable prediction flicker. In deployment, the system remains causal:
\[
(I_t^L, I_t^R)
\rightarrow
\mathrm{S2M2}
\rightarrow
D_t^{raw}
\rightarrow
\mathrm{TemporalRefiner}
\rightarrow
D_t^{refined}.
\]

\paragraph{Role of StereoAnyVideo and oracle distillation.}
StereoAnyVideo is not intended to be run online in the final lightweight
pipeline. Instead, it is used as an offline teacher. The selected-clip oracle
analysis showed that StereoAnyVideo is selected in a large fraction of pixels and
provides most of the additional all-source oracle headroom over raw S2M2,
especially in difficult failure modes such as catastrophic geometry, motion
mismatch, and low-valid stress cases.

However, the teacher/oracle clips are too few to train a general model from
scratch. Therefore, the proposed strategy is:
\[
\text{pretrain on full GT targets}
\quad \rightarrow \quad
\text{fine-tune or distill on selected oracle/SAV targets}.
\]
The full-GT stage provides broad generalization, while the selected
teacher/oracle stage provides failure-mode-specific supervision.

A hybrid objective can be formulated as:
\[
\mathcal{L}
=
\mathcal{L}_{GT}^{full}
+
\lambda_{teacher}\mathcal{L}_{oracle}^{selected}
+
\lambda_{temp}\mathcal{L}_{temporal},
\]
where:
\[
\mathcal{L}_{GT}^{full}
\]
is computed over the full 20k-frame S2M2--GT target dataset,
\[
\mathcal{L}_{oracle}^{selected}
\]
is computed over the selected clips using the all-available oracle residual,
and
\[
\mathcal{L}_{temporal}
\]
encourages frame-to-frame stability.

The selected oracle residual is:
\[
\Delta D_t^{oracle}
=
D_t^{oracle,all} - D_t^{raw},
\]
where \(D_t^{oracle,all}\) is computed from the local best candidate among raw
S2M2, fixed EMA, adaptive no-RAFT, RAFT-Small, and StereoAnyVideo:
\[
D_t^{oracle,all}(p)
=
D_t^{(k^*)}(p),
\quad
k^*
=
\arg\min_k
\left|
D_t^{(k)}(p) - D_t^{GT}(p)
\right|.
\]
This allows the lightweight model to learn when to behave like raw S2M2, when to
apply a small geometric correction, and when to imitate the behavior of a
stronger temporal/video teacher.

\paragraph{Experimental roadmap.}
The refinement experiments are organized as a sequence of controlled ablations:
\[
\begin{array}{ll}
\textbf{v1:} &
D_t^{raw} \rightarrow \Delta D_t^{GT}
\quad
\text{frame-wise geometry refiner}, \\
\textbf{v2:} &
D_{t-k:t}^{raw} \rightarrow \Delta D_t^{GT}
\quad
\text{causal temporal context refiner}, \\
\textbf{v3:} &
\text{v2} + \mathcal{L}_{oracle}^{selected}
\quad
\text{SAV/oracle teacher distillation}, \\
\textbf{v4:} &
\text{v3} + F_t^{DINO}
\quad
\text{semantic/artifact-aware refinement if useful}.
\end{array}
\]

The first stage establishes whether a lightweight residual model can improve
S2M2 on held-out surgical sequences. The temporal stage then evaluates whether
compressed causal context reduces flicker and motion mismatch. Finally, the
SAV/oracle distillation stage tests whether heavy offline video teachers can
improve difficult failure modes without being deployed at inference time.

\subsection{Intermediate Findings and Motivation for Temporal Refinement}

The full-scale S2M2--GT target dataset enabled the first generalizable training
experiments for a lightweight residual refiner. This dataset contains 27
sequences and 20,621 valid frames, stored as compact low-resolution sequence
shards. Each target contains the raw S2M2 disparity, the downsampled
ground-truth disparity, the valid mask, and the residual target
\[
\Delta D_t^{GT} = D_t^{GT} - D_t^{raw}.
\]
This setup allows the refiner to be trained without rerunning S2M2, RAFT,
StereoAnyVideo, or DINO during training.

The first full-GT refiner was an unconstrained residual model:
\[
D_t^{refined} = D_t^{raw} + f_{\theta}(D_t^{raw}, M_t).
\]
Although this model reduced the training error, it did not generalize reliably.
On the held-out test split, the improvement was marginal:
\[
4.335~\mathrm{px} \rightarrow 4.228~\mathrm{px},
\]
while on the validation split the model degraded the raw prediction:
\[
4.677~\mathrm{px} \rightarrow 5.280~\mathrm{px}.
\]
The Bad-3px rate also worsened substantially on validation:
\[
44.42\% \rightarrow 56.16\%.
\]
This indicates that an unconstrained residual model tends to overcorrect regions
where S2M2 is already accurate.

To address this, we evaluated a conservative gated formulation:
\[
D_t^{refined}
=
D_t^{raw}
+
G_t \cdot R_t,
\]
where \(G_t \in [0,1]\) is a learned correction gate and \(R_t\) is a bounded
residual proposal. The goal was to preserve the identity mapping where raw S2M2
is reliable and only apply corrections in high-error regions.

A first gated model initialized the gate close to zero, but during training the
gate progressively opened on both good and bad pixels. This caused validation
degradation, suggesting that gate sparsity alone is insufficient. A second
safe-gated variant introduced explicit gate supervision using the raw S2M2 error:
\[
G_t^{target}(p)
=
\begin{cases}
0, & |D_t^{raw}(p) - D_t^{GT}(p)| < 1~\mathrm{px}, \\
1, & |D_t^{raw}(p) - D_t^{GT}(p)| \geq 3~\mathrm{px}.
\end{cases}
\]
This made the model safer: at the first epoch it remained almost identical to
raw S2M2,
\[
4.700959~\mathrm{px} \rightarrow 4.701217~\mathrm{px},
\]
but subsequent epochs again opened the gate and degraded validation performance.
The gate classifier achieved only weak discrimination between good and bad
regions, with AUC values around \(0.54\)--\(0.59\).

These experiments suggest that frame-wise raw disparity alone is not sufficient
to robustly predict where corrections should be applied. In other words, the
model cannot reliably infer from \(D_t^{raw}\) alone whether a local structure is
a true geometric surface, a temporal flicker artifact, a boundary error, or a
failure mode caused by low valid-mask support. This motivates moving from
frame-wise refinement to temporal-context refinement.

\paragraph{Conclusion from v1 experiments.}
The v1 experiments establish three important findings:
\[
\begin{aligned}
&\text{(1) A lightweight refiner can technically learn residual corrections,} \\
&\text{(2) unconstrained correction risks damaging already-correct pixels,} \\
&\text{(3) single-frame disparity is insufficient for reliable correction gating.}
\end{aligned}
\]
Therefore, the next refinement stage must provide the model with temporal
context.

\subsection{Temporal-Context Refiner}

The next model uses a causal window of recent S2M2 predictions:
\[
\mathcal{C}_t =
\left[
D_{t-3}^{raw},
D_{t-2}^{raw},
D_{t-1}^{raw},
D_t^{raw}
\right].
\]
From this window, the model can derive temporal consistency cues:
\[
\begin{aligned}
\mu_t &= \mathrm{mean}(\mathcal{C}_t), \\
m_t &= \mathrm{median}(\mathcal{C}_t), \\
\sigma_t^2 &= \mathrm{var}(\mathcal{C}_t), \\
\Delta_t^{prev} &= |D_t^{raw} - D_{t-1}^{raw}|, \\
\Delta_t^{median} &= |D_t^{raw} - m_t|.
\end{aligned}
\]
The temporal refiner is then formulated as:
\[
(\widehat{\Delta D}_t, G_t)
=
f_{\theta}
\left(
D_t^{raw},
D_{t-1}^{raw},
D_{t-2}^{raw},
D_{t-3}^{raw},
\mu_t,
m_t,
\sigma_t^2,
\Delta_t^{prev},
\Delta_t^{median},
M_t
\right),
\]
with final prediction:
\[
D_t^{refined}
=
D_t^{raw}
+
G_t \cdot \widehat{\Delta D}_t.
\]

This formulation gives the model access to the information that was missing from
the frame-wise variants. If a region is stable across recent frames, the gate can
remain closed and preserve the raw S2M2 prediction. If a region flickers,
collapses, or deviates strongly from the recent temporal context, the gate can
open and apply a correction.

This stage directly addresses the temporal nature of the surgical stereo problem
while remaining lightweight and causal. Heavy video teachers such as
StereoAnyVideo are still not used online. Instead, they remain offline teachers
for a later distillation stage after the temporal refiner has established a
generalizable baseline.

\subsection{Safe Refinement via Abstention and Hard-Pixel Crop Training}

The previous frame-wise and temporal-context residual refiners showed that dense
correction is unsafe. Even when a learned gate is introduced, the model tends to
open the gate on both good and bad pixels, degrading regions where raw S2M2 is
already accurate. A failure analysis confirmed that the main limitation is not
model capacity alone, but unsafe correction. In particular, pixels with raw error
above 20 px represent only 8.9\% of valid pixels but contribute 51.4\% of the
total MAE. As a result, full-frame residual training is dominated by outliers and
encourages corrections that are too aggressive for easier validation sequences.

To address this, we introduce a conservative abstention-based refiner. Instead
of always modifying the raw S2M2 prediction, the model predicts both a bad-pixel
probability and a residual proposal:
\[
(P_t^{bad}, \widehat{\Delta D}_t)
=
f_{\theta}
\left(
D_{t-3:t}^{raw},
M_t,
\phi_t
\right),
\]
where \(D_{t-3:t}^{raw}\) is a causal temporal context window of raw S2M2
predictions, \(M_t\) is the valid mask, and \(\phi_t\) contains additional
low-cost features such as temporal variance, temporal difference, and spatial
disparity gradients.

The model output is interpreted as a correction policy rather than a dense
refinement field. The default action is to preserve the raw S2M2 prediction:
\[
D_t^{refined}(p) = D_t^{raw}(p).
\]
A residual correction is applied only when the predicted probability of raw S2M2
failure exceeds a threshold:
\[
D_t^{refined}(p)
=
\begin{cases}
D_t^{raw}(p) + \widehat{\Delta D}_t(p),
&
P_t^{bad}(p) \geq \tau, \\
D_t^{raw}(p),
&
P_t^{bad}(p) < \tau.
\end{cases}
\]
A soft variant is also available during training:
\[
D_t^{soft}(p)
=
D_t^{raw}(p)
+
P_t^{bad}(p)\,\widehat{\Delta D}_t(p).
\]

The residual proposal is bounded to avoid large destructive corrections:
\[
\widehat{\Delta D}_t
=
s \cdot \tanh(R_t),
\]
where \(s\) is a fixed residual scale. This prevents the model from producing
arbitrarily large disparity shifts.

\paragraph{Balanced hard-pixel crop training.}
Rather than training only on full frames, which are dominated by a small number
of large-error outliers, the abstention model is trained on balanced crops. Crops
are sampled from multiple categories:
\[
\begin{aligned}
&\text{hard/outlier crops}, \\
&\text{good/stable crops}, \\
&\text{boundary/high-gradient crops}, \\
&\text{temporal-unstable crops}, \\
&\text{random crops}.
\end{aligned}
\]
This ensures that the model sees difficult regions often enough to learn a
correction policy, while also seeing enough good regions to learn when to
abstain.

The bad-pixel target is derived directly from the raw S2M2 error:
\[
Y_t^{bad}(p)
=
\mathbb{1}
\left[
|D_t^{raw}(p) - D_t^{GT}(p)| \geq 3~\mathrm{px}
\right].
\]
Good pixels are defined as:
\[
Y_t^{good}(p)
=
\mathbb{1}
\left[
|D_t^{raw}(p) - D_t^{GT}(p)| < 1~\mathrm{px}
\right].
\]

The training objective combines bad-pixel detection, residual correction, and
preservation of already-correct regions:
\[
\mathcal{L}
=
\lambda_{det}\mathcal{L}_{det}
+
\lambda_{res}\mathcal{L}_{res}
+
\lambda_{ref}\mathcal{L}_{ref}
+
\lambda_{pres}\mathcal{L}_{pres}
+
\lambda_{sparse}\mathcal{L}_{sparse}.
\]
Here, \(\mathcal{L}_{det}\) is a BCE or focal loss for bad-pixel detection,
\(\mathcal{L}_{res}\) supervises the residual proposal on hard pixels,
\(\mathcal{L}_{ref}\) compares the soft refined prediction to ground truth,
\(\mathcal{L}_{pres}\) penalizes changes on raw-good pixels, and
\(\mathcal{L}_{sparse}\) discourages modifying too many pixels.

\paragraph{Preliminary result.}
The abstention model uses a 4-frame causal context and 16 input channels. The
current implementation contains 194,818 parameters and is trained on balanced
\(96 \times 96\) crops. Unlike previous residual refiners, this model can improve
validation MAE and Bad-3px while modifying only a small fraction of pixels.

At epoch 5, the validation result was:
\[
\begin{aligned}
\mathrm{MAE}_{raw} &= 5.1996~\mathrm{px}, \\
\mathrm{MAE}_{refined} &= 5.1538~\mathrm{px}, \\
\mathrm{Bad}\text{-}3\mathrm{px}_{raw} &= 38.878\%, \\
\mathrm{Bad}\text{-}3\mathrm{px}_{refined} &= 38.449\%.
\end{aligned}
\]
Only 3.73\% of pixels were modified at the selected threshold. This is the first
refiner variant that improves both MAE and Bad-3px without broadly damaging
already-correct S2M2 regions.

This result supports the hypothesis that surgical stereo refinement should be
formulated as a conservative correction policy rather than a dense residual
prediction problem. The refiner should abstain by default and intervene only in
regions that are likely to contain S2M2 failure.

\subsection{Staged Abstention as Safe Refinement Policy}

The previous refinement experiments showed that improving S2M2 predictions is not
simply a matter of increasing model capacity. Dense residual refinement and
jointly trained gated refinement both suffered from unsafe correction behavior:
the model learned to reduce large training errors, but also modified regions
where raw S2M2 was already accurate. This resulted in overcorrection and
introduced new Bad-3px errors on reliable pixels.

A failure analysis confirmed this behavior. The full target dataset is strongly
affected by high-error outliers: pixels with raw S2M2 error above 20 px account
for only 8.9\% of valid pixels, but contribute 51.4\% of the total MAE. As a
result, standard full-frame residual training is dominated by a small fraction of
extreme-error pixels. The model is encouraged to apply broad corrections, even
though most pixels should remain close to the raw S2M2 prediction.

This motivated a change in formulation. Instead of treating refinement as dense
residual prediction, we formulate it as a staged detection-and-correction policy.
The model should first identify where S2M2 is likely to fail, and only then apply
a bounded correction. The default behavior is abstention, i.e. preserving the raw
S2M2 prediction.

The final policy is:
\[
D_t^{refined}(p)
=
\begin{cases}
D_t^{raw}(p) + \widehat{\Delta D}_t(p),
&
P_t^{bad}(p) \geq \tau, \\
D_t^{raw}(p),
&
P_t^{bad}(p) < \tau,
\end{cases}
\]
where \(P_t^{bad}(p)\) is the predicted probability that raw S2M2 is erroneous at
pixel \(p\), \(\widehat{\Delta D}_t(p)\) is the predicted residual correction,
and \(\tau\) is a validation-calibrated threshold. The residual is bounded:
\[
\widehat{\Delta D}_t = s \cdot \tanh(R_t),
\]
where \(s\) is a fixed residual scale. This prevents the refiner from applying
arbitrarily large disparity shifts.

\paragraph{Staged training.}
Unlike the previous jointly trained gated models, the abstention refiner is
trained in two stages.

First, we train a bad-pixel detector:
\[
P_t^{bad}
=
f_{\theta}^{det}
\left(
D_{t-3:t}^{raw}, M_t, \phi_t
\right),
\]
where \(D_{t-3:t}^{raw}\) is a causal temporal context window, \(M_t\) is the
valid mask, and \(\phi_t\) contains low-cost temporal and spatial features such
as temporal variance, temporal difference, and disparity gradients. The detector
target is defined directly from the raw S2M2 error:
\[
Y_t^{bad}(p)
=
\mathbb{1}
\left[
|D_t^{raw}(p) - D_t^{GT}(p)| \geq 3~\mathrm{px}
\right].
\]
The detector is optimized independently using a balanced crop sampling strategy
and a BCE/focal detection loss.

Second, we train the residual correction head:
\[
\widehat{\Delta D}_t
=
f_{\theta}^{res}
\left(
D_{t-3:t}^{raw}, M_t, \phi_t
\right),
\]
using the ground-truth residual:
\[
\Delta D_t^{GT} = D_t^{GT} - D_t^{raw}.
\]
During this stage, the detector is initially frozen, so that residual learning
does not destabilize the correction policy. The residual loss is concentrated on
hard pixels, while a strong preservation loss discourages modifications on
already-correct pixels:
\[
\mathcal{L}_{pres}
=
\sum_{p:\ |D_t^{raw}(p)-D_t^{GT}(p)|<1}
\left|
D_t^{refined}(p) - D_t^{raw}(p)
\right|.
\]
This explicitly encourages identity behavior where S2M2 is reliable.

\paragraph{Balanced crop training.}
The staged abstention model is trained on balanced crops rather than only full
frames. Crops are sampled from hard-error regions, good/stable regions,
high-gradient boundary regions, temporally unstable regions, and random regions.
This prevents a small number of large-error pixels from dominating the objective,
while still exposing the model to the failure cases that require correction.

\paragraph{Result.}
The staged abstention model is the first refiner variant that improves both
validation and held-out test performance while preserving reliable S2M2 regions.
The model uses a 4-frame causal temporal context, 16 input channels, and 194,818
parameters. The best residual epoch was epoch 7, with a calibrated correction
threshold of \(\tau=0.5\).

On the validation split:
\[
\begin{aligned}
\mathrm{MAE}_{raw} &= 5.1996~\mathrm{px}, \\
\mathrm{MAE}_{refined} &= 5.1167~\mathrm{px}, \\
\mathrm{Bad}\text{-}3\mathrm{px}_{raw} &= 38.878\%, \\
\mathrm{Bad}\text{-}3\mathrm{px}_{refined} &= 38.479\%.
\end{aligned}
\]

On the held-out test split:
\[
\begin{aligned}
\mathrm{MAE}_{raw} &= 4.6690~\mathrm{px}, \\
\mathrm{MAE}_{refined} &= 4.5637~\mathrm{px}, \\
\mathrm{Bad}\text{-}3\mathrm{px}_{raw} &= 33.536\%, \\
\mathrm{Bad}\text{-}3\mathrm{px}_{refined} &= 33.230\%.
\end{aligned}
\]

The test MAE improvement is:
\[
4.6690 - 4.5637 = 0.1053~\mathrm{px},
\]
corresponding to a relative improvement of:
\[
\frac{0.1053}{4.6690} \times 100 \approx 2.25\%.
\]
The Bad-3px rate improves by:
\[
33.536 - 33.230 = 0.306
\]
percentage points.

Although the absolute improvement is modest, the safety behavior is substantially
better than previous refiners. On the test split, the model modifies 21.45\% of
pixels and introduces new Bad-3px errors on only 0.151\% of raw-good pixels. The
bad-pixel detector also generalizes better than the jointly trained gates,
achieving:
\[
\mathrm{AUC}_{test} = 0.768,
\qquad
\mathrm{AP}_{test} = 0.693.
\]

\paragraph{Interpretation.}
This result supports the central hypothesis that surgical stereo refinement
should not be formulated as unconstrained dense residual prediction. Since raw
S2M2 is already reliable in many regions, the refiner must be conservative. The
successful formulation is therefore:
\[
\text{detect failure}
\quad \rightarrow \quad
\text{correct only when necessary}
\quad \rightarrow \quad
\text{otherwise preserve S2M2}.
\]

Compared with dense residual and jointly trained gated variants, staged
abstention provides a safer refinement mechanism. It reduces MAE and Bad-3px on
held-out sequences while avoiding widespread degradation of good raw predictions.
This makes it a stronger foundation for subsequent extensions, including larger
abstention backbones, teacher/oracle fine-tuning with StereoAnyVideo, and
semantic or artifact-aware features.

\paragraph{Final staged-abstention analysis.}
We further analyzed the best staged-abstention refiner on validation and held-out
test sequences. The final model contains 194,818 parameters and uses a 4-frame
causal temporal context with 16 input channels. The refiner is extremely light:
a refiner-only benchmark reports approximately 1.08 ms per frame on CUDA with
batch size 64.

On the held-out test split, the staged abstention policy improves both MAE and
Bad-3px:
\[
\begin{aligned}
\mathrm{MAE}_{raw} &= 4.6690~\mathrm{px}, \\
\mathrm{MAE}_{refined} &= 4.5637~\mathrm{px}, \\
\mathrm{Bad}\text{-}3\mathrm{px}_{raw} &= 33.536\%, \\
\mathrm{Bad}\text{-}3\mathrm{px}_{refined} &= 33.230\%.
\end{aligned}
\]
The refiner modifies 21.46\% of valid pixels. Importantly, it modifies 40.33\%
of raw bad pixels but only 8.74\% of raw good pixels, indicating that the learned
policy is selective rather than densely corrective. The fraction of raw-good
pixels converted into new Bad-3px errors is only 0.152\%, showing that the model
largely preserves reliable S2M2 regions.

The detector also generalizes to the held-out test split, achieving:
\[
\mathrm{AUC}_{test} = 0.768,
\qquad
\mathrm{AP}_{test} = 0.693.
\]
Across validation and test sequences, 4 sequences improve, 3 show mixed behavior,
and only 1 degrades. This confirms that the staged abstention strategy is more
robust than previous dense residual or jointly gated refiners.

These results support the hypothesis that surgical stereo refinement should be
formulated as a conservative detection-and-correction policy. Since raw S2M2 is
already accurate in many regions, the safest online strategy is not to predict a
dense correction everywhere, but to detect likely failure regions and apply a
bounded residual only there:
\[
D_t^{refined}(p)
=
\begin{cases}
D_t^{raw}(p) + \widehat{\Delta D}_t(p),
&
P_t^{bad}(p) \geq \tau, \\
D_t^{raw}(p),
&
P_t^{bad}(p) < \tau.
\end{cases}
\]
This preserves the real-time character of the original S2M2 pipeline while
adding a lightweight, temporally informed correction module.

\subsection{Hybrid Oracle Fine-Tuning}

The staged-abstention refiner provided the first safe improvement over raw S2M2.
However, its evaluation on the selected teacher/oracle clips showed that most of
the available oracle headroom remained unused. On the selected 502-frame subset,
raw S2M2 obtained an MAE of 11.3231 px, while the v3.1 staged-abstention refiner
reduced it to 11.0421 px. The all-available oracle reached 6.8062 px, meaning
that v3.1 recovered only 6.22\% of the available oracle gap.

At the same time, the detector remained strong on these difficult clips, with
AUC/AP of 0.842/0.845. This suggests that the model can identify many unreliable
S2M2 regions, but that its residual head does not yet exploit the teacher/oracle
signal sufficiently. We therefore introduced a hybrid oracle fine-tuning stage.

The goal of this stage is not to imitate StereoAnyVideo globally. Instead, we use
the all-available oracle as a local teacher only where it is clearly better than
raw S2M2. The oracle selects among the available candidates, including raw S2M2,
simple temporal baselines, RAFT-Small, and StereoAnyVideo. This prevents the
model from treating any single teacher as universally correct.

Let \(D^{raw}\) be the raw S2M2 disparity, \(D^{GT}\) the ground-truth disparity,
and \(D^{oracle}\) the all-available oracle disparity. Oracle supervision is
activated only on valid hard pixels where the oracle improves over raw S2M2:
\[
\mathcal{M}^{oracle}(p)
=
\mathbb{1}
\left[
|D^{raw}(p)-D^{GT}(p)|
-
|D^{oracle}(p)-D^{GT}(p)|
>
\epsilon
\right]
\cdot
\mathbb{1}
\left[
|D^{raw}(p)-D^{GT}(p)| \geq 3
\right].
\]
The oracle residual target is:
\[
\Delta D^{oracle}(p)
=
D^{oracle}(p) - D^{raw}(p).
\]

During hybrid fine-tuning, training batches are mixed between the full-GT target
set and the selected oracle clips:
\[
\mathcal{B}
=
0.75 \, \mathcal{B}_{GT}
+
0.25 \, \mathcal{B}_{oracle}.
\]
The full-GT batches preserve the general safety behavior of the v3.1 model,
while the oracle batches expose the residual head to difficult failure modes with
large teacher headroom.

The hybrid objective combines GT supervision, oracle residual supervision,
good-pixel preservation, sparsity, and a new-Bad3 safety term:
\[
\mathcal{L}
=
\lambda_{GT}\mathcal{L}_{GT}
+
\lambda_{oracle}\mathcal{L}_{oracle}
+
\lambda_{pres}\mathcal{L}_{pres}
+
\lambda_{sparse}\mathcal{L}_{sparse}
+
\lambda_{newbad}\mathcal{L}_{newbad}.
\]
The preservation term discourages changes on pixels where raw S2M2 is already
accurate:
\[
\mathcal{L}_{pres}
=
\sum_{p:\ |D^{raw}(p)-D^{GT}(p)| < 1}
\left|
D^{refined}(p) - D^{raw}(p)
\right|.
\]
The new-Bad3 penalty discourages the model from converting raw-good pixels into
new Bad-3px errors:
\[
\mathcal{L}_{newbad}
=
\sum_{p:\ |D^{raw}(p)-D^{GT}(p)| < 3}
\max
\left(
0,
|D^{refined}(p)-D^{GT}(p)| - 3
\right).
\]

\paragraph{Unfrozen hybrid fine-tuning.}
The first hybrid variant, v3.2, fine-tuned both the detector and the residual
head. This made the model substantially more conservative, but it also produced
an identity-collapse failure mode. After the first epochs, the preservation,
sparsity, and new-Bad3 losses pushed the detector probabilities downward. The
validation threshold then moved to a very conservative value, causing almost no
pixels to be modified. As a result, v3.2 was safer than v3.1 but recovered less
oracle headroom.

On the selected clips, v3.2 reduced the frame-mean MAE to 11.0780 px, recovering
5.43\% of the oracle gap. This is below the v3.1 gap recovery of 6.22\%.
However, v3.2 reduced the new-Bad3 rate from raw-good pixels from 4.79\% to
2.33\%, and reduced the modified-pixel fraction from 52.27\% to 14.27\%. This
confirmed that the hybrid objective improved safety, but was too conservative
when the detector was allowed to adapt.

\paragraph{Frozen-detector hybrid fine-tuning.}
To prevent detector collapse, we froze the v3.1 detector and fine-tuned only the
residual correction pathway. This produces the v3.2b/v3.2c frozen-detector hybrid
variants. The key idea is to preserve the already well-calibrated detection
policy and adapt only the correction magnitude:
\[
\text{keep where-to-correct fixed}
\quad \rightarrow \quad
\text{improve how-to-correct}.
\]

In v3.2b, freezing the detector prevented the identity collapse observed in v3.2.
The correction threshold remained stable at \(\tau=0.7\), and the modified-pixel
fraction remained approximately constant throughout training. This yielded a
better trade-off than both v3.1 and v3.2: v3.2b improved the selected-frame MAE
to 11.0075 px and recovered 6.99\% of the oracle gap, while modifying only
18.43\% of pixels.

We then extended the same recipe with a longer cosine-decay schedule, producing
v3.2c. The model was trained for 32 epochs with the detector frozen, residual
learning rate \(3\times10^{-4}\), sparsity weight 0.02, and the same 75/25
full-GT/oracle crop mixture. The best checkpoint was selected at epoch 29.
The improvement over v3.2b was small but consistent, indicating that this recipe
had reached a plateau.

The final selected-frame metrics are:
\[
\begin{aligned}
\mathrm{MAE}_{raw} &= 11.3231~\mathrm{px}, \\
\mathrm{MAE}_{v3.1} &= 11.0421~\mathrm{px}, \\
\mathrm{MAE}_{v3.2} &= 11.0780~\mathrm{px}, \\
\mathrm{MAE}_{v3.2b} &= 11.0075~\mathrm{px}, \\
\mathrm{MAE}_{v3.2c} &= 11.0054~\mathrm{px}.
\end{aligned}
\]

The corresponding oracle-gap recovery is:
\[
\begin{aligned}
\mathrm{Gap}_{v3.1} &= 6.22\%, \\
\mathrm{Gap}_{v3.2} &= 5.43\%, \\
\mathrm{Gap}_{v3.2b} &= 6.99\%, \\
\mathrm{Gap}_{v3.2c} &= 7.03\%.
\end{aligned}
\]

On the full-GT held-out test split, v3.2c remains better than raw S2M2:
\[
\mathrm{MAE}_{raw} = 4.6690~\mathrm{px},
\qquad
\mathrm{MAE}_{v3.2c} = 4.6145~\mathrm{px}.
\]
The full-GT Bad-3px rate also improves slightly:
\[
\mathrm{Bad}\text{-}3_{raw} = 33.536\%,
\qquad
\mathrm{Bad}\text{-}3_{v3.2c} = 33.441\%.
\]

\paragraph{Comparison.}
Table~\ref{tab:hybrid_oracle_refinement} summarizes the selected-clip comparison.
All selected-clip values are reported as frame means to allow fair comparison
between v3.1, v3.2, v3.2b, and v3.2c.

\begin{table}[t]
\centering
\caption{Hybrid oracle fine-tuning on selected teacher/oracle clips. All selected
metrics are reported as frame means.}
\label{tab:hybrid_oracle_refinement}
\begin{tabular}{lccccc}
\toprule
Method & MAE $\downarrow$ & Oracle gap $\uparrow$ & Bad-3 $\downarrow$
& New Bad-3 $\downarrow$ & Modified $\downarrow$ \\
\midrule
v3.1 staged abstention
& 11.0421 & 6.22\% & 51.280 & 4.79\% & 52.27\% \\
v3.2 hybrid oracle
& 11.0780 & 5.43\% & 51.032 & 2.33\% & 14.27\% \\
v3.2b frozen-detector hybrid
& 11.0075 & 6.99\% & 50.71 & 5.04\% & 18.43\% \\
v3.2c frozen-detector hybrid, long
& \textbf{11.0054} & \textbf{7.03\%} & \textbf{50.698}
& 5.36\% & 18.43\% \\
\bottomrule
\end{tabular}
\end{table}

\paragraph{Interpretation.}
The hybrid oracle experiments show that the detector is no longer the main
bottleneck. When the detector is allowed to adapt during hybrid fine-tuning, the
model becomes overly conservative and collapses toward the identity solution.
Freezing the detector preserves the useful staged-abstention policy and allows
the residual head to absorb a small amount of oracle information.

The final v3.2c model is the best practical hybrid variant. It recovers more
oracle headroom than v3.1 while modifying substantially fewer pixels. However,
the gain over v3.2b is marginal, and the total recovered oracle gap remains only
7.03\%. This indicates that the current residual head and loss formulation still
cannot exploit most of the available oracle/SAV headroom. Further progress will
likely require a different residual-learning strategy, stronger failure-mode
conditioning, or targeted handling of the pathological clips that dominate the
frame-mean new-Bad3 metric.


\subsection{Final Hybrid-Oracle Refiner}

The final hybrid-oracle variant, v3.2c, is the strongest practical model in the
selected failure-clip setting. It builds on the staged-abstention v3.1 baseline
and fine-tunes only the residual correction pathway while keeping the v3.1
detector frozen. This prevents the identity-collapse failure mode observed in
v3.2, where preservation and sparsity losses pushed the detector probabilities
down and caused the model to stop correcting.

On the selected teacher/oracle clips, v3.2c improves over both raw S2M2 and the
v3.1 staged-abstention baseline. Raw S2M2 obtains a frame-mean MAE of
11.3231 px, v3.1 reduces it to 11.0421 px, and v3.2c further reduces it to
11.0054 px. The all-available oracle achieves 6.8062 px, so v3.2c recovers
7.03\% of the available oracle gap, compared with 6.22\% for v3.1:
\[
\begin{aligned}
\mathrm{MAE}_{raw} &= 11.3231~\mathrm{px}, \\
\mathrm{MAE}_{v3.1} &= 11.0421~\mathrm{px}, \\
\mathrm{MAE}_{v3.2c} &= 11.0054~\mathrm{px}, \\
\mathrm{MAE}_{oracle} &= 6.8062~\mathrm{px}.
\end{aligned}
\]

At the same time, v3.2c is substantially more selective than v3.1. The modified
pixel fraction decreases from 52.27\% in v3.1 to 18.43\% in v3.2c. This shows
that the frozen-detector hybrid strategy recovers more oracle headroom while
applying corrections to a much smaller subset of pixels.

On the full-GT held-out test split, v3.2c remains better than raw S2M2:
\[
\mathrm{MAE}_{raw} = 4.6690~\mathrm{px},
\qquad
\mathrm{MAE}_{v3.2c} = 4.6145~\mathrm{px}.
\]
The Bad-3px rate also improves slightly:
\[
\mathrm{Bad}\text{-}3_{raw} = 33.536\%,
\qquad
\mathrm{Bad}\text{-}3_{v3.2c} = 33.441\%.
\]
However, v3.2c is slightly weaker than the v3.1 staged-abstention baseline on
the full-GT test split, where v3.1 reaches 4.5637 px MAE. This indicates a
trade-off: hybrid oracle fine-tuning improves selected failure clips and
selectivity, but slightly sacrifices general full-GT refinement performance.

\paragraph{Failure-mode limitation.}
A pathological-clip analysis shows that the remaining safety issue is not a
simple denominator artifact. The correlation between raw-good pixel count and
new-Bad3 percentage is weak (\(-0.321\)), whereas the correlation between
modified-pixel percentage and new-Bad3 percentage is strong (\(0.784\)). This
indicates that elevated frame-mean new-Bad3 is caused by genuine over-correction
on specific content, rather than by noisy ratios on frames with few raw-good
pixels.

The issue is concentrated in two selected clips: one dominated by high temporal
flicker and one dominated by high boundary error. On these two clips, the
abstention gate opens on 39.6\% of pixels on average, compared with 9.3\% on the
other four clips. The corresponding new-Bad3 rate is 15.7\% on the two
pathological clips and only 0.89\% on the remaining clips. Among the 50 worst
frames by new-Bad3, 39 still show a net MAE improvement, while 11 are net
regressions. Thus, the model is not globally unsafe, but it is too aggressive on
specific failure modes.

\paragraph{Interpretation.}
v3.2c should therefore be considered the current best hybrid-oracle refiner. It
improves selected failure clips more than v3.1, recovers more oracle headroom,
and modifies far fewer pixels. Its limitation is mode-dependent safety: high
temporal flicker and high boundary error can cause the abstention gate to open
too broadly. The detector is no longer the main bottleneck; the remaining
bottleneck is failure-mode-conditioned residual correction under abstention
constraints.
\paragraph{Modern damping refiner: v4-tiny.}
We first evaluated whether a moderately larger modern encoder--decoder could
improve over the tiny abstention refiners while explicitly controlling correction
strength. The resulting v4-tiny model contains 942,867 parameters and uses the
same 16-channel input representation as the previous refiners. In addition to a
bad-pixel confidence head and a bounded residual head, v4-tiny introduces a
learned damping head. The final correction is therefore modulated as
\[
D^{refined} = D^{raw} + P^{bad} \cdot \alpha \cdot \widehat{\Delta D},
\]
where \(P^{bad}\) is the learned bad-pixel probability, \(\alpha \in [0,1]\) is
a per-pixel damping factor, and \(\widehat{\Delta D}\) is the bounded residual.
The motivation is that the failure mode observed in v3.2c is not only where to
apply a correction, but how aggressively to apply it once the abstention gate is
open.

v4-tiny strongly validates the damping mechanism. On the pathological selected
clips, the new-Bad3 rate from raw-good pixels drops from 15.77\% for v3.2c to
0.33\%, and the global frame-mean new-Bad3 drops from 5.36\% to 0.55\%. However,
this safety comes at the cost of reduced accuracy. Selected-clip MAE worsens from
11.0054 px for v3.2c to 11.0669 px, and oracle-gap recovery drops from 7.03\% to
5.67\%. Moreover, v4-tiny regresses on the full-GT test split, increasing MAE
from the raw S2M2 value of 4.6690 px to 4.7763 px. Thus, v4-tiny is not a new
best model, but it demonstrates that learned damping is an effective mechanism
for suppressing pathological overcorrection.
\paragraph{Lightweight suppression controller: SOG.}
We then tested whether the safety behavior of v4-tiny could be recovered without
retraining the full refiner. The SOG branch freezes the v3.2c refiner and learns
only a lightweight suppression/gating controller on top of its outputs. This
branch has only 71k trainable parameters and adds a small runtime overhead,
resulting in a combined runtime of approximately 1.65 ms/frame. Importantly, the
model is initialized to reproduce v3.2c exactly, so the experiment starts from
the current best practical operating point and only learns where to suppress
harmful corrections.

SOG identifies a new useful operating point in the accuracy--safety trade-off.
Compared with v3.2c, it reduces pathological new-Bad3 from 15.77\% to 5.77\%,
while still improving over raw S2M2 on the full-GT test split
(\(4.6690 \rightarrow 4.6221\) px). This is the first lightweight controller in
the v3--v4 line to combine pathological safety below 8\% with a full-GT test
improvement over raw S2M2. However, this safety improvement still reduces
oracle-gap recovery from 7.03\% to 5.14\%, and selected MAE worsens from
11.0054 px to 11.0909 px. Therefore, SOG is best interpreted as a practical
safety controller rather than a true oracle-gap closing model.

\paragraph{Comparison of lightweight safety variants.}
The recent lightweight variants reveal a consistent trend. v3.2c remains the
best accuracy-oriented lightweight refiner, with selected MAE 11.0054 px and
7.03\% oracle-gap recovery, but it suffers from severe pathological
overcorrection, with 15.77\% new-Bad3 on the two pathological selected clips.
v4-tiny nearly eliminates this pathological failure mode, reducing pathological
new-Bad3 to 0.33\%, but becomes too conservative and regresses on full-GT test
performance. SOG provides an intermediate operating point: it reduces
pathological new-Bad3 to 5.77\% and still beats raw S2M2 on full-GT test, but it
does not preserve the oracle-gap recovery of v3.2c.

These results indicate that the available low-level refiner features do contain
signal for distinguishing beneficial and harmful corrections. However, small
models and lightweight controllers primarily improve safety; they do not recover
substantially more of the oracle/SAV gap. Closing that gap likely requires either
higher-capacity refinement, richer temporal modeling, richer visual input, or
stronger teacher distillation.

\begin{table}[t]
\centering
\caption{Comparison of recent lightweight refiner variants on selected clips and
full-GT test. Selected metrics are frame-mean over the selected oracle clips.}
\label{tab:recent_refiner_variants}
\begin{tabular}{lcccc}
\toprule
Metric & v3.2c & v4-tiny & v3.3 threshold & SOG \\
\midrule
Selected MAE $\downarrow$ & 11.0054 & 11.0669 & 11.1062 & 11.0909 \\
Oracle gap recovered $\uparrow$ & 7.03\% & 5.67\% & 4.80\% & 5.14\% \\
Pathological new-Bad3 $\downarrow$ & 15.77\% & 0.33\% & 6.69\% & 5.77\% \\
Clean new-Bad3 $\downarrow$ & 0.89\% & 0.64\% & 0.89\% & -- \\
Full-GT test MAE $\downarrow$ & 4.6145 & 4.7763 & -- & 4.6221 \\
Trainable parameters & -- & 943k & -- & 71k \\
Runtime & 1.07 ms & 2.74 ms & -- & 1.65 ms \\
\bottomrule
\end{tabular}
\end{table}

\paragraph{Event-Gated Boundary Memory Refiner.}
To test whether additional model capacity and explicit temporal reasoning can
break the lightweight-refiner plateau, we also introduced an experimental
Event-Gated Boundary Memory Refiner (EGBM). Unlike the previous variants, EGBM is
not a simple dense residual head. It combines an event-gated temporal memory, a
boundary-aware correction field, dynamic abstention, a mixture of correction
experts, and an iterative energy-style update. The model has approximately 4.04M
parameters in its standard configuration and remains within the online refiner
budget.

The key design idea is to explicitly represent the two dominant failure modes
identified in v3.2c: temporal flicker and boundary overcorrection. Temporal
disparity inconsistencies are converted into event-like features and processed
through a low-resolution recurrent memory. Boundary confidence is used to damp or
stop correction propagation across likely disparity edges. A dynamic abstention
module predicts local threshold offsets, while a mixture-of-experts decoder
allocates pixels among identity preservation, smooth correction, boundary
correction, and larger residual correction modes. The final update is applied in
a small number of damped refinement steps rather than as a single unconstrained
residual addition.

Early training logs suggest that EGBM has substantially stronger capacity than
the lightweight variants. During detector pretraining, validation AUC peaks
around 0.87, substantially above the final detector quality of v4-tiny. During
the initial residual/damping stage, validation MAE improves from 5.1996 px to
5.1315 px with only 0.074\% new-Bad3 and approximately 28.36\% modified pixels.
This indicates that a higher-capacity architecture can modify a larger fraction
of pixels while keeping new errors low, a behavior not observed in the smaller
models. The final selected-clip and full-GT test results are still required to
determine whether EGBM becomes a new main-branch candidate.

\paragraph{Motivation for higher-capacity refiners.}
The objective of this line of work is not merely to build a lightweight
post-processing module, but to approach the behavior of stronger offline
video-level teachers such as SAV/oracle refinements. The results so far suggest
that very small models can improve safety and reduce pathological
overcorrection, but they are unlikely to close the teacher gap. In particular,
v3.2c recovers only 7.03\% of the selected-clip oracle gap, while v4-tiny and SOG
improve safety at the cost of even lower gap recovery.

This motivates a second tier of refiners: medium-capacity, online-compatible
models with richer temporal and boundary reasoning. Since S2M2 runs at roughly
58--62 ms/frame and the practical system budget is approximately 100 ms/frame, a
refiner in the range of 10--35 ms/frame remains compatible with online use. This
budget allows models in the 5M--30M parameter range, substantially larger than
the 195k-parameter v3 refiner or the 943k-parameter v4-tiny model. Such models
may provide enough capacity to recover more of the SAV/oracle gap while still
preserving the safety constraints learned from the lightweight experiments.

\paragraph{Current interpretation.}
Overall, the recent experiments separate the refinement problem into two regimes.
Small refiners and frozen-controller variants are effective for safety: they can
reduce pathological new-Bad3 and prevent harmful corrections on raw-good pixels.
However, they do not substantially increase oracle-gap recovery. In contrast,
the early EGBM results suggest that additional capacity and explicit temporal
memory may be necessary to move beyond the v3.2c plateau. The current best
practical baseline remains v3.2c, SOG provides the best lightweight
safety-oriented operating point, and EGBM represents the first serious
higher-capacity attempt to close the gap toward SAV-like behavior.

\paragraph{Event-Gated Boundary Memory Refiner.}
The final refiner is an Event-Gated Boundary Memory Refiner (EGBM), a
failure-aware post-hoc controller designed to refine frozen S2M2 disparities in
surgical stereo sequences. Unlike the previous tiny refiners, which predict a
single gated residual, EGBM decomposes refinement into several coupled decisions:
whether a pixel should be corrected, how strongly the correction should be
applied, which correction mode should be used, and whether temporal or boundary
evidence should suppress the update.

EGBM combines four components. First, an event-gated temporal memory processes
temporal disparity inconsistencies and provides a low-resolution memory of
unstable regions. Second, a boundary-aware branch predicts edge-sensitive
features to reduce residual bleeding across depth discontinuities. Third, a
dynamic abstention module locally modulates the correction policy. Finally, a
mixture-of-experts decoder combines smooth, boundary-aware, strong-correction,
and identity-preserving experts. The resulting update is modulated by a learned
damping factor, allowing the model to attenuate harmful corrections while
preserving strong corrections on oracle-beneficial pixels.

\paragraph{Final EGBM results.}
EGBM is the first refiner to substantially break the v3.2c plateau. On the
selected oracle clips, raw S2M2 obtains a frame-mean MAE of 11.3231 px, while
EGBM reduces this to 10.4032 px. This corresponds to 20.37\% oracle-gap recovery,
compared with 7.03\% for v3.2c. At the same time, EGBM reduces pathological
new-Bad3 from 15.77\% for v3.2c to 1.30\%, while keeping clean-clip new-Bad3 at
0.81\%. The global frame-mean new-Bad3 is 0.96\%, and the pixel-weighted
new-Bad3 is only 0.21\%.

Importantly, the improvement is not limited to the selected oracle clips. On the
full-GT test split, EGBM improves raw S2M2 from 4.6690 px to 4.5226 px and
reduces Bad-3 from 33.536\% to 32.941\%. It also outperforms the previous best
v3.2c test MAE of 4.6145 px. With 4.04M parameters and a refiner runtime of
approximately 6.25 ms/frame, the estimated S2M2+EGBM system runtime is
68.25 ms/frame, remaining within the 100 ms online budget.

\paragraph{Damping interpretability.}
The learned damping behavior provides evidence that EGBM does not simply apply a
globally aggressive correction. On the high-temporal-flicker pathological clip,
the mean damping is 0.387 on hard-negative pixels and 0.733 on hard-positive
pixels. On the high-boundary-error clip, this separation is even stronger:
damping is 0.186 on hard negatives and 0.792 on hard positives. Thus, EGBM
learns to suppress correction strength on pixels where refinement would create
new Bad-3 errors, while preserving high correction strength on pixels where the
oracle indicates that correction is beneficial.

\begin{table}[t]
\centering
\caption{Final comparison of the main refiner variants. Selected metrics are
frame-mean over the selected oracle clips. Full-GT metrics are reported on the
held-out test split.}
\label{tab:final_refiner_comparison}
\begin{tabular}{lcccc}
\toprule
Metric & v3.2c & v4-tiny & SOG & EGBM \\
\midrule
Selected MAE $\downarrow$ & 11.0054 & 11.0669 & 11.0909 & \textbf{10.4032} \\
Oracle gap recovered $\uparrow$ & 7.03\% & 5.67\% & 5.14\% & \textbf{20.37\%} \\
Pathological new-Bad3 $\downarrow$ & 15.77\% & \textbf{0.33\%} & 5.77\% & 1.30\% \\
Clean new-Bad3 $\downarrow$ & 0.89\% & 0.64\% & -- & 0.81\% \\
Global frame-mean new-Bad3 $\downarrow$ & 5.36\% & \textbf{0.55\%} & -- & 0.96\% \\
Pixel-weighted new-Bad3 $\downarrow$ & 0.71\% & -- & -- & \textbf{0.21\%} \\
Full-GT test MAE $\downarrow$ & 4.6145 & 4.7763 & 4.6221 & \textbf{4.5226} \\
Full-GT test Bad-3 $\downarrow$ & -- & -- & -- & \textbf{32.941\%} \\
Parameters & 195k & 943k & 71k trainable & 4.04M \\
Runtime & 1.07 ms & 2.74 ms & 1.65 ms & 6.25 ms \\
\bottomrule
\end{tabular}
\end{table}

\paragraph{Position with respect to prior work.}
Temporal aggregation, edge-aware refinement, confidence estimation, and
mixture-of-experts have all been explored in stereo, depth estimation, or 3D
reconstruction. Therefore, the novelty of EGBM is not any single component in
isolation. Instead, EGBM contributes a failure-aware post-hoc refinement
formulation for surgical stereo. Unlike video stereo methods that directly
estimate temporally consistent disparity, EGBM operates on top of a frozen stereo
matcher and explicitly targets the domain-specific failure modes observed in
surgical scenes: temporal flicker and boundary overcorrection. The model
integrates temporal event memory, boundary-aware damping, dynamic abstention, and
correction experts to close the teacher/oracle gap while suppressing
safety-critical overcorrection.

\paragraph{Current conclusion.}
The progression from v3.2c to v4-tiny, SOG, and EGBM reveals two distinct
regimes. Lightweight controllers are effective at reducing pathological
overcorrection, but they do not substantially increase oracle-gap recovery. EGBM
is the first model to improve both sides of the trade-off: it recovers a much
larger fraction of the oracle gap while reducing pathological new-Bad3 and
improving the full-GT test split. We therefore treat EGBM as the new main branch
for online surgical stereo refinement.


\subsection{EGBM-v2-CARE: Change-Aware Reliability Encoding}

Building on the EGBM-v1 refiner, we introduce EGBM-v2-CARE, a
change-aware extension designed to improve temporal flicker handling without
sacrificing online applicability. The motivation is that temporal inconsistency
in surgical stereo is inherently ambiguous. A large change between consecutive
frames may indicate a stereo artifact, such as disparity flicker or boundary
overcorrection, but it may also correspond to a real scene change, such as tool
insertion, tissue deformation, camera motion, occlusion, or specular/liquid
appearance changes. Therefore, treating every temporal prediction error as an
artifact can suppress legitimate surgical dynamics.

EGBM-v2-CARE augments the EGBM-v1 correction controller with a predictive
temporal memory and a change-aware reliability head. Instead of only accumulating
recent disparity changes, the model predicts the expected reliability features of
the current frame from the past and interprets the resulting prediction error as
a temporal surprise signal. CARE then classifies this surprise into lightweight
change-type representations that modulate damping, dynamic abstention, expert
routing, and bad-pixel confidence.

\paragraph{Predictive temporal reliability memory.}
For each frame $t$, the current raw stereo state is compressed into a
low-resolution reliability embedding,
\begin{equation}
    z_t = f(D_t, V_t, \nabla D_t, \Delta D_t),
\end{equation}
where $D_t$ is the raw S2M2 disparity, $V_t$ is the validity mask,
$\nabla D_t$ denotes disparity-gradient features, and $\Delta D_t$ represents
local temporal disparity change. A causal memory state $M_{t-1}$ summarizes past
stereo behavior. From this state, the model predicts the expected current-frame
embedding:
\begin{equation}
    \hat{z}_t = g(M_{t-1}).
\end{equation}
The discrepancy between the predicted and observed embeddings defines a temporal
surprise signal:
\begin{equation}
    e_t = \rho(z_t - \hat{z}_t),
\end{equation}
where $\rho(\cdot)$ is a robust absolute-error operator. This signal is more
informative than a raw frame-to-frame disparity difference: predictable camera or
tissue motion can produce large disparity changes but low prediction error,
whereas unstable stereo flicker tends to produce high prediction error.

\paragraph{Change-aware reliability encoding.}
The CARE module interprets temporal surprise rather than treating it as an
artifact by default. Given the observed embedding $z_t$, predicted embedding
$\hat{z}_t$, surprise $e_t$, and memory state $M_{t-1}$, CARE predicts
low-resolution change-type probabilities:
\begin{equation}
    c_t =
    \operatorname{CARE}(z_t, \hat{z}_t, e_t, M_{t-1}),
\end{equation}
where $c_t$ contains five soft change categories:
\begin{equation}
    c_t =
    \{c_t^{\mathrm{art}},
      c_t^{\mathrm{real}},
      c_t^{\mathrm{tool}},
      c_t^{\mathrm{bdry}},
      c_t^{\mathrm{stable}}\}.
\end{equation}
These correspond respectively to artifact-like change, real scene change,
tool/occlusion-like change, boundary change, and stable predictable motion. The
goal is not to obtain semantic tool labels, but to provide the correction
controller with a reliability representation that separates stereo failure modes
from legitimate scene dynamics.

\paragraph{CARE-modulated correction.}
The original EGBM controller predicts a bad-pixel probability, a bounded residual
field, a damping factor, a local abstention threshold, and a mixture-of-experts
routing distribution. EGBM-v2-CARE injects the memory, surprise, and change-type
features into these heads. The refined disparity is computed as
\begin{equation}
    D_t^{\mathrm{ref}} =
    D_t^{\mathrm{raw}} +
    G(p_t^{\mathrm{bad}}, \tau_t)
    \cdot
    \alpha_t
    \cdot
    R_t,
\end{equation}
where $G(\cdot)$ is the dynamic abstention gate, $p_t^{\mathrm{bad}}$ is the
predicted bad-pixel probability, $\tau_t$ is the local threshold, $\alpha_t$ is
the learned damping factor, and $R_t$ is the mixture-of-experts residual. In
EGBM-v2-CARE, these quantities are conditioned on the CARE context:
\begin{equation}
    (\alpha_t, \tau_t, R_t)
    =
    h(\Phi_t, M_{t-1}, e_t, c_t),
\end{equation}
where $\Phi_t$ denotes the standard EGBM spatial, temporal, and boundary
features. Artifact-like surprise encourages lower damping, higher abstention,
and stronger identity routing. Real scene or tool-like change encourages memory
adaptation and avoids blindly suppressing the correction. Boundary-change
features reinforce edge-aware damping to prevent residual bleeding across depth
discontinuities.

\paragraph{Causal memory update.}
After refining the current frame, the memory is updated causally:
\begin{equation}
    M_t = U(M_{t-1}, z_t, e_t, c_t).
\end{equation}
In the current implementation, the memory is replayed causally over the existing
short temporal context window, preserving compatibility with the EGBM-v1
training pipeline and avoiding future-frame access. The recurrence itself is
state-compatible and can be extended to a persistent streaming state over full
video sequences. This makes the model online-compatible: at frame $t$, it uses
only the current frame and the compressed history up to $t-1$.

\paragraph{Proxy supervision for CARE.}
Because explicit labels for artifacts, tools, occlusions, and real scene changes
are not available, CARE is trained using lightweight proxy supervision derived
from existing signals. Hard-negative pixels, defined as raw-good pixels that
would become Bad-3 after refinement, encourage high artifact-like scores. Oracle
beneficial pixels encourage lower artifact-like scores and higher correction
trust. Boundary-change proxies are obtained from regions with both strong
disparity gradients and temporal disparity changes. Coherent large temporal
changes are used as weak proxies for real scene or occlusion-like change, while
low temporal-change regions encourage the stable-predictable category.

The auxiliary loss is intentionally kept small so that the main refinement and
safety losses dominate:
\begin{equation}
\begin{aligned}
    \mathcal{L}_{\mathrm{CARE}} =
    &\lambda_{\mathrm{pred}}\mathcal{L}_{\mathrm{pred}}
    + \lambda_{\mathrm{art}}\mathcal{L}_{\mathrm{art}}
    + \lambda_{\mathrm{bdry}}\mathcal{L}_{\mathrm{bdry}} \\
    &+ \lambda_{\mathrm{real}}\mathcal{L}_{\mathrm{real}}
    + \lambda_{\mathrm{stable}}\mathcal{L}_{\mathrm{stable}}
    - \lambda_{\mathrm{ent}}\mathcal{H}(c_t).
\end{aligned}
\end{equation}
Here, $\mathcal{L}_{\mathrm{pred}}$ trains the next-feature predictor,
$\mathcal{L}_{\mathrm{art}}$ separates hard negatives from hard positives,
$\mathcal{L}_{\mathrm{bdry}}$ encourages boundary-change responses near moving
depth discontinuities, $\mathcal{L}_{\mathrm{real}}$ encourages real-change
responses on coherent large temporal changes, and $\mathcal{H}(c_t)$ is a small
entropy regularizer used to prevent collapse to a single change type.

\paragraph{Warm-start and computational cost.}
EGBM-v2-CARE is initialized from the EGBM-v1 checkpoint. All shape-compatible
parameters are loaded directly, while the widened CARE-conditioned layers are
initialized such that their additional CARE input channels initially contribute
zero. This preserves the EGBM-v1 behavior exactly at initialization, ensuring
that any improvement is attributable to the learned CARE pathway rather than to a
different random initialization.

The resulting model adds only a small overhead. EGBM-v1 contains 4.04M
parameters and runs in approximately 6.26 ms/frame on an H100 GPU. EGBM-v2-CARE
contains 4.35M parameters and runs in approximately 6.65 ms/frame, adding only
0.39 ms/frame. Assuming the frozen S2M2 matcher runs at approximately
58--62 ms/frame, the full S2M2+CARE pipeline remains below the 100 ms online
budget.

\begin{table}[t]
\centering
\caption{Runtime and parameter comparison between EGBM-v1 and EGBM-v2-CARE.
Both models are benchmarked at $256\times320$ resolution on an NVIDIA H100 GPU
in FP32.}
\label{tab:care_runtime}
\begin{tabular}{lccc}
\toprule
Model & Parameters & Runtime & Peak VRAM \\
\midrule
EGBM-v1 & 4.04M & 6.26 ms/frame & 8.9 GB \\
EGBM-v2-CARE & 4.35M & 6.65 ms/frame & 9.2 GB \\
\midrule
CARE overhead & +0.31M & +0.39 ms/frame & +0.3 GB \\
\bottomrule
\end{tabular}
\end{table}

\paragraph{CARE diagnostics.}
In addition to the standard selected-oracle and full-GT metrics, we evaluate
whether CARE learns the intended reliability structure. We report temporal
surprise statistics by failure mode, CARE change-type scores on hard negatives
and hard positives, damping separation on pathological clips, and entropy of the
change-type distribution. In particular, the high-temporal-flicker clip is used
to test the core hypothesis: EGBM-v1 has weaker damping separation on this mode
than on boundary errors. CARE should reduce damping on temporal hard negatives
while preserving high damping on oracle-beneficial hard positives.

\paragraph{Toward stateful streaming CARE.}
The current EGBM-v2-CARE implementation replays the predictive memory over the
same short temporal window used by EGBM-v1. This provides a controlled
architecture test while preserving the existing training and evaluation
protocol. However, the CARE recurrence naturally supports a stronger stateful
online formulation. In a streaming version, the memory state would persist across
the full video sequence:
\begin{equation}
    M_0 = 0, \qquad
    (D_t^{\mathrm{ref}}, M_t) =
    \operatorname{CAREStep}(D_t, V_t, M_{t-1}).
\end{equation}
This would allow the model to summarize arbitrary-length past stereo behavior
into a fixed-size causal reliability state. The computational cost per frame
would remain constant, while the effective temporal context would grow with the
sequence length. Such a stateful formulation is particularly attractive for
robotic surgery, where regions may remain unstable over long periods and where
new instruments or occlusions must be incorporated into the scene memory rather
than suppressed indefinitely.

\paragraph{Potential advantage of persistent memory.}
A persistent CARE state would allow the model to distinguish transient stereo
flicker from stable scene changes over longer horizons. For example, an entering
tool may initially produce high temporal surprise, but if the change remains
spatially coherent and temporally persistent, the memory update gate can adapt
and treat it as part of the new scene state. Conversely, temporally alternating
or spatially incoherent disparity changes can be retained as artifact-like
signals and suppressed through damping and abstention. This gives the refiner an
online reliability memory whose size is independent of the sequence length.

\paragraph{Interpretation if CARE improves over EGBM-v1.}
CARE improves over EGBM-v1 by resolving a key ambiguity in temporal refinement:
not all temporal inconsistency is an artifact. By predicting the expected
current-frame reliability features and interpreting the resulting surprise as
artifact-like, boundary-like, stable, or real-scene-like change, CARE enables the
controller to suppress harmful flicker without blindly damping legitimate
surgical dynamics. This makes EGBM-v2-CARE the new main branch for online
surgical stereo refinement.

\paragraph{Interpretation if CARE does not improve aggregate performance.}
Even when CARE does not immediately improve the aggregate oracle-gap metric, its
diagnostics remain informative. A reduction in hard-negative damping on
high-temporal-flicker regions, together with non-collapsed change-type scores,
would indicate that predictive change-aware reliability contains useful signal
for a future stateful streaming formulation. In this case, CARE should be
interpreted as a temporal reliability controller whose full benefit may require
persistent sequence-level memory rather than short-window replay.

\subsection{EGBM-v3-CARE-S: Stateful Change-Aware Reliability Refinement}

EGBM-v3-CARE-S extends the EGBM family with a predictive, change-aware
temporal memory designed to improve correction safety in temporally unstable
surgical stereo sequences. The model is motivated by a limitation observed in
EGBM-v1: while the original event-gated memory and boundary-aware controller
substantially improved oracle-gap recovery, temporal changes were still treated
mainly as generic instability. This is problematic because temporal disparity
variation is intrinsically ambiguous. A large change may correspond to stereo
flicker, boundary overcorrection, or invalid matching, but it may also represent
a legitimate scene transition caused by camera motion, tissue deformation,
occlusion, tool insertion, or specular appearance changes.

The goal of EGBM-v3-CARE-S is therefore not simply to retain more frames, but to
predict what the current temporal state should look like and to interpret the
difference between predicted and observed behavior. This predictive
change-aware representation is then used to control damping, abstention, expert
routing, and memory updates.

\paragraph{Predictive reliability encoding.}
At each frame $t$, the model compresses the current stereo state into a
low-resolution reliability embedding:
\begin{equation}
    z_t = f(D_t, V_t, \nabla D_t, \Delta D_t),
\end{equation}
where $D_t$ is the raw disparity map, $V_t$ is the validity mask,
$\nabla D_t$ denotes spatial disparity gradients, and $\Delta D_t$ represents
short-term temporal disparity variation.

A causal temporal state $M_{t-1}$ summarizes the preceding sequence context.
From this state, the model predicts the expected current embedding:
\begin{equation}
    \hat{z}_t = g(M_{t-1}).
\end{equation}
The temporal surprise signal is then defined as
\begin{equation}
    e_t = \left|z_t - \hat{z}_t\right|.
\end{equation}

Unlike a simple frame-to-frame difference, temporal surprise can remain low
during predictable motion and become high when the observed stereo behavior
deviates from the learned temporal dynamics. This enables the model to
distinguish predictable scene evolution from unstable stereo artifacts.

\paragraph{Change-aware reliability representation.}
The CARE head processes the observed embedding, predicted embedding, surprise,
and temporal state:
\begin{equation}
    c_t =
    \operatorname{CARE}
    \left(
        z_t,
        \hat{z}_t,
        e_t,
        M_{t-1}
    \right).
\end{equation}
The output $c_t$ is a soft change-type representation containing five channels:
\begin{equation}
    c_t =
    \left\{
        c_t^{\mathrm{artifact}},
        c_t^{\mathrm{real}},
        c_t^{\mathrm{tool}},
        c_t^{\mathrm{boundary}},
        c_t^{\mathrm{stable}}
    \right\}.
\end{equation}
These channels represent artifact-like change, real scene change,
tool/occlusion-like change, boundary change, and stable predictable motion.

The objective is not semantic scene understanding. Instead, CARE provides a
failure-aware reliability representation that helps the controller distinguish
between changes that should suppress refinement and changes that should be
accepted as legitimate scene dynamics.

\paragraph{Correction controller.}
As in EGBM-v1, the model predicts a bad-pixel probability, a mixture of bounded
residual corrections, an identity expert, a dynamic abstention threshold, a
boundary-aware correction field, and a learned damping factor. The final
refinement is expressed as
\begin{equation}
    D_t^{\mathrm{ref}} =
    D_t^{\mathrm{raw}}
    +
    G
    \left(
        p_t^{\mathrm{bad}},
        \tau_t
    \right)
    \cdot
    \alpha_t
    \cdot
    R_t,
\end{equation}
where $G(\cdot)$ is the dynamic correction gate,
$p_t^{\mathrm{bad}}$ is the predicted bad-pixel confidence,
$\tau_t$ is the local abstention threshold,
$\alpha_t$ is the damping factor, and
$R_t$ is the mixture-of-experts residual.

In EGBM-v3-CARE-S, these quantities are conditioned on the CARE context:
\begin{equation}
    \left(
        \alpha_t,
        \tau_t,
        R_t,
        r_t
    \right)
    =
    h
    \left(
        \Phi_t,
        M_{t-1},
        e_t,
        c_t
    \right),
\end{equation}
where $\Phi_t$ denotes the standard spatial, temporal, and boundary features, and
$r_t$ is the expert-routing distribution.

Artifact-like regions encourage reduced damping, stronger abstention, and
identity-expert routing. Boundary changes reinforce edge-aware correction, while
coherent real-scene changes are allowed to modify the temporal state rather than
being automatically suppressed.

\paragraph{Stateful streaming memory.}
EGBM-v3-CARE-S introduces an explicit persistent state that is updated
frame-by-frame:
\begin{equation}
    S_t =
    \left\{
        M_t^{\mathrm{EGBM}},
        M_t^{\mathrm{CARE}},
        \hat{z}_{t+1},
        D_t^{\mathrm{LR}},
        V_t^{\mathrm{LR}}
    \right\}.
\end{equation}
The state includes the original EGBM temporal memory, the predictive CARE
memory, the next predicted embedding, and the previous low-resolution disparity
and validity maps.

Streaming inference follows
\begin{equation}
    \left(
        D_t^{\mathrm{ref}},
        S_t
    \right)
    =
    \operatorname{CAREStep}
    \left(
        X_t,
        S_{t-1}
    \right),
\end{equation}
where $X_t$ contains the current feature stack.

The temporal state has fixed spatial and channel dimensions. Therefore, both the
memory footprint and computational cost per frame remain constant with respect
to the sequence length:
\begin{equation}
    \mathcal{O}(\text{memory}) = \mathcal{O}(1),
    \qquad
    \mathcal{O}(\text{cost per frame}) = \mathcal{O}(1).
\end{equation}

\paragraph{Change-aware memory preservation.}
The streaming formulation includes a learned local memory-preservation gate.
Given the current embedding, surprise, CARE probabilities, and previous memory,
the model predicts
\begin{equation}
    k_t =
    \sigma
    \left(
        q
        \left(
            z_t,
            e_t,
            c_t,
            M_{t-1}
        \right)
    \right),
\end{equation}
where high $k_t$ values preserve the previous memory and low values allow
stronger adaptation to the current observation. The state is updated as
\begin{equation}
    M_t =
    k_t \odot M_{t-1}
    +
    \left(
        1-k_t
    \right)
    \odot
    \widetilde{M}_t.
\end{equation}

This mechanism is intended to prevent temporally unstable artifact-like
observations from contaminating the persistent state, while still allowing rapid
adaptation to coherent scene changes.

\paragraph{Window and streaming modes.}
The architecture supports two inference configurations.

In \emph{window mode}, the temporal memories are reconstructed causally from a
four-frame local context for each target frame:
\begin{equation}
    \left\{
        X_{t-3},
        X_{t-2},
        X_{t-1},
        X_t
    \right\}
    \rightarrow
    D_t^{\mathrm{ref}}.
\end{equation}

In \emph{streaming mode}, the state persists across the full sequence:
\begin{equation}
    S_0 = 0,
    \qquad
    \left(
        D_t^{\mathrm{ref}},
        S_t
    \right)
    =
    \operatorname{CAREStep}
    \left(
        X_t,
        S_{t-1}
    \right).
\end{equation}

This dual formulation enables a controlled comparison between predictive
change-aware encoding with short temporal replay and predictive encoding with
long-horizon persistent memory.

\paragraph{Training with truncated backpropagation through time.}
The model is trained on contiguous temporal chunks rather than independent
frames. For each sequence chunk of length $L$, the state is initialized at the
chunk boundary and propagated causally:
\begin{equation}
    S_0 = 0,
    \qquad
    S_t =
    U
    \left(
        S_{t-1},
        X_t
    \right),
    \quad
    t=1,\ldots,L.
\end{equation}
Losses are accumulated across the chunk and the state is detached between
chunks, implementing truncated backpropagation through time. The default chunk
length is 16 frames, with an initial burn-in period whose loss contribution is
downweighted.

The training protocol follows three stages.

\begin{enumerate}
    \item \textbf{Stage 1: predictive detector adaptation.}
    The backbone is initialized from EGBM-v1, while the bad-pixel detector and
    CARE pathway are trained using detector supervision and predictive-memory
    loss.

    \item \textbf{Stage 2: full-GT streaming refinement.}
    The model is trained on full-GT temporal chunks using disparity reconstruction,
    preservation, new-Bad3, damping, detector, and temporal prediction objectives.

    \item \textbf{Stage 3: oracle and hard-negative refinement.}
    Training mixes full-GT chunks, oracle-beneficial chunks, and pathological
    hard-negative chunks. This stage explicitly penalizes unsafe corrections,
    encourages identity routing on vulnerable pixels, and suppresses memory
    contamination by hard negatives.
\end{enumerate}

\paragraph{CARE supervision.}
Explicit labels for artifact, tool, occlusion, and scene-change classes are not
available. CARE therefore uses weak proxy supervision derived from available
signals. Hard negatives encourage artifact-like responses, while
oracle-beneficial pixels discourage artifact classification. Boundary-change
responses are encouraged in regions with both high spatial gradients and high
temporal variation. Coherent large temporal changes act as weak proxies for
real-scene or occlusion-like transitions, while low temporal variation supports
the stable-predictable category.

The CARE auxiliary objective is
\begin{equation}
\begin{aligned}
    \mathcal{L}_{\mathrm{CARE}}
    =
    &\lambda_{\mathrm{pred}}
    \mathcal{L}_{\mathrm{pred}}
    +
    \lambda_{\mathrm{artifact}}
    \mathcal{L}_{\mathrm{artifact}}
    +
    \lambda_{\mathrm{boundary}}
    \mathcal{L}_{\mathrm{boundary}}
    \\
    &+
    \lambda_{\mathrm{real}}
    \mathcal{L}_{\mathrm{real}}
    +
    \lambda_{\mathrm{stable}}
    \mathcal{L}_{\mathrm{stable}}
    -
    \lambda_{\mathrm{entropy}}
    \mathcal{H}(c_t).
\end{aligned}
\end{equation}

The auxiliary losses are kept small so that the main refinement and safety
objectives remain dominant.

\paragraph{Safety-aware refinement objective.}
The complete training objective combines full-GT reconstruction, detector
supervision, residual regression, raw-good preservation, new-Bad3 suppression,
damping regularization, hard-negative preservation, identity routing, and CARE
auxiliary losses:
\begin{equation}
\begin{aligned}
    \mathcal{L}
    =
    &\lambda_{\mathrm{full}}
    \mathcal{L}_{\mathrm{full}}
    +
    \lambda_{\mathrm{det}}
    \mathcal{L}_{\mathrm{det}}
    +
    \lambda_{\mathrm{res}}
    \mathcal{L}_{\mathrm{res}}
    \\
    &+
    \lambda_{\mathrm{pres}}
    \mathcal{L}_{\mathrm{pres}}
    +
    \lambda_{\mathrm{nB3}}
    \mathcal{L}_{\mathrm{nB3}}
    +
    \lambda_{\mathrm{damp}}
    \mathcal{L}_{\mathrm{damp}}
    \\
    &+
    \lambda_{\mathrm{id}}
    \mathcal{L}_{\mathrm{id}}
    +
    \lambda_{\mathrm{memory}}
    \mathcal{L}_{\mathrm{memory}}
    +
    \mathcal{L}_{\mathrm{CARE}}.
\end{aligned}
\end{equation}

The hard-negative terms target pixels that are correct in the raw disparity but
would become erroneous after refinement. This directly trains the model to
suppress unsafe corrections.

\paragraph{Initialization and computational cost.}
EGBM-v3-CARE-S is warm-started from EGBM-v1. All compatible parameters are
loaded directly, while newly introduced CARE and streaming modules are
initialized safely. The widened decoder and damping layers copy the original
EGBM-v1 weights into their existing input channels, while the additional CARE
channels are initialized to zero. Consequently, the initial window-mode
behavior matches the original EGBM policy before CARE influence is learned.

The final model contains approximately 4.35M parameters. On an NVIDIA H100 GPU
at $256\times320$ resolution, the window and streaming configurations run at
approximately 7.41 ms/frame and 7.13 ms/frame, respectively. Combined with the
frozen S2M2 matcher, the estimated full system latency is approximately
69 ms/frame, remaining below the 100 ms online budget.

\begin{table}[t]
\centering
\caption{Comparison of the main temporal refinement variants. Selected metrics
are frame-mean over 502 frames from six selected oracle clips. Full-GT metrics
are reported on the held-out test split.}
\label{tab:egbm_v3_comparison}
\begin{tabular}{lccccc}
\toprule
Model &
Selected MAE $\downarrow$ &
Oracle gap $\uparrow$ &
Patho nB3 $\downarrow$ &
Clean nB3 $\downarrow$ &
Full-GT test $\downarrow$ \\
\midrule
v3.2c &
11.0054 &
7.03\% &
15.77\% &
0.89\% &
4.6145 \\
v4-tiny &
11.0669 &
5.67\% &
0.33\% &
0.64\% &
4.7763 \\
SOG &
11.0909 &
5.14\% &
5.77\% &
-- &
4.6221 \\
EGBM-v1 &
10.4032 &
20.37\% &
1.30\% &
0.81\% &
4.5226 \\
EGBM-v2-CARE &
\textbf{10.3866} &
\textbf{20.73\%} &
2.89\% &
0.80\% &
\textbf{4.5090} \\
EGBM-v3-CARE-S window &
10.4066 &
20.29\% &
0.67\% &
\textbf{0.50\%} &
4.5753 \\
EGBM-v3-CARE-S streaming &
10.4125 &
20.16\% &
\textbf{0.47\%} &
\textbf{0.50\%} &
4.5880 \\
\bottomrule
\end{tabular}
\end{table}

\paragraph{Accuracy--safety trade-off.}
EGBM-v3-CARE-S operates in essentially the same aggregate accuracy regime as
EGBM-v1. In window mode, selected MAE changes from 10.4032 px to 10.4066 px and
oracle-gap recovery changes from 20.37\% to 20.29\%. In streaming mode, the model
obtains 10.4125 px and 20.16\%, respectively.

Although these differences are small, the safety metrics improve substantially.
Pathological new-Bad3 decreases from 1.30\% in EGBM-v1 to 0.67\% in window mode
and 0.47\% in streaming mode. Clean-clip new-Bad3 decreases from 0.81\% to
0.50\% in both EGBM-v3 configurations.

Thus, the main benefit of EGBM-v3-CARE-S is not increased aggregate disparity
accuracy, but reduced unsafe correction behavior.

\paragraph{Damping behavior on temporal flicker.}
The strongest evidence for the intended CARE mechanism is the damping response
on the high-temporal-flicker clip. EGBM-v1 applies a mean damping of 0.387 on
hard-negative pixels and 0.733 on hard-positive pixels. EGBM-v3-CARE-S reduces
hard-negative damping to approximately 0.127 while preserving high damping on
oracle-beneficial hard positives.

The hard-negative versus hard-positive damping separation therefore increases
from approximately
\begin{equation}
    0.733 - 0.387 = 0.346
\end{equation}
for EGBM-v1 to approximately
\begin{equation}
    0.690
    -
    0.127
    =
    0.563
\end{equation}
for EGBM-v3-CARE-S.

This indicates that CARE learns a more selective correction policy: harmful
updates are strongly attenuated while beneficial corrections remain active.

\paragraph{Source of the safety improvement.}
The safety improvement should not be attributed to persistent long-horizon
memory. The windowed model, which reconstructs temporal context from four local
frames, achieves nearly the same safety gains as the streaming model. Moreover,
the window configuration does not use the persistent sequence-level memory gate
during inference.

The improvement therefore originates primarily from two factors:
\begin{equation}
    \text{predictive change-aware local encoding}
\end{equation}
and
\begin{equation}
    \text{safety-aware sequential training}.
\end{equation}

CARE enables the controller to interpret short-term temporal variation rather
than treating all changes as generic instability. Chunk-based training then
converts this representation into a conservative correction policy through
hard-negative preservation, new-Bad3 penalties, identity routing, damping
supervision, and causal temporal optimization.

The resulting interpretation is
\begin{equation}
    \text{Safety gain}
    \approx
    \text{change-aware short context}
    +
    \text{sequential safety training},
\end{equation}
rather than
\begin{equation}
    \text{Safety gain}
    \approx
    \text{arbitrarily long temporal memory}.
\end{equation}

\paragraph{Long-horizon memory ablation.}
Persistent streaming inference does not outperform short-window replay. The
windowed model obtains a selected MAE of 10.4066 px and recovers 20.29\% of the
oracle gap, compared with 10.4125 px and 20.16\% for streaming inference. The
same trend appears on the full-GT test split, where window mode achieves
4.5753 px and streaming mode achieves 4.5880 px.

Streaming also fails to improve the clean or pathological safety metrics enough
to justify the additional state-management complexity. Although pathological
new-Bad3 is slightly lower in streaming mode, the gain is small relative to the
added implementation burden and the mild degradation in aggregate accuracy.

These results do not indicate that temporal reasoning is unnecessary. Rather,
they show that the relevant temporal evidence in this dataset is mostly captured
within a short local horizon. The primary gain comes from learning how to
interpret recent change, not from preserving an arbitrarily long history.

\paragraph{Why persistent memory did not help.}
Several explanations are consistent with the observed result. First, the selected
validation clips may not contain temporal dependencies that require context
beyond a few frames. Second, the predictive state may already summarize most
useful information within the four-frame local window. Third, the weak proxy
supervision used for CARE may be sufficient for local reliability
classification but not strong enough to train meaningful long-horizon scene
adaptation. Finally, persistent memory may introduce mild temporal smoothing or
state inertia that reduces responsiveness to rapidly changing surgical content.

Importantly, empirical analysis did not support a simple long-sequence drift
failure. Both window and streaming modes showed similar degradation on difficult
sequence segments, suggesting that some regions are intrinsically harder rather
than being degraded by persistent-state accumulation.

\paragraph{Practical deployment recommendation.}
For deployment, the window-mode EGBM-v3-CARE-S configuration offers the best
practical trade-off. It retains almost the same disparity accuracy and oracle-gap
recovery as EGBM-v1, while approximately halving pathological new-Bad3 and
reducing clean-clip new-Bad3.

The streaming configuration provides slightly stronger pathological safety, but
does not improve aggregate accuracy, full-GT generalization, or runtime enough to
justify persistent state management. The recommended operational configuration
is therefore:
\begin{equation}
    \boxed{
    \text{EGBM-v3-CARE-S in window mode}
    }
\end{equation}
for safety-oriented robotic deployment.

The EGBM-v2-CARE model remains the accuracy-oriented variant, as it achieves the
best selected MAE, oracle-gap recovery, and full-GT test MAE, but its higher
pathological new-Bad3 makes it less suitable when unsafe overcorrection is the
primary concern.

\paragraph{Scientific interpretation.}
The architectural progression supports three conclusions.

First, increased capacity alone is insufficient. Larger models and candidate
fusion did not consistently outperform EGBM-v1.

Second, failure-aware temporal interpretation is beneficial. CARE substantially
improves the separation between harmful and useful corrections, especially on
temporal flicker.

Third, long-horizon persistent memory is not required for the observed gain.
Short-window predictive encoding combined with safety-aware sequential training
captures most of the useful temporal structure.

The central finding is therefore:
\begin{quote}
The dominant benefit does not arise from retaining a longer temporal history,
but from learning how to interpret short-term temporal change and suppress
unsafe corrections.
\end{quote}

\paragraph{Final conclusion.}
EGBM-v3-CARE-S is not a clean accuracy winner over EGBM-v1, but it establishes a
substantially safer refinement policy at almost unchanged disparity accuracy.
The model reduces pathological new-Bad3 from 1.30\% to 0.67\% in window mode and
0.47\% in streaming mode, while maintaining approximately 20\% oracle-gap
recovery.

The long-context hypothesis is not confirmed: persistent sequence-level memory
does not outperform short-window replay. However, the predictive CARE
representation and sequence-aware safety training clearly improve correction
selectivity. EGBM-v3-CARE-S should therefore be interpreted as a safety-oriented
temporal refinement architecture whose main contribution is more reliable
correction control rather than longer temporal memory.

% Report 1's \appendix switched \thesection to \Alph, and every report pasted in after it
% inherited that: past the twenty-sixth section LaTeX raises a fatal "Counter too large",
% which it did thirty-two times. Numbering returns to arabic here, where report 2 begins.
\renewcommand{\thesection}{\arabic{section}}
\setcounter{section}{0}

\section*{ARGOS Temporal Stereo Refinement: From Universal Refinement to
Lightweight Domain Adaptation}

% This report was pasted in with its own preamble, which silently broke the diary:
% a second \documentclass mid-document is a fatal error, so nothing compiled from
% here on. Its definitions moved to the preamble above and the title became a
% section, which is what it always was.


\section{Original Objective}

The original objective was to develop a lightweight temporal refinement module on
top of a frozen stereo model, initially \smodel, to:

\begin{itemize}[leftmargin=*]
    \item reduce frame-to-frame disparity flickering;
    \item improve temporal coherence;
    \item correct large geometric errors;
    \item preserve already-correct raw predictions;
    \item operate causally and in real time;
    \item generalize to unseen surgical domains without retraining.
\end{itemize}

The intended system was:

\begin{equation}
    \left\{
    D_{t-k}^{\mathrm{raw}}, \ldots, D_t^{\mathrm{raw}}
    \right\}
    \xrightarrow{\text{temporal refiner}}
    D_t^{\mathrm{refined}},
\end{equation}

where the stereo backbone remained frozen and only the temporal refinement module
was trained.

\section{Initial Experimental Findings}

\subsection{Refinement on SCARED}

The developed refiners successfully improved \smodel{} on the SCARED domain.
The strongest conservative variants achieved a favorable compromise between
accuracy and safety, while remaining lightweight and fast.

\begin{table}[h]
\centering
\caption{Summary of representative refiner variants on the primary SCARED setup.}
\small
\begin{tabularx}{\textwidth}{
    >{\raggedright\arraybackslash}X
    r
    r
    r
    r
}
\toprule
\textbf{Method}
& \textbf{Oracle gap}
& \textbf{Pathological new-Bad3}
& \textbf{Full-GT MAE}
& \textbf{Runtime} \\
\midrule
Raw \smodel
& 0\%
& 0\%
& 4.6690
& -- \\

EGBM-v1
& 20.37\%
& 1.30\%
& 4.5226
& 6.25 ms \\

EGBM-v2 CARE
& 20.73\%
& 2.89\%
& \textbf{4.5090}
& 6.65 ms \\

EGBM-v3 CARE window
& 20.29\%
& 0.67\%
& 4.5753
& 7.41 ms \\

EGBM-v3 CARE streaming
& 20.16\%
& \textbf{0.47\%}
& 4.5880
& 7.13 ms \\
\bottomrule
\end{tabularx}
\end{table}

The short causal window was at least as effective as persistent long-term state.
This suggests that the useful temporal signal is mostly local and that long memory
is not automatically beneficial for surgical stereo refinement.

\subsection{Oracle-Gap Analysis}

A detailed oracle analysis showed that the main limitation was not detecting where
the raw prediction was wrong, nor predicting the direction of the correction.

Approximately:

\begin{itemize}[leftmargin=*]
    \item 95\% of the remaining oracle gap required corrections larger than 3 px;
    \item more than 92\% required corrections larger than 6 px;
    \item approximately 86\% required corrections larger than 12 px;
    \item correction support was approximately 76\%;
    \item correction sign accuracy was approximately 94--95\%;
    \item applied correction magnitude was only approximately 15--17\% of the
    oracle magnitude when the sign was correct.
\end{itemize}

The main bottleneck was therefore correction magnitude.

\subsection{Large-Magnitude Proposal Branch}

The MPC, CPV, and ASSR branches were introduced to allow large corrections.

\begin{table}[h]
\centering
\caption{Large-correction experiments.}
\small
\begin{tabularx}{\textwidth}{
    >{\raggedright\arraybackslash}X
    r
    r
    r
    r
}
\toprule
\textbf{Method}
& \textbf{Selected MAE}
& \textbf{Oracle gap}
& \textbf{Pathological new-Bad3}
& \textbf{Full-GT MAE} \\
\midrule
MPC
& 9.1788
& 47.47\%
& 7.33\%
& 4.6386 \\

CPV
& 9.2977
& 44.84\%
& 5.05\%
& 4.5939 \\

ASSR
& 9.2431
& 46.05\%
& 6.02\%
& 4.5997 \\
\bottomrule
\end{tabularx}
\end{table}

These results confirmed that large-magnitude corrections could close a substantial
part of the oracle gap. However, they also exposed a strong safety trade-off:

\begin{quote}
Increasing correction magnitude improved difficult frames, but caused
over-correction of predictions that were already geometrically accurate.
\end{quote}

A conservative post-hoc configuration reduced pathological errors to approximately
0.43\%, but also collapsed oracle-gap recovery to approximately 12.48\%.

\section{What Did Not Work}

\subsection{Universal Zero-Shot Refinement}

The original assumption was that a temporal refiner trained on SCARED could be
placed on top of \smodel{} and improve its outputs in other surgical domains.

This was disproven by the SERV-CT zero-shot evaluation.

\begin{table}[h]
\centering
\caption{Zero-shot behavior on SERV-CT. Positive MAE change indicates degradation.}
\small
\begin{tabularx}{\textwidth}{
    >{\raggedright\arraybackslash}X
    r
    r
    r
}
\toprule
\textbf{Method}
& $\Delta$ \textbf{MAE}
& \textbf{new-Bad3}
& \textbf{Harmful correction rate} \\
\midrule
EGBM-v1
& +0.53
& 8.5\%
& 0.70 \\

EGBM-v2
& +0.60
& 8.5\%
& 0.68 \\

v3.2c
& +0.54
& 13.3\%
& 0.47 \\

MPC
& +5.35
& 44.2\%
& 0.83 \\

CPV
& +5.06
& 39.3\%
& 0.82 \\
\bottomrule
\end{tabularx}
\end{table}

On SERV-CT, raw \smodel{} already achieved a MAE of approximately 1.28 px.
The refiner had been trained in a regime where significant corrections were often
required. When transferred to a domain where the raw matcher was already accurate,
the learned correction policy continued to activate and degraded the result.

\subsection{Interpretation}

The failure is not simply caused by insufficient model capacity.

A disparity-only refiner observes patterns such as:

\begin{itemize}[leftmargin=*]
    \item disparity gradients;
    \item temporal differences;
    \item boundaries;
    \item validity masks;
    \item local instability.
\end{itemize}

However, the same pattern may correspond to:

\begin{itemize}[leftmargin=*]
    \item a real geometric discontinuity;
    \item a stereo-matching failure;
    \item true object or tissue motion;
    \item temporal flicker;
    \item a prediction that is already correct.
\end{itemize}

Without access to the original matching evidence, cost volume, photometric
consistency, or calibrated matcher confidence, the problem may be partially
non-identifiable.

\subsection{Why More SCARED Data Alone Is Not Sufficient}

Even a large number of correlated frames does not guarantee domain coverage.

SCARED contains approximately 20,000 frames but only a limited number of independent
sequences and acquisition conditions. Consecutive frames have high correlation, so
the effective number of independent examples is much lower than the raw frame count.

More frames from the same error distribution do not necessarily teach the refiner
how to behave when:

\begin{itemize}[leftmargin=*]
    \item the camera changes;
    \item calibration changes;
    \item the stereo backbone changes;
    \item raw predictions become much more accurate;
    \item the distribution of occlusions, specularities, and boundaries changes.
\end{itemize}

\section{Revised Scope}

\subsection{New Central Idea}

The refiner is no longer presented as a universal zero-shot correction module.

The revised concept is:

\begin{quote}
A lightweight causal temporal adapter that specializes a frozen stereo backbone
to the error dynamics of a target domain using a limited amount of target data.
\end{quote}

The full system becomes:

\begin{equation}
    \underbrace{
    (I_t^L, I_t^R)
    \xrightarrow{\text{frozen stereo backbone}}
    D_t^{\mathrm{raw}}
    }_{\text{general stereo estimation}}
    \xrightarrow{
    \substack{
    \text{lightweight temporal}\\
    \text{domain adapter}
    }}
    D_t^{\mathrm{refined}}.
\end{equation}

Only the temporal adapter is trained or fine-tuned for the target domain.

\subsection{Main Claim}

\begin{quote}
A small causal temporal adapter, with approximately four million trainable
parameters, can efficiently specialize frozen stereo backbones to domain-specific
error dynamics, improving geometric accuracy and temporal stability without
retraining a large stereo network.
\end{quote}

\subsection{Claims Explicitly Avoided}

The project will not claim that:

\begin{itemize}[leftmargin=*]
    \item one checkpoint improves every stereo backbone;
    \item one checkpoint improves every unseen domain zero-shot;
    \item the complete system contains only four million parameters;
    \item the method universally outperforms all large stereo models;
    \item sparse D4D keyframes provide dense temporal ground truth.
\end{itemize}

The approximately four million parameters refer only to the additional trainable
adapter.

\section{Research Questions}

\begin{description}[leftmargin=1.3cm,style=nextline]
    \item[RQ1]
    Can a lightweight causal temporal adapter improve \smodel{} on SCARED in both
    geometric accuracy and temporal consistency?

    \item[RQ2]
    Can the same adapter architecture be trained successfully on top of different
    frozen stereo backbones?

    \item[RQ3]
    How much target-domain data is required before the adapter becomes effective?

    \item[RQ4]
    Does pretraining the adapter on a source domain reduce the target data required
    compared with training from scratch?

    \item[RQ5]
    Can an efficient stereo backbone plus a small adapter approach or outperform
    substantially larger online stereo systems on a target domain?

    \item[RQ6]
    Does the adapter improve true temporal stability, or does it merely learn a
    static dataset-specific geometric correction?
\end{description}

\section{Dataset Strategy}

\subsection{SCARED}

SCARED remains the primary dataset.

Its role is:

\begin{itemize}[leftmargin=*]
    \item dense supervised training;
    \item dense temporal evaluation;
    \item oracle analysis;
    \item architectural ablations;
    \item flicker and jitter evaluation;
    \item boundary stability;
    \item temporal safety analysis.
\end{itemize}

SCARED is the dataset that directly supports the temporal-refinement claim.

\subsection{SERV-CT}

SERV-CT contains only a small collection of non-consecutive stereo images and is
therefore not suitable for dense temporal-consistency evaluation.

Its role is restricted to:

\begin{itemize}[leftmargin=*]
    \item static cross-domain geometric evaluation;
    \item raw-good preservation;
    \item analysis of zero-shot false activation;
    \item optional few-shot geometric recalibration.
\end{itemize}

SERV-CT should not be used to support claims about flicker reduction.

\subsection{D4D}

D4D contains long stereo sequences and sparse Zivid structured-light geometry.

The raw organization is:

\begin{verbatim}
dataset/D4D/raw/extracted/
    specimen_N/
        session_timestamp/
            left_images/
            right_images/
            camera_info/
            depth_images/
            pointcloud/
            color_images/
            snr_images/
            masks/
            tf/
            clips.json
\end{verbatim}

The stereo images are not originally rectified. Each session contains approximately
500 stereo frames, while Zivid depth and point clouds are available only at sparse
start/end geometry anchors.

D4D cannot provide exact per-frame dense temporal ground truth because the scene is
non-rigid.

Its scientifically valid role is:

\begin{itemize}[leftmargin=*]
    \item independent sparse-keyframe geometric evaluation;
    \item structured-light target-domain supervision;
    \item specimen-disjoint adaptation experiments;
    \item session-level few-shot adaptation;
    \item cross-backbone evaluation at real geometric anchors.
\end{itemize}

\subsection{D4D Transform Chain}

The Zivid geometry and the stereo endoscope are connected through optical tracking.
For specimen~1, the validated marker-bridge transformation is:

\begin{equation}
T_{\mathrm{cam}\leftarrow\mathrm{zivid}}
=
T_{\mathrm{ps}\leftarrow\mathrm{cam}}^{-1}
T_{\mathrm{ps}\leftarrow\mathrm{MiRe45}}
T_{\mathrm{polaris}\leftarrow\mathrm{MiRe45}}^{-1}
T_{\mathrm{polaris}\leftarrow\mathrm{zivid}}.
\end{equation}

The full-dataset analysis revealed that transform conventions differ between
specimens:

\begin{itemize}[leftmargin=*]
    \item \texttt{mire45\_bridge}: used for specimen~1;
    \item \texttt{direct\_ps}: camera and Zivid both represented in
    \texttt{polaris\_spectra};
    \item \texttt{direct\_polaris}: camera and Zivid directly represented in
    \texttt{polaris}.
\end{itemize}

The convention is detected per session rather than forced globally.

\subsection{D4D Benchmark Status}

At the current frozen state:

\begin{table}[h]
\centering
\caption{Current D4D sparse-keyframe benchmark.}
\small
\begin{tabularx}{\textwidth}{
    >{\raggedright\arraybackslash}X
    r
    r
    >{\raggedright\arraybackslash}X
}
\toprule
\textbf{Specimen}
& \textbf{Candidate anchors}
& \textbf{Valid anchors}
& \textbf{TF convention} \\
\midrule
Specimen 1
& 82
& 60
& MiRe45 bridge and direct Polaris \\

Specimen 2
& 39
& 31
& Direct Polaris Spectra \\

Specimen 3
& 76
& 51
& Direct Polaris Spectra and direct Polaris \\
\midrule
Total
& 197
& 142 valid + 24 warning
& 166 usable \\
\bottomrule
\end{tabularx}
\end{table}

Additional properties:

\begin{itemize}[leftmargin=*]
    \item structured-light GT coverage: approximately 19--34\%, median 27\%;
    \item stereo--Zivid synchronization: approximately 1--55 ms, median 20 ms;
    \item exact depth--disparity numerical round trip;
    \item per-session calibration loading;
    \item focal length around 798 px for specimen~1 and 834 px for specimen~2;
    \item stereo baseline around 4.24 mm for specimen~1 and 3.97 mm for specimen~2;
    \item deterministic, resumable, and idempotent processing;
    \item session-disjoint, specimen-disjoint, leave-one-specimen-out, and few-shot
    splits;
    \item 47 automated validation checks passing.
\end{itemize}

Specimen~4 is pending complete extraction. Specimen~5 is currently blocked by a
corrupted archive.

\section{Proposed Model Scope}

\subsection{Main Architecture}

The main candidate is a consolidated version of \method{} with a short causal
window.

The architecture contains:

\begin{itemize}[leftmargin=*]
    \item causal temporal context;
    \item residual disparity prediction;
    \item boundary-aware processing;
    \item identity/no-op expert;
    \item correction damping;
    \item change-aware reliability;
    \item safety-oriented gating;
    \item approximately four million trainable parameters;
    \item approximately seven milliseconds of additional runtime.
\end{itemize}

\subsection{Role of MPC and Large-Correction Models}

MPC, CPV, and ASSR are not the main deployable systems.

They are retained as analytical variants showing that:

\begin{itemize}[leftmargin=*]
    \item correction magnitude is the main oracle-gap bottleneck;
    \item large corrections can recover substantially more oracle error;
    \item unrestricted correction magnitude causes severe over-correction;
    \item safety and accuracy must be balanced explicitly.
\end{itemize}

\section{Backbone Strategy}

The adapter should be evaluated on multiple frozen stereo backbones.

A practical initial set is:

\begin{enumerate}[leftmargin=*]
    \item \smodel{} as the primary efficient backbone;
    \item RAFT-Stereo or CREStereo as an architecturally different matcher;
    \item Fast-FoundationStereo or FoundationStereo as a large modern reference.
\end{enumerate}

For each backbone, three settings are evaluated.

\subsection{Zero-Shot Cross-Backbone Transfer}

\begin{equation}
    \text{adapter trained on \smodel}
    \longrightarrow
    \text{different frozen backbone}.
\end{equation}

This tests whether the learned correction policy transfers directly.

A negative result is still informative because it quantifies backbone-specific
error distributions.

\subsection{Backbone-Specific Adapter Training}

The same adapter architecture is trained from scratch on predictions generated by
the new frozen backbone.

This tests whether the architecture itself is reusable across matchers.

\subsection{Lightweight Fine-Tuning}

A source-pretrained adapter is fine-tuned using a small amount of target data.

This is the main proposed operating mode.

\section{Data-Efficiency Study}

The target-domain sweep should be performed at the level of complete sessions or
sequences, not individual randomly sampled frames.

Suggested fractions are:

\begin{equation}
    0,\; 1,\; 2.5,\; 5,\; 10,\; 25,\; 50,\; 100\%.
\end{equation}

For D4D, more interpretable units are:

\begin{equation}
    1,\; 2,\; 4,\; 8 \text{ sessions},
\end{equation}

in addition to percentage-based subsets.

For every subset size, compare:

\begin{itemize}[leftmargin=*]
    \item raw frozen backbone;
    \item simple residual refiner;
    \item adapter trained from scratch;
    \item source-pretrained adapter;
    \item source-pretrained adapter fine-tuned on target data;
    \item full-data target adapter;
    \item large stereo baseline.
\end{itemize}

At small subset sizes, multiple deterministic seeds should be used.

\section{Experimental Protocol}

\subsection{Experiment A: SCARED Main Evaluation}

Evaluate:

\begin{itemize}[leftmargin=*]
    \item raw \smodel;
    \item classical temporal filters;
    \item simple convolutional refiner;
    \item EGBM-v1;
    \item EGBM-v2 CARE;
    \item EGBM-v3 CARE window;
    \item streaming variant;
    \item published temporal-refinement baseline, where compatible.
\end{itemize}

\subsection{Experiment B: Cross-Backbone Adaptation}

For every frozen backbone:

\begin{enumerate}[leftmargin=*]
    \item generate raw disparities on the same sequences;
    \item build backbone-specific oracle targets;
    \item evaluate source-trained adapter zero-shot;
    \item train the adapter from scratch;
    \item fine-tune a pretrained adapter;
    \item compare sample-efficiency curves.
\end{enumerate}

\subsection{Experiment C: D4D Few-Shot Adaptation}

The temporal adapter receives causal frame context, but supervision is applied only
at the sparse Zivid keyframe:

\begin{equation}
\left(
D_{t-3}^{\mathrm{raw}},
D_{t-2}^{\mathrm{raw}},
D_{t-1}^{\mathrm{raw}},
D_t^{\mathrm{raw}}
\right)
\xrightarrow{\text{adapter}}
D_t^{\mathrm{refined}},
\end{equation}

with loss:

\begin{equation}
\mathcal{L}_t =
\mathcal{L}_{\mathrm{geometry}}
\left(
D_t^{\mathrm{refined}},
D_t^{\mathrm{Zivid}}
\right)
\end{equation}

only at valid Zivid anchors.

This provides sparse but highly accurate target-domain supervision.

\subsection{Experiment D: SERV-CT Static OOD Evaluation}

SERV-CT is used to measure:

\begin{itemize}[leftmargin=*]
    \item raw-preservation;
    \item geometric MAE;
    \item Bad-1 and Bad-3;
    \item new-Bad3;
    \item harmful correction rate;
    \item effect of limited target calibration.
\end{itemize}

No temporal consistency claim is made from SERV-CT.

\section{Evaluation Metrics}

\subsection{Geometric Metrics}

\begin{itemize}[leftmargin=*]
    \item disparity MAE;
    \item depth MAE;
    \item Bad-1;
    \item Bad-3;
    \item Bad-5;
    \item new-Bad3;
    \item raw-good preservation;
    \item oracle-gap recovery;
    \item boundary and interior error;
    \item SNR-stratified error on D4D.
\end{itemize}

\subsection{Temporal Metrics}

\begin{itemize}[leftmargin=*]
    \item temporal disparity jitter;
    \item temporal error jitter;
    \item high-frequency temporal energy;
    \item motion-compensated consistency;
    \item boundary jitter;
    \item residual flicker;
    \item 3D point or surface jitter;
    \item motion lag;
    \item oversmoothing;
    \item drift over a sequence.
\end{itemize}

\subsection{Safety Metrics}

\begin{itemize}[leftmargin=*]
    \item harmful correction rate;
    \item pathological new-Bad3;
    \item clean-region degradation;
    \item fraction of already-correct pixels modified;
    \item modified-pixel ratio;
    \item correction magnitude distribution.
\end{itemize}

\subsection{Efficiency Metrics}

\begin{itemize}[leftmargin=*]
    \item total parameters;
    \item trainable parameters;
    \item frozen parameters;
    \item runtime;
    \item peak VRAM;
    \item adaptation GPU-hours;
    \item amount of target data;
    \item number of target sessions;
    \item performance per trainable parameter.
\end{itemize}

\section{Required Baselines}

A strong evaluation requires:

\begin{itemize}[leftmargin=*]
    \item identity/no-refinement baseline;
    \item exponential temporal averaging;
    \item temporal median filtering;
    \item edge-aware temporal filtering;
    \item small generic U-Net or residual CNN;
    \item EGBM variants;
    \item at least one existing temporal depth/stereo refiner;
    \item raw frozen backbone;
    \item large stereo or video-stereo reference model.
\end{itemize}

The generic small refiner is particularly important. Without it, it is impossible
to determine whether the proposed architecture is genuinely more sample-efficient
than a standard lightweight network.

\section{Main Paper Narrative}

The proposed paper narrative is:

\begin{enumerate}[leftmargin=*]
    \item Modern stereo models provide strong per-frame geometry but may produce
    temporally unstable predictions in surgical video.

    \item A lightweight temporal refiner can improve the target surgical domain
    without modifying the main stereo backbone.

    \item Oracle analysis shows that conservative refiners are limited mainly by
    correction magnitude, while aggressive refiners introduce unsafe
    over-correction.

    \item Zero-shot tests demonstrate that learned correction policies do not
    automatically transfer across domains or stereo backbones.

    \item Instead of claiming universal refinement, the method is reframed as a
    lightweight temporal domain adapter.

    \item The adapter is fine-tuned using only a small fraction of target-domain
    sequences while the large stereo backbone remains frozen.

    \item SCARED validates dense temporal improvement, D4D validates sparse
    structured-light cross-domain adaptation, and SERV-CT measures geometric
    raw-preservation under domain shift.

    \item The resulting system offers a favorable trade-off between adaptation
    cost, runtime, geometric accuracy, and temporal stability compared with
    retraining or replacing a large stereo network.
\end{enumerate}

\section{Expected Contributions}

\begin{enumerate}[leftmargin=*]
    \item A lightweight causal temporal adapter for frozen stereo backbones.

    \item An oracle-driven analysis of correction support, sign, magnitude, and
    safety in learned stereo refinement.

    \item A data-efficiency study across target-domain subset sizes and stereo
    backbones.

    \item A processed D4D sparse-keyframe benchmark with structured-light ground
    truth, session/specimen-safe splits, and per-session transform conventions.

    \item A systematic evaluation separating geometric accuracy, temporal
    consistency, correction safety, and adaptation cost.

    \item Evidence that domain-specific lightweight adaptation is a practical
    alternative to retraining large stereo models.
\end{enumerate}

\section{Minimum Viable ICRA Paper}

The minimum viable paper requires:

\begin{itemize}[leftmargin=*]
    \item clear improvement of \smodel{} on SCARED;
    \item direct improvement in temporal metrics;
    \item controlled new-error and harmful-correction rates;
    \item successful training of the same adapter architecture on at least one
    additional stereo backbone;
    \item a target-data learning curve;
    \item evidence that pretraining or the proposed architecture improves
    sample efficiency;
    \item D4D sparse-keyframe external evaluation;
    \item runtime and trainable-parameter comparison;
    \item at least one strong external baseline.
\end{itemize}

\section{Strong Paper Criteria}

A stronger submission would additionally show:

\begin{itemize}[leftmargin=*]
    \item two or three stereo backbones;
    \item effective adaptation using only 5--10\% of target data;
    \item specimen-disjoint D4D adaptation;
    \item improvement over a generic lightweight refiner;
    \item favorable comparison with a large stereo/video model;
    \item 3D surface-stability improvement;
    \item a multi-backbone adapter or leave-one-backbone-out experiment.
\end{itemize}

\section{Immediate Experimental Plan}

\begin{enumerate}[leftmargin=*]
    \item Complete D4D specimen~4 extraction and conversion.

    \item Freeze D4D benchmark version, quality thresholds, manifests, and splits.

    \item Freeze the main SCARED architecture, provisionally EGBM-v3 CARE window.

    \item Finalize dense temporal metrics on SCARED.

    \item Select one second stereo backbone already supported by the ARGOS
    benchmark pipeline.

    \item Generate its raw predictions on SCARED and D4D.

    \item Train the same adapter architecture from scratch on the second backbone.

    \item Test the \smodel-trained adapter zero-shot on the second backbone.

    \item Run a reduced pilot adaptation sweep:
    \[
        0,\; 5,\; 10,\; 25,\; 100\%.
    \]

    \item If the pilot succeeds, expand to:
    \[
        0,\; 1,\; 2.5,\; 5,\; 10,\; 25,\; 50,\; 100\%.
    \]

    \item Compare from-scratch training against source-pretrained fine-tuning.

    \item Run D4D specimen/session-disjoint few-shot adaptation.

    \item Add a third backbone only after confirming the hypothesis on the second.

    \item Stop developing unrelated new architectures after the cross-backbone
    pilot.
\end{enumerate}

\section{Provisional Title Options}

\begin{itemize}[leftmargin=*]
    \item \textbf{Lightweight Temporal Adaptation for Frozen Stereo Backbones}

    \item \textbf{Sample-Efficient Temporal Refinement of Frozen Stereo Models
    for Surgical Reconstruction}

    \item \textbf{Adapting Stereo Foundation Models with Lightweight Temporal
    Refiners}

    \item \textbf{Sparse-Supervised Temporal Adaptation for Surgical Stereo Depth}

    \item \textbf{A Lightweight Temporal Adapter for Domain-Specific Stereo
    Reconstruction}
\end{itemize}

\section{Final Scope Statement}

\begin{quote}
We propose a lightweight causal temporal adapter for frozen stereo backbones.
Rather than assuming that one learned correction policy transfers universally
across domains and matchers, the adapter is efficiently specialized using a small
amount of target-domain sequential data. The method is evaluated through dense
temporal supervision on SCARED, independent sparse structured-light validation and
few-shot adaptation on D4D, and static geometric domain-shift analysis on SERV-CT.
The goal is to obtain target-specific geometric and temporal improvements while
training only approximately four million additional parameters and preserving the
large stereo backbone unchanged.
\end{quote}

% ============================================================
% FAILURE ANALYSIS
% ============================================================

\section{Failure Analysis and Scope Revision}

The initial objective was to develop a lightweight temporal refiner capable of
improving a frozen stereo backbone across surgical domains.

The experiments showed that the current model does not satisfy this objective.
Although it improves S2M2 on SCARED, the learned behavior does not transfer to
SERV-CT or D4D and does not provide measurable temporal stabilization on unseen
D4D sequences.

The main failure modes are summarized below.

\subsection{Domain-Specific Overfitting}

On SCARED, the proposed EGBM/CARE refiners improve the raw S2M2 predictions and
recover approximately 20\% of the available oracle gap while maintaining a
relatively low rate of newly introduced large errors.

However, the same SCARED-trained checkpoints degrade the raw predictions when
applied zero-shot to other surgical datasets.

\begin{table}[h]
\centering
\caption{Zero-shot behavior across surgical domains. Positive $\Delta$MAE denotes
improvement; negative values denote degradation.}
\small
\begin{tabular}{lccc}
\toprule
\textbf{Dataset}
& \textbf{Raw S2M2 MAE}
& \textbf{Best refiner $\Delta$MAE}
& \textbf{Outcome} \\
\midrule
SCARED
& 4.67 px
& positive
& improvement \\

SERV-CT
& 1.28 px
& approximately $-0.5$ px
& degradation \\

D4D
& 4.09 px
& approximately $-0.37$ px
& degradation \\
\bottomrule
\end{tabular}
\end{table}

This indicates that the refiner does not learn a domain-independent correction
rule. Instead, it learns the error distribution of the specific combination of:

\begin{itemize}[leftmargin=*]
    \item SCARED imagery;
    \item S2M2 prediction statistics;
    \item SCARED boundary and occlusion patterns;
    \item the frequency and magnitude of large disparity errors;
    \item the spatial priors present in the training sequences.
\end{itemize}

When the raw-error distribution changes, the correction policy continues to
activate even when the raw prediction is already accurate.

\subsection{False Activation on Raw-Good Pixels}

The D4D raw-error stratification exposes the failure mechanism directly.

For EGBM-v3-CARE-S:

\begin{table}[h]
\centering
\caption{Zero-shot D4D performance stratified by raw S2M2 error.
Positive $\Delta$MAE indicates improvement.}
\small
\begin{tabular}{lccc}
\toprule
\textbf{Raw error bin}
& \textbf{Raw MAE}
& \textbf{Refined MAE}
& $\boldsymbol{\Delta}$\textbf{MAE} \\
\midrule
$<1$ px
& 0.476
& 1.186
& $-0.710$ \\

$1$--$3$ px
& 1.721
& 1.966
& $-0.245$ \\

$3$--$6$ px
& 4.266
& 4.361
& $-0.095$ \\

$6$--$12$ px
& 8.197
& 7.785
& $+0.412$ \\

$>12$ px
& 26.369
& 24.726
& $+1.643$ \\
\bottomrule
\end{tabular}
\end{table}

The model retains some useful correction capability in the large-error regime.
However, it applies corrections too frequently to raw-good pixels.

Since most D4D pixels fall in the low- and medium-error regimes, the damage
introduced on already accurate predictions dominates the improvements obtained on
large errors.

The problem is therefore not only correction quality, but correction selectivity:

\begin{quote}
The model can improve large errors, but it cannot reliably determine when a
correction should be applied outside the source domain.
\end{quote}

\subsection{Calibration Recovers Safety but Collapses to Identity}

A few-shot adaptation experiment was performed on D4D using only a small calibration
subset of the network.

The calibration-only configuration updated approximately 10.9k parameters,
corresponding to roughly 0.25\% of EGBM-v3-CARE-S.

\begin{table}[h]
\centering
\caption{D4D few-shot adaptation results. The frozen raw S2M2 test MAE is 3.50 px.}
\small
\begin{tabular}{lrrrrr}
\toprule
\textbf{Configuration}
& \textbf{MAE}
& \textbf{new-Bad3}
& \textbf{Harmful}
& $\boldsymbol{\Delta}_{<1\mathrm{px}}$
& $\boldsymbol{\Delta}_{>6\mathrm{px}}$ \\
\midrule
Zero-shot
& 4.72
& 23.5
& 0.83
& $-1.44$
& $+0.86$ \\

Calibration-only, 2 sessions
& 3.66
& 1.9
& 0.34
& $-0.05$
& $+0.08$ \\

Calibration-only, 8 sessions
& 3.50
& 0.0
& 0.04
& $\approx 0$
& $\approx 0$ \\

Full fine-tuning, 4 sessions
& 4.09
& 3.4
& 0.61
& $-0.23$
& $+0.58$ \\

From scratch, 4 sessions
& 3.47
& 1.2
& 0.13
& $-0.03$
& $+0.14$ \\
\bottomrule
\end{tabular}
\end{table}

Calibration-only adaptation successfully suppresses false activation:

\begin{itemize}[leftmargin=*]
    \item raw-good degradation is reduced by approximately 96\%;
    \item new-Bad3 falls from 23.5\% to almost zero;
    \item the harmful correction rate decreases substantially.
\end{itemize}

However, this safety improvement is obtained by suppressing almost all corrections.

At eight sessions, the calibrated model approximately reproduces the raw S2M2
output:

\begin{equation}
D_t^{\mathrm{refined}}
\approx
D_t^{\mathrm{raw}}.
\end{equation}

The model becomes safe, but no longer provides useful refinement.

Therefore, lightweight calibration transfers abstention, not beneficial correction.

\subsection{Limited Evidence of Transferable Correction Skill}

Full fine-tuning preserves more of the large-error correction ability:

\begin{equation}
\Delta_{\mathrm{MAE},\,>6\mathrm{px}} = +0.58\ \mathrm{px}.
\end{equation}

However, with only a few target anchors it still over-corrects raw-good regions and
does not outperform the raw stereo backbone in aggregate MAE.

The from-scratch model reaches approximately the raw performance, while the
pretrained model does not demonstrate a clear sample-efficiency advantage.

This weakens the original hypothesis that SCARED pretraining learns a transferable
temporal correction skill.

The current evidence supports only the following weaker statement:

\begin{quote}
A small target-domain calibration set can rapidly suppress unsafe corrections, but
more data or a different objective is required to obtain net-beneficial refinement.
\end{quote}

\subsection{No Temporal-Stability Improvement on D4D}

The most important negative result is that the refiner does not improve temporal
stability on full D4D clips.

The evaluation used real consecutive stereo frames and motion-compensated temporal
diagnostics based on optical flow and forward--backward occlusion filtering.

\begin{table}[h]
\centering
\caption{D4D temporal evaluation. The temporal metrics are prediction-space or
motion-compensated diagnostics, since dense temporal GT is unavailable.}
\small
\begin{tabular}{lrrrr}
\toprule
\textbf{Configuration}
& \textbf{Anchor MAE}
& \textbf{MC inconsistency}
& \textbf{HF energy}
& \textbf{Depth MC [mm]} \\
\midrule
Raw S2M2
& 3.50
& \textbf{0.390}
& \textbf{0.507}
& \textbf{0.691} \\

Zero-shot
& 4.72
& 0.397
& 0.525
& 0.717 \\

Calibration-only, 8 sessions
& 3.50
& 0.400
& 0.525
& 0.723 \\

Full fine-tuning, 4 sessions
& 4.09
& 0.394
& 0.521
& 0.706 \\

From scratch, 4 sessions
& 3.47
& \textbf{0.390}
& \textbf{0.507}
& \textbf{0.691} \\
\bottomrule
\end{tabular}
\end{table}

The temporal axis is effectively flat:

\begin{itemize}[leftmargin=*]
    \item motion-compensated inconsistency changes by only a few percent;
    \item high-frequency temporal energy slightly increases;
    \item motion-compensated depth inconsistency slightly increases;
    \item pretrained full fine-tuning is not more stable than the scratch identity
    solution;
    \item calibration-only adaptation reduces correction magnitude but does not
    reduce temporal instability.
\end{itemize}

The refiner therefore does not act as a temporal stabilizer on D4D.

Instead, it behaves primarily as a frame-wise spatial residual corrector.

\subsection{The Architecture Is Not Forced to Use Time}

The current architecture receives temporal inputs, but it can minimize the training
objective without learning a genuinely temporal function.

Its inputs include:

\begin{itemize}[leftmargin=*]
    \item current disparity;
    \item previous disparities;
    \item temporal differences;
    \item spatial gradients;
    \item edge features;
    \item local statistics;
    \item validity masks.
\end{itemize}

However, the network can still rely mainly on the current-frame disparity and
dataset-specific spatial patterns.

The training target is the full residual:

\begin{equation}
R_t^{\star}
=
D_t^{\mathrm{GT}}
-
D_t^{\mathrm{raw}},
\end{equation}

which contains both:

\begin{itemize}[leftmargin=*]
    \item static spatial error;
    \item temporal flicker;
    \item bias;
    \item boundary mismatch;
    \item large geometric failure.
\end{itemize}

Nothing in the objective explicitly forces the model to separate temporal
inconsistency from generic per-frame error.

As a result, the easiest solution is a domain-specific spatial correction policy.

\subsection{Longer Temporal Memory Did Not Solve the Problem}

A persistent streaming variant was compared with a short causal replay window.

The longer-memory model did not improve over the four-frame window.

This suggests that the main limitation is not insufficient temporal context.

The failure is more fundamental:

\begin{quote}
The model is not learning the correct temporal task, rather than merely lacking
enough history.
\end{quote}

\subsection{Large-Correction Models Expose a Safety Trade-Off}

The MPC, CPV, and ASSR variants confirmed that the correction-magnitude bottleneck
is real.

Allowing larger corrections increased oracle-gap recovery from approximately 20\%
to approximately 45--47\%.

However, this also caused:

\begin{itemize}[leftmargin=*]
    \item higher new-Bad3;
    \item more harmful corrections;
    \item worse full-GT performance;
    \item severe degradation under domain shift.
\end{itemize}

The resulting Pareto frontier is:

\begin{equation}
\text{large correction capacity}
\quad\Longleftrightarrow\quad
\text{higher risk of unsafe over-correction}.
\end{equation}

Post-hoc damping restored safety only by reducing the model toward a near-identity
solution.

\subsection{Resolution Sensitivity of Large-Proposal Branches}

MPC and CPV also exposed an implementation-level limitation.

Their large-context path assumed feature-map dimensions compatible with the SCARED
training resolution.

On D4D, the odd target-grid height caused incompatible tensor sizes in the
multi-scale fusion path.

Although this issue is technically fixable through padding or explicit alignment,
it reveals that the large-proposal branch was not yet robust to arbitrary deployment
resolutions.

These branches are therefore retained only as analytical evidence for the
magnitude bottleneck, not as deployable models.

\subsection{Summary of Negative Findings}

The current experiments support the following conclusions:

\begin{enumerate}[leftmargin=*]
    \item The refiner improves S2M2 on SCARED but overfits strongly to the SCARED
    error distribution.

    \item Zero-shot transfer fails on both SERV-CT and D4D.

    \item The model helps large errors but damages already accurate predictions.

    \item Minimal calibration restores safety mainly by suppressing the correction
    toward identity.

    \item Full fine-tuning retains some large-error capability but does not yet
    produce a net target-domain benefit.

    \item SCARED pretraining does not provide a measurable temporal-stability
    advantage on D4D.

    \item The current architecture behaves primarily as a spatial residual
    corrector, despite receiving temporal inputs.

    \item Longer temporal memory does not solve the problem.

    \item Increasing correction magnitude improves oracle recovery but worsens
    safety and OOD behavior.

    \item The current evidence does not support the claim of a transferable
    temporal refiner.
\end{enumerate}

\subsection{Implication for the Next Model}

The failure analysis motivates a strict redesign.

The next model should not predict an unrestricted per-frame residual.

Instead, it should be explicitly constrained to operate on the motion-aligned
temporal inconsistency:

\begin{equation}
R_t^{\mathrm{temp}}
=
D_t^{\mathrm{raw}}
-
\mathcal{W}
\left(
D_{t-1}^{\mathrm{raw}},
F_{t-1\rightarrow t}
\right),
\end{equation}

where $\mathcal{W}$ denotes warping according to image motion and
$F_{t-1\rightarrow t}$ is the estimated optical flow.

The refined disparity should take the form:

\begin{equation}
D_t^{\mathrm{refined}}
=
D_t^{\mathrm{raw}}
-
G_t
\odot
R_t^{\mathrm{temp}},
\end{equation}

where $G_t$ is a bounded confidence gate.

This formulation restricts the model to removing temporally unexplained components
instead of learning arbitrary spatial corrections.

The next architecture must also be validated with mandatory temporal ablations:

\begin{itemize}[leftmargin=*]
    \item full temporal model;
    \item current-frame-only model;
    \item shuffled previous-frame model;
    \item warped-history model without learning;
    \item identity baseline.
\end{itemize}

If removing or shuffling the temporal context does not reduce performance, the model
cannot be considered a temporal refiner.

% ============================================================
% TEMPORAL SUPERVISION CONFIRMATION STUDY
% ============================================================

\section{Does Explicit Temporal Supervision Induce Genuine Temporal Reasoning?}

After observing that the original EGBM-based refiners behaved primarily as
domain-specific spatial residual correctors, we investigated whether their failure
was caused by the absence of explicit temporal supervision.

To isolate this question, we implemented a lightweight causal baseline inspired by
video depth post-processing methods. The model consisted of a shared geometric
encoder followed by a causal ConvGRU and a bounded residual head. The stereo
backbone remained frozen, and the refiner received only disparity-derived inputs.

The main hypothesis was:

\begin{quote}
If the model is trained with dense temporal-gradient supervision, then it should
learn to exploit causal history and improve temporal consistency beyond what can
be achieved by a frame-wise regularizer.
\end{quote}

The initial pilot appeared to support this hypothesis. However, its ablation design
was later found to be confounded: only the full temporal model was trained with the
temporal loss, while the current-frame and shuffled-history variants were trained
without it. Therefore, the observed separation could have originated from the loss
rather than from genuine use of temporal history.

A dedicated confirmation study was consequently conducted.

% ------------------------------------------------------------
\subsection{Temporal Gradient Matching}

The temporal supervision was implemented using a Temporal Gradient Matching
loss.

For two consecutive refined disparity maps
$\hat{D}_{t-1}$ and $\hat{D}_{t}$, and the corresponding ground-truth maps
$D^{GT}_{t-1}$ and $D^{GT}_{t}$, the loss is defined as

\begin{equation}
\mathcal{L}_{\mathrm{TGM}}
=
\frac{1}{|\Omega_t|}
\sum_{\mathbf{x}\in\Omega_t}
\left|
\left(
\hat{D}_{t}(\mathbf{x})
-
\hat{D}_{t-1}(\mathbf{x})
\right)
-
\left(
D^{GT}_{t}(\mathbf{x})
-
D^{GT}_{t-1}(\mathbf{x})
\right)
\right|,
\end{equation}

where $\Omega_t$ denotes pixels with valid ground truth in both frames.

The complete objective was

\begin{equation}
\mathcal{L}
=
\mathcal{L}_{\mathrm{spatial}}
+
\lambda_{\mathrm{TGM}}
\mathcal{L}_{\mathrm{TGM}}
+
\lambda_{\mathrm{safe}}
\mathcal{L}_{\mathrm{safe}},
\end{equation}

where

\begin{itemize}[leftmargin=*]
    \item $\mathcal{L}_{\mathrm{spatial}}$ supervises per-frame disparity accuracy;
    \item $\mathcal{L}_{\mathrm{TGM}}$ supervises frame-to-frame disparity changes;
    \item $\mathcal{L}_{\mathrm{safe}}$ discourages unnecessary modification of
    raw-good pixels.
\end{itemize}

The refined disparity was predicted through a bounded residual:

\begin{equation}
\hat{D}_t
=
D_t^{\mathrm{raw}}
+
\alpha
\tanh
\left(
R_t
\right),
\end{equation}

where $\alpha$ limits the maximum correction magnitude.

% ------------------------------------------------------------
\subsection{Factorial Experimental Design}

To separate temporal input structure from loss supervision, we decoupled two
factors:

\begin{enumerate}[leftmargin=*]
    \item temporal input mode;
    \item loss mode.
\end{enumerate}

The temporal input modes were:

\begin{itemize}[leftmargin=*]
    \item \textbf{full history}: the ConvGRU receives the correct causal sequence;
    \item \textbf{current frame only}: the recurrent state is reset at every frame;
    \item \textbf{shuffled history}: the current frame remains the target, while the
    historical frames are randomly reordered.
\end{itemize}

The loss modes were:

\begin{itemize}[leftmargin=*]
    \item spatial supervision only;
    \item spatial supervision plus TGM.
\end{itemize}

The core factorial comparison was therefore

\begin{table}[h]
\centering
\caption{Confirmation-study matrix. The decisive comparisons use the same TGM loss,
same architecture, and same training budget.}
\small
\begin{tabular}{lll}
\toprule
\textbf{Configuration}
& \textbf{Temporal input}
& \textbf{Loss} \\
\midrule
Full history + spatial
& Correct causal history
& Spatial only \\

Full history + TGM
& Correct causal history
& Spatial + TGM \\

Current frame only + TGM
& No temporal memory
& Spatial + TGM \\

Shuffled history + TGM
& Corrupted history order
& Spatial + TGM \\
\bottomrule
\end{tabular}
\end{table}

All configurations used:

\begin{itemize}[leftmargin=*]
    \item the same causal ConvGRU architecture;
    \item clip length of eight frames;
    \item frozen S2M2 stereo predictions;
    \item the same SCARED sequence-disjoint train/validation/test splits;
    \item identical optimizer and training duration;
    \item three independent random seeds.
\end{itemize}

This design tests whether the improvement is due to

\begin{equation}
\text{temporal supervision},
\end{equation}

or specifically due to

\begin{equation}
\text{temporal supervision}
+
\text{correct causal history}.
\end{equation}

% ------------------------------------------------------------
\subsection{SCARED Results}

Table~\ref{tab:tgm_confirmation} reports the temporal-gradient error on the frozen
SCARED test split.

\begin{table}[h]
\centering
\caption{SCARED temporal-gradient error across three seeds.
Lower is better.}
\label{tab:tgm_confirmation}
\small
\begin{tabular}{lccc}
\toprule
\textbf{Configuration}
& \textbf{Mean}
& \textbf{Std.}
& \textbf{Seed values} \\
\midrule
Full history + spatial
& 1.7023
& 0.0117
& 1.6934,\;1.6946,\;1.7189 \\

Full history + TGM
& \textbf{1.6456}
& 0.0534
& 1.6121,\;1.7210,\;1.6038 \\

Current frame only + TGM
& 1.6816
& 0.0033
& 1.6781,\;1.6861,\;1.6806 \\

Shuffled history + TGM
& 1.6501
& 0.0107
& 1.6380,\;1.6641,\;1.6482 \\
\bottomrule
\end{tabular}
\end{table}

The first comparison confirms that TGM is beneficial:

\begin{equation}
1.7023
\rightarrow
1.6456.
\end{equation}

Thus, adding dense temporal-gradient supervision reduces the average temporal error.

However, the comparisons against the temporal ablations are not robust.

The full-history model only marginally outperforms shuffled history:

\begin{equation}
1.6456
\quad \text{vs.} \quad
1.6501,
\end{equation}

corresponding to an average difference of approximately $0.3\%$.

Moreover, the seed distributions overlap substantially:

\begin{equation}
\max
\left(
E_{\mathrm{full}}
\right)
=
1.7210
>
\min
\left(
E_{\mathrm{shuffled}}
\right)
=
1.6380.
\end{equation}

The worst full-history run is also worse than every current-frame-only run.

Therefore, the group mean follows the expected direction, but the result is not
stable enough to support the claim that the ConvGRU reliably exploits causal history.

% ------------------------------------------------------------
\subsection{Geometric Accuracy and Safety}

The temporal objective did not produce a clear geometric separation between the
three TGM configurations.

\begin{table}[h]
\centering
\caption{Mean geometric behavior of the three TGM configurations on SCARED.}
\small
\begin{tabular}{lcc}
\toprule
\textbf{Configuration}
& \textbf{MAE [px]}
& \textbf{Interpretation} \\
\midrule
Full history + TGM
& 5.130
& no clear advantage \\

Current frame only + TGM
& 5.131
& effectively equivalent \\

Shuffled history + TGM
& 5.150
& within experimental noise \\
\bottomrule
\end{tabular}
\end{table}

The harmful-correction rate also showed high variance across seeds and no stable
ordering between the configurations.

This indicates that the full-history architecture did not learn a reliably superior
correction policy.

% ------------------------------------------------------------
\subsection{TGM Weight Sweep}

To investigate whether the high variance was caused by an overly aggressive
temporal weight, we tested

\begin{equation}
\lambda_{\mathrm{TGM}}
\in
\{0.2,\;0.5,\;1.0\}.
\end{equation}

\begin{table}[h]
\centering
\caption{Effect of the TGM weight for the full-history model.}
\small
\begin{tabular}{lccc}
\toprule
$\boldsymbol{\lambda_{\mathrm{TGM}}}$
& \textbf{TGM error}
& \textbf{MAE [px]}
& \textbf{Harmful rate} \\
\midrule
0.2
& $1.6887 \pm 0.0147$
& $5.118 \pm 0.009$
& $0.303 \pm 0.068$ \\

0.5
& $1.6619 \pm 0.0081$
& $5.137 \pm 0.050$
& $0.327 \pm 0.042$ \\

1.0
& $1.6456 \pm 0.0534$
& $5.130 \pm 0.079$
& $0.234 \pm 0.167$ \\
\bottomrule
\end{tabular}
\end{table}

The highest weight produced the best mean temporal score but also the highest
variance.

Lower weights were more stable, suggesting a trade-off between nominal temporal
improvement and optimization robustness.

Importantly, none of the tested weights produced a clear and stable separation from
the current-frame or shuffled-history variants.

% ------------------------------------------------------------
\subsection{D4D Zero-Shot Evaluation}

The temporal-usage hypothesis was further tested on D4D.

Because D4D does not provide dense temporal ground truth, the evaluation used
prediction-space motion-compensated diagnostics and sparse Zivid ground-truth
anchors.

The primary temporal diagnostic was motion-compensated disparity inconsistency.

For each configuration and seed, we measured

\begin{equation}
\Delta E_{\mathrm{MC}}
=
E_{\mathrm{MC}}^{\mathrm{refined}}
-
E_{\mathrm{MC}}^{\mathrm{raw}},
\end{equation}

where negative values indicate improvement.

\begin{table}[h]
\centering
\caption{D4D motion-compensated inconsistency relative to raw S2M2.
Values are mean paired differences with bootstrap confidence intervals.}
\small
\begin{tabular}{lccc}
\toprule
\textbf{Configuration}
& \textbf{Seed 0}
& \textbf{Seed 1}
& \textbf{Seed 2} \\
\midrule
Full history
& $+0.0001$
& $0.0000$
& $+0.0029$ \\

Current frame only
& $-0.0020$
& $-0.0020$
& $-0.0038$ \\

Shuffled history
& $-0.0006$
& $+0.0018$
& $+0.0001$ \\
\bottomrule
\end{tabular}
\end{table}

The full-history model does not show a stable improvement.

Its effect is approximately zero or slightly negative in all three seeds.

Unexpectedly, the current-frame-only model shows the most consistent reduction in
motion-compensated inconsistency.

This strongly suggests that the small temporal improvement observed on SCARED is
not caused by robust exploitation of causal history.

Instead, TGM appears to regularize the output even when no temporal memory is
available.

% ------------------------------------------------------------
\subsection{Sparse Geometric Transfer on D4D}

All three TGM variants remained geometrically safe on the sparse Zivid anchors.

The raw S2M2 MAE was approximately

\begin{equation}
4.090\ \mathrm{px}.
\end{equation}

The different configurations produced small improvements of similar magnitude.

\begin{table}[h]
\centering
\caption{D4D sparse-anchor zero-shot behavior. Positive $\Delta$MAE indicates
improvement over raw S2M2.}
\small
\begin{tabular}{lccc}
\toprule
\textbf{Configuration}
& $\boldsymbol{\Delta}$\textbf{MAE [px]}
& \textbf{Harmful rate}
& \textbf{Modified ratio} \\
\midrule
Full history
& $0.00$--$0.04$
& $0.00$--$0.02$
& $0.00$--$0.02$ \\

Current frame only
& $0.04$--$0.05$
& $0.00$--$0.01$
& $0.01$--$0.02$ \\

Shuffled history
& $0.03$--$0.05$
& $0.01$--$0.08$
& $0.01$--$0.02$ \\
\bottomrule
\end{tabular}
\end{table}

The variants are all conservative and do not exhibit the catastrophic
false-activation observed with EGBM.

However, their similar behavior again suggests that the transferable component is
a conservative correction prior rather than a learned causal temporal mechanism.

% ------------------------------------------------------------
\subsection{Five-Gate Decision}

The confirmation study used five predefined gates.

\begin{table}[h]
\centering
\caption{Temporal-usage decision gates.}
\small
\begin{tabular}{lp{7cm}c}
\toprule
\textbf{Gate}
& \textbf{Criterion}
& \textbf{Outcome} \\
\midrule
1
& Full history + TGM beats full history + spatial-only
& Pass \\

2
& Full history + TGM robustly beats current-frame-only + TGM
& Marginal \\

3
& Full history + TGM robustly beats shuffled-history + TGM
& Marginal \\

4
& Full-history improvement transfers to D4D with stable sign
& Fail \\

5
& No identity collapse or unsafe false activation
& Pass \\
\bottomrule
\end{tabular}
\end{table}

The resulting verdict is:

\begin{quote}
\textbf{Marginal --- genuine causal-history usage is not confirmed.}
\end{quote}

% ------------------------------------------------------------
\subsection{Interpretation}

The study supports two distinct conclusions.

First, explicit temporal supervision is useful:

\begin{equation}
\mathcal{L}_{\mathrm{TGM}}
\quad
\text{reduces temporal-gradient error}.
\end{equation}

Second, the causal ConvGRU does not demonstrate a reliable ability to exploit the
correct temporal history:

\begin{equation}
E_{\mathrm{full+TGM}}
\approx
E_{\mathrm{shuffled+TGM}}
\approx
E_{\mathrm{current+TGM}}.
\end{equation}

The most plausible interpretation is that TGM acts as an output regularizer.

It encourages conservative frame-to-frame changes, but the network is not forced to
build accurate temporal correspondences.

The current architecture can therefore satisfy the objective through a mostly
frame-wise correction policy.

% ------------------------------------------------------------
\subsection{Implication for the Model Design}

The result clarifies the next methodological step.

TGM should be retained as a training objective, but the disparity-only ConvGRU
should not be used as the final temporal architecture.

The next refiner should explicitly construct correspondences between the target
frame and causal reference frames.

A suitable formulation is:

\begin{equation}
\mathbf{q}_t
=
\phi
\left(
I_t,
D_t^{\mathrm{raw}}
\right),
\end{equation}

\begin{equation}
\mathbf{k}_{t-\tau},
\mathbf{v}_{t-\tau}
=
\phi
\left(
I_{t-\tau},
D_{t-\tau}^{\mathrm{raw}}
\right),
\qquad
\tau \in \{1,2,3\},
\end{equation}

followed by local target-to-history matching:

\begin{equation}
\mathbf{z}_t
=
\mathrm{CrossAttention}
\left(
\mathbf{q}_t,
\{\mathbf{k}_{t-\tau}\},
\{\mathbf{v}_{t-\tau}\}
\right).
\end{equation}

The final correction should remain bounded:

\begin{equation}
\hat{D}_t
=
D_t^{\mathrm{raw}}
+
G_t
\odot
\tanh
\left(
R_t
\right),
\end{equation}

where $G_t$ is a learned confidence gate.

The training objective should combine:

\begin{equation}
\mathcal{L}
=
\mathcal{L}_{\mathrm{spatial}}
+
\lambda_{\mathrm{TGM}}
\mathcal{L}_{\mathrm{TGM}}
+
\lambda_{\mathrm{warp}}
\mathcal{L}_{\mathrm{warp}}
+
\lambda_{\mathrm{safe}}
\mathcal{L}_{\mathrm{safe}}.
\end{equation}

Here, $\mathcal{L}_{\mathrm{warp}}$ should supervise motion-aligned temporal
consistency:

\begin{equation}
\mathcal{L}_{\mathrm{warp}}
=
\frac{1}{|\Omega_t^{\mathrm{flow}}|}
\sum_{\mathbf{x}\in\Omega_t^{\mathrm{flow}}}
\left|
\hat{D}_t(\mathbf{x})
-
\mathcal{W}
\left(
\hat{D}_{t-1},
F_{t-1\rightarrow t}
\right)
(\mathbf{x})
\right|,
\end{equation}

where

\begin{itemize}[leftmargin=*]
    \item $\mathcal{W}$ denotes warping;
    \item $F_{t-1\rightarrow t}$ is optical flow;
    \item $\Omega_t^{\mathrm{flow}}$ excludes occlusions and unreliable flow.
\end{itemize}

This design makes temporal matching explicit rather than expecting it to emerge
implicitly from recurrent state.

% ------------------------------------------------------------
\subsection{Final Conclusion}

The confirmation study revises the interpretation of the original pilot.

The initial result did not demonstrate robust temporal reasoning.

Instead, the experiments show that:

\begin{enumerate}[leftmargin=*]
    \item TGM is a useful temporal regularizer;
    \item the current ConvGRU does not reliably exploit causal history;
    \item shuffled and current-frame variants perform similarly to the full model;
    \item D4D does not show stable temporal transfer for the full-history model;
    \item the conservative geometric behavior transfers better than the temporal
    mechanism;
    \item the next model requires explicit target-to-history correspondence.
\end{enumerate}

Therefore, the disparity-only ConvGRU branch is retained as a diagnostic baseline,
while the next model will combine explicit temporal matching, motion-aware
supervision, bounded correction, and surgical-domain visual features.

\section{ARGOS v2}
\label{sec:argos_v2}

ARGOS v2 extends the original frame-wise stereo evaluation into a unified framework for
temporally consistent surgical depth estimation. The objective is not only to reduce
frame-to-frame flickering, but also to preserve or improve geometric accuracy under camera
motion, tissue deformation, occlusions, specularities, and instrument interaction.

A central limitation of existing surgical stereo datasets is that no single benchmark
simultaneously provides long stereo sequences, dense metric depth independently measured at
every frame, realistic tissue deformation, and sufficient scale for supervised temporal
training. ARGOS v2 therefore adopts a complementary multi-dataset protocol in which each
dataset is assigned a clearly defined role.

\subsection{Dataset Roles}

\paragraph{SCARED.}
The original SCARED dataset provides 45 structured-light keyframes with directly measured
geometric ground truth. These keyframes are used for static geometric evaluation. Since
datasets 4 and 5 contain known calibration issues, the main clean evaluation protocol uses
the remaining 35 keyframes, while the complete 45-keyframe benchmark is reported separately.
The intermediate stereo video frames are retained only for no-reference temporal analysis,
as their original depth annotations were propagated using inaccurate robot kinematics.

\paragraph{SCARED-C.}
SCARED-C replaces the original kinematics-based camera poses with poses estimated using
Structure-from-Motion and subsequently aligned to metric scale using the structured-light
keyframe depth. The resulting temporal depth is therefore not independently measured at
every frame, but is derived by reprojecting the structured-light keyframe geometry through
the corrected camera poses.

The official release contains 17,135 RGB-D pairs. After sequence-level quality control,
the ARGOS v2 curated subset contains 16,921 temporally aligned stereo frames with metric
pseudo-ground truth. SCARED-C is consequently used as the primary dataset for supervised
training, validation, and in-domain evaluation of the temporal refinement module.

\paragraph{D4D.}
D4D provides sparse but high-quality Zivid structured-light geometry anchors within long
sequences containing substantial non-rigid deformation and camera motion. It is used as an
out-of-distribution benchmark for evaluating geometric fidelity at measured anchor frames
and temporal robustness between anchors.

\paragraph{SERV-CT.}
SERV-CT contains 16 stereo pairs with geometry derived from registered CT data. Although too
small for training, it provides an independent static out-of-distribution test of geometric
accuracy and consistency.

\paragraph{StereoMIS.}
StereoMIS provides long in-vivo stereo sequences containing respiration, tissue deformation,
specularities, occlusions, and surgical instruments, but no dense per-frame depth ground
truth. It is therefore used for label-free domain adaptation and no-reference temporal
evaluation. To preserve a meaningful in-vivo generalization test, adaptation and evaluation
are separated at the subject level.

\begin{table}[t]
    \centering
    \caption{Role of each dataset in the ARGOS v2 evaluation protocol.}
    \label{tab:argos_v2_datasets}
    \resizebox{\columnwidth}{!}{
    \begin{tabular}{lcccc}
        \toprule
        Dataset &
        Geometric reference &
        Temporal sequences &
        Dense temporal GT &
        ARGOS v2 role \\
        \midrule
        SCARED &
        Structured light at 45 keyframes &
        Yes &
        No reliable per-frame GT &
        Static geometry and no-reference temporal test \\
        
        SCARED-C &
        Structured-light-derived metric pseudo-GT &
        Yes &
        16,921 curated frames &
        Supervised temporal train/validation/test \\
        
        D4D &
        Zivid structured-light anchors &
        Yes &
        Sparse anchors only &
        OOD geometry and deformation test \\
        
        SERV-CT &
        CT-derived geometry on 16 stereo pairs &
        No &
        No &
        Static OOD geometry test \\
        
        StereoMIS &
        No dense depth GT &
        Yes, in vivo &
        No &
        Unsupervised adaptation and held-out in-vivo test \\
        \bottomrule
    \end{tabular}
    }
\end{table}

\subsection{Multi-Backbone Temporal Refinement}

The temporal module is trained as a post-processing refiner over predictions produced by
multiple frozen stereo backbones. For each SCARED-C training clip, the input disparity
sequence is sampled from a heterogeneous pool of frame-wise stereo models. The refiner
therefore learns from different error distributions rather than specializing to a single
backbone.

Given the current stereo pair $(I_t^L, I_t^R)$, the raw backbone disparity
$\hat{d}_t^{\,\mathrm{raw}}$, and the previous refined state, the temporal module predicts a
residual correction $\Delta d_t$ and an update gate $g_t$:

\begin{equation}
    \hat{d}_t^{\,\mathrm{ref}}
    =
    \hat{d}_t^{\,\mathrm{raw}}
    +
    g_t \odot \Delta d_t,
\end{equation}

where $g_t \in [0,1]$ controls how strongly temporal information is allowed to modify the
current frame. This formulation enables the model to preserve an already accurate
frame-wise estimate while correcting unstable or temporally inconsistent regions.

The main training objective combines geometric accuracy, temporal consistency, edge
preservation, and regression safety:

\begin{equation}
    \mathcal{L}
    =
    \lambda_{\mathrm{geo}} \mathcal{L}_{\mathrm{geo}}
    +
    \lambda_{\mathrm{temp}} \mathcal{L}_{\mathrm{temp}}
    +
    \lambda_{\mathrm{edge}} \mathcal{L}_{\mathrm{edge}}
    +
    \lambda_{\mathrm{safe}} \mathcal{L}_{\mathrm{safe}}
    +
    \lambda_{\mathrm{upd}} \mathcal{L}_{\mathrm{upd}}.
\end{equation}

The safety term explicitly penalizes cases in which temporal refinement makes the prediction
less accurate than the original backbone output:

\begin{equation}
    \mathcal{L}_{\mathrm{safe}}
    =
    \max
    \left(
        0,
        \left|
        \hat{d}_t^{\,\mathrm{ref}} - d_t^{*}
        \right|
        -
        \left|
        \hat{d}_t^{\,\mathrm{raw}} - d_t^{*}
        \right|
        -
        m
    \right),
\end{equation}

where $d_t^{*}$ denotes the SCARED-C metric pseudo-ground truth and $m$ is a small tolerance
margin.

\subsection{Training and Evaluation Protocol}

The complete ARGOS v2 protocol is divided into four stages:

\begin{enumerate}
    \item \textbf{Supervised temporal training.}
    The refiner is trained on SCARED-C using cached disparity predictions from multiple
    frozen stereo backbones.

    \item \textbf{In-domain model selection.}
    Architecture selection, loss weighting, checkpoint selection, and early stopping are
    performed exclusively on a sequence-disjoint SCARED-C validation split.

    \item \textbf{Label-free in-vivo adaptation.}
    The trained refiner may be adapted on selected StereoMIS subjects using stereo
    photometric consistency, left--right consistency, occlusion-aware temporal consistency,
    and teacher--student regularization. The frame-wise stereo backbone remains frozen.

    \item \textbf{Frozen external evaluation.}
    The final model is evaluated without additional tuning on held-out SCARED-C sequences,
    held-out StereoMIS subjects, SCARED structured-light keyframes, SERV-CT, and D4D.
\end{enumerate}

\subsection{Evaluation Questions}

ARGOS v2 is designed to answer the following questions:

\begin{enumerate}
    \item Does temporal refinement reduce flickering while preserving or improving metric
    geometric accuracy?

    \item Does a refiner trained with predictions from multiple stereo backbones generalize
    to an unseen backbone?

    \item Does the learned temporal prior transfer zero-shot from SCARED-C to independent
    surgical datasets?

    \item Can label-free adaptation on in-vivo StereoMIS sequences improve temporal behavior
    without causing geometric degradation on SCARED, SERV-CT, or D4D?

    \item Does the method remain stable around depth discontinuities, occlusions, surgical
    instruments, and non-rigid tissue deformation?
\end{enumerate}

\subsection{Evaluation Metrics}

On datasets providing geometric reference, the evaluation reports disparity EPE, Bad-3,
depth MAE, AbsRel, RMSE, and $\delta_1$. Temporal performance is evaluated using
motion-compensated disparity consistency, temporal error variance, the fraction of frames
improved over the raw backbone, and the frequency and magnitude of geometric regressions.

For datasets without dense depth ground truth, temporal evaluation relies on stereo
photometric consistency, left--right consistency, motion-compensated temporal residuals,
edge stability, and qualitative analysis of ghosting, occlusions, and newly visible regions.

The primary claim of ARGOS v2 is therefore not merely that the refined depth is smoother,
but that temporal refinement produces more stable predictions without sacrificing geometric
fidelity, and that this behavior generalizes across stereo backbones and surgical domains.

\section{Positioning ARGOS v2 Against the Temporal Stereo State of the Art}
\label{sec:argos_v2_positioning}

The recent literature on temporally consistent depth estimation contains several strong and partially overlapping lines of research. However, the existing approaches address different subsets of the problem targeted by ARGOS v2. In particular, three recent methods are especially relevant: EndoStreamDepth, BiDAStabilizer, and PPMStereo.

\subsection{EndoStreamDepth: Causal Hierarchical Temporal Modeling for Endoscopy}

EndoStreamDepth is a streaming monocular depth estimation framework designed specifically for endoscopic videos. Its architecture combines a Depth Anything V2-based encoder, a DPT-style decoder, and hierarchical temporal Mamba modules inserted at multiple decoder resolutions. Each decoder level maintains a recurrent hidden state that is propagated sequentially across frames, allowing the model to process arbitrarily long video streams in a causal manner.

The main architectural contribution is therefore a multi-scale temporal memory:

\begin{equation}
\mathbf{H}_{t}^{(l)}
=
\mathcal{M}^{(l)}
\left(
\mathbf{F}_{t}^{(l)},
\mathbf{H}_{t-1}^{(l)}
\right),
\end{equation}

where $\mathbf{F}_{t}^{(l)}$ denotes the current feature representation at scale $l$, and $\mathbf{H}_{t-1}^{(l)}$ is the temporal hidden state propagated from the previous frame.

This hierarchical structure allows the network to model both fine anatomical boundaries and coarse global geometry. EndoStreamDepth further introduces endoscopy-specific transformations, including rotations, blur, defocus, brightness and gamma perturbations, fog, and motion artifacts, in order to improve robustness to surgical imaging conditions.

The method is highly relevant to ARGOS v2 because it demonstrates that:

\begin{itemize}
    \item causal streaming inference is feasible for endoscopic depth estimation;
    \item temporal memory should operate at multiple spatial resolutions;
    \item endoscopy-specific data augmentation can substantially improve robustness;
    \item temporal consistency and anatomical boundary preservation must be optimized jointly.
\end{itemize}

However, EndoStreamDepth is not a plug-in temporal refiner. It modifies the full monocular depth network and relies on backbone-specific encoder and decoder features. Consequently, it cannot be directly applied to arbitrary stereo predictors such as S2M2, RAFT-Stereo, StereoAnywhere, or FoundationStereo. In addition, it does not explicitly include an identity-preserving mechanism that prevents modifications when the input prediction is already correct.

Its temporal regularization is based on direct frame-to-frame variation of normalized depth:

\begin{equation}
\mathcal{L}_{\mathrm{temp}}
=
\frac{1}{(T-1)N}
\sum_{t=1}^{T-1}
\sum_{i=1}^{N}
\left|
\bar{D}_{t+1}(i)
-
\bar{D}_{t}(i)
\right|,
\end{equation}

which may encourage smoothness, but does not explicitly compensate for camera motion, tissue deformation, or newly appearing structures.

\subsection{BiDAStabilizer: Plug-and-Play Stereo Stabilization}

BiDAStabilizer is the closest existing method to the ARGOS v2 plug-in formulation. It receives disparity predictions from a frozen image-based stereo network and transforms them into temporally consistent disparity sequences without retraining the original stereo matcher.

The pipeline can be summarized as follows:

\begin{equation}
\left\{
d_{t-1},
d_t,
d_{t+1}
\right\}
\rightarrow
\text{flow-based alignment}
\rightarrow
\text{temporal propagation}
\rightarrow
\Delta d_t,
\end{equation}

with final prediction

\begin{equation}
d_t^{\mathrm{out}}
=
d_t^{\mathrm{raw}}
+
\Delta d_t.
\end{equation}

The method first estimates bidirectional optical flow, then aligns neighboring disparity maps toward the central frame. A forward hidden state and a backward hidden state are propagated across the sequence and fused to predict the final disparity residual.

BiDAStabilizer establishes several important principles:

\begin{itemize}
    \item temporal aggregation should operate on motion-aligned disparity information;
    \item a residual plug-in can improve existing frozen stereo models;
    \item local alignment and global propagation provide complementary benefits;
    \item temporal processing can improve both accuracy and consistency across datasets.
\end{itemize}

However, BiDAStabilizer is fundamentally non-causal. It uses future disparity $d_{t+1}$ and a backward propagation pass from the end of the sequence toward the beginning. Therefore, it is unsuitable for strict online deployment in a robotic control loop.

Furthermore, BiDAStabilizer predicts an unconstrained residual and does not include:

\begin{itemize}
    \item an explicit error-detection gate;
    \item a bounded correction magnitude;
    \item an identity-preserving initialization;
    \item a dedicated penalty on predictions that were already correct;
    \item an unseen-backbone evaluation protocol;
    \item surgical cross-domain validation.
\end{itemize}

These limitations are particularly important in light of the ARGOS v1 results, where temporal refiners trained on SCARED-specific error distributions generated false positive corrections on SERV-CT and strongly degraded already accurate predictions.

\subsection{PPMStereo: Reliability-Aware Long-Term Memory}

PPMStereo introduces a Pick-and-Play Memory mechanism for dynamic stereo matching. Instead of using only adjacent frames or storing the complete sequence indiscriminately, the model maintains a memory buffer and dynamically selects a compact subset of high-quality references.

The memory selection process combines three signals:

\begin{equation}
S_t
=
S_t^{\mathrm{conf}}
+
S_t^{\mathrm{rel}},
\end{equation}

where $S_t^{\mathrm{conf}}$ represents confidence and $S_t^{\mathrm{rel}}$ combines semantic similarity and redundancy awareness.

The top-$K$ memory entries are retained:

\begin{equation}
\mathcal{I}_t
=
\left\{
i
\mid
\operatorname{rank}(S_t[i])
\leq K
\right\},
\end{equation}

and then adaptively weighted during the play stage:

\begin{equation}
\tilde{S}_t[i]
=
\frac{S_t[i]}
{\sum_{j \in \mathcal{I}_t} S_t[j]}.
\end{equation}

An attention-based read-out aggregates the selected memory values:

\begin{equation}
\mathbf{F}_{t}^{\mathrm{agg}}
=
\mathbf{F}_{t}^{\mathrm{cost}}
+
\alpha
\operatorname{Softmax}
\left(
\frac{
\mathbf{q}_t
\mathbf{k}_m^{\prime\top}
}{
\sqrt{D_k}
}
\right)
\mathbf{v}_m^{\prime},
\end{equation}

where $\alpha$ is initialized to zero.

The main lesson from PPMStereo is that memory quantity alone is not sufficient. Naively retaining all frames, the latest frames, or random frames performs worse than selecting references based on reliability, relevance, and redundancy. The reported ablations indicate that performance saturates around $K=5$, suggesting that a compact, high-quality memory is preferable to unrestricted temporal accumulation.

However, PPMStereo is not a plug-in refiner. It relies on:

\begin{itemize}
    \item backbone-specific feature maps;
    \item internal cost volumes;
    \item context encoders;
    \item iterative GRU-based disparity updates;
    \item a DynamicStereo-style full stereo architecture.
\end{itemize}

It therefore cannot be directly applied to cached disparity predictions from heterogeneous frozen stereo networks. Moreover, its formulation is not a pure online causal system, because the vanilla memory is constructed over the full sequence. Finally, it does not explicitly include a mechanism that guarantees identity preservation on already correct predictions.

\subsection{Combined Comparison}

Table~\ref{tab:temporal_sota_positioning} summarizes the main differences between these methods and the target formulation of ARGOS v2.

\begin{table*}[t]
\centering
\caption{Positioning of ARGOS v2 relative to closely related temporal depth and stereo methods.}
\label{tab:temporal_sota_positioning}
\resizebox{\textwidth}{!}{
\begin{tabular}{lcccccc}
\toprule
\textbf{Method}
&
\textbf{Plug-in}
&
\textbf{Causal}
&
\textbf{Long-term memory}
&
\textbf{Backbone-agnostic}
&
\textbf{Safety on clean inputs}
&
\textbf{Surgical validation}
\\
\midrule
EndoStreamDepth
&
No
&
Yes
&
Hierarchical Mamba
&
No
&
No explicit mechanism
&
Yes
\\

BiDAStabilizer
&
Yes
&
No
&
Bidirectional propagation
&
Partially
&
No explicit mechanism
&
No
\\

PPMStereo
&
No
&
No
&
Selective memory buffer
&
No
&
No explicit mechanism
&
No
\\

\textbf{ARGOS v2}
&
\textbf{Yes}
&
\textbf{Yes}
&
\textbf{Selective causal memory}
&
\textbf{Yes}
&
\textbf{Explicit identity preservation}
&
\textbf{Yes}
\\
\bottomrule
\end{tabular}
}
\end{table*}

\subsection{Resulting Design Space for ARGOS v2}

The analysis above suggests that ARGOS v2 should not be designed as a simple convolutional smoothing module or as a larger ConvGRU. Instead, it should combine the strongest principles from the existing literature while addressing the limitations that remain unsolved.

The proposed formulation is a:

\begin{quote}
\textbf{Safe, causal, backbone-agnostic temporal stereo refiner with reliability-aware selective memory.}
\end{quote}

The architecture should contain five main components.

\subsubsection{Universal Evidence Encoder}

The refiner must rely only on signals that are available for any stereo backbone:

\begin{equation}
\mathbf{x}_t
=
\left[
I_t^L,
I_t^R,
d_t^{\mathrm{raw}},
r_t^{\mathrm{stereo}},
\nabla d_t^{\mathrm{raw}},
\tilde{d}_{t-1}^{\mathrm{ref}},
d_t^{\mathrm{raw}}
-
\tilde{d}_{t-1}^{\mathrm{ref}}
\right],
\end{equation}

where:

\begin{itemize}
    \item $I_t^L$ and $I_t^R$ are the left and right images;
    \item $d_t^{\mathrm{raw}}$ is the disparity predicted by the frozen backbone;
    \item $r_t^{\mathrm{stereo}}$ is a stereo photometric residual;
    \item $\tilde{d}_{t-1}^{\mathrm{ref}}$ is the previous refined disparity aligned to the current frame.
\end{itemize}

No internal backbone-specific feature map, correlation volume, or confidence representation should be required.

\subsubsection{Causal Motion Alignment}

Following the strongest result from BiDAStabilizer, memory should be aligned before temporal fusion. For a previous disparity $d_{t-k}^{\mathrm{ref}}$, the aligned memory is:

\begin{equation}
\tilde{d}_{t-k \rightarrow t}^{\mathrm{ref}}
=
\mathcal{W}
\left(
d_{t-k}^{\mathrm{ref}},
f_{t-k \rightarrow t}
\right),
\end{equation}

where $\mathcal{W}$ denotes differentiable warping and $f_{t-k \rightarrow t}$ is a motion field.

Only past frames are used:

\begin{equation}
\mathcal{M}_t
=
\left\{
m_{t-1},
m_{t-2},
\dots,
m_{t-K}
\right\},
\end{equation}

ensuring strict causal inference.

\subsubsection{Reliability-Aware Pick-and-Play Memory}

Inspired by PPMStereo, memory entries should be ranked according to reliability rather than temporal proximity alone.

A generic memory score may be defined as:

\begin{equation}
s_{t,k}
=
\lambda_q q_{t,k}
+
\lambda_s s_{t,k}^{\mathrm{sim}}
-
\lambda_r r_{t,k}^{\mathrm{red}}
-
\lambda_o o_{t,k}
-
\lambda_a a_{t,k},
\end{equation}

where:

\begin{itemize}
    \item $q_{t,k}$ is estimated memory quality;
    \item $s_{t,k}^{\mathrm{sim}}$ is similarity to the current frame;
    \item $r_{t,k}^{\mathrm{red}}$ measures redundancy;
    \item $o_{t,k}$ measures occlusion or warp invalidity;
    \item $a_{t,k}$ encodes memory age.
\end{itemize}

The top-$K$ memories are selected and adaptively weighted before aggregation.

\subsubsection{Hierarchical Temporal State}

Following EndoStreamDepth, temporal information should be represented at multiple spatial scales:

\begin{equation}
\mathbf{H}_t
=
\left\{
\mathbf{H}_t^{1/4},
\mathbf{H}_t^{1/8},
\mathbf{H}_t^{1/16}
\right\}.
\end{equation}

The finest scale captures local corrections and disparity boundaries, while coarser scales preserve global shape, motion context, and long-term structural information.

The temporal operator may be implemented using Mamba blocks, recurrent gated units, or a hybrid design. The scientific contribution should not be framed merely as the use of Mamba, since hierarchical temporal Mamba has already been introduced for endoscopic depth estimation.

\subsubsection{Safe Identity-Preserving Output}

The final correction should be explicitly bounded:

\begin{equation}
d_t^{\mathrm{ref}}
=
d_t^{\mathrm{raw}}
+
g_t
c_t
\tau_t
\tanh
\left(
\Delta d_t
\right),
\label{eq:safe_update}
\end{equation}

where:

\begin{itemize}
    \item $g_t \in [0,1]$ is the probability that the raw disparity is incorrect;
    \item $c_t \in [0,1]$ is the confidence in the temporal correction;
    \item $\tau_t$ is a local maximum update magnitude;
    \item $\Delta d_t$ is the proposed residual correction.
\end{itemize}

The network should be initialized such that:

\begin{equation}
g_t \approx 0,
\qquad
\Delta d_t \approx 0,
\end{equation}

and therefore:

\begin{equation}
d_t^{\mathrm{ref}}
\approx
d_t^{\mathrm{raw}}
\end{equation}

at the beginning of training.

This identity-preserving property must be structural rather than being left entirely to the optimization objective.

\subsection{Novelty Boundary}

The novelty of ARGOS v2 cannot be stated as temporal memory alone, since selective long-term memory has already been introduced by PPMStereo. It cannot be stated as hierarchical Mamba modeling for endoscopy, since this is already present in EndoStreamDepth. It also cannot be stated merely as a stereo stabilization plug-in, since BiDAStabilizer already establishes that paradigm.

The contribution instead lies in the combination of properties that has not been demonstrated jointly:

\begin{itemize}
    \item plug-in refinement of frozen stereo backbones;
    \item strict causal and streaming operation;
    \item backbone-independent input representation;
    \item supervised training on multiple heterogeneous stereo predictors;
    \item evaluation on completely unseen stereo backbones;
    \item reliability-aware selective temporal memory;
    \item explicit identity preservation on already accurate predictions;
    \item zero-shot evaluation across surgical datasets;
    \item validation on both geometric ground truth and in-vivo unlabeled sequences.
\end{itemize}

The resulting central claim of ARGOS v2 is therefore:

\begin{quote}
\textit{ARGOS v2 introduces a causal, safe, and backbone-agnostic temporal stereo refiner trained across heterogeneous frozen stereo predictors. The model uses reliability-aware selective memory to improve temporal consistency while preserving accurate geometry, and is evaluated under unseen-backbone and surgical cross-domain shifts.}
\end{quote}

\subsection{Core Scientific Gap}

The central unresolved problem is not simply how to smooth temporally inconsistent stereo predictions. The key challenge is determining when temporal correction is justified.

A successful temporal refiner must distinguish between:

\begin{equation}
\text{prediction variation caused by stereo flicker}
\end{equation}

and

\begin{equation}
\text{prediction variation caused by real scene motion, deformation, or occlusion}.
\end{equation}

The failure mode observed in ARGOS v1 demonstrates that a refiner trained on one backbone and one dataset can learn a specific error signature rather than a transferable temporal prior. When evaluated on an out-of-domain dataset in which the raw stereo prediction is already accurate, the refiner may trigger false positive corrections and significantly degrade the result.

ARGOS v2 is therefore motivated by a stricter principle:

\begin{quote}
\textbf{The temporal refiner must earn the right to modify the raw prediction.}
\end{quote}

This principle motivates the combination of multi-backbone training, selective memory, uncertainty-aware gating, bounded updates, and explicit safety evaluation on clean predictions.

\subsection{Experimental Consequences}

The proposed formulation requires evaluation along three independent axes.

\paragraph{Seen-backbone performance.}
The refiner must improve or preserve all stereo backbones used during training.

\paragraph{Unseen-backbone generalization.}
A backbone that is completely excluded from training and model selection must be evaluated on the same test domain.

\paragraph{Cross-domain generalization.}
The frozen refiner must then be evaluated on unseen surgical datasets, including SERV-CT and D4D, and on unlabeled in-vivo sequences such as StereoMIS.

Particular attention should be given to already accurate predictions. Let:

\begin{equation}
\Omega_{\mathrm{clean}}
=
\left\{
i
\mid
\left|
d_i^{\mathrm{raw}}
-
d_i^{\mathrm{gt}}
\right|
<
\epsilon
\right\}.
\end{equation}

The safety error on clean pixels can be defined as:

\begin{equation}
\mathcal{L}_{\mathrm{clean}}
=
\frac{1}
{
|\Omega_{\mathrm{clean}}|
}
\sum_{i \in \Omega_{\mathrm{clean}}}
\left|
d_i^{\mathrm{ref}}
-
d_i^{\mathrm{raw}}
\right|.
\end{equation}

Additional safety metrics should include:

\begin{itemize}
    \item new-Bad3 introduced by the refiner;
    \item percentage of clean pixels degraded;
    \item false update rate;
    \item mean update magnitude on clean predictions;
    \item fraction of frames with worse EPE after refinement.
\end{itemize}

These metrics are necessary to verify that the model improves temporal consistency without reproducing the out-of-domain degradation observed in ARGOS v1.

\subsection{BiDA-Style Motion Alignment as the First Validated Building Block}

The first validated component of ARGOS v2 is a causal BiDA-style motion-alignment module. 
For each current frame $t$, a frozen optical-flow network estimates the temporal motion between the current and previous left images. The previous disparity prediction is then warped into the current-frame coordinate system:

\[
\hat{d}_{t-1 \rightarrow t}
=
\mathcal{W}
\left(
d_{t-1},
F_{t \rightarrow t-1}
\right),
\]

where $d_{t-1}$ is the previous frame-wise stereo prediction, $F_{t \rightarrow t-1}$ is the current-to-previous optical flow, and $\mathcal{W}$ denotes target-to-source bilinear warping using the BiDA convention.

The validation shows that motion alignment alone is not sufficient to improve the current disparity through direct replacement, fixed blending, or a hand-crafted forward--backward and photometric gate. However, the aligned temporal memory contains complementary geometric evidence. In particular, a per-pixel oracle that selects the lower-error value between the current raw disparity and the aligned previous disparity improves the result across all evaluated sequences and stereo backbones.

This leads to the following interpretation:

\[
\text{Optical flow}
\neq
\text{automatic disparity improvement},
\]

but rather

\[
\text{Optical flow}
\rightarrow
\text{temporally aligned evidence}.
\]

The remaining scientific problem is therefore not how to transport information from the past, but how to determine whether the current prediction or the aligned temporal memory is more reliable at each spatial location.

Accordingly, the BiDA-style alignment block is retained as the first architectural building block of ARGOS v2:

\[
\left(
I_{t-1}^{L},
I_t^{L},
d_{t-1},
d_t
\right)
\rightarrow
\left(
\hat{d}_{t-1 \rightarrow t},
m_{\mathrm{warp}},
e_{\mathrm{FB}},
e_{\mathrm{photo}},
\Delta d_t
\right),
\]

where:

\begin{itemize}
    \item $\hat{d}_{t-1 \rightarrow t}$ is the aligned previous disparity;
    \item $m_{\mathrm{warp}}$ is the valid warp-support mask;
    \item $e_{\mathrm{FB}}$ is the forward--backward flow-consistency error;
    \item $e_{\mathrm{photo}}$ is the photometric residual;
    \item $\Delta d_t = d_t - \hat{d}_{t-1 \rightarrow t}$ is the temporal disparity disagreement.
\end{itemize}

These signals will be passed to a learned, identity-preserving error detector and refinement module. The intended output remains:

\[
d_t^{\mathrm{ref}}
=
d_t
+
g_{\mathrm{error}}
\cdot
c_{\mathrm{memory}}
\cdot
\tau
\cdot
\tanh(\delta_t),
\]

where the learned gate must determine when the aligned memory provides trustworthy corrective evidence.

The current experimental conclusion is therefore:

\begin{quote}
BiDA-style causal alignment is a \textbf{GO} as a universal temporal evidence module, while direct memory replacement, fixed blending, and heuristic gating remain a \textbf{NO-GO}.
\end{quote}

\subsection{Milestone: Validation of Causal BiDA-Style Temporal Alignment}

\textbf{Status: Major Scientific Milestone Achieved}

The first major milestone of ARGOS v2 was successfully completed.

Rather than directly developing a new temporal refiner, we first isolated and validated the core temporal alignment mechanism inspired by BiDAStabilizer in order to determine whether motion-compensated temporal evidence actually contains useful corrective information for surgical stereo.

The study produced three important findings:

\begin{itemize}
    \item A causal adaptation of the BiDA alignment module successfully transfers temporal information from previous frames using only past observations, making it compatible with real-time surgical deployment.
    
    \item Oracle experiments demonstrate that the aligned previous disparity consistently contains complementary geometric information capable of improving the current prediction across multiple frozen stereo backbones.
    
    \item Simple heuristic selectors based on forward-backward consistency and photometric consistency are insufficient to exploit this information safely, motivating the need for a learned refinement strategy.
\end{itemize}

This represents an important scientific result because it demonstrates that temporal alignment itself is valuable, while identifying the true research problem:

\begin{quote}
\emph{The challenge is not obtaining temporal information, but learning when that information should be trusted.}
\end{quote}

Building on this observation, the first learned identity-preserving refiner was developed.

The learned bounded residual model:

\begin{itemize}
    \item improves all three training backbones;
    \item generalizes without tuning to the unseen Fast-FoundationStereo backbone;
    \item recovers a substantial fraction of the oracle improvement;
    \item remains lightweight (approximately 39k parameters).
\end{itemize}

Although the current selector is not yet sufficiently reliable to satisfy the desired safety constraints (clean-input preservation and low false-update rate), these experiments validate the central hypothesis of ARGOS v2:

\begin{quote}
\emph{Causally aligned temporal evidence can improve frozen stereo predictors in a backbone-agnostic manner; the remaining challenge is reliable intervention rather than temporal information extraction.}
\end{quote}

This milestone establishes the first building block of the final ARGOS v2 architecture and justifies subsequent research on learned confidence estimation, selective memory, and safety-aware temporal refinement.

\subsection{Milestone: Long-Range Temporal Evidence and the Limits of Explicit Pick-and-Play Memory}

\textbf{Status: Long-range temporal evidence validated; current PPMStereo-style selector not promoted.}

After validating causal BiDA-style alignment and a learned bounded residual refiner based on the immediately preceding frame, we investigated whether longer temporal memory could provide additional corrective information.

A causal memory bank was constructed using BiDA-aligned disparity predictions from:

\[
t-1,\quad t-2,\quad t-4,\quad t-8.
\]

The study compared the existing learned \(t-1\) refiner against deterministic and learned PPMStereo-inspired memory-selection strategies.

At cache-grid coverage \(0.50\), the following aggregate results were obtained over the three seen stereo backbones:

\[
\begin{array}{l|c}
\textbf{Method} & \textbf{EPE} \\
\hline
\text{Raw disparity} & 0.64384 \\
\text{Oracle using } t-1 & 0.57495 \\
\text{Oracle using any memory} & 0.47672 \\
\text{Learned BiDA } t-1 \text{ refiner} & 0.61971 \\
\text{Learned long-memory refiner} & 0.62951
\end{array}
\]

The multi-memory oracle therefore provides an additional improvement of:

\[
0.57495 - 0.47672 = 0.09823 \ \mathrm{px}
\]

beyond the \(t-1\) oracle.

This demonstrates that older aligned memories contain substantial complementary geometric evidence. In particular, frames at \(t-2\), \(t-4\), and \(t-8\) repeatedly contain correct disparity information that is unavailable in the current prediction and in the immediately preceding frame.

However, the tested PPMStereo-style selectors were unable to exploit this additional evidence safely.

The learned long-memory model remained worse than the validated \(t-1\) refiner on all evaluated backbones, including the unseen Fast-FoundationStereo model:

\[
\begin{array}{l|ccc}
\textbf{Backbone} &
\textbf{Raw EPE} &
\textbf{Learned } t-1 &
\textbf{Long Memory} \\
\hline
\text{S2M2-S} & 0.69611 & 0.67702 & 0.67913 \\
\text{RAFT-Stereo} & 0.61744 & 0.59070 & 0.60420 \\
\text{StereoAnywhere} & 0.61795 & 0.59143 & 0.60520 \\
\text{Fast-FoundationStereo} & 0.64515 & 0.62332 & 0.63578
\end{array}
\]

The learned selector assigned approximately \(72\%\) of its total mass to temporal memory and distributed this mass almost uniformly across memory ages. It therefore failed to perform reliable selective memory access or abstention.

The resulting safety behaviour was insufficient:

\begin{itemize}
    \item clean-pixel degradation of approximately \(36\%\);
    \item false-update rate above \(41\%\);
    \item approximately \(39\%-41\%\) of frames worsened;
    \item degraded performance relative to the learned \(t-1\) refiner on both seen and unseen backbones.
\end{itemize}

The correct interpretation is therefore not that long temporal memory is uninformative. On the contrary, the oracle results demonstrate a substantially higher temporal correction ceiling.

The current limitation is the inability of universal disparity- and flow-based features to identify which memory is trustworthy at each spatial location.

The resulting decision is:

\begin{quote}
\textbf{GO:} long-range aligned memory as a source of complementary temporal evidence.

\textbf{NO-GO:} the current deterministic and learned PPMStereo-style top-\(K\) selectors for inclusion in the first ARGOS v2 refiner.
\end{quote}

This milestone reveals that the bottleneck is no longer the availability of temporal information, but its representation and safe selection.

The PPMStereo memory-bank infrastructure is retained for future experiments, while the first ARGOS v2 refiner remains based on the validated causal BiDA \(t-1\) pathway.

\subsection{Milestone 4 -- Evaluation of Foundation Visual Features (DINOv3)}

A natural hypothesis was that richer semantic visual representations could substantially improve temporal memory selection.

To test this, we integrated the official frozen DINOv3 ViT-L/16 foundation model as an alternative feature extractor for the temporal selector while keeping the remainder of the pipeline unchanged. Multiple feature configurations were evaluated, including intermediate transformer layers and multi-layer feature fusion.

The objective was to determine whether high-level semantic descriptors provide better evidence for deciding when historical temporal information should be trusted compared with purely geometric cues.

The controlled experiments led to a clear negative result.

Although DINOv3 slightly increased the ranking AUROC, all practically relevant decision metrics deteriorated:

\[
\text{Geometry-only selector}
>
\text{DINOv3-enhanced selector}
\]

The DINO-based selector showed:

\begin{itemize}
    \item lower Average Precision,
    \item lower Top-1 retrieval accuracy,
    \item higher decision regret,
    \item no measurable improvement in temporal refinement quality.
\end{itemize}

Consequently, DINOv3 was not promoted to the main ARGOS v2 architecture.

This experiment provides an important scientific insight: for plug-and-play temporal stereo refinement, deciding whether historical disparity should be trusted appears to depend primarily on geometric consistency rather than on high-level semantic image understanding.

As a result, ARGOS v2 continues to rely on universal geometric evidence (disparity disagreement, optical-flow alignment, validity masks and temporal consistency) instead of expensive frozen foundation features.

This milestone substantially narrows the architectural search space and suggests that future improvements should focus on better temporal reasoning mechanisms rather than richer image representations.

\subsection{Milestone: Cross-Memory Consensus Reveals Correlated Temporal Errors}

After observing that long-range aligned memories contain substantial oracle improvement, we investigated whether reliable temporal corrections could be identified without learning a per-memory selector.

We introduced a zero-parameter causal baseline, termed \emph{Cross-Memory Consensus Correction} (CMC). Given the BiDA-aligned disparities from previous frames

\[
\mathcal{M}_t =
\left\{
\tilde d_{t-1\rightarrow t},
\tilde d_{t-2\rightarrow t},
\tilde d_{t-4\rightarrow t},
\tilde d_{t-8\rightarrow t}
\right\},
\]

CMC computes a validity-aware median:

\[
d_{\mathrm{med}}(p)
=
\operatorname{median}
\left\{
\tilde d_{t-k\rightarrow t}(p)
\mid v_{t-k}(p)=1
\right\},
\]

and a robust dispersion estimate using the median absolute deviation:

\[
s_{\mathrm{MAD}}(p)
=
\operatorname{median}
\left\{
\left|
\tilde d_{t-k\rightarrow t}(p)
-
d_{\mathrm{med}}(p)
\right|
\mid v_{t-k}(p)=1
\right\}.
\]

The current raw disparity is corrected only when:

\[
n_{\mathrm{valid}}(p) \geq n_{\min},
\]

\[
s_{\mathrm{MAD}}(p) \leq \tau_s,
\]

and

\[
\left|
d_{\mathrm{raw}}(p)-d_{\mathrm{med}}(p)
\right|
\geq
\tau_d+\kappa s_{\mathrm{MAD}}(p).
\]

The final correction is bounded:

\[
d_{\mathrm{CMC}}(p)
=
d_{\mathrm{raw}}(p)
+
g_{\mathrm{CMC}}(p)
\operatorname{clip}
\left(
d_{\mathrm{med}}(p)-d_{\mathrm{raw}}(p),
-B,
B
\right),
\]

with \(B=3\) cache-grid pixels.

The hypothesis was that several temporally separated memories agreeing tightly with each other while disagreeing with the current prediction would indicate that the current prediction is an outlier.

A predeclared sweep over 36 configurations was performed exclusively on the training sequences and three seen stereo backbones, covering approximately 30 million evaluated pixels.

The predefined promotion criteria required:

\[
\Delta\mathrm{EPE} \geq 0.005 \ \mathrm{px},
\]

\[
\mathrm{FalseUpdateRate} \leq 20\%,
\]

and

\[
\mathrm{CleanDegradation} \leq 15\%.
\]

CMC did not satisfy these criteria. The best configuration achieved only:

\[
\Delta\mathrm{EPE}=0.0013 \ \mathrm{px},
\]

with:

\[
\mathrm{FalseUpdateRate}=22.8\%.
\]

Aggressive configurations reached false-update rates as high as \(86\%\).

However, the mechanism remained highly conservative on already accurate predictions, with a clean-pixel degradation rate of approximately:

\[
0.04\%.
\]

The most important result is therefore not the failure of the median-based correction itself, but the structure of the failure.

When all temporal memories agree against the current prediction, they are frequently wrong together. Persistent stereo artifacts, specularities, occlusion errors, and backbone-specific failure modes remain correlated across multiple memory ages, even after optical-flow alignment.

The CMC mechanism oracle obtained a maximum gain of only:

\[
0.0297 \ \mathrm{px},
\]

whereas the unrestricted multi-memory oracle previously achieved approximately:

\[
0.098 \ \mathrm{px}
\]

beyond the \(t-1\) oracle.

Therefore, consensus captures only about:

\[
\frac{0.0297}{0.098}
\approx 30\%
\]

of the available long-memory improvement.

The remaining \(70\%\) is concentrated in cases where only one temporal age contains the correct evidence, while the remaining memories and the current prediction may be jointly incorrect.

This leads to the following conclusion:

\begin{quote}
Temporal agreement is not equivalent to temporal correctness. Long-range stereo errors are strongly correlated across memory ages, and the useful evidence is frequently contained in a minority temporal candidate rather than in the consensus.
\end{quote}

CMC is therefore not promoted as the main ARGOS v2 correction mechanism. It is retained as:

\begin{itemize}
    \item a zero-parameter temporal baseline;
    \item an ultra-conservative safety prior;
    \item a source of high-precision pseudo-labels;
    \item a hard-positive mining mechanism for future quality estimation.
\end{itemize}

Together with the PPMStereo, DINOv3, and latent-state experiments, this result motivates explicit prediction of candidate error and uncertainty rather than similarity, consensus, or latent-state aggregation alone.

\subsection{Milestone: Absolute Quality Estimation Does Not Enable Reliable Memory Selection}

We investigated whether safe temporal refinement could be obtained by explicitly predicting the expected disparity error of the current prediction and four causally aligned temporal candidates:

\[
\mathcal{D}_t =
\left\{
d_{\mathrm{raw}},
\tilde d_{t-1\rightarrow t},
\tilde d_{t-2\rightarrow t},
\tilde d_{t-4\rightarrow t},
\tilde d_{t-8\rightarrow t}
\right\}.
\]

For each candidate \(i\), the target quality was defined as its absolute disparity error:

\[
e_i(p)=\left|d_i(p)-d_{\mathrm{GT}}(p)\right|.
\]

A lightweight backbone-agnostic Quality Prediction Module estimated the expected error and uncertainty of every candidate:

\[
\hat e_i(p), \qquad \hat{\sigma}_i(p).
\]

The model used only universal geometric and temporal signals. No RGB features, backbone identity, stereo cost volumes, internal hidden states, or future frames were provided.

The quality predictor learned a meaningful measure of absolute scene difficulty. Pearson correlations between predicted and true candidate errors ranged approximately from \(0.60\) to \(0.70\), while uncertainty remained positively correlated with actual prediction error.

However, this signal did not translate into reliable relative candidate ranking. The main results were:

\[
\begin{array}{lccc}
\toprule
\text{Method} &
\text{Top-1 Accuracy} &
\text{Diagnostic Regret} &
\text{Raw-vs-Memory AUROC} \\
\midrule
\text{Constant baseline} & 27.9\% & 0.150 & 0.506 \\
\text{PPMStereo selector} & 28.3\% & 0.119 & 0.422 \\
\text{CMC spread heuristic} & 11.2\% & 0.154 & 0.679 \\
\text{Q0 quality predictor} & 19.7\% & 0.153 & 0.514 \\
\bottomrule
\end{array}
\]

The proposed quality predictor therefore performed worse than both the constant baseline and the previous PPMStereo-style selector for candidate ranking.

The failure was particularly evident where a minority temporal candidate contained the correct evidence:

\[
R_{\mathrm{minority}}^{Q0}=0.280,
\qquad
R_{\mathrm{minority}}^{PPM}=0.194.
\]

The model predominantly learned a shared difficulty component:

\[
\hat e_{\mathrm{raw}}
\approx
\hat e_{t-1}
\approx
\hat e_{t-2}
\approx
\hat e_{t-4}
\approx
\hat e_{t-8},
\]

rather than identifying which specific candidate retained correct geometry.

Risk--coverage analysis did not provide a useful intervention regime. At \(10\%\) coverage, regret decreased to approximately \(0.049\), but \(98.1\%\) of the selected pixels were already raw-clean and intervention precision was only \(1.33\%\).

The central finding is therefore:

\begin{quote}
Absolute prediction difficulty is learnable, whereas relative temporal candidate correctness is not reliably recoverable from the tested universal evidence.
\end{quote}

Together with the PPMStereo, DINOv3, latent-state, and cross-memory consensus experiments, this result motivates abandoning multi-age candidate selection for the first ARGOS v2 model.

The next architecture therefore focuses on the more tractable safety question:

\[
\textit{Is the current raw disparity sufficiently unreliable to justify intervention?}
\]

The resulting design direction is:

\[
\boxed{
\text{causal BiDA }t-1
+
\text{raw-error detection}
+
\text{abstention}
+
\text{bounded residual correction}
}
\]

\subsection{Milestone: Safe Intervention via Raw Error Detection and Abstention}

The central objective of ARGOS v2 is not to maximize temporal smoothing, but to determine when temporal correction is actually justified.

Following the failure of multiple temporal memory selection strategies (PPMStereo-inspired memory selection, DINOv3 features, latent temporal state, cross-memory consensus, and explicit quality prediction), we reformulated the problem as a binary safety decision:

\[
\textit{Is the current raw disparity sufficiently unreliable to justify any temporal correction?}
\]

Instead of learning which temporal candidate should replace the current prediction, ARGOS v2 separates temporal refinement into two independent stages:

\[
\boxed{
\text{Correction Proposal}
\rightarrow
\text{Intervention Authorization}
}
\]

The correction proposal is generated by the previously validated bounded BiDA \(t-1\) residual refiner,

\[
d_{\mathrm{proposal}}
=
d_{\mathrm{raw}}
+
g
\cdot
c
\cdot
\tau
\tanh(\Delta),
\]

while a lightweight Raw Error Detector predicts

\[
P(\text{raw is wrong}),
\qquad
\mu_{\mathrm{raw}},
\qquad
\sigma_{\mathrm{raw}},
\]

where

\(\mu_{\mathrm{raw}}\) is the expected raw disparity error and
\(\sigma_{\mathrm{raw}}\) represents predictive uncertainty.

The detector authorizes refinement only when the estimated error is sufficiently large and the uncertainty remains below a calibrated threshold. Otherwise the model abstains:

\[
d_{\mathrm{out}}
=
d_{\mathrm{raw}}.
\]

Importantly, the detector contains only 1,107 trainable parameters and uses exclusively backbone-independent geometric evidence.

On held-out SCARED-C sequences, the resulting system achieved:

\[
\begin{array}{lcccc}
\toprule
Method &
EPE &
False\ Update &
Clean\ Degradation &
Coverage\\
\midrule
Raw &
0.9616 &
0\% &
0\% &
0\%\\
A2 &
0.9190 &
35.1\% &
13.2\% &
40.0\%\\
Authorized\ A2 &
0.9347 &
1.9\% &
0.97\% &
5.6\%\\
\bottomrule
\end{array}
\]

Although intervention coverage is intentionally conservative, the detector preserves approximately

\[
63.2\%
\]

of the original A2 refinement gain while reducing false updates by nearly an order of magnitude.

Most importantly, after freezing all model parameters and operating thresholds, the detector generalized without retraining to the unseen Fast-FoundationStereo backbone:

\[
0.9757
\rightarrow
0.9523
\]

while maintaining virtually identical safety metrics.

This experiment validates the central hypothesis of ARGOS v2:

\begin{quote}
Safe temporal stereo refinement is better formulated as an intervention authorization problem than as a direct temporal replacement problem.
\end{quote}

The resulting architecture becomes

\[
\boxed{
\text{BiDA causal alignment}
+
\text{bounded residual proposal}
+
\text{raw-error detector}
+
\text{calibrated abstention}
}
\]

which constitutes the first promoted ARGOS v2 architecture.

\subsection{Milestone: Backbone Transfer Succeeds, but Cross-Domain Authorization Remains Unsolved}

After promoting the balanced ARGOS v2 authorization policy, we froze the complete inference system before observing any out-of-distribution result.

The frozen system comprised:

\[
\text{SEA-RAFT alignment}
+
\text{bounded BiDA } t-1 \text{ proposal}
+
\text{Raw Error Detector}
+
\text{calibrated abstention}.
\]

All model parameters, detector temperature, intervention thresholds, and source files were frozen and recorded through SHA-256 hashes. No optimization or post-hoc tuning was performed on any OOD dataset.

The evaluation therefore tested two distinct forms of generalization:

\begin{enumerate}
    \item transfer to an unseen stereo backbone within the SCARED-C domain;
    \item transfer across unseen surgical datasets and acquisition conditions.
\end{enumerate}

The main results were:

\[
\begin{array}{lrrrrl}
\toprule
\text{Dataset / protocol}
&
\text{Raw EPE}
&
\text{Refined EPE}
&
\Delta\text{EPE}
&
\text{False Update}
&
\text{Verdict}
\\
\midrule
\text{CREStereo, unseen backbone}
& 1.1003 & 1.0625 & -0.0378 & 2.08\% & \text{GO}
\\
\text{SERV-CT}
& 1.0454 & 1.1015 & +0.0561 & 38.41\% & \text{NO-GO}
\\
\text{SCARED structured-light}
& 2.9026 & 2.9132 & +0.0106 & 0.71\% & \text{Preservation}
\\
\text{D4D anchors}
& 3.7498 & 3.7004 & -0.0494 & 29.15\% & \text{NO-GO}
\\
\bottomrule
\end{array}
\]

Here, negative \(\Delta\mathrm{EPE}\) denotes improvement.

On the second unseen stereo backbone, CREStereo, ARGOS v2 reduced EPE by:

\[
3.43\%,
\]

while maintaining:

\[
\mathrm{FalseUpdate}=2.08\%,
\qquad
\mathrm{CleanDegradation}=1.00\%,
\]

and an intervention precision of:

\[
81.67\%.
\]

This confirms that the promoted refiner does not merely memorize the error distribution of the three training backbones. The combination of universal evidence, bounded correction, and intervention authorization transfers to a previously unseen stereo architecture within the same surgical domain.

However, cross-dataset evaluation exposed a different limitation.

On SERV-CT, the frozen detector authorized incorrect updates too frequently:

\[
\mathrm{FalseUpdate}=38.41\%,
\qquad
\mathrm{CleanDegradation}=22.69\%,
\]

causing EPE to increase from:

\[
1.0454
\rightarrow
1.1015.
\]

On D4D structured-light anchors, mean EPE improved slightly:

\[
3.7498
\rightarrow
3.7004,
\]

but safety remained unacceptable:

\[
\mathrm{FalseUpdate}=29.15\%,
\qquad
\mathrm{CleanDegradation}=17.95\%,
\]

and Bad3 increased despite the lower average EPE.

This demonstrates that average geometric improvement is insufficient evidence of safe refinement.

The temporal evaluations further confirmed that improved temporal smoothness does not necessarily imply improved geometry. D4D temporal disagreement decreased, while anchor-based safety metrics worsened. Similarly, StereoMIS showed sparse bounded interventions, but aggregate no-reference temporal consistency did not improve:

\[
\Delta E_{\mathrm{temporal}} = +0.00149.
\]

The principal finding is therefore:

\begin{quote}
ARGOS v2 generalizes across unseen stereo backbones within the training domain, but the intervention authorization policy remains sensitive to cross-dataset calibration shift.
\end{quote}

The failure is localized primarily in the authorization stage rather than in the causal alignment or bounded proposal mechanism:

\[
\text{backbone transfer} \quad \checkmark
\]

\[
\text{bounded causal proposal} \quad \checkmark
\]

\[
\text{within-domain authorization} \quad \checkmark
\]

\[
\text{cross-domain authorization calibration} \quad \times
\]

This distinction is central to the ARGOS v2 contribution. The model does not fail because temporal evidence is absent; it fails because the detector assigns excessive confidence to intervention decisions under unseen acquisition distributions.

The defensible conclusion is therefore:

\begin{quote}
Safe temporal refinement requires both error detection and awareness of whether the detector itself remains within its calibrated operating domain.
\end{quote}

Consequently, the next ARGOS v2 milestone is not a more complex temporal refiner, but a domain-shift-aware authorization mechanism capable of abstaining when its own evidence distribution is unreliable.

\subsection{Milestone: Cross-Domain Failure Is Explained by Feature-Support Shift}

After validating safe authorization on held-out SCARED-C sequences and two unseen stereo backbones, we investigated why the same frozen authorization policy failed on SERV-CT and D4D.

The complete ARGOS v2 pipeline was kept frozen:

\[
\text{SEA-RAFT alignment}
+
\text{bounded A2 proposal}
+
\text{Raw Error Detector}
+
\text{balanced abstention policy}.
\]

No retraining, threshold tuning, or architectural modification was performed. The detector, A2 checkpoint, BiDA implementation, temperature, and operating thresholds were verified through frozen SHA-256 hashes before inference.

The analysis compared the detector evidence distribution of each test domain against held-out SCARED-C using:

\begin{itemize}
    \item nearest-neighbour support overlap;
    \item Mahalanobis distance;
    \item detector calibration;
    \item raw-error discrimination;
    \item authorization frequency;
    \item clean-pixel safety.
\end{itemize}

The results revealed a clear separation between unseen-backbone transfer and unseen-domain transfer:

\[
\begin{array}{lrrrr}
\toprule
\text{Dataset}
&
\text{NN Support Overlap}
&
\text{Mean Mahalanobis}
&
\text{AUROC}
&
\text{ECE}
\\
\midrule
\text{SCARED-C reference}
& \text{--} & 22.4 & 0.843 & 0.059
\\
\text{CREStereo}
& 92.8\% & 21.3 & 0.859 & 0.056
\\
\text{SERV-CT}
& 0.0\% & 581.8 & 0.487 & 0.409
\\
\text{D4D}
& 7.8\% & 168.8 & 0.553 & 0.290
\\
\text{StereoMIS}
& 42.8\% & 66.3 & \text{n/a} & \text{n/a}
\\
\bottomrule
\end{array}
\]

CREStereo, despite being an unseen stereo architecture, remained strongly inside the calibrated SCARED-C feature support. Its detector discrimination and calibration remained comparable to the reference distribution, explaining the successful transfer observed previously.

In contrast, SERV-CT was entirely outside the training support:

\[
\mathrm{SupportOverlap}_{\mathrm{SERV\text{-}CT}}=0\%,
\]

\[
D_{\mathrm{Mahalanobis}}=581.8.
\]

Under this shift, the Raw Error Detector became effectively non-discriminative:

\[
\mathrm{AUROC}=0.487,
\]

while remaining strongly overconfident:

\[
\mathrm{ECE}=0.409.
\]

It authorized intervention on approximately:

\[
48.3\%
\]

of evaluated pixels, and \(64.2\%\) of GT-paired authorizations were diagnostically incorrect.

D4D exhibited the same failure mode at lower severity:

\[
\mathrm{SupportOverlap}_{\mathrm{D4D}}=7.8\%,
\qquad
D_{\mathrm{Mahalanobis}}=168.8,
\]

with:

\[
\mathrm{AUROC}=0.553,
\qquad
\mathrm{ECE}=0.290.
\]

A stricter global authorization threshold was insufficient. The previously frozen ultra-safe operating point continued to authorize approximately:

\[
43.3\%
\]

of SERV-CT pixels and:

\[
28.5\%
\]

of D4D pixels.

Therefore, the failure cannot be explained solely by scalar miscalibration. The detector is confidently applying an in-domain decision rule to evidence lying outside its learned support.

The central finding is:

\begin{quote}
ARGOS v2 transfers across unseen stereo backbones when the universal evidence remains inside the calibrated training support, but authorization becomes unsafe when the evidence distribution undergoes substantial cross-domain shift.
\end{quote}

This establishes the distinction:

\[
\boxed{
\text{unseen backbone}
\neq
\text{unseen evidence distribution}
}
\]

and localizes the remaining limitation to support awareness rather than correction capacity.

The promoted architecture currently performs:

\[
\text{raw-error authorization}
\]

but does not yet answer:

\[
\textit{Is the authorization evidence itself within a region where the detector is reliable?}
\]

The next ARGOS v2 mechanism is therefore a second abstention layer:

\[
a_{\mathrm{final}}
=
a_{\mathrm{support}}
\cdot
a_{\mathrm{error}},
\]

where:

\[
a_{\mathrm{support}}
=
\mathbb{1}
\left[
s_{\mathrm{OOD}} < \tau_{\mathrm{support}}
\right].
\]

When evidence lies outside the calibrated feature support, the complete system must preserve the raw disparity:

\[
a_{\mathrm{support}}=0
\quad\Longrightarrow\quad
d_{\mathrm{out}}=d_{\mathrm{raw}}.
\]

This motivates a feature-support-aware authorization policy rather than a larger refiner or further temporal-memory complexity.

\subsection{Milestone: Supervised Multi-Domain Exposure Does Not Yield Third-Domain Authorization Robustness}

After identifying cross-domain calibration shift as the main limitation of the ARGOS v2 Raw Error Detector, we tested whether supervised exposure to an additional surgical domain was sufficient to produce domain-general authorization.

The correction pipeline was kept unchanged:

\[
\text{frozen SEA-RAFT}
+
\text{causal BiDA alignment}
+
\text{frozen A2 bounded proposal}
+
\text{trainable Raw Error Detector}
+
\text{bit-exact abstention}.
\]

Only the Raw Error Detector and its calibration policy were retrained. No new architecture, temporal module, support detector, or feature extractor was introduced.

Two leave-one-domain-out experiments were conducted:

\[
\mathrm{M1}:
\quad
\text{SCARED-C}+\text{D4D}
\rightarrow
\text{SERV-CT},
\]

\[
\mathrm{M2}:
\quad
\text{SCARED-C}+\text{SERV-CT}
\rightarrow
\text{D4D}.
\]

All final-domain evaluations were performed only after freezing the detector, calibration parameters, operating thresholds, and artifact hashes.

At GT coverage \(0.50\), the main results were:

\[
\begin{array}{llrrrrr}
\toprule
\text{Experiment}
&
\text{Dataset}
&
\text{Raw EPE}
&
\text{Output EPE}
&
\text{Gain}
&
\text{False Update}
&
\text{Coverage}
\\
\midrule
\mathrm{M1}
&
\text{SCARED-C}
&
0.9616
&
0.9491
&
+0.0125
&
0.88\%
&
2.33\%
\\
\mathrm{M1}
&
\text{Fast-FoundationStereo}
&
0.9757
&
0.9656
&
+0.0101
&
0.79\%
&
2.19\%
\\
\mathrm{M1}
&
\text{CREStereo}
&
1.1003
&
1.0862
&
+0.0141
&
0.89\%
&
2.52\%
\\
\mathrm{M1}
&
\text{SERV-CT unseen}
&
1.0463
&
1.4140
&
-0.3678
&
49.08\%
&
51.02\%
\\
\midrule
\mathrm{M2}
&
\text{SCARED-C}
&
0.9616
&
0.9421
&
+0.0195
&
0.25\%
&
1.97\%
\\
\mathrm{M2}
&
\text{Fast-FoundationStereo}
&
0.9757
&
0.9581
&
+0.0177
&
0.19\%
&
1.97\%
\\
\mathrm{M2}
&
\text{CREStereo}
&
1.1003
&
1.0752
&
+0.0251
&
0.20\%
&
2.08\%
\\
\mathrm{M2}
&
\text{D4D unseen}
&
2.9055
&
2.9001
&
+0.0054
&
1.80\%
&
3.31\%
\\
\bottomrule
\end{array}
\]

The primary hypothesis was rejected.

In M1, adding D4D supervision did not improve transfer to SERV-CT. Instead, the detector failed more severely than the original SCARED-C-only detector:

\[
\mathrm{FalseUpdate}_{\mathrm{M1, SERV}}
=
49.08\%,
\]

compared with approximately \(38.3\%\) for the previous SCARED-C-only model.

The SERV-CT EPE increased by:

\[
1.4140 - 1.0463 = 0.3678\ \text{px},
\]

and \(78.6\%\) of evaluated frames were degraded.

The detector remained effectively non-discriminative:

\[
\mathrm{AUROC}_{\mathrm{M1, SERV}}=0.468.
\]

Therefore, exposure to SCARED-C and D4D did not produce a transferable definition of raw stereo error for SERV-CT.

M2 behaved more conservatively. Training with SCARED-C and SERV-CT reduced the risk on unseen D4D:

\[
\mathrm{FalseUpdate}_{\mathrm{M2, D4D}}=1.80\%,
\]

with a small aggregate gain:

\[
\Delta\mathrm{EPE}=+0.0054\ \text{px}.
\]

However, this behaviour was obtained at only \(3.31\%\) intervention coverage and was not consistent across specimens:

\[
\begin{array}{lr}
\toprule
\text{D4D specimen} & \text{Gain}\\
\midrule
1 & +0.00765\ \text{px}\\
2 & -0.00036\ \text{px}\\
3 & -0.01269\ \text{px}\\
\bottomrule
\end{array}
\]

Thus, M2 primarily learned conservative abstention rather than a useful and repeatable cross-domain authorization rule.

The main scientific conclusion is:

\begin{quote}
Supervised exposure to two heterogeneous surgical domains is not sufficient for the current Raw Error Detector to learn an authorization rule that transfers safely to a third domain.
\end{quote}

This result rejects the simple hypothesis:

\[
\text{more supervised domains}
\quad\Longrightarrow\quad
\text{domain-invariant raw-error detection}.
\]

Instead, the meaning of ``raw prediction is wrong'' remains strongly dependent on the acquisition and error distribution of each domain.

The experiments indicate that the current detector learns:

\[
P
\left(
e_{\mathrm{raw}}>\epsilon
\mid
\text{training-domain evidence}
\right),
\]

rather than a universal intervention rule.

This explains why:

\begin{itemize}
    \item it generalizes across unseen stereo backbones inside SCARED-C;
    \item it preserves reasonable performance on domains represented during training;
    \item it can remain confidently wrong on a third unseen domain;
    \item adding domains may reduce coverage without improving decision quality.
\end{itemize}

The limitation is therefore not solved by simply concatenating supervised datasets.

The next design step should stop estimating raw error as the final authorization target. Instead, authorization should be conditioned on the specific correction proposal:

\[
\textit{Will applying this bounded proposal improve the current prediction?}
\]

Define the proposal utility:

\[
u(p)
=
e_{\mathrm{raw}}(p)
-
e_{\mathrm{proposal}}(p).
\]

A positive value indicates that the proposed correction is beneficial:

\[
u(p)>0
\quad\Longrightarrow\quad
e_{\mathrm{proposal}}(p)<e_{\mathrm{raw}}(p).
\]

The next ARGOS v2 decision module should estimate either:

\[
\hat u(p),
\]

or the probability:

\[
P
\left(
u(p)>\epsilon
\right),
\]

with an explicit neutral region and conservative abstention.

This changes the decision problem from:

\[
\textit{Is the raw disparity wrong?}
\]

to:

\[
\boxed{
\textit{Is this particular bounded temporal correction expected to help?}
}
\]

This proposal-conditioned authorization is more directly aligned with the actual safety objective of ARGOS v2.

\subsection{Milestone: Proposal Utility Becomes Effective When Used as a Conditional Safety Veto}

The dual-stage authorization experiment evaluated whether the frozen P4 proposal-applicability model could improve the safety of the existing Raw Error authorization policy.

The cascade was defined as:

\[
a_{\mathrm{final}}(p)
=
a_{\mathrm{raw}}(p)
\left(
1-v_{\mathrm{P4}}(p)
\right),
\]

where \(a_{\mathrm{raw}}\) is the frozen Raw Error authorization and \(v_{\mathrm{P4}}\) can only reject an already authorized A2 proposal.

The final output remained structurally safe:

\[
d_{\mathrm{out}}(p)
=
\begin{cases}
d_{\mathrm{A2}}(p), & a_{\mathrm{final}}(p)=1,\\
d_{\mathrm{raw}}(p), & a_{\mathrm{final}}(p)=0.
\end{cases}
\]

Therefore, abstention returned the raw disparity bit-exactly, while accepted interventions returned the frozen A2 proposal bit-exactly.

The selected cascade used the rule:

\[
v_{\mathrm{P4}}(p)
=
\mathbb{1}
\left[
\hat u_{\mathrm{P4}}(p)\leq 0
\right],
\]

where \(\hat u_{\mathrm{P4}}\) is the predicted proposal utility.

The final held-out SCARED-C results were:

\[
\begin{array}{lrrr}
\toprule
\text{Metric}
&
\text{Raw Error authorization}
&
\text{P4 cascade}
&
\text{Target}
\\
\midrule
\text{EPE}
& 0.934660 & 0.934845 & \approx \text{baseline}
\\
\text{Gain}
& 0.026946 & 0.026760 & \geq 80\%
\\
\text{Gain retained}
& \text{--} & 99.31\% & \geq 80\%
\\
\text{False update}
& 1.90\% & 1.53\% & <1.25\%
\\
\text{Clean degradation}
& 0.98\% & 0.73\% & <0.60\%
\\
\text{Intervention precision}
& 75.63\% & 79.46\% & >80\%
\\
\text{Coverage}
& 5.62\% & 4.72\% & >0
\\
\bottomrule
\end{array}
\]

The cascade retained:

\[
99.31\%
\]

of the original authorized-A2 gain, while reducing false updates by:

\[
\frac{1.90-1.53}{1.90}
=
19.3\%,
\]

and clean-pixel degradation by:

\[
\frac{0.98-0.73}{0.98}
=
25.5\%.
\]

The EPE penalty relative to the original Raw Error authorization was only:

\[
0.934845-0.934660
=
0.000185\ \mathrm{px}.
\]

A rejection-rate-matched random veto produced an EPE of:

\[
0.940068,
\]

demonstrating that the P4 veto contains genuine information and that the safety improvement is not explained solely by reduced intervention coverage.

All three seen stereo backbones remained improved relative to their raw predictions:

\[
\begin{array}{lrrr}
\toprule
\text{Backbone}
&
\text{Raw}
&
\text{Raw Error authorization}
&
\text{P4 cascade}
\\
\midrule
\text{S2M2-S}
& 1.028950 & 1.009592 & 1.009830
\\
\text{RAFT-Stereo}
& 0.925869 & 0.897324 & 0.897360
\\
\text{StereoAnywhere}
& 0.929996 & 0.897062 & 0.897347
\\
\bottomrule
\end{array}
\]

Although the predeclared promotion criteria were narrowly missed, the experiment demonstrates that raw-error authorization and proposal-utility estimation provide complementary decision signals.

The conditional audit revealed that, among proposals already authorized by the Raw Error Detector:

\[
56.91\%
\]

were helpful,

\[
32.00\%
\]

were harmful, and:

\[
11.09\%
\]

were indifferent.

This conditional distribution is substantially different from the global proposal-utility distribution, where harmful proposals represented only approximately \(1.93\%\) of all pixels.

P4 was originally trained on the global distribution, but was deployed as a veto on the restricted distribution:

\[
P
\left(
\text{proposal harmful}
\mid
a_{\mathrm{raw}}=1
\right).
\]

Despite this training--deployment mismatch, P4 retained \(94.68\%\) of useful proposals on validation and \(91.52\%\) on the final test, while rejecting \(27.42\%\) of harmful proposals.

The central finding is therefore:

\begin{quote}
Proposal-utility prediction provides a useful safety signal when conditioned on proposals already authorized by the Raw Error Detector, but a globally trained applicability model is not optimally matched to this conditional veto task.
\end{quote}

This motivates a dedicated Conditional Harmfulness Veto trained exclusively on the decision subset produced by the frozen Raw Error authorization policy.

For each pixel satisfying:

\[
a_{\mathrm{raw}}(p)=1,
\]

the conditional target is:

\[
y_{\mathrm{harm}}(p)
=
\mathbb{1}
\left[
e_{\mathrm{A2}}(p)
>
e_{\mathrm{raw}}(p)+\epsilon
\right].
\]

The resulting architecture remains:

\[
\boxed{
\text{Raw Error authorization}
\rightarrow
\text{Conditional Harmfulness Veto}
\rightarrow
\text{frozen bounded A2 proposal}
}
\]

with:

\[
a_{\mathrm{final}}(p)
=
a_{\mathrm{raw}}(p)
\left(
1-\hat y_{\mathrm{harm}}(p)
\right).
\]

The new veto is not required to discover useful interventions or estimate raw disparity error. Its only task is to identify dangerous corrections among proposals that the first-stage policy has already decided to apply.

This conditional decomposition targets the remaining safety gap directly while preserving the gain and bit-exact abstention properties of the existing ARGOS v2 pipeline.

\subsection{Current Status of ARGOS v2}

At this stage, the core architecture of ARGOS v2 has stabilized around a clear design principle:

\begin{center}
\emph{Temporal evidence is useful; the fundamental challenge is deciding when it is safe to use it.}
\end{center}

Extensive controlled investigations have been conducted to evaluate multiple temporal refinement strategies.

\paragraph{Validated components (GO).}

\begin{itemize}
    \item \textbf{Causal BiDA alignment} successfully provides temporally aligned disparity evidence without introducing future information.
    \item \textbf{Bounded residual proposal (A2)} consistently improves disparity estimation using only geometric cues while guaranteeing bounded corrections.
    \item \textbf{Raw Error Detector} is the first promoted authorization module, reducing false updates from approximately $35\%$ to below $2\%$ while preserving more than $60\%$ of the refinement gain.
    \item The complete pipeline generalizes across multiple unseen stereo backbones (Fast-FoundationStereo and CREStereo) without retraining.
\end{itemize}

\paragraph{Rejected research directions (NO-GO).}

The following hypotheses were experimentally validated and rejected:

\begin{itemize}
    \item Long-term memory (PPMStereo-inspired)
    \item Latent recurrent state (EndoStreamDepth-inspired)
    \item Frozen semantic features (DINOv3)
    \item Cross-memory consensus correction
    \item Quality prediction for memory ranking
    \item Proposal-only authorization
    \item Lightweight support guard
    \item Multi-domain supervised training
\end{itemize}

Although several approaches demonstrated useful signals, none produced a safer or more general solution than the current bounded proposal plus authorization strategy.

\paragraph{Current scientific understanding.}

The accumulated evidence consistently indicates that temporal refinement itself is no longer the limiting factor.

Instead, the dominant challenge is reliable \emph{decision making}:

\[
\boxed{
\text{Should this particular correction be applied?}
}
\]

rather than

\[
\boxed{
\text{How should the correction be computed?}
}
\]

This shifts the contribution of ARGOS v2 from a temporal refinement network toward a \emph{safe decision framework} for plug-and-play stereo refinement.

\paragraph{Current architecture.}

\[
\text{Frozen Stereo}
\rightarrow
\text{SEA-RAFT}
\rightarrow
\text{Causal BiDA Alignment}
\rightarrow
\text{Bounded Proposal (A2)}
\rightarrow
\text{Raw Error Authorization}
\rightarrow
\text{Abstention}
\]

The proposal-conditioned veto remains an optional safety-enhanced variant. Although it preserves over $99\%$ of the authorized gain while reducing false updates and clean degradation, it narrowly misses the pre-registered promotion criteria and is therefore retained as a positive ablation rather than the primary model.

\paragraph{Current research direction.}

The next phase no longer focuses on inventing new temporal modules.

Instead, the objective is to understand and improve the authorization mechanism itself, investigating whether harmful proposals can be identified more accurately without sacrificing the gain already achieved by the bounded residual refiner.

\subsection{From Temporal Oracle Gain to Deployable Selective Refinement}
\label{sec:temporal_selection_findings}

Our experiments show that causally propagated disparity memory contains substantial complementary geometric information. Given the current stereo estimate \(d_t^{\mathrm{raw}}\) and the temporally aligned memory estimate \(d_t^{\mathrm{mem}}\), a ground-truth oracle that selects the better candidate independently at each pixel consistently outperforms the raw stereo prediction. Across the stereo backbones used during development, the oracle reduces the mean end-point error (EPE) from \(0.2666\) to \(0.2210\), corresponding to a relative improvement of \(17.1\%\). The gain remains positive on unseen backbones, where the oracle provides a relative EPE reduction of \(10.5\%\). These results indicate that the aligned temporal memory is not merely temporally smoother, but contains genuine corrections to errors made by the current-frame stereo estimator.

However, the existence of a strong oracle ceiling does not imply that this information can be exploited reliably without ground truth. A broad set of deterministic and learned selection strategies was investigated, including direct replacement, fixed blending, forward--backward flow consistency, photometric gating, PPMStereo-inspired cues, DINO features, recurrent state representations, consensus signals, quality prediction, support guards, and multi-domain selector training. None of these approaches recovered the oracle gain with sufficient reliability across sequence-disjoint evaluation splits.

The final selector was evaluated using three independent training seeds and a strict sequence-disjoint protocol. The raw stereo baseline achieved an EPE of
\[
\mathrm{EPE}_{\mathrm{raw}} = 0.5337,
\]
whereas the selective raw-versus-memory policy obtained
\[
\mathrm{EPE}_{\mathrm{sel}} = 0.5218 \pm 0.0031.
\]
This corresponds to an absolute improvement of
\[
\Delta \mathrm{EPE} = 0.0119 \pm 0.0031.
\]
Although the average change is positive, the selector recovers only
\[
22.7\% \pm 5.8\%
\]
of the oracle improvement. Moreover, the effective intervention coverage is only
\[
0.97\% \pm 1.07\%,
\]
while
\[
9.0\% \pm 5.5\%
\]
of the evaluated frames are degraded relative to the raw baseline.

These results lead to a clear interpretation. The selector is conservative enough to remain beneficial on average, but it intervenes too rarely and captures too little of the available oracle gain. It therefore does not satisfy our predefined requirement of recovering at least \(50\%\) of the oracle improvement. Although all three development backbones improve slightly, the magnitude of the gain is insufficient to justify scaling the method to additional backbones as a final stereo-improvement result.

Earlier positive results obtained with smaller development protocols should consequently be interpreted as exploratory evidence rather than confirmatory results. Those experiments were useful for identifying promising signals and failure modes, but they are superseded for final claims by the stricter sequence-disjoint, multi-seed evaluation.

\subsubsection{Interpretation of the Bottleneck}

The main bottleneck is not the existence of useful temporal information, nor does the evidence suggest that optical-flow alignment is the dominant limitation. Instead, the unresolved problem is decision-making: determining, from observable inference-time cues, whether \(d_t^{\mathrm{mem}}\) is more accurate than \(d_t^{\mathrm{raw}}\) at a given pixel or region.

Many investigated cues identify ambiguity, disagreement, poor flow support, or possible inconsistency. However, such signals do not necessarily reveal the sign of the relative error
\[
\Delta e_t(x)
=
e\!\left(d_t^{\mathrm{raw}}(x)\right)
-
e\!\left(d_t^{\mathrm{mem}}(x)\right),
\]
where \(e(\cdot)\) denotes the unknown geometric error with respect to ground truth. In particular, disagreement between the two candidates may indicate that one estimate is wrong, but does not identify which estimate is more accurate. This distinction explains why a selector may achieve non-trivial ranking performance while still producing limited or even negative geometric utility at an operational threshold.

The out-of-distribution experiments further indicate that this limitation cannot be reduced to backbone shift alone. CREStereo remained relatively close to the training support, whereas SERV-CT and D4D exhibited considerably stronger distribution shifts. Neither global confidence thresholds nor feature-space guards based on Mahalanobis distance or \(k\)-nearest neighbours produced a satisfactory trade-off: conservative guards reduced the intervention coverage almost to zero, while less restrictive settings retained harmful failures.

\subsubsection{Relation to Prior Work and Novelty}

The underlying candidate-selection problem is closely related to CODD, which already considers causal fusion between a current stereo estimate and an aligned temporal estimate, and supervises its fusion weights using the relative errors of the two candidates. Consequently, the generic idea of learning whether current or temporally propagated disparity is preferable cannot by itself constitute the main novelty of ARGOS v2.

A defensible contribution would instead require an explicit selective-decision formulation with a keep-current default, together with calibrated control of harmful interventions. Relevant distinctions from conventional temporal fusion include:
\begin{itemize}
    \item an explicit decision between using temporal memory and retaining the current estimate;
    \item abstention when the predicted benefit is insufficiently reliable;
    \item optimization for net geometric utility rather than temporal smoothness alone;
    \item evaluation through intervention coverage, harm rate, oracle recovery, and risk--coverage behaviour;
    \item strict sequence-disjoint and cross-backbone validation.
\end{itemize}

The present results do not yet establish such a method. They instead identify a robust scientific gap between the amount of complementary information available in temporal memory and the fraction of that information that can be safely recovered by an inference-time selector.

\subsubsection{Current Conclusion}

The findings can be summarized as follows:
\begin{enumerate}
    \item BiDA-aligned temporal memory contains genuine and substantial complementary geometric information.
    \item Blind replacement or blending of the raw stereo estimate is not reliable.
    \item The tested selectors do not yet recover a sufficient fraction of the oracle gain.
    \item The current ARGOS v2 pipeline is therefore not ready to be presented as a general stereo-improvement method.
    \item The central open problem is safe prediction of the relative utility of current and propagated disparity estimates.
\end{enumerate}

Accordingly, the most accurate current framing is not that temporal refinement has been solved, but that temporal memory exhibits a strong and reproducible oracle advantage while safe per-pixel exploitation of that advantage remains unresolved. Future work should therefore focus on determining whether the limitation arises from target construction, sampling and calibration, or from an intrinsic lack of sufficiently informative observable cues. Until this gap is closed under the strict evaluation protocol, further architectural scaling or additional out-of-distribution experiments would not provide a reliable basis for a final method claim.

\subsection{Relation to Selective Prediction and Temporal Memory Arbitration}
\label{sec:related_selective_temporal}

The problem studied in ARGOS v2 is related to several partially overlapping research areas: temporal stereo fusion, confidence estimation, learning to defer, selective prediction, and quality-aware temporal memory. However, these areas differ substantially in the quantity being predicted and in the operational objective.

Classical confidence and uncertainty estimation methods associate a reliability score with a single prediction. Such a score can identify potentially inaccurate disparity estimates, but it does not necessarily determine which of two competing candidates is preferable. In our setting, the relevant quantity is instead the signed difference between the errors of the current and propagated estimates,
\begin{equation}
    \Delta e_t(x)
    =
    \left|d_t^{\mathrm{raw}}(x)-d_t^{\mathrm{gt}}(x)\right|
    -
    \left|d_t^{\mathrm{mem}}(x)-d_t^{\mathrm{gt}}(x)\right|.
\end{equation}
A positive value indicates that replacing the current estimate with temporal memory would be beneficial, whereas a negative value indicates a harmful intervention. Consequently, generic confidence ranking is insufficient unless it is also informative about this comparative error.

This distinction is consistent with results from learning-to-defer and cascade-deferral theory. Optimal arbitration between two predictors generally depends on their relative expected losses rather than on the confidence of either predictor in isolation \citep{okati2021triage,jitkrittum2023confidence}. This is particularly relevant when the alternative predictor behaves as a specialist: temporally propagated disparity can be highly accurate in static and well-aligned regions, yet unreliable around occlusions, newly visible surfaces, motion boundaries, or previously corrupted estimates.

Selective prediction further introduces a reject option, allowing the system to avoid an intervention when its expected benefit is uncertain \citep{geifman2017selective,geifman2019selectivenet,zaoui2020regression}. In ARGOS v2, rejection is naturally interpreted as retaining the current stereo estimate. The objective is therefore not merely to identify uncertain pixels, but to maximize geometric benefit while controlling the risk that accepted temporal updates worsen the baseline.

\subsubsection{Closest Prior Art: CODD}

The closest prior work is CODD, which performs causal online depth estimation by combining a current stereo disparity with an aligned temporal estimate \citep{li2023codd}. CODD already contains several elements that closely match the problem considered here:

\begin{itemize}
    \item a current-frame stereo candidate and an aligned temporal candidate;
    \item causal per-pixel arbitration;
    \item supervision based on the relative errors of the two candidates;
    \item margin regions in which the candidates are treated as similarly accurate;
    \item a structural fallback to the current estimate;
    \item an empirical oracle demonstrating the potential gain from selecting the better candidate.
\end{itemize}

Its fused prediction can be expressed as
\begin{equation}
    d_t^{\mathrm{fused}}
    =
    \left(1-w_t^{\mathrm{reset}}w_t^{\mathrm{fusion}}\right)
    d_t^{\mathrm{raw}}
    +
    w_t^{\mathrm{reset}}w_t^{\mathrm{fusion}}
    d_t^{\mathrm{mem}}.
\end{equation}
Thus, the general concept of learning whether temporal memory is preferable to a current stereo estimate cannot be considered novel by itself.

There are nevertheless important differences between CODD and the intended ARGOS v2 formulation. CODD performs continuous soft fusion and is primarily motivated by temporally consistent online depth. It uses a comparatively rich and tightly integrated evidence set, including learned scene motion, per-pixel transformations, internal stereo features, appearance correlations, visibility, semantic information, and previously fused memory. By contrast, ARGOS v2 investigates whether the relative utility of memory can be estimated from lightweight, frozen, and largely backbone-independent signals.

Moreover, CODD does not formulate deployment as an explicit risk-controlled switch with a reject option, nor does it report intervention coverage, harmful-update rate, oracle recovery, or risk--coverage curves against a never-update baseline. The potential contribution space is therefore not generic temporal candidate selection, but the narrower problem of safe and calibrated temporal intervention under limited observable evidence.

\subsubsection{Related Gating Mechanisms in Other Video Tasks}

Several adjacent fields contain mechanisms that reject or attenuate temporal information. Dynamic video segmentation methods choose between warped previous predictions and current-frame computation \citep{xu2018dvsnet,nilsson2018grfp,jain2019accel}. Video object segmentation methods regulate which observations are written to memory or restrict the memory bank to reduce degradation from stale information \citep{liu2022qdmn,cheng2022xmem,zhou2024rmem}. Tracking systems similarly learn whether an appearance model should be updated, since incorrect updates can cause irreversible drift \citep{zhang2019updatenet,dai2020ltmu}.

These methods establish that temporal memory can be conditionally harmful and that learned gating is useful. However, most predict absolute quality, optimize the downstream task implicitly, or operate at frame or region level. They generally do not learn an explicit pixel-wise comparative error between two finished dense-regression candidates.

Video restoration provides additional evidence that propagation can amplify previous errors. Recurrent super-resolution systems exploit temporal recurrence effectively, while later work explicitly observes that real-world degradations and artifacts may accumulate through propagation \citep{chan2021basicvsr,chan2022basicvsrpp,chan2022realbasicvsr}. These works motivate memory rejection, but their mitigation strategies typically involve learned reconstruction, pre-cleaning, or recurrent feature refinement rather than selective replacement between two disparity maps.

\subsubsection{Risk and Net-Utility Evaluation}

The final operational value of a selector cannot be inferred from AUROC or AUPRC alone. A ranking model may separate useful and harmful updates reasonably well, yet still yield negative or negligible benefit at the threshold used for deployment. This discrepancy is closely related to decision-curve analysis, in which prediction models are evaluated by their net benefit relative to default policies rather than only by discrimination metrics \citep{vickers2006decision,vickers2016netbenefit}.

For temporal disparity selection, the natural default policy is
\begin{equation}
    \pi_0(x) = d_t^{\mathrm{raw}}(x),
\end{equation}
corresponding to never applying temporal memory. A selective policy with threshold \(\tau\) is
\begin{equation}
    \pi_{\tau}(x)
    =
    \begin{cases}
        d_t^{\mathrm{mem}}(x), & s_t(x) \geq \tau,\\
        d_t^{\mathrm{raw}}(x), & s_t(x) < \tau,
    \end{cases}
\end{equation}
where \(s_t(x)\) estimates the expected relative utility of the temporal candidate.

Evaluation should therefore report, as a function of \(\tau\),
\begin{itemize}
    \item intervention coverage;
    \item mean geometric gain;
    \item harmful-update rate;
    \item fraction of degraded frames;
    \item recovery of oracle improvement;
    \item risk--coverage or utility--coverage behaviour.
\end{itemize}

Distribution-free calibration approaches such as Learn-then-Test, conformal risk control, and conformal decision theory provide possible tools for selecting operating points under explicit risk constraints \citep{angelopoulos2021learntest,angelopoulos2022crc,lekeufack2024cdt}. Their direct application to dense, temporally correlated regression is not immediate, however, because the calibration unit, dependence structure, and definition of acceptable geometric risk must first be specified.

\subsubsection{Novelty Positioning}

Based on the reviewed literature, the following claims would not be defensible:

\begin{itemize}
    \item that ARGOS v2 introduces the first learned selection between current and temporally propagated disparity;
    \item that pairwise relative-error supervision is itself novel;
    \item that per-pixel disparity candidate selection is new;
    \item that temporal memory rejection has not been previously studied.
\end{itemize}

A more defensible contribution would require the combination of:
\begin{enumerate}
    \item a hard keep-current-versus-memory decision;
    \item explicit abstention to the current estimate;
    \item optimization and calibration for net geometric utility;
    \item direct measurement of harmful temporal interventions;
    \item risk--coverage evaluation under sequence-disjoint testing;
    \item lightweight and backbone-independent inference signals.
\end{enumerate}

The experiments presented in this work do not yet establish such a successful method. Instead, they reveal a reproducible gap between the oracle utility contained in temporal memory and the amount of utility recoverable from lightweight causal signals.

\subsection{Capacity and Observability as Competing Explanations}
\label{sec:capacity_observability}

The failure of the lightweight selector admits two fundamentally different explanations. First, the available signals may contain sufficient information, but the selected architecture, loss, sampling strategy, or calibration procedure may fail to extract it. Second, the winner between the current and temporal candidates may be only weakly identifiable from the available causal inputs.

To distinguish these cases, a capacity probe should preserve the candidates, input channels, target construction, data split, and calibration protocol, while increasing only model capacity and spatial receptive field. The interpretation is then:

\begin{itemize}
    \item If a larger model substantially increases oracle recovery on sequence-disjoint test data, the limitation lies primarily in representation capacity, optimization, or calibration.
    \item If the model fits the training set but fails to generalize, the observable signals are sequence- or domain-specific, or the comparative target is excessively noisy.
    \item If even training performance remains weak, the current input representation is insufficient to identify relative utility.
    \item If ranking metrics improve while deployment utility remains low, the main limitation lies in thresholding and cost-sensitive decision-making rather than representation.
\end{itemize}

This test cannot establish that temporal candidate selection is impossible in general. It can only establish whether the currently available causal, backbone-independent signals contain sufficient generalizable information. A negative result would imply that further progress requires additional evidence, such as internal stereo matching features, cost-volume information, learned 3D motion, additional viewpoints, or stronger geometric supervision, rather than merely a larger selector.

\subsection{Capacity Scaling Does Not Close the Temporal Oracle Gap}
\label{sec:capacity_probe}

The limited oracle recovery of the utility selector may arise from two fundamentally
different causes. First, the available causal signals may contain sufficient information,
while the selector lacks the capacity or spatial context required to extract it. Second,
the relative correctness of the current and propagated disparities may be only weakly
identifiable from the available backbone-independent observations.

To distinguish these explanations, we conducted a controlled capacity probe. The data,
candidate disparities, thirteen input signals, target construction, sampling strategy,
training objective, sequence-disjoint split, checkpoint selection, and calibration
protocol were kept fixed. Only the capacity and receptive field of the selector were
increased.

The original selector used \(64\) channels and four residual blocks, corresponding to
\(303{,}747\) parameters and an approximate receptive field of \(19\times19\) pixels.
The capacity probe increased the model to \(128\) channels and eight residual blocks,
corresponding to \(2{,}379{,}011\) parameters and a receptive field of \(35\times35\)
pixels. Both configurations were evaluated using three independent seeds on the same
held-out sequence-disjoint test set.

\begin{table}[t]
    \centering
    \caption{Capacity-only comparison of the raw-versus-memory utility selector.
    The larger model uses the same observable evidence and differs only in capacity
    and receptive field. Values are aggregated over three seeds.}
    \label{tab:capacity_probe}
    \begin{tabular}{lcc}
        \toprule
        Metric & \(64\times4\) selector & \(128\times8\) selector \\
        \midrule
        Parameters
            & \(0.304\,\mathrm{M}\)
            & \(2.379\,\mathrm{M}\) \\
        Receptive field
            & \(19\times19\)
            & \(35\times35\) \\
        EPE gain
            & \(0.01189 \pm 0.00306\)
            & \(0.01411 \pm 0.00242\) \\
        Oracle recovery
            & \(22.7\% \pm 5.8\%\)
            & \(27.0\% \pm 4.6\%\) \\
        Intervention coverage
            & \(0.97\%\)
            & \(0.99\%\) \\
        False-update rate
            & \(0.60\%\)
            & \(0.55\%\) \\
        Clean-pixel degradation
            & \(0.289\%\)
            & \(0.286\%\) \\
        Intervention precision
            & \(39.4\%\)
            & \(47.3\%\) \\
        Degraded frames
            & \(9.0\%\)
            & \(10.7\%\) \\
        Worst-frame degradation
            & \(0.445\ \mathrm{px}\)
            & \(1.065\ \mathrm{px}\) \\
        \bottomrule
    \end{tabular}
\end{table}

Increasing the parameter count by approximately a factor of eight produces a measurable
but limited improvement. Mean oracle recovery increases by only \(4.3\) percentage
points, from \(22.7\%\) to \(27.0\%\), while the EPE gain increases by approximately
\(0.00222\) pixels. The larger selector improves all three evaluated stereo backbones
across all three seeds, indicating that additional capacity can extract some further
signal from the existing inputs.

However, the operational behaviour remains essentially unchanged. Intervention coverage
stays below \(1\%\), oracle recovery remains far below the predefined \(50\%\) promotion
requirement, and every seed continues to fail the promotion gate. Moreover, although
average false-update and clean-pixel degradation rates remain low, the worst-frame
degradation increases from \(0.445\) to \(1.065\) pixels. Thus, greater capacity does not
eliminate the harmful tail of the intervention distribution.

These results indicate that model capacity is not the dominant bottleneck. The available
raw-disparity, propagated-memory, optical-flow, support, and consistency signals contain
some predictive information, as demonstrated by the improvement of the larger model.
Nevertheless, they do not contain enough generalizable evidence to reliably determine
which candidate is geometrically more accurate over a sufficiently large fraction of
pixels.

The result should not be interpreted as proving that temporal candidate selection is
impossible. Rather, it establishes a narrower observability result:

\begin{quote}
Given the currently available causal and backbone-independent signals, scaling the
selector alone does not recover a substantial fraction of the temporal oracle advantage.
\end{quote}

Further progress therefore requires an information probe rather than another capacity
increase. In particular, the selector should be provided with additional evidence that
is directly related to stereo matching ambiguity, such as candidate-conditioned matching
costs, local cost-volume shape, confidence margins, internal stereo features, or an
independent geometric signal. A controlled information ablation can then determine
whether the remaining oracle gap is caused by missing observable evidence or by a deeper
generalization limitation in relative utility prediction.


\subsection{Milestone: CODD-Style Fusion Recovers More Than Half of the Temporal Oracle}
\label{sec:codd_style_milestone}

The preceding experiments established that the aligned temporal memory contains genuine
complementary geometric information, but that this information cannot be recovered
effectively through lightweight post-hoc selection alone. Increasing the selector capacity
from \(0.304\) to \(2.379\) million parameters raised oracle recovery only from
\(22.7\%\) to approximately \(27.0\%\). Similarly, augmenting the selector with compact
candidate-conditioned census evidence produced no material improvement.

These negative probes motivated a more substantial change in how temporal evidence is
represented and consumed. We therefore implemented a CODD-style causal fusion module
inspired by \citet{li2023codd}, while retaining the existing BiDA-aligned temporal memory.
The fused disparity is defined as
\begin{equation}
    d_t^{F}
    =
    \left(
        1 -
        w_t^{\mathrm{reset}}
        w_t^{\mathrm{fusion}}
    \right)
    d_t^{S}
    +
    w_t^{\mathrm{reset}}
    w_t^{\mathrm{fusion}}
    d_t^{M},
    \label{eq:codd_style_fusion}
\end{equation}
where \(d_t^{S}\) is the current stereo prediction, \(d_t^{M}\) is the aligned temporal
memory, and both weights are bounded to \([0,1]\).

The reset weight is supervised to reject the temporal candidate when its error is
substantially larger than that of the current estimate. The fusion weight is supervised
to favour the lower-error candidate outside a dead-band and to approach equal averaging
when the two candidates have similar errors. In addition, the final fused disparity is
directly supervised against ground truth using a robust regression loss.

Unlike the previous hard selector, the CODD-style module receives learned stereo-matching
and spatio-temporal evidence, including candidate-conditioned left--right feature
distances, local disparity and appearance self-correlations, cross-frame correlations,
motion magnitude, flow validity, visibility, and support information. Training is
performed through a four-frame causal unrolling procedure in which subsequent frames
consume the previous fused estimate, thereby reducing the mismatch between training and
online inference.

\begin{table}[t]
    \centering
    \caption{Strict three-seed result of the CODD-style fusion module on the held-out
    dataset-7 split. The raw-versus-memory oracle is computed using the same candidate
    pair and evaluation support.}
    \label{tab:codd_style_main_result}
    \begin{tabular}{lc}
        \toprule
        Metric & CODD-style BiDA fusion \\
        \midrule
        Raw stereo EPE
            & \(0.53317\) \\
        Fused EPE
            & \(0.50357 \pm 0.00176\) \\
        Absolute EPE gain
            & \(0.02960 \pm 0.00176\) \\
        Raw-versus-memory oracle gain
            & \(0.05233\) \\
        Oracle recovery
            & \(56.56\% \pm 3.36\%\) \\
        Harmful-update rate
            & \(0.658\% \pm 0.074\%\) \\
        Clean-pixel degradation
            & \(3.41\% \pm 0.28\%\) \\
        Degraded frames
            & \(27.4\%\) \\
        Worst-frame degradation
            & \(0.681 \pm 0.327\ \mathrm{px}\) \\
        \bottomrule
    \end{tabular}
\end{table}

The CODD-style configuration reduces EPE from \(0.53317\) to
\(0.50357 \pm 0.00176\), corresponding to an average gain of
\(0.02960 \pm 0.00176\) pixels. Most importantly, it recovers
\(56.56\% \pm 3.36\%\) of the raw-versus-memory oracle advantage, exceeding the
predefined promotion requirement of \(50\%\). The improvement is observed across all
three evaluated stereo backbones, indicating that it is not produced by a single
backbone-specific failure mode.

\begin{table}[t]
    \centering
    \caption{Progressive recovery of the temporal oracle under the strict
    sequence-disjoint protocol.}
    \label{tab:oracle_recovery_progression}
    \begin{tabular}{lc}
        \toprule
        Configuration & Oracle recovery \\
        \midrule
        Lightweight \(64\times4\) hard selector
            & \(22.72\% \pm 5.84\%\) \\
        Capacity-scaled \(128\times8\) hard selector
            & \(26.96\% \pm 4.63\%\) \\
        \(128\times8\) selector with local census evidence
            & \(27.24\% \pm 1.85\%\) \\
        CODD-style recurrent soft fusion
            & \(\mathbf{56.56\% \pm 3.36\%}\) \\
        \bottomrule
    \end{tabular}
\end{table}

This progression changes the interpretation of the previous negative results. The
raw-versus-memory decision is not fundamentally unobservable. Rather, it is poorly
identified by simple post-hoc cues and hard selective replacement. The substantial
increase obtained by the CODD-style configuration indicates that the missing ingredients
were a combination of richer learned matching evidence, explicit spatial and temporal
correlations, recurrent training with fused memory, direct supervision of the fused
disparity, and a continuous fusion mechanism.

The result represents the first ARGOS v2 configuration to convert the strong BiDA oracle
ceiling into a substantial deployable improvement. It therefore promotes temporal fusion
from an unresolved diagnostic problem to a promising method direction.

\paragraph{Difference from CODD.}
The experiment should be described as CODD-style rather than as a faithful reproduction
of CODD. The temporal candidate is aligned using the existing BiDA and optical-flow
pipeline rather than a learned per-pixel SE(3) motion network. Moreover, because the
cached stereo backbones do not expose a common internal feature representation, the
fusion module uses a shared frozen feature encoder instead of the native stereo features
used by CODD. The experiment therefore isolates whether CODD's fusion principles transfer
to the ARGOS setting without simultaneously introducing its learned 3D motion module.

\paragraph{Remaining safety limitation.}
Passing the oracle-recovery gate does not imply that the method is fully ready for
deployment. Although the average geometric gain is substantial, approximately
\(27.4\%\) of frames are degraded, and the worst-frame degradation remains non-negligible.
The next stage must therefore focus on reproducibility, common-support verification,
component ablation, cross-backbone transfer, and risk-controlled operation rather than
immediately adding a more complex motion model.

\paragraph{Milestone conclusion.}
The principal result of this phase is that BiDA memory can be exploited effectively when
it is embedded in a CODD-style recurrent fusion architecture. ARGOS v2 therefore has a
credible temporal stereo improvement mechanism for the first time. The central open
question shifts from whether the temporal oracle is recoverable to whether the recovered
gain can be made robust across frames, backbones, and deployment conditions.

\subsection{CODD-Style Recurrent Fusion with BiDA-Aligned Memory}
\label{sec:codd_style_configuration}

The preceding experiments demonstrated that causally aligned temporal disparity contains
substantial complementary geometric information, but that lightweight post-hoc selectors
recover only a limited fraction of this potential. In particular, increasing the capacity
of the hard selector and adding compact census-based candidate evidence both resulted in
approximately \(27\%\) recovery of the historical raw-versus-memory oracle.

To determine whether this limitation originated from the fusion formulation and observable
evidence rather than from the temporal candidate itself, we implemented a causal
CODD-style fusion module inspired by \citet{li2023codd}. The configuration does not
retrain the stereo predictors and does not modify the validated BiDA/SEA-RAFT alignment
pipeline. Instead, it replaces the previous hard raw-versus-memory decision with a
continuous, evidence-conditioned fusion of the current stereo estimate and a causally
propagated temporal state.

\subsubsection{Current and Temporal Disparity Candidates}

Let
\begin{equation}
    d_t^{S}
\end{equation}
denote the current raw disparity produced by a frozen stereo backbone at time \(t\).
The temporal state used at the previous time step is defined as
\begin{equation}
    d_{t-1}^{\mathrm{state}}
    =
    \begin{cases}
        d_{t-1}^{S},
        & \text{at the first step of a training clip},\\[2mm]
        d_{t-1}^{F},
        & \text{at subsequent steps},
    \end{cases}
    \label{eq:codd_state_definition}
\end{equation}
where \(d_{t-1}^{F}\) is the preceding fused prediction.

SEA-RAFT estimates target-to-source optical flow between the current and previous left
images. The validated causal BiDA operator then aligns the preceding state with the
current frame:
\begin{equation}
    d_t^{M}
    =
    \operatorname{BiDAWarp}
    \left(
        d_{t-1}^{\mathrm{state}},
        F_{t\rightarrow t-1}
    \right).
    \label{eq:codd_bida_memory}
\end{equation}

The aligned memory is accompanied by warp support, disparity validity,
forward--backward flow consistency, and visibility information. No future frame is used.

\subsubsection{Dual-Weight Soft Fusion}

The fusion module predicts two spatially varying weights,
\begin{equation}
    w_t^{\mathrm{reset}},\;
    w_t^{\mathrm{fusion}}
    \in [0,1]^{H\times W}.
\end{equation}

The effective temporal contribution is
\begin{equation}
    w_t^{T}
    =
    w_t^{\mathrm{reset}}
    w_t^{\mathrm{fusion}}.
    \label{eq:effective_temporal_weight}
\end{equation}

The output disparity follows the CODD fusion equation:
\begin{equation}
    d_t^{F}
    =
    \left(1-w_t^{T}\right)d_t^{S}
    +
    w_t^{T}d_t^{M}.
    \label{eq:argos_codd_fusion}
\end{equation}

Therefore,
\begin{equation}
    w_t^{T}=0
    \quad\Longrightarrow\quad
    d_t^{F}=d_t^{S},
\end{equation}
whereas
\begin{equation}
    w_t^{T}=1
    \quad\Longrightarrow\quad
    d_t^{F}=d_t^{M}.
\end{equation}

Intermediate values produce a convex interpolation between the current stereo prediction
and the aligned temporal memory. Unlike the previous ARGOS hard selector, the primary
configuration does not threshold \(w_t^{T}\) during inference.

\subsubsection{CODD-Style Comparative Supervision}

Let the current and temporal disparity errors be
\begin{equation}
    e_t^{S}
    =
    \left|d_t^{S}-d_t^{GT}\right|,
    \qquad
    e_t^{M}
    =
    \left|d_t^{M}-d_t^{GT}\right|,
\end{equation}
and define their relative difference as
\begin{equation}
    \Delta e_t
    =
    e_t^{M}-e_t^{S}.
\end{equation}

The reset head is trained to suppress temporal memory when it is substantially worse than
the current prediction:
\begin{equation}
    \mathcal{L}_{\mathrm{reset}}
    =
    \begin{cases}
        w_t^{\mathrm{reset}},
        & \Delta e_t > \tau_{\mathrm{reset}},\\[1mm]
        1-w_t^{\mathrm{reset}},
        & \Delta e_t < -\tau_{\mathrm{reset}},\\[1mm]
        0,
        & \text{otherwise}.
    \end{cases}
    \label{eq:argos_reset_loss}
\end{equation}

The fusion head uses a narrower dead-band:
\begin{equation}
    \mathcal{L}_{\mathrm{fusion}}
    =
    \begin{cases}
        w_t^{\mathrm{fusion}},
        & \Delta e_t > \tau_{\mathrm{fusion}},\\[1mm]
        1-w_t^{\mathrm{fusion}},
        & \Delta e_t < -\tau_{\mathrm{fusion}},\\[1mm]
        \alpha_{\mathrm{reg}}
        \left|w_t^{\mathrm{fusion}}-0.5\right|,
        & \left|\Delta e_t\right|
          \leq \tau_{\mathrm{fusion}}.
    \end{cases}
    \label{eq:argos_fusion_loss}
\end{equation}

CODD defines
\(\tau_{\mathrm{reset}}=5\) and
\(\tau_{\mathrm{fusion}}=1\) in native-resolution disparity pixels.
The SCARED-C images have width \(1280\), whereas the cached disparity maps have width
\(180\). The thresholds were consequently converted while preserving their image-space
meaning:
\begin{align}
    \tau_{\mathrm{reset}}
    &=
    5\frac{180}{1280}
    =
    0.703125,
    \\
    \tau_{\mathrm{fusion}}
    &=
    1\frac{180}{1280}
    =
    0.140625.
\end{align}

The fused output is additionally supervised using a Huber loss:
\begin{equation}
    \mathcal{L}_{\mathrm{disp}}
    =
    \operatorname{Huber}
    \left(
        d_t^{F},
        d_t^{GT}
    \right).
\end{equation}

The complete objective is
\begin{equation}
    \mathcal{L}
    =
    \alpha_{\mathrm{disp}}\mathcal{L}_{\mathrm{disp}}
    +
    \alpha_{\mathrm{reset}}\mathcal{L}_{\mathrm{reset}}
    +
    \alpha_{\mathrm{fusion}}\mathcal{L}_{\mathrm{fusion}},
    \label{eq:argos_codd_total_loss}
\end{equation}
with
\begin{equation}
    \alpha_{\mathrm{disp}}
    =
    \alpha_{\mathrm{reset}}
    =
    \alpha_{\mathrm{fusion}}
    =
    1,
    \qquad
    \alpha_{\mathrm{reg}}=0.2.
\end{equation}

The two heads therefore perform related but distinct functions. The reset head learns
coarse outlier rejection, whereas the fusion head estimates how strongly temporal
information should be incorporated when memory is not clearly invalid.

\subsubsection{Candidate-Conditioned Stereo Evidence}

The fusion network receives \(142\) input channels. These include the thirteen universal
geometric channels used by the previous ARGOS selectors, together with explicit stereo
matching, spatial correlation, temporal correlation, and contextual evidence.

\paragraph{Stereo matching at the candidate disparities.}

A frozen ImageNet-pretrained ResNet-18 layer-1 encoder extracts feature maps from the
current rectified left and right images. For both \(d_t^{S}\) and \(d_t^{M}\), the right
feature map is sampled at the corresponding disparity location. Matching evidence is also
evaluated at offsets
\begin{equation}
    \{-1,0,+1\}
\end{equation}
in the feature grid.

The candidate-conditioned matching cost is based on the channel-wise feature distance:
\begin{equation}
    C(d)
    =
    \frac{1}{C}
    \left\|
        \phi_L(x,y)
        -
        \phi_R(x-d,y)
    \right\|_1.
    \label{eq:candidate_matching_cost}
\end{equation}

Because the ResNet feature map operates at one-quarter resolution, an offset of one
corresponds to one feature-grid pixel rather than one cached-disparity pixel.

This evidence allows the network to observe not only that the raw and temporal
disparities disagree, but also which candidate is better supported by the current
left--right stereo correspondence.

\paragraph{Local disparity correlations.}

Pixel-to-patch correlations are computed in \(3\times3\) neighborhoods with dilation \(2\)
for both current and temporal disparity. The similarity between a central disparity and a
neighboring value is represented as
\begin{equation}
    S_d(d_c,d_n)
    =
    1-
    \frac{
        \min\left(
            |d_n-d_c|,
            4
        \right)
    }{4}.
    \label{eq:disparity_similarity}
\end{equation}

These features describe local continuity, boundaries, isolated disparity errors, and
structural disagreement between the two candidates.

\paragraph{Appearance and cross-frame correlations.}

The input also includes:

\begin{itemize}
    \item current-frame feature self-correlation;
    \item aligned previous-frame feature self-correlation;
    \item cross-frame correlation between current features and BiDA-warped past features;
    \item local cross-correlation between current and memory disparity structures.
\end{itemize}

The cross-frame appearance features use the same validated target-to-source BiDA warp as
the temporal disparity. No alternative flow convention is introduced.

\paragraph{Geometric and contextual cues.}

The remaining channels include:

\begin{itemize}
    \item current raw disparity \(d_t^{S}\);
    \item aligned memory disparity \(d_t^{M}\);
    \item signed and absolute raw--memory disagreement;
    \item horizontal and vertical optical flow;
    \item flow magnitude;
    \item forward--backward flow confidence;
    \item warp support and aligned-disparity validity;
    \item current RGB;
    \item \(64\) frozen ResNet context channels.
\end{itemize}

The principal normalization ranges are:
\begin{align}
    C(d) &\in [0,1],\\
    S_{\mathrm{cosine}} &\in [-1,1],\\
    F_x,F_y &\in [-32,32],\\
    d &\in [0,64],\\
    I_{\mathrm{RGB}} &\in [-1,1].
\end{align}

The ResNet feature extractor contains \(157{,}504\) frozen parameters. Only the fusion
head is optimized.

\subsubsection{Fusion Architecture}

The trainable fusion module contains \(177{,}338\) parameters, approximately matching the
small scale of the original CODD fusion component.

Its fine-resolution branch consists of:
\begin{equation}
    \operatorname{Conv}_{3\times3}(142,24)
    \rightarrow
    \operatorname{GroupNorm}
    \rightarrow
    \operatorname{SiLU}
    \rightarrow
    \text{residual block}.
\end{equation}

A coarse branch performs spatial downsampling and increases the representation to
\(48\) channels:
\begin{equation}
    \operatorname{AvgPool}
    \rightarrow
    \operatorname{Conv}_{s=2}
    \rightarrow
    3\times\text{residual block}.
\end{equation}

The fusion head predicts \(w_t^{\mathrm{fusion}}\) at reduced resolution and upsamples it
bilinearly. The reset head combines fine features with upsampled coarse context and predicts
\(w_t^{\mathrm{reset}}\) directly on the cached disparity grid.

The reset-head bias is initialized to \(-2\), giving
\begin{equation}
    w_t^{\mathrm{reset}}
    \approx
    \sigma(-2)
    \approx
    0.119.
\end{equation}

The fusion-head bias is initialized to zero:
\begin{equation}
    w_t^{\mathrm{fusion}}
    =
    \sigma(0)
    =
    0.5.
\end{equation}

The initial temporal contribution is therefore approximately
\begin{equation}
    w_t^{T}
    \approx
    0.119\times0.5
    \approx
    0.0595.
\end{equation}

The network consequently begins close to the raw prediction rather than initially replacing
it with temporal memory.

\subsubsection{Four-Frame Causal Training}

Training samples consist of four-frame causal clips. At the first step, temporal memory is
constructed from the previous raw disparity:
\begin{equation}
    d_{t-1}^{\mathrm{state}}
    =
    d_{t-1}^{S}.
\end{equation}

At subsequent steps, the fused output becomes the propagated state:
\begin{equation}
    d_t^{\mathrm{state}}
    =
    d_t^{F}.
\end{equation}

The model therefore consumes its own preceding fused prediction during the second, third,
and fourth steps of the clip. Gradients propagate through the complete four-step unroll,
and the clip loss is the mean of the individual frame losses:
\begin{equation}
    \mathcal{L}_{\mathrm{clip}}
    =
    \frac{1}{4}
    \sum_{k=1}^{4}
    \mathcal{L}_{t+k}.
\end{equation}

The state is reset at clip boundaries. Thus, the current experiment establishes bounded
four-frame recurrence rather than arbitrary long-sequence streaming.

\subsubsection{Training and Evaluation Protocol}

The experiment uses a strictly dataset-ID-disjoint protocol:
\begin{table}[t]
    \centering
    \caption{Dataset roles for the CODD-style Phase-1 experiment.}
    \label{tab:codd_style_split}
    \begin{tabular}{ll}
        \toprule
        Role & Dataset IDs \\
        \midrule
        Training & \(1,3,6\) \\
        Validation and checkpoint selection & \(2\) \\
        Final test & \(7\) \\
        \bottomrule
    \end{tabular}
\end{table}

The seen stereo backbones are S2M2-S, RAFT-Stereo, and StereoAnywhere. The model is
trained for \(12\) epochs with AdamW, learning rate \(2\times10^{-4}\), weight decay
\(10^{-4}\), batch size \(4\) clips, automatic mixed precision, and gradient clipping at
\(1.0\). Checkpoints are selected exclusively by validation fused EPE.

All comparisons use the common support
\begin{equation}
    \mathcal{M}
    =
    \mathcal{M}_{GT}
    \cap
    \mathcal{M}_{\mathrm{raw}}
    \cap
    \mathcal{M}_{\mathrm{aligned}}
    \cap
    \mathcal{M}_{\mathrm{warp}}.
    \label{eq:codd_common_support}
\end{equation}

Raw disparity, temporal memory, fused disparity, and oracle metrics are therefore evaluated
on the same valid pixels.

\subsubsection{Main Result}

\begin{table}[t]
    \centering
    \caption{Three-seed CODD-style Phase-1 result on the held-out dataset-7 split.}
    \label{tab:codd_style_configuration_result}
    \begin{tabular}{lc}
        \toprule
        Metric & Result \\
        \midrule
        Raw EPE
            & \(0.53317\) \\
        Fused EPE
            & \(0.50357 \pm 0.00176\) \\
        Absolute EPE gain
            & \(0.02960 \pm 0.00176\) \\
        Historical raw-versus-memory oracle gain
            & \(0.05233\) \\
        Historical normalized gain
            & \(56.56\% \pm 3.36\%\) \\
        Harmful-update rate
            & \(0.658\% \pm 0.074\%\) \\
        Clean-pixel degradation
            & \(3.41\% \pm 0.28\%\) \\
        Frames with increased EPE
            & \(27.4\%\) \\
        Worst-frame degradation
            & \(0.681 \pm 0.327\ \mathrm{px}\) \\
        \bottomrule
    \end{tabular}
\end{table}

The CODD-style module almost doubles the normalized gain obtained by the previous hard
selectors:
\begin{equation}
    \sim27\%
    \longrightarrow
    56.56\%.
\end{equation}

The improvement is also observed in temporal metrics. Raw warped disparity variation is
approximately \(0.10696\), whereas the three fused models obtain \(0.06459\),
\(0.06236\), and \(0.06040\). Similarly, raw TEPE is \(0.10883\), compared with
\(0.07483\), \(0.07126\), and \(0.07037\) for the three seeds. The reduction in temporal
variation is therefore accompanied by an improvement in geometric EPE rather than being
produced solely by blind smoothing.

\subsubsection{Why the Configuration Works}

The improvement cannot yet be attributed to one isolated component because the
configuration changes several aspects of the previous selector simultaneously. Nevertheless,
the results support several complementary mechanisms.

\paragraph{Candidate quality becomes directly observable.}

The former selector mainly observed disparity disagreement, flow, and support. These cues
can reveal that the candidates differ, but not necessarily which candidate is geometrically
supported by the current stereo pair. Candidate-conditioned left--right feature costs give
the model direct evidence about whether \(d_t^{S}\) or \(d_t^{M}\) corresponds to a plausible
stereo match.

\paragraph{Spatial correlation distinguishes noise from structure.}

Disparity disagreement alone cannot distinguish a stereo error from a real depth boundary.
Local disparity and appearance correlations provide the context required to recognize
smooth surfaces, discontinuities, isolated errors, and occlusion boundaries.

\paragraph{Cross-frame evidence tests the alignment itself.}

A propagated disparity can be plausible geometrically but originate from an incorrect
temporal correspondence. Cross-frame appearance and disparity correlations indicate whether
the aligned past observation remains compatible with the current image.

\paragraph{Reset and fusion solve different subproblems.}

A single hard score must simultaneously detect catastrophic memory failures and estimate
the preferred candidate. The dual-head design separates these tasks. The reset head removes
large temporal outliers, while the fusion head controls aggregation in less decisive regions.

\paragraph{Soft fusion can correct both candidates.}

A hard selector is restricted to
\begin{equation}
    d_t^{S}
    \quad\text{or}\quad
    d_t^{M}.
\end{equation}

The CODD-style module can instead predict any point on the segment joining them. When
\(d_t^{GT}\) lies between the two candidates, a convex combination can outperform both.
Part of the observed gain may therefore originate from interpolation rather than from
correct endpoint selection.

\paragraph{Direct disparity supervision optimizes the actual output.}

The previous selector was primarily optimized to classify relative candidate utility.
The CODD-style objective also supervises the final \(d_t^{F}\). This permits the network
to learn a weight that minimizes geometric error directly rather than only reproducing a
binary candidate label.

\paragraph{Recurrent training reduces train--inference mismatch.}

During deployment, temporal memory may contain an earlier fused output rather than an
independent raw prediction. Four-frame causal unrolling exposes the model to this feedback
during training and allows earlier corrections to propagate into subsequent states.

The most defensible current conclusion is therefore:
\begin{quote}
The failure of the previous hard selectors was not caused by an absolute absence of causal
temporal information. A compact recurrent fusion model can exploit BiDA-aligned memory when
it receives explicit stereo-matching and spatio-temporal correlation evidence, is trained
with comparative reset/fusion supervision, and is allowed to perform continuous convex
correction.
\end{quote}

\subsubsection{Interpretive Caveats}

The current configuration should be described as \emph{CODD-style}, not as a complete
reproduction of CODD. CODD uses a learned per-pixel SE(3) motion module based on RAFT3D
and native stereo-backbone features. ARGOS v2 instead uses the existing SEA-RAFT/BiDA
alignment and a common frozen ResNet-18 feature space.

The reported \(56.56\%\) value also uses a historical oracle based on
\begin{equation}
    \operatorname{warp}\!\left(d_{t-1}^{S}\right),
\end{equation}
whereas the model uses
\begin{equation}
    \operatorname{warp}\!\left(d_{t-1}^{F}\right)
\end{equation}
after the first step. Furthermore, the historical oracle selects the better endpoint,
while the model can interpolate between the endpoints. The value should therefore be
reported as a
\emph{historical-selection normalized gain}, not as the exact fraction of the full
convex-fusion ceiling recovered.

Finally, the continuous sigmoid weights do not provide bit-exact abstention. Approximately
\(27.4\%\) of frames are worsened despite the strong mean improvement. The experiment is
therefore a strong validation of the fusion mechanism, but not yet a demonstration of a
clinically safe operating policy.

\paragraph{Current milestone.}

This configuration constitutes the first robust ARGOS v2 method that converts
BiDA-aligned temporal information into substantial improvements in both geometric and
temporal accuracy. The next analysis must determine how much of the gain arises from
learned stereo evidence, convex interpolation, and recurrent candidate improvement, and
whether the four-frame state remains stable during continuous sequence-level streaming.

\subsection{Bounded Recurrence Reveals a Temporal-Memory Provenance Problem}
\label{sec:bounded_memory_finding}

The preceding CODD-style experiment established that BiDA-aligned temporal information
can improve frame-wise stereo disparity. However, the initial experiment did not determine
whether the improvement originated from recurrent state propagation, learned stereo
evidence, soft interpolation, or the periodic four-frame reset used during training and
evaluation.

We therefore conducted a controlled bounded-memory study in which the principal components
were independently removed or modified while preserving the frozen stereo predictions,
SEA-RAFT/BiDA alignment, common-support masks, optimization protocol, and causal
four-frame construction.

\subsubsection{Component Ablations}

Table~\ref{tab:bounded_memory_ablation} reports the results on the held-out dataset-7
split.

\begin{table}[t]
    \centering
    \caption{Bounded-memory ablations on the held-out dataset-7 split.
    Gain is computed relative to the corresponding raw stereo prediction.
    All disparity values are expressed in cache-grid pixels.}
    \label{tab:bounded_memory_ablation}
    \resizebox{\columnwidth}{!}{
    \begin{tabular}{lccccc}
        \toprule
        Configuration
        & EPE
        & Gain
        & Harmful updates
        & Clean degradation
        & Worst frame \\
        \midrule
        Full, fixed \(H=4\)
        & \(0.50613\)
        & \(+0.02755\)
        & \(2.85\%\)
        & \(3.72\%\)
        & \(0.41\) \\

        No recurrent fused state
        & \(0.51211\)
        & \(+0.02157\)
        & \(1.80\%\)
        & \(2.28\%\)
        & \(2.17\) \\

        No learned stereo evidence
        & \(\mathbf{0.50348}\)
        & \(\mathbf{+0.03020}\)
        & \(2.59\%\)
        & \(3.10\%\)
        & \(0.80\) \\

        Hard endpoint, \(H=4\)
        & \(0.50887\)
        & \(+0.02482\)
        & \(5.55\%\)
        & \(7.91\%\)
        & \(0.41\) \\

        Adaptive hybrid, \(H_{\max}=8\)
        & \(0.51160\)
        & \(+0.02208\)
        & \(3.85\%\)
        & \(4.80\%\)
        & \(0.44\) \\
        \bottomrule
    \end{tabular}}
\end{table}

The full bounded configuration improves the raw prediction by
\begin{equation}
    G_{\mathrm{full}}
    =
    0.02755.
\end{equation}

When the preceding fused output is never propagated and temporal memory is always
constructed from the previous raw prediction,
\begin{equation}
    d_t^{M}
    =
    \operatorname{BiDAWarp}
    \left(
        d_{t-1}^{S}
    \right),
\end{equation}
the gain remains positive:
\begin{equation}
    G_{\mathrm{no\text{-}rec}}
    =
    0.02157.
\end{equation}

The additional benefit attributable to bounded fused-state recurrence is therefore
approximately
\begin{equation}
    G_{\mathrm{full}}
    -
    G_{\mathrm{no\text{-}rec}}
    =
    0.00598.
    \label{eq:bounded_recurrence_gain}
\end{equation}

Thus, recurrence is useful but is not the sole source of temporal improvement.

\subsubsection{Short Recurrence Helps, Indefinite Recurrence Fails}

The bounded model propagates the preceding fused prediction according to
\begin{equation}
    d_t^{M}
    =
    \operatorname{BiDAWarp}
    \left(
        d_{t-1}^{F}
    \right),
\end{equation}
and produces
\begin{equation}
    d_t^{F}
    =
    \left(1-w_t\right)d_t^{S}
    +
    w_t d_t^{M}.
    \label{eq:bounded_soft_fusion}
\end{equation}

With a reset every four frames, this recurrent correction remains beneficial:
\begin{equation}
    \operatorname{EPE}_{H=4}
    =
    0.50613.
\end{equation}

When the same state is propagated continuously without periodic re-anchoring, performance
collapses:
\begin{align}
    \operatorname{EPE}_{\mathrm{continuous}}
    &= 1.6259,\\
    G_{\mathrm{continuous}}
    &= -1.0922.
\end{align}

In continuous streaming, \(63.5\%\) of the evaluated frames are degraded and the
worst-frame degradation reaches \(22.81\) EPE. The mean geometric correction increases
from approximately \(0.06\) pixels under bounded recurrence to approximately
\(1.50\) pixels under unrestricted recurrence.

This behavior indicates a recurrent feedback process:
\begin{equation}
    \text{state error}
    \rightarrow
    \text{temporal warp}
    \rightarrow
    \text{incorrect fusion}
    \rightarrow
    \text{larger state error}.
\end{equation}

The preceding fused prediction is therefore useful as a short-lived temporal candidate,
but it is unsafe as an indefinitely persistent representation of the complete history.

\subsubsection{Learned Stereo Evidence Is Not the Dominant Factor}

The no-evidence ablation removes:

\begin{itemize}
    \item candidate-conditioned learned left--right matching costs;
    \item learned appearance self-correlation;
    \item learned cross-frame appearance correlation;
    \item frozen ResNet context features.
\end{itemize}

It retains disparity disagreement, flow, support, validity, motion information, dual
reset/fusion supervision, bounded recurrence, and soft disparity fusion.

This reduced configuration obtains
\begin{equation}
    \operatorname{EPE}_{\mathrm{no\text{-}features}}
    =
    0.50348,
\end{equation}
corresponding to a gain of
\begin{equation}
    G_{\mathrm{no\text{-}features}}
    =
    0.03020.
\end{equation}

In the canonical run, removing the learned stereo evidence therefore does not reduce
performance. This result does not prove that learned stereo features are universally
harmful, because the ablation was not replicated across multiple random initializations.
It does establish that they are not necessary for obtaining the current temporal gain.

The evidence consequently suggests that the principal bottleneck is not solely feature
availability. The lifecycle and provenance of the temporal state are at least as important
as the representation used to score it.

\subsubsection{Soft Correction Is Preferable to Hard Replacement}

The hard endpoint variant restricts the output to
\begin{equation}
    d_t^{H}
    \in
    \left\{
        d_t^{S},
        d_t^{M}
    \right\}.
\end{equation}

Although it preserves a positive gain of \(0.02482\), its harmful-update rate increases
to \(5.55\%\) and its clean-pixel degradation increases to \(7.91\%\).

The soft model instead performs a bounded correction:
\begin{equation}
    d_t^{F}
    =
    d_t^{S}
    +
    w_t
    \left(
        d_t^{M}-d_t^{S}
    \right).
\end{equation}

The benefit of soft fusion is not primarily due to predicting values that outperform both
endpoints. The measured aggregate gain from such cases is only approximately
\(0.00072\) EPE. Rather, the continuous weight allows the model to move conservatively
toward a useful temporal candidate without fully replacing the current stereo estimate.

\subsubsection{Fixed Re-Anchoring Generalizes Better than the Adaptive Policy}

A validation-only policy search evaluated fixed horizons
\begin{equation}
    H
    \in
    \left\{
        1,2,4,6,8,12,16
    \right\},
\end{equation}
as well as cumulative-update, evidence-based, and hybrid reset policies.

The hybrid policy with maximum age \(H_{\max}=8\) was preferred on the validation split.
After being frozen and evaluated on dataset 7, however, it produced
\begin{equation}
    \operatorname{EPE}_{\mathrm{hybrid}}
    =
    0.51160,
\end{equation}
which is worse than the fixed \(H=4\) result.

Thus, the available reset cues do not reliably detect state corruption before it affects
the output. Fixed four-frame re-anchoring remains the most defensible operating policy.

\subsubsection{Temporal-Memory Interpretation}

The combined ablations reveal a temporal-memory provenance problem.

A raw disparity prediction
\begin{equation}
    d_{t-k}^{S}
\end{equation}
is an independent observation produced directly by the stereo backbone. A corrected
prediction
\begin{equation}
    d_{t-k}^{F}
\end{equation}
may be more accurate locally, but it also contains the accumulated effect of preceding
warps and fusions.

The two types of memory therefore have different properties:

\begin{align}
    d^{F}
    &: \text{high short-term utility, increasing recurrent risk},\\
    d^{S}
    &: \text{lower correction level, independent long-term anchor}.
\end{align}

This motivates a provenance-typed temporal memory:
\begin{equation}
    \mathcal{M}_t
    =
    \mathcal{M}^{\mathrm{corr}}_t
    \cup
    \mathcal{M}^{\mathrm{raw}}_t
    \cup
    \mathcal{B}^{\mathrm{emb}}_t,
    \label{eq:hybrid_temporal_memory}
\end{equation}
where
\begin{equation}
    \mathcal{M}^{\mathrm{corr}}_t
    =
    \left\{
        d_{t-1}^{F},
        d_{t-2}^{F}
    \right\}
\end{equation}
contains short-lived corrected states,
\begin{equation}
    \mathcal{M}^{\mathrm{raw}}_t
    =
    \left\{
        d_{t-1}^{S},
        d_{t-2}^{S},
        d_{t-4}^{S},
        d_{t-8}^{S}
    \right\}
\end{equation}
contains immutable raw anchors, and
\(\mathcal{B}^{\mathrm{emb}}_t\) stores compact temporal evidence for consensus,
retrieval, flicker detection, and memory-confidence estimation.

Under this formulation, corrected predictions are never promoted to permanent long-term
memory. They expire after a short temporal lifetime, whereas raw observations preserve
access to a longer causal history without accumulating fusion drift.

\subsubsection{Validated Scope}

The reported results establish held-out generalization to dataset ID 7 within the
SCARED-C experimental protocol. They do not yet establish external out-of-distribution
generalization.

The refiner was trained and evaluated using predictions from S2M2-S, RAFT-Stereo, and
StereoAnywhere. Consequently, the experiment demonstrates compatibility across multiple
heterogeneous but seen stereo backbones. It does not establish transfer to a stereo
backbone excluded from refiner training.

Similarly, dataset ID 7 belongs to the same surgical dataset family used for training and
validation. No external surgical dataset, different acquisition system, or unseen
backbone has yet been evaluated.

The current claim is therefore:
\begin{quote}
ARGOS v2 provides a bounded-horizon causal temporal refinement mechanism that generalizes
to held-out SCARED-C sequences and multiple seen stereo backbones. Backbone-agnostic
transfer and external-dataset robustness remain unverified.
\end{quote}

\paragraph{Main finding.}

The principal requirement for stable temporal refinement is not an indefinitely expressive
recurrent state. It is a controlled memory lifecycle: corrected states provide additional
short-term gain, raw predictions provide stable re-anchoring, and soft fusion prevents
aggressive endpoint replacement. This finding motivates a hybrid causal memory in which
short-lived corrected estimates and immutable raw anchors are stored separately.

\subsection{Long-Range Causal History Is Preserved by Immutable Raw Anchors}
\label{sec:raw_anchor_memory}

The preceding bounded-recurrence experiments showed that short-term propagation of
corrected disparity improves temporal refinement, while unrestricted recurrent propagation
causes severe long-horizon drift. This raised a separate question: whether distant causal
history contains useful geometric information that is lost by the current single-state
representation.

We therefore performed a post-hoc multi-anchor oracle audit using direct current-to-anchor
BiDA alignment. No model was trained during this experiment. At frame \(t\), the candidate
bank contained the current stereo estimate
\begin{equation}
    d_t^{S},
\end{equation}
short-term corrected candidates
\begin{equation}
    \mathcal{M}^{F}_t
    =
    \left\{
        \hat d_{t-1}^{F},
        \hat d_{t-2}^{F}
    \right\},
\end{equation}
and immutable raw anchors
\begin{equation}
    \mathcal{M}^{S}_t
    =
    \left\{
        \hat d_{t-1}^{S},
        \hat d_{t-2}^{S},
        \hat d_{t-4}^{S},
        \hat d_{t-8}^{S}
    \right\},
\end{equation}
where every temporal candidate was independently warped into the current frame using
direct current-to-anchor optical flow.

\subsubsection{Multi-Anchor Oracle Ceiling}

Table~\ref{tab:hybrid_memory_oracle} reports the availability-aware oracle results on
dataset 7. The evaluation contains \(28{,}031{,}346\) valid pixels.

\begin{table}[t]
    \centering
    \caption{Post-hoc temporal-memory oracle audit on dataset 7.
    The oracle selects the candidate with the smallest disparity error independently at
    each valid pixel. These results measure information availability and do not represent
    the performance of a deployable model.}
    \label{tab:hybrid_memory_oracle}
    \resizebox{\columnwidth}{!}{
    \begin{tabular}{lcc}
        \toprule
        Candidate bank
        & Oracle EPE
        & Gain versus raw \\
        \midrule
        Current raw
        & \(0.54773\)
        & -- \\

        Current \(+\) CS1
        & \(0.49537\)
        & \(0.05236\) \\

        Current \(+\) CF1/CF2
        & \(0.45401\)
        & \(0.09372\) \\

        Current \(+\) CS1/CS2/CS4/CS8
        & \(0.40821\)
        & \(0.13952\) \\

        Full hybrid
        & \(0.40331\)
        & \(0.14442\) \\

        Full hybrid, convex oracle
        & \(\mathbf{0.38879}\)
        & \(\mathbf{0.15894}\) \\
        \bottomrule
    \end{tabular}}
\end{table}

The immutable raw-anchor bank reduces the oracle EPE from
\begin{equation}
    0.54773
\end{equation}
to
\begin{equation}
    0.40821,
\end{equation}
corresponding to a ceiling gain of
\begin{equation}
    G_{\mathrm{raw\ bank}}
    =
    0.13952.
\end{equation}

Adding the corrected short-term candidates improves the full hybrid ceiling to
\begin{equation}
    G_{\mathrm{hybrid}}
    =
    0.14442.
\end{equation}

Consequently, the raw anchors explain
\begin{equation}
    \frac{
        G_{\mathrm{raw\ bank}}
    }{
        G_{\mathrm{hybrid}}
    }
    =
    96.6\%
\end{equation}
of the full hybrid selection ceiling. The two corrected candidates contribute only
\begin{equation}
    G_{\mathrm{CF1+CF2}\mid\mathrm{raw\ bank}}
    =
    0.00489,
\end{equation}
or approximately \(3.4\%\) of the complete gain.

\subsubsection{Incremental Value of Distant Raw Anchors}

The raw-first incremental audit produced the following gains:

\begin{table}[t]
    \centering
    \caption{Incremental oracle gain when raw anchors are added in increasing temporal
    order.}
    \label{tab:incremental_raw_anchor_gain}
    \begin{tabular}{lc}
        \toprule
        Candidate added & Incremental gain \\
        \midrule
        CS1 & \(0.05236\) \\
        CS2 & \(0.02841\) \\
        CS4 & \(0.02787\) \\
        CS8 & \(0.03088\) \\
        CF1 & \(0.00274\) \\
        CF2 & \(0.00215\) \\
        \bottomrule
    \end{tabular}
\end{table}

The gain does not saturate after the immediately preceding frame. In particular, the
\(t-8\) raw anchor contributes an additional
\begin{equation}
    0.03088
\end{equation}
EPE of oracle gain after the inclusion of CS1, CS2, and CS4.

Although CS8 has a direct candidate EPE of
\begin{equation}
    \operatorname{EPE}_{\mathrm{CS8}}
    =
    0.61433,
\end{equation}
which is worse than the current raw prediction, it provides a selective oracle gain of
\begin{equation}
    G_{\mathrm{CS8}}
    =
    0.09137.
\end{equation}

CS8 outperforms the raw estimate on \(41.91\%\) of its valid pixels and contributes to
\(82.20\%\) of evaluated frames. It is therefore not a reliable global prediction, but a
highly informative local specialist.

\subsubsection{Raw and Corrected Memory Have Asymmetric Complementarity}

The provenance-complementarity audit found that:

\begin{itemize}
    \item a raw anchor is useful while both corrected candidates fail on \(14.34\%\) of
    evaluated pixels;
    \item CS4 or CS8 recovers a failure of short corrected memory on \(9.08\%\);
    \item a corrected candidate is useful while all raw anchors fail on only \(0.85\%\);
    \item both provenance families contain some useful candidate on \(60.24\%\);
    \item the current raw prediction remains optimal on \(24.58\%\).
\end{itemize}

The complementarity is therefore strongly asymmetric. Corrected states provide a small
residual benefit, whereas immutable raw observations carry nearly all of the long-range
temporal signal.

\subsubsection{The Gain Is Not Caused by Support Differences}

Temporal support decreases only moderately with anchor age:

\begin{align}
    \operatorname{coverage}(\mathrm{CS1})
    &= 99.06\%,\\
    \operatorname{coverage}(\mathrm{CS2})
    &= 98.73\%,\\
    \operatorname{coverage}(\mathrm{CS4})
    &= 97.99\%,\\
    \operatorname{coverage}(\mathrm{CS8})
    &= 96.54\%.
\end{align}

Under strict-intersection support, where every compared candidate must be valid, coverage
remains \(96.40\%\). The full hybrid selection oracle still achieves
\begin{equation}
    \operatorname{EPE}_{\mathrm{strict}}
    =
    0.39393
\end{equation}
and a gain of
\begin{equation}
    G_{\mathrm{strict}}
    =
    0.14396.
\end{equation}

The observed long-range signal therefore cannot be explained by incompatible masks or
candidate-specific support advantages.

\subsubsection{Consistency Across Backbones and Sequences}

The additional full-hybrid gain over the best short-term family is positive for every
evaluated backbone:

\begin{align}
    G_{\mathrm{RAFT\text{-}Stereo}}
    &= 0.05847,\\
    G_{\mathrm{S2M2\text{-}S}}
    &= 0.04469,\\
    G_{\mathrm{StereoAnywhere}}
    &= 0.04893.
\end{align}

It is also positive for each dataset-7 sequence, ranging from
\begin{equation}
    0.02666
\end{equation}
to
\begin{equation}
    0.07624.
\end{equation}

Thus, the signal is not confined to a single seen stereo backbone or individual sequence.

\subsubsection{Memory-Failure Detection}

Compact temporal statistics were evaluated as diagnostic predictors of raw error,
temporal-memory failure, and harmful corrected updates. An Isolation Forest trained only
on training datasets 1, 3, and 6 achieved:

\begin{align}
    \operatorname{AUROC}_{\mathrm{memory\ failure}}
    &= 0.812,\\
    \operatorname{AUROC}_{\mathrm{harmful\ corrected}}
    &= 0.828.
\end{align}

These results indicate that candidate failure is partially observable from compact temporal
statistics. However, the anomaly detector remains a validation-only diagnostic and is not
used as a fusion or authorization policy.

\subsubsection{Architectural Implication}

The audit supports a non-recurrent causal memory bank:
\begin{equation}
    \mathcal{M}^{S}_t
    =
    \left\{
        \hat d_{t-1}^{S},
        \hat d_{t-2}^{S},
        \hat d_{t-4}^{S},
        \hat d_{t-8}^{S}
    \right\}.
\end{equation}

A deployable model should retrieve one candidate
\begin{equation}
    k^\star
    =
    \operatorname*{arg\,max}_{k\in\{1,2,4,8\}}
    s_t^{(k)}
\end{equation}
and perform a bounded soft correction:
\begin{equation}
    d_t^{\mathrm{out}}
    =
    \left(1-w_t\right)d_t^{S}
    +
    w_t\hat d_{t-k^\star}^{S}.
\end{equation}

When no temporal candidate is sufficiently reliable, the model must return
\begin{equation}
    d_t^{\mathrm{out}}
    =
    d_t^{S}
\end{equation}
with a bit-exact raw fallback.

Corrected candidates may be added as an explicitly typed memory with maximum lifetime
two frames:
\begin{equation}
    \mathcal{M}^{F}_t
    =
    \left\{
        \hat d_{t-1}^{F},
        \hat d_{t-2}^{F}
    \right\},
\end{equation}
but they should be treated as a controlled ablation rather than as the primary memory
representation.

\paragraph{Main finding.}

Long-range causal history is highly informative for surgical stereo, but its utility is
preserved by independently anchored raw observations rather than by recurrent propagation
of corrected disparity. This motivates a causal multi-anchor retrieval and soft-fusion
adapter, not an unrestricted recurrent or generic embedded memory bank.

\paragraph{Scope.}

The result is an oracle ceiling analysis rather than deployable model performance. It is
validated on held-out SCARED-C sequences and three seen stereo backbones. Transfer to
unseen stereo estimators and external surgical datasets remains unverified.

\subsection{Multi-Anchor Raw Memory Outperforms Bounded Recurrent Fusion}
\label{sec:multi_anchor_realizable}

The oracle analysis established that long-range causal history contains substantial
geometric information, primarily in the form of independently generated raw disparity
anchors. We therefore evaluated whether this signal could be recovered by a deployable
model without recurrently propagating refined disparity.

At frame \(t\), the proposed adapter constructs the immutable raw-anchor bank
\begin{equation}
    \mathcal{M}^{S}_{t}
    =
    \left\{
        \hat d^{S}_{t-1},
        \hat d^{S}_{t-2},
        \hat d^{S}_{t-4},
        \hat d^{S}_{t-8}
    \right\},
\end{equation}
where every past raw disparity is independently aligned to the current frame using
direct current-to-anchor SEA-RAFT flow and a causal BiDA-style pull warp.

A shared candidate-wise module predicts a utility score and a fusion weight for every
valid anchor:
\begin{equation}
    \left(s_t^{(k)}, w_t^{(k)}\right)
    =
    f_{\theta}\left(x_t^{(k)}\right),
    \qquad
    k\in\{1,2,4,8\}.
\end{equation}

The highest-scoring valid candidate is retrieved:
\begin{equation}
    k^\star
    =
    \operatorname*{arg\,max}_{k\in\{1,2,4,8\}}
    s_t^{(k)},
\end{equation}
and softly fused with the current raw stereo prediction:
\begin{equation}
    d_t^{\mathrm{out}}
    =
    \left(1-w_t^{(k^\star)}\right)d_t^{S}
    +
    w_t^{(k^\star)}\hat d^{S}_{t-k^\star}.
\end{equation}

The refined output is never written back into the raw-anchor bank. Consequently, the
method accesses long-range causal history without introducing a recurrent disparity
state.

\subsubsection{Frozen Dataset-7 Results}

Table~\ref{tab:multi_anchor_realizable} compares the learned multi-anchor adapter with
the one-step reference, hard candidate retrieval, and a deterministic temporal-consensus
baseline. All comparisons use the same evaluation support.

\begin{table}[t]
    \centering
    \caption{Frozen dataset-7 results for the realizable raw multi-anchor experiment.
    Gain versus \(H=4\) is measured against the bounded CODD-style recurrent reference
    on common support.}
    \label{tab:multi_anchor_realizable}
    \resizebox{\columnwidth}{!}{
    \begin{tabular}{lccc}
        \toprule
        Configuration
        & EPE
        & Gain vs.\ raw
        & Gain vs.\ \(H=4\) \\
        \midrule
        CS1 soft
        & \(0.53240\)
        & \(0.01552\)
        & \(-0.01515\) \\

        Hard multi-anchor retrieval
        & \(0.51420\)
        & \(0.03372\)
        & \(+0.00306\) \\

        \textbf{Soft multi-anchor fusion}
        & \(\mathbf{0.51047}\)
        & \(\mathbf{0.03745}\)
        & \(\mathbf{+0.00679}\) \\

        Temporal consensus
        & \(0.54100\)
        & \(0.00693\)
        & \(-0.02374\) \\
        \bottomrule
    \end{tabular}}
\end{table}

The soft multi-anchor model improves upon the bounded recurrent \(H=4\) reference by
\begin{equation}
    G_{\mathrm{multi\text{-}anchor}\;vs.\;H4}
    =
    0.00679.
\end{equation}

The improvement is positive for all three evaluated stereo backbones and all four
dataset-7 sequences. Moreover, the learned model recovers
\begin{equation}
    26.8\%
\end{equation}
of the raw-bank selection-oracle gain and
\begin{equation}
    24.2\%
\end{equation}
of the corresponding convex-oracle gain.

The one-step CS1 model does not match the \(H=4\) reference, while both hard and soft
multi-anchor models outperform it. This confirms that the improvement is not explained
by the immediately preceding frame alone. The additional ages
\begin{equation}
    \{t-2,t-4,t-8\}
\end{equation}
provide realizable complementary information.

Soft fusion also outperforms hard retrieval:
\begin{equation}
    0.51047 < 0.51420.
\end{equation}
Thus, selecting an informative temporal anchor is useful, but retaining a bounded convex
correction toward the current raw estimate remains beneficial.

\subsubsection{Selective-Intervention Limitation}

Despite its geometric improvement, the current policy does not yet provide sufficiently
safe intervention behavior. On dataset 7, it produces:

\begin{align}
    \operatorname{HarmfulUpdateRate}
    &= 35.1\%,\\
    \operatorname{CleanDegradation}
    &= 4.86\%,\\
    \operatorname{InterventionPrecision}
    &= 49.1\%,\\
    \operatorname{DegradedFrameFraction}
    &= 26.2\%.
\end{align}

The model therefore identifies enough useful temporal corrections to reduce aggregate
EPE, but approximately half of its accepted interventions are not beneficial. The
remaining limitation is not the availability of temporal information, but the reliability
of the intervention-authorization policy.

\paragraph{Finding.}

Immutable multi-age raw disparity anchors provide a realizable improvement over bounded
recurrent CODD-style fusion. A lightweight retrieval and pairwise soft-fusion adapter
recovers a meaningful fraction of the multi-anchor oracle ceiling and improves every
evaluated backbone and sequence. However, the current candidate acceptance policy remains
insufficiently selective, motivating an explicit risk-aware reject option before deployment.

\paragraph{Scope.}

The result demonstrates multi-backbone compatibility across three stereo estimators seen
during training. It does not yet establish unseen-backbone transfer, external-dataset
generalization, or safety for clinical deployment.

\subsection{Milestone: Transfer to Unseen Stereo Backbones}
\label{sec:h4_unseen_backbone_transfer}

A central objective of ARGOS v2 is to develop a temporal refinement module that
operates independently of the internal architecture of the underlying stereo
matcher. To test this property, we evaluated the frozen bounded CODD-style
\(H=4\) refiner on stereo backbones that were completely excluded from refiner
training.

\paragraph{Training and evaluation protocol.}

The temporal refiner was trained using disparity predictions produced by three
heterogeneous frozen stereo estimators:
\begin{equation}
    \mathcal{B}_{\mathrm{train}}
    =
    \left\{
        \text{S2M2-S},
        \text{RAFT-Stereo},
        \text{StereoAnywhere}
    \right\}.
\end{equation}

After model selection, the canonical temporal checkpoint was frozen. No
backbone-specific fine-tuning, threshold calibration, feature normalization
adaptation, or parameter update was performed during the transfer experiment.

The frozen refiner was then applied to:
\begin{equation}
    \mathcal{B}_{\mathrm{unseen}}
    =
    \left\{
        \text{CREStereo},
        \text{Fast-FoundationStereo}
    \right\},
\end{equation}
neither of which contributed predictions during temporal-refiner training.

For every backbone, the temporal module retains the same bounded correction rule:
\begin{equation}
    \alpha_t
    =
    w_t^{\mathrm{reset}}
    w_t^{\mathrm{fusion}},
\end{equation}
\begin{equation}
    d_t^{F}
    =
    \left(1-\alpha_t\right)d_t^{S}
    +
    \alpha_t d_t^{M},
\end{equation}
where \(d_t^{S}\) is the current raw stereo disparity and \(d_t^{M}\) is the
causally aligned temporal memory. The recurrent state is re-anchored after a
fixed horizon of four frames to prevent indefinite error accumulation.

\paragraph{Canonical reproduction.}

Before opening the unseen-backbone results, the original evaluation path was
reproduced using the canonical checkpoint. The like-for-like reproduction
differed from the validated reference by at most:
\begin{equation}
    \epsilon_{\mathrm{reproduction}}
    =
    2.72\times10^{-6},
\end{equation}
confirming that the transfer experiment reused the unchanged model,
feature-construction path, fusion rule, and reset protocol.

\paragraph{Unseen-backbone results.}

Table~\ref{tab:h4_unseen_backbone} reports the frozen SCARED-C results under
strict common support.

\begin{table}[t]
    \centering
    \caption{Frozen transfer of the bounded CODD-style \(H=4\) refiner to
    stereo backbones excluded from temporal-refiner training. Lower EPE is
    better. Gain is computed relative to the corresponding raw backbone.}
    \label{tab:h4_unseen_backbone}
    \resizebox{\columnwidth}{!}{
    \begin{tabular}{lccc}
        \toprule
        Backbone
        & Raw EPE
        & Refined EPE
        & Gain \\
        \midrule
        CREStereo
        & \(0.57046\)
        & \(0.54823\)
        & \(+0.02224\) \\

        Fast-FoundationStereo
        & \(0.53227\)
        & \(0.51754\)
        & \(+0.01473\) \\
        \bottomrule
    \end{tabular}}
\end{table}

The frozen refiner improves both unseen stereo estimators:
\begin{align}
    G_{\mathrm{CREStereo}}
    &= 0.02224,\\
    G_{\mathrm{FastFoundationStereo}}
    &= 0.01473.
\end{align}

For context, the same checkpoint also improves all three stereo estimators
represented during training:
\begin{align}
    G_{\mathrm{S2M2-S}}
    &= 0.01941,\\
    G_{\mathrm{RAFT-Stereo}}
    &= 0.02928,\\
    G_{\mathrm{StereoAnywhere}}
    &= 0.03397.
\end{align}

The resulting evaluation therefore covers five substantially different stereo
architectures:
\begin{equation}
    \left\{
    \begin{array}{c}
        \text{S2M2-S},\,
        \text{RAFT-Stereo},\,
        \text{StereoAnywhere},\\
        \text{CREStereo},\,
        \text{Fast-FoundationStereo}
    \end{array}
    \right\}.
\end{equation}

\paragraph{Backbone-independent interface.}

The temporal refiner does not require internal backbone-specific representations
such as stereo cost volumes, recurrent hidden states, matching-network features,
or proprietary confidence outputs. Instead, it operates through an external
interface constructed from signals available for heterogeneous stereo matchers:
\begin{equation}
    x_t
    =
    \left[
        d_t^{S},
        d_t^{M},
        d_t^{S}-d_t^{M},
        F_{t\rightarrow t-1},
        C_t^{\mathrm{FB}},
        V_t,
        S_t,
        I_t^{L},
        I_t^{R}
    \right],
\end{equation}
where \(F\) denotes temporal motion, \(C^{\mathrm{FB}}\) forward--backward
confidence, \(V\) validity, and \(S\) warp support.

This interface allows the same temporal checkpoint to operate on predictions
from stereo estimators with substantially different architectures.

\paragraph{Finding.}

The experiment provides direct evidence that the bounded CODD-style ARGOS v2
refiner is backbone-agnostic across the evaluated stereo estimators within the
SCARED-C domain. A single frozen checkpoint improves three training-seen
backbones and two stereo backbones excluded from training, without
backbone-specific adaptation.

More precisely, the result demonstrates:
\begin{equation}
    \boxed{
    \text{same-domain transfer to unseen stereo backbones}
    }.
\end{equation}

It does not yet demonstrate:
\begin{equation}
    \boxed{
    \text{joint unseen-backbone and external-domain transfer}
    }.
\end{equation}

The term \emph{backbone-agnostic} is therefore used to describe the architecture
and its verified transfer across the five evaluated stereo estimators, rather
than as a universal claim covering every possible matcher and acquisition
domain.

\paragraph{Research significance.}

This result establishes the bounded CODD-style model as a robust architectural
foundation for ARGOS v2. Subsequent work can focus on improving correction
selection, temporal-memory utilization, and intervention safety without
requiring a redesign tied to a specific stereo backbone.

The immediate research path is:
\begin{equation}
    \boxed{
    \begin{array}{c}
    \text{validated backbone-independent temporal fusion}\\
    \Downarrow\\
    \text{safety-aware intervention selection}\\
    \Downarrow\\
    \text{stronger training and multi-anchor refinement}
    \end{array}}
\end{equation}

If the current spatial safety selector successfully identifies harmful temporal
updates, it can be integrated as an abstention layer over the multi-anchor
architecture. If it does not, the bounded CODD-style \(H=4\) system remains a
validated and transferable base architecture that can be further improved
through larger-scale training, stronger safety supervision, and controlled
architectural refinement.

\paragraph{Scope.}

The current evidence supports backbone-independent transfer within SCARED-C.
External-domain geometric generalization remains unresolved because the
available D4D structured-light reference does not currently satisfy the required
stereo-disparity compatibility contract. SERV-CT does not contain temporally
consecutive frames and therefore cannot support an \(H=4\) temporal evaluation.
No claim of universal cross-domain or clinical robustness is made.

\subsection{Conditional Milestone: Spatially Selective Multi-Anchor Refinement}
\label{sec:spatial_safety_critic}

The immutable multi-anchor temporal refiner demonstrated a realizable geometric
advantage over bounded recurrent fusion, but its ungated intervention policy
introduced an excessive number of harmful updates. We therefore investigated
whether a lightweight spatial comparative critic could identify when the
multi-anchor proposal should replace the current raw stereo estimate.

\paragraph{Frozen proposal and selective intervention.}

Let \(d_t^{S}\) denote the current raw stereo disparity and
\(d_t^{P}\) the frozen multi-anchor proposal. The critic does not modify the
proposal, the selected anchor, or the fusion weight. It only predicts whether
the proposal should be accepted:
\begin{equation}
    d_t^{\mathrm{out}}
    =
    \begin{cases}
        d_t^{P}, & a_t = 1,\\
        d_t^{S}, & a_t = 0.
    \end{cases}
\end{equation}

The fallback is implemented through exact tensor selection, ensuring that
rejected interventions reproduce the raw disparity bit-exactly.

The critic predicts separate heteroscedastic error estimates for the raw and
proposed disparities:
\begin{equation}
    \left(\mu_t^{S}, \sigma_t^{S}\right),
    \qquad
    \left(\mu_t^{P}, \sigma_t^{P}\right),
\end{equation}
together with a direct harmful-update probability \(p_t^{\mathrm{harm}}\).

The predicted proposal gain is:
\begin{equation}
    \widehat{\Delta}_t
    =
    \mu_t^{S}-\mu_t^{P},
\end{equation}
with comparative uncertainty:
\begin{equation}
    \sigma_t^{\Delta}
    =
    \sqrt{
        \left(\sigma_t^{S}\right)^2+
        \left(\sigma_t^{P}\right)^2
    }.
\end{equation}

A conservative lower confidence bound is defined as:
\begin{equation}
    \mathrm{LCB}_{t}^{\Delta}
    =
    \widehat{\Delta}_t
    -
    \lambda_{\mathrm{unc}}
    \sigma_t^{\Delta}.
\end{equation}

The frozen intervention rule is:
\begin{equation}
    a_t
    =
    \mathbbm{1}
    \left[
        \mathrm{LCB}_{t}^{\Delta}>\tau_{\Delta}
        \land
        p_t^{\mathrm{harm}}<\tau_{\mathrm{harm}}
        \land
        V_t=1
    \right],
\end{equation}
where \(V_t\) denotes proposal validity.

\paragraph{Spatial comparative critic.}

The selected critic is a lightweight fully convolutional network using
Group Normalization, SiLU activations, and dilated convolutions with dilation
pattern:
\begin{equation}
    [1,2,4,2,1].
\end{equation}

The resulting receptive field is approximately \(43\) pixels and the selected
stereo-evidence configuration contains approximately \(0.40\) million trainable
parameters.

The critic receives spatially structured evidence describing:

\begin{itemize}
    \item current raw, selected-anchor, and proposed disparities;
    \item raw--anchor and raw--proposal disagreement;
    \item proposal fusion weight and validity;
    \item temporal motion and forward--backward consistency;
    \item local disparity gradients and spatial consensus;
    \item temporal photometric residuals;
    \item stereo photometric residuals for both the raw and proposed hypotheses.
\end{itemize}

The multi-anchor proposal generator remains completely frozen throughout critic
training.

\paragraph{Protocol.}

Training, validation, and testing use dataset-disjoint SCARED-C splits:
\begin{equation}
    \mathcal{D}_{\mathrm{train}}=\{1,3,6\},
    \qquad
    \mathcal{D}_{\mathrm{val}}=\{2\},
    \qquad
    \mathcal{D}_{\mathrm{test}}=\{7\}.
\end{equation}

The critic checkpoint and intervention thresholds were selected exclusively on
dataset \(2\). The complete checkpoint, feature schema, thresholds, and policy
were then frozen and hashed before the single evaluation on dataset \(7\).

No test-time calibration, threshold adjustment, or model selection was
performed.

\paragraph{Frozen test results.}

Table~\ref{tab:spatial_critic_main} compares the principal methods on the
held-out dataset-7 test split.

\begin{table}[t]
    \centering
    \caption{Frozen dataset-7 evaluation of the spatial safety critic.
    Lower EPE is better.}
    \label{tab:spatial_critic_main}
    \resizebox{\columnwidth}{!}{
    \begin{tabular}{lcc}
        \toprule
        Method & EPE & Gain over raw \\
        \midrule
        Raw stereo
        & \(0.5479\)
        & \(0.0000\) \\

        Bounded CODD-style \(H=4\)
        & \(0.5173\)
        & \(0.0306\) \\

        Ungated immutable multi-anchor
        & \(0.5105\)
        & \(0.0374\) \\

        Multi-anchor + spatial critic
        & \(\mathbf{0.5072}\)
        & \(\mathbf{0.0407}\) \\
        \bottomrule
    \end{tabular}}
\end{table}

The frozen critic improves over the bounded CODD-style baseline by:
\begin{equation}
    0.5173-0.5072
    =
    0.0101\ \mathrm{px}.
\end{equation}

It also improves over the ungated multi-anchor proposal:
\begin{equation}
    0.5105-0.5072
    =
    0.0033\ \mathrm{px}.
\end{equation}

Consequently, the selective policy retains more than the complete geometric
gain of the ungated proposal:
\begin{equation}
    \mathrm{GainRetained}
    =
    108.8\%.
\end{equation}

This indicates that rejecting harmful proposals increases the net geometric
benefit rather than merely trading accuracy for safety.

\paragraph{Safety behavior.}

Table~\ref{tab:spatial_critic_safety} compares the ungated and selectively
gated multi-anchor variants.

\begin{table}[t]
    \centering
    \caption{Safety and intervention behavior on frozen dataset 7.}
    \label{tab:spatial_critic_safety}
    \resizebox{\columnwidth}{!}{
    \begin{tabular}{lcc}
        \toprule
        Metric
        & Ungated
        & Spatial critic \\
        \midrule
        Harmful interventions
        & \(35.1\%\)
        & \(24.7\%\) \\

        Intervention precision
        & \(49\%\)
        & \(62\%\) \\

        Clean-pixel degradation
        & \(4.9\%\)
        & \(2.1\%\) \\

        Degraded frames
        & \(26.0\%\)
        & \(16.7\%\) \\

        Intervention coverage
        & \(100\%\)
        & \(7.9\%\) \\
        \bottomrule
    \end{tabular}}
\end{table}

The critic satisfies the predefined clean-pixel and degraded-frame constraints:
\begin{equation}
    \mathrm{CleanDegradation}=2.1\% < 3\%,
\end{equation}
\begin{equation}
    \mathrm{DegradedFrames}=16.7\% < 25\%.
\end{equation}

However, the accepted harmful-update rate remains:
\begin{equation}
    \mathrm{HarmRate}=24.7\%,
\end{equation}
which exceeds the predefined \(10\%\) target.

The policy therefore improves safety substantially but does not yet provide
fully risk-controlled intervention.

\paragraph{Validation-to-test transfer.}

A central limitation of the previous scalar gate was the collapse of its
validation-selected operating point on the test split. The spatial critic does
not exhibit the same failure.

The harmful-update rate changes only slightly from validation to test:
\begin{equation}
    26\% \rightarrow 24.7\%,
\end{equation}
while intervention coverage remains non-trivial:
\begin{equation}
    \mathrm{Coverage}_{\mathrm{test}}=7.9\%.
\end{equation}

The critic improves all three evaluated stereo backbones and all four held-out
dataset-7 sequences. Approximately \(51\%\) of accepted interventions originate
from the distant \(CS4\) and \(CS8\) raw anchors, demonstrating that the policy
does not reduce the multi-anchor model to an adjacent-frame refiner.

The harm-discrimination AUROC improves from approximately \(0.57\) for the
previous scalar gate to approximately:
\begin{equation}
    \mathrm{AUROC}_{\mathrm{harm}}\approx0.69
\end{equation}
on the frozen test split.

\paragraph{Finding.}

The experiment establishes that spatial comparative evidence provides a
substantially stronger intervention signal than pointwise scalar statistics.
The frozen critic:

\begin{itemize}
    \item transfers from dataset 2 to dataset 7 without coverage collapse;
    \item improves the raw prediction and the bounded CODD-style baseline;
    \item improves over the ungated multi-anchor proposal;
    \item reduces harmful updates and clean-region degradation;
    \item preserves useful long-range \(CS4/CS8\) interventions;
    \item improves every evaluated backbone and sequence.
\end{itemize}

Nevertheless, the critic does not meet the primary accepted-harm target.
The resulting verdict is therefore:
\begin{equation}
    \boxed{
        \text{Geometry GO}
        \quad+\quad
        \text{Safety CONDITIONAL}
        \quad=\quad
        \text{Overall CONDITIONAL GO}
    }.
\end{equation}

\paragraph{Scientific interpretation.}

The remaining limitation is no longer the availability of useful temporal
geometry, validation-to-test transfer, or coverage collapse. Instead, the
dominant bottleneck is the fine-grained discrimination and calibration of
harmful interventions.

The current result supports the following architectural progression:
\begin{equation}
    \boxed{
    \begin{array}{c}
        \text{immutable raw multi-anchor retrieval}\\
        \Downarrow\\
        \text{pairwise soft temporal fusion}\\
        \Downarrow\\
        \text{spatial comparative error prediction}\\
        \Downarrow\\
        \text{exact raw abstention}
    \end{array}}
\end{equation}

The post-hoc veto formulation improves both accuracy and safety, but a fully
risk-controlled system will likely require the predicted intervention utility
to influence anchor retrieval and fusion magnitude during training rather than
only accepting or rejecting an already constructed proposal.

\paragraph{Scope.}

These results demonstrate held-out same-domain generalization across the three
stereo backbones represented during critic training. They do not establish
external-dataset robustness, unseen-backbone transfer of the safety critic,
clinical safety, or formal statistical coverage guarantees.


% Required packages:
% \usepackage{amsmath,amssymb,booktabs}
% \usepackage{algorithm}
% \usepackage{algpseudocode}
% \usepackage{graphicx}
% \usepackage{xcolor}

\section{Immutable Raw Multi-Anchor Temporal Refinement}
\label{sec:method}

\subsection{Overview}

We consider causal temporal refinement of a frozen stereo estimator.
At time \(t\), a stereo backbone \(S\) predicts a raw disparity map

\begin{equation}
d_t^{\mathrm{raw}}
=
S\left(I_t^L,I_t^R\right),
\end{equation}

where \(I_t^L\) and \(I_t^R\) are the current rectified stereo images.
The temporal module is required to operate without future frames and
without accessing backbone-specific cost volumes, hidden states, or
internal feature representations.

Our final geometric architecture is composed of five stages:

\begin{equation}
\boxed{
\begin{aligned}
&\text{frozen stereo prediction}
\\[-1mm]
&\rightarrow
\text{direct causal motion alignment}
\\[-1mm]
&\rightarrow
\text{immutable multi-age raw-anchor bank}
\\[-1mm]
&\rightarrow
\text{learned candidate retrieval}
\\[-1mm]
&\rightarrow
\text{pairwise soft fusion with exact raw fallback}.
\end{aligned}
}
\end{equation}

The central design choice is that long-term memory contains only
independently generated raw stereo predictions.
Temporally fused predictions are never written back into the long-term
bank.

% ---------------------------------------------------------------------
\subsection{Motivation: Bounded Corrected Memory}
\label{sec:bounded_memory}

We first investigated a CODD-style recurrent formulation
\cite{li2023codd}.
A previous disparity estimate is aligned to the current frame and softly
fused with the current stereo prediction:

\begin{equation}
d_t^{F}
=
\left(
1-w_t^{\mathrm{reset}}w_t^{\mathrm{fusion}}
\right)d_t^{\mathrm{raw}}
+
w_t^{\mathrm{reset}}w_t^{\mathrm{fusion}}
d_t^{\mathrm{mem}}.
\label{eq:codd_bounded}
\end{equation}

This formulation is effective when the corrected state is maintained
only over a short horizon.
In our experiments, a four-frame memory horizon, denoted \(H=4\),
provides stable improvements.
In contrast, indefinite recurrence causes prediction errors to be
repeatedly aligned, fused, and propagated, leading to severe drift.

This result motivates a separation between:

\begin{itemize}
    \item \textbf{short-lived corrected memory}, which can be useful
    locally but must be bounded; and
    \item \textbf{long-lived raw memory}, which should preserve
    independent observations without recursive contamination.
\end{itemize}

The proposed architecture therefore removes long-term corrected-state
recurrence entirely.

% ---------------------------------------------------------------------
\subsection{Immutable Raw Anchor Bank}
\label{sec:anchor_bank}

At time \(t\), we maintain a causal bank of raw disparity predictions
generated at different temporal ages:

\begin{equation}
\mathcal{B}_t
=
\left\{
d_{t-1}^{\mathrm{raw}},
d_{t-2}^{\mathrm{raw}},
d_{t-4}^{\mathrm{raw}},
d_{t-8}^{\mathrm{raw}}
\right\}.
\label{eq:raw_bank}
\end{equation}

We denote these candidates by

\begin{equation}
\mathrm{CS1},\quad
\mathrm{CS2},\quad
\mathrm{CS4},\quad
\mathrm{CS8}.
\end{equation}

Each anchor satisfies the following properties:

\begin{enumerate}
    \item it is generated independently by the frozen stereo backbone;
    \item it is stored without modification;
    \item it is never replaced by a fused prediction;
    \item it is aligned directly to the current frame;
    \item it has explicit temporal age and validity provenance.
\end{enumerate}

The bank update is therefore

\begin{equation}
\mathcal{B}_{t+1}
=
\operatorname{Update}
\left(
\mathcal{B}_t,
d_t^{\mathrm{raw}}
\right),
\label{eq:bank_update}
\end{equation}

and not

\begin{equation}
\mathcal{B}_{t+1}
\neq
\operatorname{Update}
\left(
\mathcal{B}_t,
d_t^{\mathrm{out}}
\right).
\end{equation}

This prevents recursively accumulated temporal errors from contaminating
future long-range memory.

% ---------------------------------------------------------------------
\subsection{Direct Causal Alignment}
\label{sec:direct_alignment}

For every available anchor of age \(k\in\{1,2,4,8\}\), a frozen
SEA-RAFT model \cite{wang2024searaft} estimates target-to-source optical
flow:

\begin{equation}
F_{t\rightarrow t-k}
=
M\left(I_t^L,I_{t-k}^L\right).
\end{equation}

The historical disparity is brought into the current coordinate system
using BiDA-style pull warping \cite{bidastereo2024}:

\begin{equation}
\widetilde d_t^{(k)}
=
\mathcal{W}
\left(
d_{t-k}^{\mathrm{raw}},
F_{t\rightarrow t-k}
\right).
\label{eq:anchor_warp}
\end{equation}

Every anchor is aligned independently from the current frame to its
original source frame.
Flow fields are not recursively composed across time.

For each aligned anchor, we construct a validity mask

\begin{equation}
V_t^{(k)}
=
V_{t,k}^{\mathrm{bounds}}
\land
V_{t,k}^{\mathrm{flow}}
\land
V_{t,k}^{\mathrm{FB}}
\land
V_{t-k}^{\mathrm{disp}},
\label{eq:validity}
\end{equation}

where:

\begin{itemize}
    \item \(V^{\mathrm{bounds}}\) indicates that the sampling location
    lies inside the source image;
    \item \(V^{\mathrm{flow}}\) indicates finite and valid optical flow;
    \item \(V^{\mathrm{FB}}\) represents forward--backward flow
    consistency;
    \item \(V^{\mathrm{disp}}\) represents valid source disparity.
\end{itemize}

% ---------------------------------------------------------------------
\subsection{Shared Candidate Representation}
\label{sec:candidate_representation}

Each aligned anchor is described by a shared candidate representation

\begin{equation}
x_t^{(k)}
=
\Phi
\left(
d_t^{\mathrm{raw}},
\widetilde d_t^{(k)},
I_t^L,
I_{t-k}^L,
F_{t\rightarrow t-k},
V_t^{(k)},
k
\right).
\end{equation}

The validated implementation uses 17 universal channels per candidate.
These channels contain only externally available evidence, including:

\begin{itemize}
    \item current raw disparity;
    \item aligned anchor disparity;
    \item signed raw--anchor disparity difference;
    \item absolute raw--anchor disparity difference;
    \item local disparity gradients;
    \item warp support and validity;
    \item forward--backward flow consistency;
    \item optical-flow magnitude;
    \item temporal photometric agreement;
    \item local robust statistics;
    \item anchor age and provenance.
\end{itemize}

The same encoder is applied to every candidate:

\begin{equation}
z_t^{(k)}
=
E_\theta
\left(
x_t^{(k)}
\right).
\end{equation}

No backbone identifier, backbone-specific feature map, cost volume, or
hidden state is provided to the model.

% ---------------------------------------------------------------------
\subsection{Candidate Retrieval}
\label{sec:retrieval}

The shared encoder predicts a candidate score

\begin{equation}
s_t^{(k)}
=
G_\theta
\left(
z_t^{(k)}
\right).
\end{equation}

Unavailable or invalid anchors are assigned

\begin{equation}
s_t^{(k)}
=
-\infty.
\end{equation}

The selected anchor age is

\begin{equation}
k^\star
=
\arg\max_{k\in\{1,2,4,8\}}
s_t^{(k)}.
\label{eq:retrieval}
\end{equation}

Only one anchor is retrieved for each output location.
The model does not perform an unconditional average over the complete
memory bank.

This distinction is important because distant anchors may be globally
less accurate than the current prediction while remaining locally useful
in regions where the current stereo backbone fails.

% ---------------------------------------------------------------------
\subsection{Pairwise Soft Fusion}
\label{sec:soft_fusion}

After retrieval, a lightweight fusion head compares the current raw
prediction and the selected aligned anchor:

\begin{equation}
w_t
=
\sigma
\left[
H_\theta
\left(
d_t^{\mathrm{raw}},
\widetilde d_t^{(k^\star)},
z_t^{(k^\star)}
\right)
\right],
\end{equation}

where \(w_t\in[0,1]\) is a spatially varying fusion weight.

The proposed disparity is

\begin{equation}
d_t^{\mathrm{prop}}
=
\left(1-w_t\right)d_t^{\mathrm{raw}}
+
w_t\widetilde d_t^{(k^\star)}.
\label{eq:soft_fusion}
\end{equation}

The retrieval and fusion network contains approximately

\begin{equation}
60\,739
\end{equation}

trainable parameters.

The stereo backbone, SEA-RAFT model, and all external feature
extractors remain frozen.

% ---------------------------------------------------------------------
\subsection{Exact Raw Fallback}
\label{sec:fallback}

If no valid temporal candidate is available, or if the frozen retrieval
policy rejects the candidate, the method returns the current raw
prediction:

\begin{equation}
d_t^{\mathrm{out}}
=
\begin{cases}
d_t^{\mathrm{prop}},
&
a_t=1,
\\[1mm]
d_t^{\mathrm{raw}},
&
a_t=0,
\end{cases}
\label{eq:final_output}
\end{equation}

where \(a_t\) denotes the frozen candidate-acceptance decision.

The fallback is implemented using exact tensor selection, rather than an
arithmetic blend with a zero weight.
Consequently, rejected output pixels are bit-identical to the original
raw stereo prediction.

\textbf{TODO:}
Insert the exact frozen probability and utility thresholds from the
validated experiment manifest.

% ---------------------------------------------------------------------
\subsection{Causal Inference Procedure}

\begin{algorithm}[t]
\caption{Immutable raw multi-anchor temporal refinement}
\label{alg:multi_anchor}
\begin{algorithmic}[1]
\Require Current stereo pair
\((I_t^L,I_t^R)\), raw anchor bank \(\mathcal{B}_t\)
\State \(d_t^{\mathrm{raw}}
       \gets S(I_t^L,I_t^R)\)
\For{\(k\in\{1,2,4,8\}\)}
    \If{\(d_{t-k}^{\mathrm{raw}}\) is available}
        \State
        \(F_{t\rightarrow t-k}
        \gets M(I_t^L,I_{t-k}^L)\)
        \State
        \(\widetilde d_t^{(k)}
        \gets
        \mathcal{W}
        (d_{t-k}^{\mathrm{raw}},
        F_{t\rightarrow t-k})\)
        \State Construct \(V_t^{(k)}\)
        \State Construct candidate representation \(x_t^{(k)}\)
        \State
        \(s_t^{(k)}
        \gets
        G_\theta(E_\theta(x_t^{(k)}))\)
    \Else
        \State \(s_t^{(k)}\gets-\infty\)
    \EndIf
\EndFor
\If{no valid candidate exists}
    \State \(d_t^{\mathrm{out}}\gets d_t^{\mathrm{raw}}\)
\Else
    \State
    \(k^\star\gets\arg\max_k s_t^{(k)}\)
    \State Predict pairwise fusion weight \(w_t\)
    \State Construct \(d_t^{\mathrm{prop}}\) using
    Eq.~\eqref{eq:soft_fusion}
    \State Apply frozen acceptance policy
    \State Select \(d_t^{\mathrm{out}}\) using
    Eq.~\eqref{eq:final_output}
\EndIf
\State Store \(d_t^{\mathrm{raw}}\) in \(\mathcal{B}_{t+1}\)
\State Never store \(d_t^{\mathrm{out}}\) in long-term memory
\State \Return \(d_t^{\mathrm{out}}\)
\end{algorithmic}
\end{algorithm}

% ---------------------------------------------------------------------
\subsection{Backbone-Independent Interface}
\label{sec:backbone_independent}

The complete temporal adapter consumes only:

\begin{itemize}
    \item rectified RGB images;
    \item raw disparity predictions;
    \item optical flow;
    \item warp support;
    \item validity masks;
    \item forward--backward consistency;
    \item externally computed local evidence.
\end{itemize}

It does not consume:

\begin{itemize}
    \item stereo cost volumes;
    \item backbone confidence heads;
    \item recurrent hidden states;
    \item backbone-specific features;
    \item backbone identifiers.
\end{itemize}

The same frozen retrieval and fusion checkpoint can therefore be applied
to heterogeneous stereo estimators without architectural modification.

% =====================================================================
\section{Transfer to Stereo Backbones Excluded from Training}
\label{sec:unseen_backbone_transfer}

We evaluate the frozen immutable raw multi-anchor refiner on two stereo
estimators excluded from temporal-refiner training:

\begin{itemize}
    \item CREStereo;
    \item Fast-FoundationStereo.
\end{itemize}

The evaluation is performed within SCARED-C using dataset 7, which is
held out from the frozen training and calibration protocol.
No model parameter, threshold, or preprocessing configuration is tuned
on these two backbones.

The canonical frozen multi-anchor checkpoint is

\begin{verbatim}
results/raw_multi_anchor_temporal_refiner/
soft_fusion/checkpoints/best_validation.pt
\end{verbatim}

with SHA-256

\begin{verbatim}
40526a32ef6e9a62a3ea2b59e6751a60
c441b8190f9b96522e3b12b35895d5cd
\end{verbatim}

The frozen SEA-RAFT checkpoint has SHA-256

\begin{verbatim}
1a21575ed6ca2c6945fb8e25c4169d24
1cf59ee5d12b8802c01c965206268cac
\end{verbatim}

\begin{table}[t]
\centering
\caption{
Same-domain transfer to stereo estimators excluded from temporal-refiner
training. Results use strict common support across the raw prediction,
bounded \(H=4\) refiner, and immutable multi-anchor refiner.
Lower EPE is better.
}
\label{tab:unseen_multi_anchor}
\resizebox{\columnwidth}{!}{
\begin{tabular}{lccccc}
\toprule
Backbone
& Raw
& \(H=4\)
& Multi-anchor
& Gain vs.\ raw
& Gain vs.\ \(H=4\)
\\
\midrule
CREStereo
& 0.5784
& 0.5555
& \textbf{0.5437}
& \textbf{+0.0347}
& \textbf{+0.0119}
\\
Fast-FoundationStereo
& 0.5388
& 0.5237
& \textbf{0.5209}
& \textbf{+0.0179}
& \textbf{+0.0028}
\\
\bottomrule
\end{tabular}
}
\end{table}

The frozen multi-anchor model improves both unseen stereo estimators
relative to their raw predictions and relative to bounded recurrent
fusion.

The improvement is consistent across sequences:
all four evaluated sequences improve for CREStereo and all four improve
for Fast-FoundationStereo.

Long-range memory remains active after transfer.
The combined fraction of selected \(\mathrm{CS4}\) and \(\mathrm{CS8}\)
anchors is

\begin{equation}
53.4\%
\end{equation}

for CREStereo and

\begin{equation}
55.8\%
\end{equation}

for Fast-FoundationStereo.
The model therefore does not collapse to an adjacent-frame-only
solution when applied to stereo estimators excluded from training.

% ---------------------------------------------------------------------
\subsection{Evaluation Support}

The comparison uses strict pixelwise common support across:

\begin{itemize}
    \item the raw stereo prediction;
    \item the bounded \(H=4\) prediction;
    \item all required multi-anchor candidates;
    \item ground-truth disparity.
\end{itemize}

The resulting common pixel support is approximately

\begin{equation}
97.31\%.
\end{equation}

Frames with no pixels in the complete strict-support intersection are
recorded explicitly and excluded from metric aggregation.
A total of \(1\,226\) such frame entries are reported in the evaluation
artifacts.

This exclusion is caused by an empty intersection among the required
temporal and \(H=4\) support masks, rather than by post-hoc error
filtering.

% ---------------------------------------------------------------------
\subsection{Supported Claim}

The experiment supports the following claim:

\begin{quote}
The frozen immutable raw multi-anchor refiner transfers to stereo
estimators excluded from its training within SCARED-C.
The same retrieval and soft-fusion module improves both raw stereo
predictions and bounded recurrent refinement on every evaluated
sequence, while retaining substantial use of long-range temporal
anchors.
\end{quote}

The result supports a backbone-independent input interface with
demonstrated same-domain transfer across the evaluated stereo
estimators.

It does not establish:

\begin{itemize}
    \item universal backbone agnosticism;
    \item external-domain generalization;
    \item clinical safety;
    \item risk-controlled intervention;
    \item real-time operation on deployment hardware.
\end{itemize}

% =====================================================================
\section{Selective Intervention as an Open Problem}
\label{sec:open_safety}

The immutable multi-anchor model provides the strongest geometric
performance among the evaluated temporal architectures.
However, an ungated temporal correction can still degrade locally
accurate raw predictions.

We investigated a post-hoc spatial comparative critic that estimates
whether a frozen multi-anchor proposal should be accepted.
The critic improves average disparity accuracy and reduces several
degradation metrics, but the strict validation target

\begin{equation}
\mathrm{harmful\ accepted\ updates}
\leq 10\%
\end{equation}

at

\begin{equation}
\mathrm{coverage}
\geq 5\%
\end{equation}

was not achieved.

Hard-negative training, calibration, and an ensemble reduced the
accepted harmful-update rate, but the resulting validation operating
point remained approximately

\begin{equation}
\mathrm{harm}=20.1\%,
\qquad
\mathrm{coverage}=2.06\%.
\end{equation}

We therefore treat strict risk-controlled authorization as an open
problem rather than as a solved component of the proposed geometric
method.

A natural extension is to jointly optimize

\begin{equation}
\boxed{
\text{anchor retrieval}
+
\text{fusion magnitude}
+
\text{harm prediction}
+
\text{abstention}
}
\end{equation}

instead of applying a binary veto to an already constructed proposal.

% =====================================================================
\section{Architecture Summary}
\label{sec:architecture_summary}

The two principal ARGOS v2 temporal designs can be summarized as:

\paragraph{Bounded corrected memory.}

\begin{equation}
\boxed{
\text{frozen stereo}
\rightarrow
\text{SEA-RAFT}
\rightarrow
\text{BiDA alignment}
\rightarrow
\text{CODD-style fusion}
\rightarrow
\text{reset every four frames}
}
\end{equation}

This design is robust and provides the controlled baseline demonstrating
that corrected state is useful only over bounded horizons.

\paragraph{Immutable raw multi-anchor memory.}

\begin{equation}
\boxed{
\text{frozen stereo}
\rightarrow
\text{directly aligned CS1/CS2/CS4/CS8}
\rightarrow
\text{shared retrieval}
\rightarrow
\text{pairwise soft fusion}
\rightarrow
\text{raw-only memory write}
}
\end{equation}

This design provides the strongest geometric performance, exploits
long-range temporal evidence, and transfers to stereo estimators
excluded from training within the tested domain.

Together, the two architectures support the following design principle:

\begin{quote}
Corrected disparity memory is effective when short-lived and bounded,
whereas long causal history is better preserved as independently
generated immutable raw observations.
\end{quote}

% ================================================================
% Required packages
% ================================================================
% \usepackage{amsmath,amssymb,bm}
% \usepackage{booktabs}
% \usepackage{graphicx}
% \usepackage{xcolor}

% \providecommand, not \newcommand: the preamble already carries \warp from ICRA.tex, and
% a mid-document \newcommand for a name that exists is a fatal error rather than a clash.
\providecommand{\clip}{\operatorname{clip}}
\providecommand{\warp}{\mathcal{W}}
\providecommand{\ind}{\mathbb{I}}

% ================================================================
\section{Immutable Raw Multi-Anchor Temporal Refinement}
\label{sec:argos_method}

\subsection{Problem Formulation}
\label{sec:problem_formulation}

Let
%
\begin{equation}
\left(I_t^L,I_t^R\right)
\end{equation}
%
denote the current rectified stereo pair at time \(t\).
An external frozen stereo estimator \(S\) produces the positive-left
disparity map and its validity mask:
%
\begin{equation}
d_t^S,\,M_t
=
S\left(I_t^L,I_t^R\right),
\label{eq:frozen_stereo}
\end{equation}
%
where
%
\begin{equation}
d_t^S\in\mathbb{R}^{B\times1\times H\times W},
\qquad
M_t\in\{0,1\}^{B\times1\times H\times W}.
\end{equation}

ARGOS v2 operates exclusively on the externally available RGB images,
raw disparity predictions, optical flow, and geometric validity
information. It does not consume backbone identifiers, stereo cost
volumes, confidence heads, hidden states, or internal backbone features.

The temporal adapter implements the following causal pipeline:
%
\begin{equation}
\boxed{
\begin{aligned}
&\text{frozen stereo prediction}
\\[-1mm]
&\rightarrow
\text{immutable multi-age raw memory}
\\[-1mm]
&\rightarrow
\text{direct current-to-anchor motion estimation}
\\[-1mm]
&\rightarrow
\text{causal pull alignment}
\\[-1mm]
&\rightarrow
\text{shared utility-aware candidate encoding}
\\[-1mm]
&\rightarrow
\text{pixelwise single-anchor retrieval}
\\[-1mm]
&\rightarrow
\text{pairwise soft correction}
\\[-1mm]
&\rightarrow
\text{geometric acceptance and exact raw fallback}.
\end{aligned}
}
\label{eq:overall_pipeline}
\end{equation}

The right image \(I_t^R\) is used by the external stereo estimator.
Once \(d_t^S\) has been produced, the temporal adapter operates on the
current left image \(I_t^L\), the raw disparity, and the causal memory
bank.

% ================================================================
\subsection{Immutable Raw Temporal Memory}
\label{sec:raw_memory}

Before processing frame \(t\), the temporal memory contains only
independently generated raw stereo records from previous frames:
%
\begin{equation}
\mathcal{M}_t
=
\left\{
R_j:
\max(0,t-8)\leq j\leq t-1
\right\},
\label{eq:memory_before_inference}
\end{equation}
%
where each record is
%
\begin{equation}
R_j
=
\left(
d_j^S,
M_j,
I_j^L,
\operatorname{id}_j,
j,
\tau_j
\right).
\label{eq:memory_record}
\end{equation}

The stored disparity, validity mask, and RGB image are detached and
cloned before insertion. Frame indices must be strictly increasing, and
valid disparity values must be finite and positive. Each record carries
the provenance label
%
\begin{equation}
\texttt{independent\_frozen\_stereo}.
\end{equation}

The canonical anchor-age set is
%
\begin{equation}
\mathcal{A}
=
\{1,2,4,8\}.
\label{eq:anchor_ages}
\end{equation}

For an age \(k\in\mathcal{A}\), the corresponding raw anchor is
%
\begin{equation}
d_{t-k}^S,
\end{equation}
%
which we denote as \(\mathrm{CS}k\). Thus, \(\mathrm{CS1}\),
\(\mathrm{CS2}\), \(\mathrm{CS4}\), and \(\mathrm{CS8}\) are raw
observations produced independently by the frozen stereo estimator.
They are not recursively fused states.

If the source frame \(t-k\) is unavailable, the corresponding candidate
is represented by a zero-supported dummy tensor and is excluded from
retrieval through explicit masking.

After the current output has been constructed, the bank is updated as
%
\begin{equation}
\mathcal{M}_{t+1}
=
\operatorname{EvictAndAppendRaw}
\left(
\mathcal{M}_t,
R_t
\right),
\label{eq:raw_memory_update}
\end{equation}
%
with
%
\begin{equation}
R_t
=
\left(
d_t^S,
M_t,
I_t^L,
\operatorname{id}_t,
t,
\tau_t
\right).
\end{equation}

Importantly,
%
\begin{equation}
d_t^{\mathrm{out}}
\notin
\mathcal{M}_{t+1}.
\label{eq:no_fused_writeback}
\end{equation}

Consequently, an error introduced by a temporal correction cannot become
the source of a later temporal correction. Long-range information is
therefore preserved without recursively propagating corrected-state
errors.

% ================================================================
\subsection{Direct Causal Motion Estimation}
\label{sec:motion_estimation}

For every available age \(k\in\mathcal{A}\), a frozen SEA-RAFT model
estimates optical flow independently in both directions:
%
\begin{align}
F_{t\rightarrow t-k}
&=
M_{\mathrm{flow}}
\left(
I_t^L,
I_{t-k}^L
\right),
\label{eq:current_to_anchor_flow}
\\
F_{t-k\rightarrow t}
&=
M_{\mathrm{flow}}
\left(
I_{t-k}^L,
I_t^L
\right).
\label{eq:anchor_to_current_flow}
\end{align}

The frozen flow model uses four recurrent iterations and returns flow in
pixel units:
%
\begin{equation}
F_{a\rightarrow b}
\in
\mathbb{R}^{B\times2\times H\times W}.
\end{equation}

All flow estimates are computed directly between the current frame and
the original source frame. No flow fields are composed through
intermediate time steps:
%
\begin{equation}
F_{t\rightarrow t-k}
\neq
F_{t\rightarrow t-1}
\circ
\cdots
\circ
F_{t-k+1\rightarrow t-k}.
\end{equation}

For \(K=4\) available anchors, the architecture requires \(2K\)
directional flow estimates per frame. These evaluations may be batched
in implementation, but they remain geometrically independent.

SEA-RAFT is a frozen external motion estimator and is not an ARGOS
contribution.

% ================================================================
\subsection{BiDA-Style Causal Pull Alignment}
\label{sec:pull_alignment}

The historical disparity is aligned to the current image grid using
target-to-source pull warping.

For current-grid coordinate \((x,y)\), define the source sampling
coordinate:
%
\begin{align}
q_x^{(k)}(x,y)
&=
x+
F_{t\rightarrow t-k}^{x}(x,y),
\\
q_y^{(k)}(x,y)
&=
y+
F_{t\rightarrow t-k}^{y}(x,y).
\end{align}

The normalized sampling coordinates used by
\texttt{grid\_sample} are
%
\begin{align}
g_x^{(k)}
&=
\frac{2q_x^{(k)}}{W-1}-1,
\\
g_y^{(k)}
&=
\frac{2q_y^{(k)}}{H-1}-1.
\end{align}

The aligned candidate disparity is
%
\begin{equation}
\widetilde d_t^{(k)}
=
\warp
\left(
d_{t-k}^S,
F_{t\rightarrow t-k}
\right),
\label{eq:aligned_candidate}
\end{equation}
%
or, pixelwise,
%
\begin{equation}
\widetilde d_t^{(k)}(x,y)
=
\operatorname{grid\_sample}
\left(
d_{t-k}^S,
g_x^{(k)}(x,y),
g_y^{(k)}(x,y)
\right).
\end{equation}

The frozen implementation uses:

\begin{itemize}
    \item bilinear interpolation;
    \item zero padding;
    \item \texttt{align\_corners=True}.
\end{itemize}

Disparity is treated as a scalar field and sampled directly. No
horizontal-flow correction is added after warping.

The geometric support mask is
%
\begin{equation}
S_t^{(k)}(x,y)
=
\ind
\left[
0\leq q_x^{(k)}(x,y)\leq W-1
\right]
\ind
\left[
0\leq q_y^{(k)}(x,y)\leq H-1
\right].
\label{eq:warp_support}
\end{equation}

The historical source-validity mask is warped with the same grid and
thresholded at \(0.999\):
%
\begin{equation}
\widetilde M_{t-k}^{(k)}
=
\ind
\left[
\warp
\left(
M_{t-k},
F_{t\rightarrow t-k}
\right)
\geq0.999
\right].
\label{eq:warped_source_validity}
\end{equation}

The candidate availability mask used by retrieval is
%
\begin{equation}
V_t^{(k)}
=
M_t
\land
\widetilde M_{t-k}^{(k)}
\land
S_t^{(k)}.
\label{eq:candidate_availability}
\end{equation}

% ================================================================
\subsection{Forward--Backward Motion Confidence}
\label{sec:fb_confidence}

The reverse flow is pulled onto the current grid:
%
\begin{equation}
\overline F_{t-k\rightarrow t}
=
\warp
\left(
F_{t-k\rightarrow t},
F_{t\rightarrow t-k}
\right).
\label{eq:aligned_reverse_flow}
\end{equation}

The forward--backward residual is
%
\begin{equation}
e_t^{(k)}
=
\left\|
F_{t\rightarrow t-k}
+
\overline F_{t-k\rightarrow t}
\right\|_2.
\label{eq:fb_residual}
\end{equation}

The adaptive consistency scale is
%
\begin{equation}
\theta_t^{(k)}
=
0.5
+
0.01
\left(
\left\|
F_{t\rightarrow t-k}
\right\|_2
+
\left\|
\overline F_{t-k\rightarrow t}
\right\|_2
\right).
\label{eq:fb_threshold}
\end{equation}

A continuous forward--backward confidence map is then computed as
%
\begin{equation}
C_t^{(k)}
=
S_t^{(k)}
\exp
\left[
-\frac{
e_t^{(k)}
}{
\max\left(\theta_t^{(k)},10^{-6}\right)
}
\right].
\label{eq:fb_confidence}
\end{equation}

The implementation also computes the binary consistency diagnostic
%
\begin{equation}
B_t^{(k)}
=
S_t^{(k)}
\land
\left[
e_t^{(k)}
\leq
\theta_t^{(k)}
\right].
\label{eq:binary_fb_consistency}
\end{equation}

However, \(B_t^{(k)}\) does not directly determine candidate
availability in Eq.~\eqref{eq:candidate_availability}. Motion
consistency is instead provided to the learned model through the
continuous confidence channel \(C_t^{(k)}\).

% ================================================================
\subsection{Robust Cross-Anchor Statistics}
\label{sec:cross_anchor_statistics}

The four aligned candidates provide both local candidate evidence and
robust bank-level context.

Let
%
\begin{equation}
n_t(x)
=
\sum_{j\in\mathcal{A}}
V_t^{(j)}(x)
\end{equation}
%
denote the number of available anchors at pixel \(x\).

The candidate median over the available anchors is denoted by
%
\begin{equation}
m_t(x)
=
\operatorname{Median}
\left\{
\widetilde d_t^{(j)}(x):
V_t^{(j)}(x)=1
\right\}.
\label{eq:candidate_median}
\end{equation}

The corresponding median absolute deviation is
%
\begin{equation}
q_t(x)
=
\operatorname{Median}
\left\{
\left|
\widetilde d_t^{(j)}(x)-m_t(x)
\right|:
V_t^{(j)}(x)=1
\right\}.
\label{eq:candidate_mad}
\end{equation}

The number of agreeing anchors is
%
\begin{equation}
A_t(x)
=
\sum_{j\in\mathcal{A}}
\ind
\left[
V_t^{(j)}(x)=1
\right]
\ind
\left[
\left|
\widetilde d_t^{(j)}(x)-m_t(x)
\right|
\leq q_t(x)+0.10
\right].
\label{eq:agreement_count}
\end{equation}

These statistics provide cross-anchor context without using temporal
attention, 3D convolutions, or recurrent hidden states.

% ================================================================
\subsection{Exact Candidate Evidence}
\label{sec:candidate_evidence}

For clarity, let
%
\begin{equation}
r_t=d_t^S,
\qquad
c_t^{(k)}=\widetilde d_t^{(k)},
\qquad
\epsilon_t^{(k)}=c_t^{(k)}-r_t.
\end{equation}

For each candidate age \(k\), the model receives an exactly ordered
17-channel evidence tensor:
%
\begin{equation}
x_t^{(k)}
=
\left[
x_{t,1}^{(k)},
\ldots,
x_{t,17}^{(k)}
\right].
\label{eq:evidence_tensor}
\end{equation}

The channels are:

\begin{table*}[t]
\centering
\caption{
Exact candidate-evidence channels used by the frozen ARGOS v2
geometry-v1 checkpoint. The ordering is part of the checkpoint
contract.
}
\label{tab:evidence_channels}
\resizebox{\textwidth}{!}{
\begin{tabular}{cll}
\toprule
Index & Channel & Definition \\
\midrule
1
& Current raw disparity
& \(\clip(r_t,0,64)/64\)
\\
2
& Aligned candidate disparity
& \(\clip(c_t^{(k)},0,64)/64\)
\\
3
& Signed raw--candidate residual
& \(\clip(\epsilon_t^{(k)},-16,16)/16\)
\\
4
& Absolute raw--candidate residual
& \(\left|\epsilon_t^{(k)}\right|/16\)
\\
5
& Local absolute-residual mean
& \(\clip\!\left(\mu_5(|\epsilon_t^{(k)}|),0,8\right)/8\)
\\
6
& Local residual standard deviation
& \(\clip\!\left(\sigma_5(\epsilon_t^{(k)}),0,8\right)/8\)
\\
7
& Normalized anchor age
& \(k/8\)
\\
8
& Candidate availability
& \(V_t^{(k)}\)
\\
9
& Geometric warp support
& \(S_t^{(k)}\)
\\
10
& Forward--backward confidence
& \(C_t^{(k)}\)
\\
11
& Valid-witness fraction
& \(n_t/|\mathcal A|\)
\\
12
& Candidate median
& \(\clip(m_t,0,64)/64\)
\\
13
& Candidate MAD
& \(\clip(q_t,0,8)/8\)
\\
14
& Candidate-to-median deviation
& \(\clip(|c_t^{(k)}-m_t|,0,8)/8\)
\\
15
& Raw-to-median deviation
& \(\clip(|r_t-m_t|,0,8)/8\)
\\
16
& Agreement fraction
& \(A_t/|\mathcal A|\)
\\
17
& Raw-anchor provenance
& \(R_t^{(k)}=0\) in the canonical pipeline
\\
\bottomrule
\end{tabular}
}
\end{table*}

The local operators \(\mu_5\) and \(\sigma_5\) use a \(5\times5\)
window, unit stride, and padding two. The local variance is clamped to
zero before applying the square root.

For \(K=4\) anchors, the complete evidence tensor has shape
%
\begin{equation}
X_t
\in
\mathbb{R}^{B\times K\times17\times H\times W}.
\label{eq:complete_evidence_shape}
\end{equation}

The provenance channel is retained for checkpoint compatibility and is
identically zero for all canonical immutable raw anchors.

The alignment utility additionally computes photometric residual,
optical-flow magnitude, disparity-disagreement, and gradient helper
quantities. These quantities are not included in the 17-channel tensor
passed to the CNN.

% ================================================================
\subsection{Shared Candidate Encoder}
\label{sec:shared_encoder}

The candidate dimension is folded into the batch dimension:
%
\begin{equation}
X_t:
[B,K,17,H,W]
\longrightarrow
[BK,17,H,W].
\end{equation}

A single shared CNN \(E_\theta\) is therefore applied independently to
every candidate age:
%
\begin{equation}
Z_t^{(k)}
=
E_\theta
\left(
x_t^{(k)}
\right).
\label{eq:shared_candidate_encoder}
\end{equation}

The candidates share all learned weights. They interact only through the
robust bank-level evidence described in
Section~\ref{sec:cross_anchor_statistics} and through the subsequent
pixelwise retrieval operation.

The exact network is reported in Table~\ref{tab:network_architecture}.

\begin{table}[t]
\centering
\caption{
Layer-by-layer architecture of the shared candidate encoder.
}
\label{tab:network_architecture}
\resizebox{\columnwidth}{!}{
\begin{tabular}{llr}
\toprule
Component & Operation & Parameters \\
\midrule
Stem convolution
& \(3\times3\), \(17\rightarrow32\), padding 1, no bias
& 4,896
\\
Stem normalization
& GroupNorm\((8,32)\)
& 64
\\
Stem activation
& SiLU
& 0
\\
Residual block 1
& Two \(3\times3\), \(32\rightarrow32\) convolutions with GN
& 18,560
\\
Residual block 2
& Two \(3\times3\), \(32\rightarrow32\) convolutions with GN
& 18,560
\\
Residual block 3
& Two \(3\times3\), \(32\rightarrow32\) convolutions with GN
& 18,560
\\
Output head
& \(1\times1\), \(32\rightarrow3\)
& 99
\\
\midrule
Total
&
&
\textbf{60,739}
\\
\bottomrule
\end{tabular}
}
\end{table}

Each residual block applies
%
\begin{equation}
h'
=
\operatorname{SiLU}
\left[
h+
\operatorname{GN}_2
\left(
\operatorname{Conv}_2
\left(
\operatorname{SiLU}
\left[
\operatorname{GN}_1
\left(
\operatorname{Conv}_1(h)
\right)
\right]
\right)
\right)
\right].
\end{equation}

The output head is initialized with zero weights and bias
%
\begin{equation}
[-2,\;0,\;-3].
\end{equation}

The theoretical local receptive field of the CNN is \(15\times15\)
pixels.

% ================================================================
\subsection{Utility-Aware Pixelwise Retrieval}
\label{sec:retrieval}

For each candidate \(k\), the three output channels are transformed into:

\begin{align}
\pi_t^{(k)}
&=
\sigma
\left(
z_{t,0}^{(k)}
\right),
\label{eq:helpfulness_probability}
\\
\widehat{\Delta}_t^{(k)}
&=
8
\tanh
\left(
z_{t,1}^{(k)}
\right),
\label{eq:predicted_improvement}
\\
w_t^{(k)}
&=
\sigma
\left(
z_{t,2}^{(k)}
\right).
\label{eq:candidate_fusion_weight}
\end{align}

Here:

\begin{itemize}
    \item \(\pi_t^{(k)}\) is a learned helpfulness output;
    \item \(\widehat{\Delta}_t^{(k)}\) is the predicted signed geometric
    improvement in pixels;
    \item \(w_t^{(k)}\in(0,1)\) is the candidate-specific fusion weight.
\end{itemize}

The retrieval utility is
%
\begin{equation}
u_t^{(k)}
=
\pi_t^{(k)}
\widehat{\Delta}_t^{(k)}.
\label{eq:retrieval_utility}
\end{equation}

Unavailable candidates are assigned
%
\begin{equation}
u_t^{(k)}
=
-\infty
\qquad
\text{where }
V_t^{(k)}=0.
\label{eq:invalid_score_masking}
\end{equation}

The selected age is determined independently for every output pixel:
%
\begin{equation}
k_t^\star(x)
=
\arg\max_{k\in\mathcal{A}}
u_t^{(k)}(x).
\label{eq:pixelwise_retrieval}
\end{equation}

The retrieval operation selects exactly one candidate per pixel.
There is no softmax over ages and no unconditional averaging of the
candidate bank.

The output \(\pi_t^{(k)}\) is not interpreted as a calibrated probability
of geometric correctness. It is a learned helpfulness quantity used
jointly with the predicted improvement.

% ================================================================
\subsection{Pairwise Soft Correction}
\label{sec:soft_correction}

Let
%
\begin{equation}
\widetilde d_t^\star
=
\widetilde d_t^{(k_t^\star)}
\end{equation}
%
and
%
\begin{equation}
w_t^\star
=
w_t^{(k_t^\star)}
\end{equation}
%
denote the selected aligned disparity and fusion weight.

The proposed corrected disparity is
%
\begin{align}
d_t^{\mathrm{prop}}
&=
d_t^S
+
w_t^\star
\left(
\widetilde d_t^\star-d_t^S
\right)
\label{eq:proposal_residual_form}
\\
&=
\left(
1-w_t^\star
\right)d_t^S
+
w_t^\star
\widetilde d_t^\star.
\label{eq:proposal_convex_form}
\end{align}

The two expressions are numerically equivalent. The implementation uses
the residual form in Eq.~\eqref{eq:proposal_residual_form}.

This operation constitutes a CODD-style pairwise soft correction, but
the complete architecture is not a full CODD model.

% ================================================================
\subsection{Geometric Acceptance and Exact Raw Fallback}
\label{sec:acceptance}

Let
%
\begin{equation}
V_t^\star
=
V_t^{(k_t^\star)},
\qquad
\pi_t^\star
=
\pi_t^{(k_t^\star)},
\qquad
u_t^\star
=
u_t^{(k_t^\star)}.
\end{equation}

The frozen acceptance mask is
%
\begin{equation}
a_t
=
V_t^\star
\land
\left[
\pi_t^\star\geq0.9
\right]
\land
\left[
u_t^\star\geq0.1
\right].
\label{eq:acceptance_mask}
\end{equation}

The final output is
%
\begin{equation}
d_t^{\mathrm{out}}
=
\operatorname{where}
\left(
a_t,
d_t^{\mathrm{prop}},
d_t^S
\right).
\label{eq:exact_raw_fallback}
\end{equation}

The effective fusion weight is therefore
%
\begin{equation}
w_t^{\mathrm{eff}}
=
a_t w_t^\star.
\label{eq:effective_weight}
\end{equation}

Equation~\eqref{eq:exact_raw_fallback} is implemented using exact tensor
selection. Rejected pixels are consequently bit-identical to the input
raw disparity.

A pixel falls back to raw when:

\begin{itemize}
    \item no temporal anchor is available;
    \item the selected candidate lacks valid geometric support;
    \item the predicted helpfulness is below \(0.9\);
    \item the predicted utility is below \(0.1\) pixels.
\end{itemize}

An accepted correction can still remain close to raw when
\(w_t^\star\) is small.

The acceptance mechanism provides geometric abstention, but it does not
guarantee that accepted corrections are safe or non-degrading.

% ================================================================
\subsection{Canonical Training Objective}
\label{sec:training_objective}

Let \(g_t\) denote the ground-truth disparity and \(G_t\) its coverage
map. For candidate \(k\), define the raw and candidate errors:
%
\begin{align}
e_t^S
&=
\left|
d_t^S-g_t
\right|,
\\
e_t^{(k)}
&=
\left|
\widetilde d_t^{(k)}-g_t
\right|.
\end{align}

The true candidate improvement is
%
\begin{equation}
\Delta_t^{(k)}
=
e_t^S-e_t^{(k)}.
\label{eq:true_candidate_improvement}
\end{equation}

The supervised-valid mask is
%
\begin{equation}
Q_t^{(k)}
=
\left[
G_t>0.5
\right]
\land
M_t
\land
V_t^{(k)}
\land
\operatorname{finite}
\left(
\Delta_t^{(k)}
\right).
\label{eq:training_validity}
\end{equation}

A candidate is considered helpful when it improves raw disparity by more
than \(0.10\) pixels:
%
\begin{equation}
h_t^{(k)}
=
Q_t^{(k)}
\land
\left[
\Delta_t^{(k)}>0.10
\right].
\label{eq:helpful_target}
\end{equation}

Let \(N_+\) and \(N_-\) denote the positive and negative supervised
counts. The positive-class weight is
%
\begin{equation}
\rho
=
\clip
\left(
\frac{N_-}{\max(N_+,1)},
1,
50
\right).
\label{eq:positive_class_weight}
\end{equation}

The helpfulness-classification loss is
%
\begin{equation}
\mathcal{L}_{\mathrm{cls}}
=
\operatorname{MaskedMean}_{Q}
\left[
\operatorname{BCEWithLogits}
\left(
z_0,
h;
\operatorname{pos\_weight}=\rho
\right)
\right].
\label{eq:classification_loss}
\end{equation}

The target improvement is clipped to
%
\begin{equation}
\overline{\Delta}_t^{(k)}
=
\clip
\left(
\Delta_t^{(k)},
-8,
8
\right).
\end{equation}

The improvement-regression loss is
%
\begin{equation}
\mathcal{L}_{\mathrm{reg}}
=
\operatorname{MaskedMean}_{Q}
\left[
\operatorname{SmoothL1}
\left(
\widehat{\Delta},
\overline{\Delta}
\right)
\operatorname{where}
\left(
h,
\rho,
1
\right)
\right].
\label{eq:regression_loss}
\end{equation}

For a valid ordered pair of candidates \((i,j)\), define
%
\begin{equation}
o_{ij}
=
\operatorname{sign}
\left(
\Delta^{(i)}-\Delta^{(j)}
\right).
\end{equation}

Pairs are used when
%
\begin{equation}
\left|
\Delta^{(i)}-\Delta^{(j)}
\right|
>
0.10.
\end{equation}

The ranking loss is
%
\begin{equation}
\mathcal{L}_{\mathrm{rank}}
=
\operatorname{MaskedMean}
\left[
\max
\left(
0,
0.10-
o_{ij}
\left(
\widehat{\Delta}^{(i)}
-
\widehat{\Delta}^{(j)}
\right)
\right)
\right].
\label{eq:ranking_loss}
\end{equation}

During training, the differentiable proposal for each candidate is
%
\begin{equation}
q_t^{(k)}
=
d_t^S
+
w_t^{(k)}
\left(
\widetilde d_t^{(k)}-d_t^S
\right).
\end{equation}

The fusion loss is
%
\begin{equation}
\mathcal{L}_{\mathrm{fuse}}
=
\operatorname{MaskedMean}_{Q}
\left[
\operatorname{SmoothL1}
\left(
q_t^{(k)},
g_t
\right)
\operatorname{where}
\left(
h_t^{(k)},
\rho,
1
\right)
\right].
\label{eq:fusion_loss}
\end{equation}

The harmful-update regularizer penalizes fusion weight on candidates that
do not improve raw by more than \(0.10\) pixels:
%
\begin{equation}
\mathcal{L}_{\mathrm{harm}}
=
\operatorname{MaskedMean}_{
Q\land[\Delta\leq0.10]
}
\left[
w
\right].
\label{eq:harm_regularizer}
\end{equation}

The complete canonical soft-training objective is
%
\begin{equation}
\boxed{
\mathcal{L}
=
1.0\mathcal{L}_{\mathrm{cls}}
+
0.5\mathcal{L}_{\mathrm{reg}}
+
0.25\mathcal{L}_{\mathrm{rank}}
+
1.0\mathcal{L}_{\mathrm{fuse}}
+
0.2\mathcal{L}_{\mathrm{harm}}
}
\label{eq:complete_training_loss}
\end{equation}

SEA-RAFT, pull alignment, frozen stereo predictions, and raw disparity
anchors are outside the learned gradient path. Only the 60,739-parameter
candidate retrieval and fusion network is optimized.

% ================================================================
\subsection{Canonical Training Configuration}
\label{sec:training_configuration}

The validated training configuration uses:

\begin{table}[t]
\centering
\caption{Canonical optimization configuration.}
\label{tab:training_configuration}
\begin{tabular}{ll}
\toprule
Setting & Value \\
\midrule
Training sequence IDs & \(1,3,6\) \\
Validation sequence ID & \(2\) \\
Training backbones
& S2M2-S, RAFT-Stereo, StereoAnywhere \\
Crop size & \(64\times80\) \\
Batch size & 12 \\
Epochs & 10 \\
Optimizer & AdamW \\
Initial learning rate & \(2\times10^{-3}\) \\
Weight decay & \(10^{-4}\) \\
Scheduler & Cosine annealing \\
Minimum learning rate & \(0.05\eta_0\) \\
Gradient clipping & 5 \\
Mixed precision & Enabled on CUDA \\
\bottomrule
\end{tabular}
\end{table}

Training samples are grouped by backbone and sequence. Each group is
shuffled deterministically using the run seed and epoch index. Shorter
groups are repeated to match the longest group and the groups are then
interleaved.

The model does not receive the identity of the stereo backbone.
Multi-backbone learning is achieved through the shared external evidence
representation.

% ================================================================
\subsection{Computational Structure}
\label{sec:complexity}

For \(K=4\) available anchors, one current frame requires:

\begin{itemize}
    \item \(K\) current-to-anchor flow estimates;
    \item \(K\) anchor-to-current flow estimates;
    \item \(K\) disparity pull warps;
    \item \(K\) reverse-flow pull warps;
    \item \(K\) 17-channel candidate evidence tensors;
    \item \(K\) shared candidate-encoder evaluations;
    \item one pixelwise candidate selection;
    \item one pairwise fusion and fallback operation.
\end{itemize}

Ignoring the external optical-flow network, evidence construction and the
learned adapter scale as
%
\begin{equation}
\mathcal{O}(KHW),
\end{equation}
%
with \(K=4\).

The learned component contains only
%
\begin{equation}
60\,739
\end{equation}
%
parameters. Temporal memory scales with the retained raw disparity,
validity, and left-image records over an eight-frame horizon.

A dedicated runtime benchmark is required before making a real-time
deployment claim.

% ================================================================
\subsection{Architectural Invariants}
\label{sec:invariants}

The frozen geometry-v1 implementation enforces the following invariants:

\begin{enumerate}
    \item only past frames are addressable;
    \item anchor ages are exactly \(\{1,2,4,8\}\);
    \item every anchor is aligned directly from the current frame;
    \item optical-flow chains are never composed;
    \item long-term memory contains only independent raw stereo outputs;
    \item fused outputs are never written back into memory;
    \item no recurrent corrected disparity state is maintained;
    \item invalid candidate scores are set to \(-\infty\);
    \item rejected corrections return the raw tensor exactly;
    \item disparity follows the positive-left convention;
    \item the temporal model receives no backbone identifier;
    \item the temporal model receives no internal stereo features;
    \item no spatial critic or safety gate belongs to the frozen geometry
    core;
    \item frozen inference is deterministic and parity-tested against the
    validated implementation.
\end{enumerate}

% ================================================================
\subsection{Relation to Bounded Corrected Memory}
\label{sec:h4_relation}

The bounded CODD-style \(H=4\) model maintains a corrected disparity
state within a short recurrent horizon. Periodic reset limits the
accumulation of corrected-state errors.

In contrast, geometry-v1 does not propagate a corrected state. It
retrieves from multiple independently generated raw anchors and applies
a single pairwise correction to the current raw prediction.

The two designs therefore represent different forms of causal memory:

\begin{equation}
\boxed{
\begin{aligned}
\text{bounded }H=4
&:
\text{short-lived corrected memory},
\\
\text{geometry-v1}
&:
\text{long-lived immutable raw memory}.
\end{aligned}
}
\end{equation}

The bounded \(H=4\) model is retained as a strong architectural baseline,
whereas immutable raw multi-anchor retrieval constitutes the principal
ARGOS v2 geometric method.

% ================================================================
\subsection{Scope and Limitations}
\label{sec:scope_limitations}

The architecture supports the following claims:

\begin{itemize}
    \item causal and online framewise operation;
    \item direct multi-age temporal alignment;
    \item immutable raw long-term memory;
    \item no recursive propagation of fused-state errors;
    \item a backbone-independent external input interface;
    \item pixelwise use of long-range temporal anchors;
    \item exact raw fallback on rejected corrections.
\end{itemize}

The current evidence does not establish:

\begin{itemize}
    \item universal backbone agnosticism;
    \item external-domain or OOD generalization;
    \item clinical safety;
    \item strict risk-controlled intervention;
    \item guaranteed non-degradation on accepted pixels;
    \item real-time deployment performance.
\end{itemize}

Exact raw fallback protects only rejected pixels. It does not imply that
accepted temporal corrections are always beneficial.

\subsection{Geometric refinement across stereo backbones}

The bounded H4 temporal refiner consistently improves geometric accuracy
over the frozen raw stereo predictions on SCARED-C D2. Across five stereo
backbones, the aggregate EPE decreases from \(0.302805\) to \(0.289806\),
corresponding to a relative reduction of approximately \(4.29\%\).

\begin{table}[t]
\centering
\caption{Geometric refinement on SCARED-C D2. Relative improvement is
computed as \((\mathrm{EPE}_{\mathrm{raw}}-\mathrm{EPE}_{\mathrm{H4}})
/\mathrm{EPE}_{\mathrm{raw}}\).}
\label{tab:h4_geometry_gain}
\begin{tabular}{lccc}
\toprule
\textbf{Backbone} &
\textbf{Raw EPE} &
\textbf{H4 EPE} &
\textbf{Relative reduction} \\
\midrule
S2M2-S                  & 0.294882 & 0.285021 & 3.34\% \\
RAFT-Stereo             & 0.302213 & 0.291390 & 3.58\% \\
StereoAnywhere          & 0.318889 & 0.296360 & 7.07\% \\
CREStereo               & 0.294341 & 0.283241 & 3.77\% \\
Fast-FoundationStereo   & 0.303701 & 0.293015 & 3.52\% \\
\midrule
\textbf{Aggregate}      & \textbf{0.302805} & \textbf{0.289806} &
\textbf{4.29\%} \\
\bottomrule
\end{tabular}
\end{table}

Importantly, the improvement is not confined to the stereo estimators
used during temporal-refiner training. H4 also improves CREStereo and
Fast-FoundationStereo, which were excluded from temporal-refiner
training, with relative EPE reductions of \(3.77\%\) and \(3.52\%\),
respectively. Overall, H4 improves \(14/15\) backbone--sequence groups,
with an aggregate frame-level win rate of \(57.62\%\) and a
sequence-level win rate of \(93.3\%\).

These results indicate that the temporal module provides a consistent
geometric refinement effect across heterogeneous frozen stereo
estimators. We therefore describe this result as \emph{backbone
transfer within SCARED-C}, rather than universal backbone agnosticism.

\paragraph{Temporal consistency remains unresolved.}
The geometric improvement does not yet translate into improved
GT-referenced temporal dynamics. Although H4 reduces the no-reference
motion-compensated temporal inconsistency (NR-TCE) across all evaluated
backbones and temporal spans, GT-referenced temporal consistency
(GT-TCE) worsens for four of the five stereo backbones. StereoAnywhere
is the only backbone for which GT-TCE improves consistently across
\(k\in\{1,2,4,8\}\).

This discrepancy suggests that improved self-consistency may partly
originate from temporal smoothing or lag rather than more faithful
tracking of the true geometric dynamics. Consequently, we consider the
geometric refinement result established, while temporal-consistency
improvement remains an open objective.

The current result can therefore be summarized as:

\[
\boxed{
\text{H4: } \approx 4.3\% \text{ mean EPE reduction}
\quad+\quad
5/5 \text{ backbone improvement}
\quad+\quad
\text{transfer to unseen stereo estimators}
}
\]

while a successful temporal refiner should ultimately satisfy both

\[
\Delta \mathrm{EPE}
=
\mathrm{EPE}_{\mathrm{raw}}
-
\mathrm{EPE}_{\mathrm{refined}}
>0
\]

and

\[
\Delta \mathrm{TCE}
=
\mathrm{TCE}_{\mathrm{raw}}
-
\mathrm{TCE}_{\mathrm{refined}}
>0.
\]

\begin{table}[t]
    \centering
    \caption{
        Summary of the current transfer evidence for the frozen H4
        temporal module.
        ``OOD geometry'' is intentionally left unclaimed where a reliable
        metric reference is unavailable.
    }
    \label{tab:argos_transfer_summary}
    \small
    \begin{tabular}{llll}
        \toprule
        Evaluation & Shift & Result & Evidence \\
        \midrule
        SCARED-C D2
            & Backbone
            & $-4.29\%$ EPE
            & 5/5 backbones improve \\
        SCARED-C D2
            & Unseen backbone
            & Positive
            & CREStereo + Fast-FoundationStereo \\
        SCARED-C D7
            & Held-out sequence
            & $-4.42\%$ EPE
            & 5/5 backbones improve \\
        D4D
            & Cross-dataset
            & $-54.3\%$ temporal inconsistency
            & No-reference \\
        SERV-CT
            & Cross-dataset static
            & Identity
            & Temporal test N/A \\
        \bottomrule
    \end{tabular}
\end{table}

% Requires:
% \usepackage{booktabs}
% \usepackage{multirow}

\subsection{Current Experimental Status}
\label{sec:current_status}

We evaluate a single frozen causal temporal refinement model, denoted
\textbf{Canonical H4}, using the same checkpoint and bounded recurrent
policy across stereo backbones and datasets.
The model was trained using disparity streams from S2M2-S,
RAFT-Stereo, and StereoAnywhere, while CREStereo and
Fast-FoundationStereo were excluded from temporal-refiner training.

The current evidence addresses three complementary questions:
(i) whether the same temporal refiner transfers across stereo
architectures;
(ii) whether it remains beneficial on held-out and external surgical
data; and
(iii) how it compares with existing temporal stabilization methods.


% ============================================================
\subsubsection{Cross-Backbone Transfer on SCARED-C}
% ============================================================

On SCARED-C D2, Canonical H4 improves disparity accuracy for all five
evaluated stereo estimators.
Importantly, the improvement extends to CREStereo and
Fast-FoundationStereo, whose predictions were never observed during
temporal-refiner training.

\begin{table}[t]
    \centering
    \caption{
        Cross-backbone transfer of Canonical H4 on SCARED-C D2.
        Values correspond to the primary macro-sequence EPE aggregation.
        ``Unseen'' denotes stereo estimators excluded from temporal-refiner
        training.
    }
    \label{tab:scared_d2_current}
    \small
    \begin{tabular}{llccc}
        \toprule
        Backbone & Status &
        Raw EPE $\downarrow$ &
        H4 EPE $\downarrow$ &
        Rel. gain \\
        \midrule
        S2M2-S
            & Seen
            & 0.270422 & 0.264544 & 2.17\% \\

        RAFT-Stereo
            & Seen
            & 0.269113 & 0.261388 & 2.87\% \\

        StereoAnywhere
            & Seen
            & 0.302181 & 0.276674 & 8.44\% \\

        CREStereo
            & Unseen
            & 0.266808 & 0.257823 & 3.37\% \\

        Fast-FoundationStereo
            & Unseen
            & 0.272485 & 0.264712 & 2.85\% \\
        \bottomrule
    \end{tabular}
\end{table}

Under the canonical pixel-weighted aggregation, the corresponding
overall D2 EPE decreases from

\[
0.302805 \rightarrow 0.289806,
\]

which corresponds to a relative reduction of approximately
\textbf{4.29\%}.


% ============================================================
\subsubsection{Held-Out In-Domain Evaluation}
% ============================================================

The same frozen checkpoint also improves all five backbones on
SCARED-C D7, without backbone-specific retraining.

\begin{table}[t]
    \centering
    \caption{
        Held-out SCARED-C D7 evaluation using the unchanged Canonical H4
        checkpoint.
    }
    \label{tab:scared_d7_current}
    \small
    \begin{tabular}{llccc}
        \toprule
        Backbone & Status &
        Raw EPE $\downarrow$ &
        H4 EPE $\downarrow$ &
        Rel. gain \\
        \midrule
        S2M2-S
            & Seen
            & 0.517635 & 0.500710 & 3.27\% \\

        RAFT-Stereo
            & Seen
            & 0.467080 & 0.438604 & 6.10\% \\

        StereoAnywhere
            & Seen
            & 0.477363 & 0.444676 & 6.85\% \\

        CREStereo
            & Unseen
            & 0.516462 & 0.495545 & 4.05\% \\

        Fast-FoundationStereo
            & Unseen
            & 0.481883 & 0.466750 & 3.14\% \\
        \bottomrule
    \end{tabular}
\end{table}

The previously established pixel-weighted aggregate decreases from
approximately

\[
0.540759 \rightarrow 0.516854,
\]

corresponding to an overall improvement of approximately
\textbf{4.42\%}.


% ============================================================
\subsubsection{External-Domain Evaluation}
% ============================================================

\paragraph{DRENDS.}
We next evaluate Canonical H4 without retraining on a chronological
DRENDS ex-vivo recording containing 1,499 synchronized frames.
The disparity reference is derived from the metric ToF reference and is
therefore not an independent stereo disparity ground truth.
Accordingly, DRENDS is treated as an external metric evaluation with this
explicit caveat.

After correcting the DRENDS depth-evaluation protocol, both disparity
EPE and depth MAE improve.

\begin{table}[t]
    \centering
    \caption{
        Zero-shot evaluation on the 1,499-frame DRENDS
        \texttt{Vid14\_Pancreas\_High} recording.
        Depth is evaluated using the corrected DRENDS protocol.
    }
    \label{tab:drends_current}
    \small
    \begin{tabular}{lccc}
        \toprule
        Metric &
        Raw &
        Canonical H4 &
        Change \\
        \midrule
        Disparity EPE $\downarrow$
            & 1.137046 px
            & 1.084080 px
            & $-4.66\%$ \\

        Depth MAE $\downarrow$
            & 14.3581 mm
            & 13.8997 mm
            & $-3.19\%$ \\

        Depth median $\downarrow$
            & 13.4207 mm
            & 13.0666 mm
            & $-2.64\%$ \\

        Depth P95 $\downarrow$
            & 28.4528 mm
            & 26.8138 mm
            & $-5.76\%$ \\

        Depth RMSE $\downarrow$
            & 16.6314 mm
            & 17.9957 mm
            & $+8.20\%$ \\
        \bottomrule
    \end{tabular}
\end{table}

The improvement in median, MAE, and P95 depth error indicates that the
benefit is not restricted to disparity space.
However, depth RMSE increases, revealing a remaining tail-risk failure
mode that is not visible from mean disparity EPE alone.
DRENDS therefore provides encouraging preliminary evidence of
cross-dataset geometric transfer, while also motivating explicit
tail-risk and safety-aware evaluation.


\paragraph{D4D.}
D4D cannot currently provide a trustworthy metric geometric comparison
because of an unresolved inconsistency between the available structured
reference and the cached stereo geometry.
We therefore use it only as a no-reference temporal-transfer experiment.

\begin{table}[t]
    \centering
    \caption{
        No-reference motion-compensated temporal inconsistency on D4D.
        These values measure prediction-space temporal stability and must
        not be interpreted as geometric ground-truth error.
    }
    \label{tab:d4d_current}
    \small
    \begin{tabular}{lccc}
        \toprule
        Backbone &
        Raw $\downarrow$ &
        H4 $\downarrow$ &
        Rel. change \\
        \midrule
        RAFT-Stereo
            & 0.509406
            & 0.245049
            & $-51.9\%$ \\

        StereoAnywhere
            & 6.927071
            & 3.151616
            & $-54.5\%$ \\
        \bottomrule
    \end{tabular}
\end{table}


% ============================================================
\subsubsection{Comparison with Existing Temporal Stabilizers}
% ============================================================

We additionally reproduced two external temporal stabilization methods
on the SCARED-C D2 development split.

\paragraph{BiDAStabilizer.}
The official BiDAStabilizer pipeline was evaluated on 4,249 D2 frames
using RAFT-Stereo predictions.
The upstream disparities were verified to be numerically identical to
the corresponding raw baseline inputs.
Its pixel-weighted EPE changes from

\[
0.302862 \rightarrow 0.302109,
\]

corresponding to only a small geometric improvement.
The method improves two of the three evaluated D2 sequences and slightly
degrades the third.

In contrast, Canonical H4 obtains a substantially larger geometric
reduction under its canonical D2 evaluation:

\[
0.302805 \rightarrow 0.289806.
\]

The exact numerical difference between the two raw EPE values reflects
their independently validated evaluation paths and should therefore not
be interpreted as a direct subtraction between methods.


\paragraph{NVDS.}
We also reproduced the official bidirectional NVDS pipeline using its
intended monocular relative-depth setting.
Because NVDS does not operate in metric stereo-disparity space, it is
not compared to H4 using disparity EPE.

Across all 4,249 D2 frames, NVDS reduces its official-style
motion-based OPW diagnostic from

\[
0.127475 \rightarrow 0.059383,
\]

corresponding to approximately a

\[
\boxed{53.4\%}
\]

reduction in relative-depth temporal inconsistency.
NVDS therefore serves as a strong temporal-stability reference, rather
than as a direct metric-stereo competitor.


\begin{table}[t]
    \centering
    \caption{
        Current positioning with respect to external refinement and
        stabilization approaches.
        Metrics are intentionally kept in their appropriate prediction
        spaces.
    }
    \label{tab:external_comparison_status}
    \small
    \begin{tabular}{llll}
        \toprule
        Method &
        Setting &
        Primary result &
        Interpretation \\
        \midrule
        Canonical H4
            & Causal stereo
            & $-4.29\%$ EPE
            & Metric refinement \\

        BiDAStabilizer
            & Bidirectional stereo
            & Small EPE gain
            & Direct temporal competitor \\

        NVDS
            & Bidirectional monocular depth
            & $-53.4\%$ OPW
            & Temporal-stability reference \\
        \bottomrule
    \end{tabular}
\end{table}


% ============================================================
\subsubsection{Failed Safety Variants and Current Model Selection}
% ============================================================

Several post-hoc hard/gated variants of Canonical H4 were explored on
the D2 development split.
Although some variants improved mean geometric metrics, none satisfied
the complete promotion criteria across geometry, safety, unseen
backbones, and external-domain behavior.

In particular, a hard threshold of $0.35$ reduced D2 mean disparity and
depth errors and also improved CREStereo and Fast-FoundationStereo.
However, it slightly increased Bad1 and produced severe near-zero
disparity outliers under DRENDS.
More conservative relative and support-aware guards also failed the
predefined D2 promotion criteria, primarily because of recurrent
degradation on StereoAnywhere.

Consequently, none of these variants replaces Canonical H4.

\[
\boxed{
\text{Canonical H4 remains the frozen reference model.}
}
\]

This negative result is informative: apparently safe frame-wise gating
rules can alter the recurrent state and create substantial downstream
tail errors.
Safety mechanisms must therefore be designed and evaluated as part of
the recurrent dynamics rather than as purely post-hoc per-frame filters.


% ============================================================
\subsubsection{Current Evidence Summary}
% ============================================================

\begin{table}[t]
    \centering
    \caption{
        Current experimental evidence for ARGOS.
        A checkmark denotes directly supported evidence; ``partial''
        indicates that the current experiment supports only part of the
        corresponding claim.
    }
    \label{tab:current_evidence_summary}
    \small
    \begin{tabular}{lll}
        \toprule
        Question & Current evidence & Status \\
        \midrule
        Improves seen stereo backbones
            & SCARED-C D2/D7
            & \checkmark \\

        Transfers to unseen backbones
            & CREStereo, Fast-FoundationStereo
            & \checkmark \\

        Held-out in-domain transfer
            & SCARED-C D7
            & \checkmark \\

        Cross-dataset temporal transfer
            & D4D
            & \checkmark \\

        Cross-dataset geometric transfer
            & DRENDS Vid14
            & Preliminary \\

        Improves metric depth OOD
            & DRENDS MAE/Median/P95
            & Preliminary \\

        Controls OOD tail risk
            & DRENDS RMSE
            & Not yet \\

        Beats direct temporal competitors
            & BiDA D2
            & Promising \\

        General temporal stability
            & NVDS comparison
            & Different prediction space \\
        \bottomrule
    \end{tabular}
\end{table}


\paragraph{Current takeaway.}
The results obtained so far indicate that temporal refinement learned on
a limited set of stereo estimators is transferable across architectures
and remains beneficial under held-out and external surgical data.
Most importantly, the same frozen checkpoint improves two stereo
backbones excluded from temporal-refiner training and shows preliminary
zero-shot geometric improvement on DRENDS.

At the same time, the DRENDS RMSE and the failed post-hoc gating
experiments show that improvements in average accuracy do not guarantee
control of rare harmful corrections.
This motivates the next stage of ARGOS: preserving the demonstrated
cross-backbone and cross-dataset gains while explicitly controlling
temporal-refinement risk and tail degradation.

\section*{STEER -- D2 Experimental Closure}
\textbf{Date:} 14 August 2026

\paragraph{Protocol.}
The experimental closure was completed on SCARED-C D2 using:
15 temporal variants, 5 stereo backbones, and 3 temporal sequences.
The final evaluation contains 225 standard evaluations plus 15 oracle
diagnostics, for a total of 240 rows.
All 241 recorded hashes were verified.
Raw and refined predictions were evaluated on identical support.

The common raw macro EPE is:
\[
\boxed{
\EPE_{\mathrm{raw}} = 0.276202
}
\]

\subsection*{Main ranking}

\begin{table}[h]
\centering
\small
\begin{tabular}{lccc}
\toprule
Method & Refined EPE & EPE reduction & Improved groups \\
\midrule
STEER $H=6$ & \textbf{0.264656} & 4.18\% & 12/15 \\
STEER $H=8$ & 0.264912 & 4.09\% & 11/15 \\
STEER $H=4$ & 0.265028 & 4.05\% & 14/15 \\
STEER $H=2$ & 0.267265 & 3.24\% & 14/15 \\
STEER $H=1$ & 0.269001 & 2.61\% & 15/15 \\
\midrule
Fixed $w=0.5$ & 0.274562 & 0.59\% & 10/15 \\
EMA-2 & 0.274562 & 0.59\% & 10/15 \\
Fixed $w=0.3$ & 0.274852 & 0.49\% & 10/15 \\
Fixed $w=0.2$ & 0.275229 & 0.35\% & 10/15 \\
EMA-3 & 0.275386 & 0.30\% & 6/15 \\
Fixed $w=0.1$ & 0.275684 & 0.19\% & 10/15 \\
Warped raw previous & 0.276298 & -0.03\% & 5/15 \\
FB confidence & 0.282182 & -2.17\% & 5/15 \\
Warped recurrent $w=1$ & 0.286035 & -3.56\% & 5/15 \\
Continuous recurrence & 0.449349 & -62.69\% & 0/15 \\
\midrule
Oracle raw/memory & 0.242223 & 12.30\% & 15/15 \\
\bottomrule
\end{tabular}
\caption{SCARED-C D2 experimental closure. Positive EPE reduction
denotes improvement relative to the common raw prediction.}
\end{table}

\paragraph{Learned head versus trivial blending.}

For the canonical $H=4$ policy:
\[
\Delta_{\mathrm{STEER}}
=
0.276202 - 0.265028
=
0.011174~\mathrm{px}.
\]

For the best fixed blend:
\[
\Delta_{\mathrm{fixed}}
=
0.276202 - 0.274562
=
0.001640~\mathrm{px}.
\]

Therefore,
\[
\boxed{
\frac{\Delta_{\mathrm{STEER}}}
{\Delta_{\mathrm{fixed}}}
\approx 6.8
}
\]

i.e.\ the learned temporal policy recovers almost seven times the gain
of the best fixed temporal blend.

The result is therefore not explained by simple averaging, EMA,
warped-memory copying, or direct use of forward--backward confidence.

\subsection*{Horizon risk--gain trade-off}

\begin{table}[h]
\centering
\small
\begin{tabular}{lccccc}
\toprule
Method & EPE & Degraded frames & HUR & BUR & $B^+/H^+$ \\
\midrule
$H=4$ & 0.265028 & 41.46\% & 0.4048 & 0.5074 & \textbf{2.53} \\
$H=6$ & \textbf{0.264656} & 42.18\% & 0.4053 & 0.5065 & 2.30 \\
$H=8$ & 0.264912 & 42.18\% & 0.4068 & 0.5047 & 2.04 \\
\bottomrule
\end{tabular}
\caption{Risk--gain comparison for the three strongest bounded horizons.}
\end{table}

The pure-EPE difference between $H=4$ and $H=6$ is:
\[
0.265028 - 0.264656
=
\boxed{0.000372~\mathrm{px}},
\]
corresponding to approximately $0.14\%$ of the $H=4$ EPE.

However, $H=4$ provides a better intervention ratio:
\[
\boxed{
B^+/H^+ = 2.53
}
\]
compared with $2.30$ for $H=6$ and $2.04$ for $H=8$.

Tail behaviour also favours $H=4$:
\[
\begin{array}{lccc}
\toprule
 & H=4 & H=6 & H=8 \\
\midrule
P95\ \mathrm{frame\ degradation} &
0.01715 &
0.02037 &
0.02441 \\
\mathrm{NewBad1} &
0.0467\% &
0.0676\% &
0.1282\% \\
\bottomrule
\end{array}
\]

Hence:
\[
\boxed{
H=6:\ \text{best development EPE}
\qquad
H=4:\ \text{selected risk--gain operating point}
}
\]

No significance claim is made for the $H=4$ versus $H=6$ EPE difference.

\subsection*{Bounded recurrence}

Continuous corrected-state recurrence fails strongly:
\[
0.276202
\rightarrow
0.449349,
\]
corresponding to a $62.69\%$ increase in disparity error.

This supports the conclusion that periodic raw re-anchoring is a
structural component of STEER rather than a minor implementation detail.

\subsection*{Cross-backbone behaviour}

\begin{table}[h]
\centering
\small
\begin{tabular}{lccc}
\toprule
Backbone & Training status & EPE reduction & Improved groups \\
\midrule
Stereo Anywhere & seen & 8.44\% & 3/3 \\
CREStereo & unseen & 3.37\% & 3/3 \\
RAFT-Stereo & seen & 2.87\% & 3/3 \\
Fast-FoundationStereo & unseen & 2.85\% & 3/3 \\
S$^2$M$^2$-S & seen & 2.17\% & 2/3 \\
\bottomrule
\end{tabular}
\caption{Canonical STEER $H=4$ performance by stereo backbone.}
\end{table}

Average gain:
\[
\boxed{
\text{seen backbones}: 4.65\%
}
\]

\[
\boxed{
\text{unseen backbones}: 3.11\%
}
\]

All six unseen-backbone sequence groups improve.

\subsection*{Sequence behaviour}

\begin{table}[h]
\centering
\small
\begin{tabular}{lcc}
\toprule
Sequence & EPE reduction & Improved backbones \\
\midrule
Keyframe 4 & 4.84\% & 5/5 \\
Keyframe 2 & 4.63\% & 5/5 \\
Keyframe 3 & 2.48\% & 4/5 \\
\bottomrule
\end{tabular}
\caption{Canonical STEER $H=4$ performance by SCARED-C D2 sequence.}
\end{table}

The only negative $H=4$ backbone--sequence group is
S$^2$M$^2$-S on keyframe 3.

\subsection*{Oracle headroom}

The pixelwise raw-versus-aligned-memory oracle obtains:
\[
0.276202
\rightarrow
0.242223,
\]
corresponding to:
\[
\boxed{
12.30\% \text{ potential EPE reduction}
}
\]

The available oracle gain is:
\[
\Delta_{\mathrm{oracle}}
=
0.033979~\mathrm{px}.
\]

Canonical STEER $H=4$ captures:
\[
\frac{0.011174}{0.033979}
=
\boxed{32.9\%}
\]
of this oracle improvement.

This indicates that motion-aligned memory contains substantially more
useful information than the current learned policy is able to exploit.

\subsection*{Current interpretation}

The D2 closure supports the following mechanism:

\begin{enumerate}
    \item motion alignment creates useful but imperfect temporal
    candidates;
    \item indiscriminate temporal fusion does not improve stereo
    substantially and can degrade it;
    \item the learned STEER head provides the majority of the observed
    temporal gain by selectively weighting aligned memory;
    \item bounded recurrence prevents error accumulation;
    \item $H=4$ sacrifices only $0.000372$ px relative to the lowest
    development EPE at $H=6$ while providing better tail and
    intervention behaviour.
\end{enumerate}

The current concise interpretation is:

\begin{quote}
\textbf{STEER does not behave as a simple temporal smoother.
Its gain arises from the combination of geometric alignment,
learned selective fusion, and bounded recurrent memory.}
\end{quote}

\paragraph{Status.}
This experiment is development-only (SCARED-C D2).
The frozen $H=4$ policy must be confirmed on held-out D7 without
changing the selection criterion.

% =====================================================================
\section*{Winning Configuration (current)}
% =====================================================================

This section is maintained as the single source of truth for which
configuration produces the paper numbers. It is updated whenever the
winning configuration changes, and the superseded one is recorded
immediately below it.

\begin{center}
\begin{tabular}{ll}
\toprule
\textbf{Item} & \textbf{Value} \\
\midrule
Temporal module   & \texttt{canonical\_h4\_masked} \\
Module code       & \texttt{model\_design/comparison/canonical\_h4\_masked.py} \\
Parent module     & \texttt{canonical\_h4}, sha256 \texttt{4148ee8b4b6d\ldots} (unmodified) \\
Checkpoint        & \texttt{codd\_style\_h4\_best\_validation.pt} \\
Checkpoint sha256 & \texttt{99c5745c164fd\ldots0d4ed4305cd10609324d725} \\
Head parameters   & 177{,}338 (unchanged) \\
Evidence channels & 142 (unchanged) \\
State lifetime    & bounded $H=4$ with raw re-anchor (unchanged) \\
Policy file       & \texttt{codd\_style\_h4\_policy.json} (unchanged) \\
Delta vs.\ parent & one line: \texttt{fused = where(support, fused, raw)} \\
\bottomrule
\end{tabular}
\end{center}

\textbf{Nothing was retrained.} The winning configuration differs from
the previously reported one only by confining the module output to the
support mask that the module already computes, which is the same
masking every matched baseline in \texttt{experimental\_policies.py}
has always applied.

\paragraph{Superseded.}
\texttt{canonical\_h4} with unmasked output. Retained for provenance
and for the published canonical-versus-masked comparison, but it must
not be used for new paper numbers.

% =====================================================================
\section*{Root Cause: Unmasked Fusion Output}
% =====================================================================

Metric depth regressed on DRENDS and D4D while disparity improved.
This was root-caused to a missing mask rather than to geometry.

\begin{enumerate}[itemsep=2pt]
    \item \texttt{causal\_warp} uses
    \texttt{F.grid\_sample(..., padding\_mode="zeros")}, so an
    out-of-bounds warp returns an aligned past disparity of
    \emph{exactly} $0.0$.
    \item \texttt{CODDStyleFusionHead.forward} computes
    $d^{\text{fused}} = (1-w)\,d^{\text{raw}} + w\,\widetilde{m}$ over
    \emph{every} pixel, with no support mask.
    \item \texttt{canonical\_h4.step} returns that value unmasked,
    while computing and reporting a \texttt{support} it never applies.
    Every baseline in \texttt{experimental\_policies.py} does apply
    \texttt{torch.where(support, fused, raw)}.
    \item Where the head places $w \approx 1$ on a zero-memory pixel,
    the fused disparity collapses to $0$ and the prediction becomes
    invalid.
    \item \texttt{drive()} feeds \texttt{result["disparity"]} back as
    the next state, so the zero propagates until the next raw
    re-anchor. The offender state-age histogram is dominated by ages
    2--4 (Vid10: 10/103/135/52 at ages 1/2/3/4), and offending frames
    arrive in bursts of at most four consecutive frames. Bounded $H=4$
    contains the damage exactly as designed, but does not prevent the
    initial injection.
    \item \texttt{MetricConfig.invalid\_penalty\_mm}$\,=10000.0$
    charges each such pixel.
\end{enumerate}

Measured: Vid10 has 300 invalid pixels out of $38{,}828{,}160$, Vid11
has 265 out of $38{,}439{,}360$. The aligned memory of \emph{every}
offender has $\min = \max = 0.0$ with zero non-finite values.
Diagnostic script:
\texttt{scripts/diagnose\_invalid\_refined.py}.

\paragraph{Corrections to two earlier readings.}
The $1000$\,px and $10000$\,mm values appearing in the tables are
\emph{error penalties} for invalid predictions, not clamped predicted
values. The regression is also \emph{not} a $Z = fB/d$ conversion
artefact; that earlier hypothesis was wrong.

% =====================================================================
\section*{Fix Verified and Adopted}
% =====================================================================

\begin{center}
\begin{tabular}{lrr}
\toprule
Recording & invalid px, canonical & invalid px, masked \\
\midrule
Vid10\_Liver\_Med  & 300 / 38{,}828{,}160 & \textbf{0} \\
Vid11\_Liver\_High & 265 / 38{,}439{,}360 & \textbf{0} \\
\bottomrule
\end{tabular}
\end{center}

DRENDS recompute, five recordings by three backbones:

\begin{center}
\begin{tabular}{lrr}
\toprule
Metric & canonical & masked \\
\midrule
depth RMSE      & regresses \textbf{14/14}, mean $+110.4\%$ & improves \textbf{14/14}, mean $-3.2\%$ \\
depth MAE       & regresses 10/14, mean $+4.2\%$            & improves \textbf{14/14}, mean $-3.7\%$ \\
InvalidRate     & up to $1.3\times10^{-5}$                  & $0$ everywhere \\
disparity EPE   & $-4.30\%$                                 & $\mathbf{-5.02\%}$ \\
disparity Bad3  & $-14.04\%$                                & $\mathbf{-18.09\%}$ \\
\bottomrule
\end{tabular}
\end{center}

Worst case: Vid10 with RAFT-Stereo, depth RMSE $69.12$ ($+256.5\%$)
becomes $18.74$ ($-3.4\%$). Across 14 combinations and 5 metrics,
\textbf{70 of 70 cells are equal or better and none regress}.
Disparity improves as well, because masking the output also stops the
zero from entering memory.

\paragraph{Consequence.}
\texttt{transfer\_closure\_v7} justified
\texttt{external\_ood\_transfer: NOT\_CONFIRMED} with ``a severe
refined tail'' and ``maxima hit 1{,}000\,px / 10{,}000\,mm''. That
tail was the missing line. The same defect drives the failing gates in
\texttt{joint\_d4d\_v5}, where CREStereo on specimen\_3 shows a depth
MAE of $23{,}060$\,mm while depth P99 stays at $53$\,mm.

\paragraph{Protocol note.}
This is a protocol change, not a silent bug fix. The
canonical-versus-masked comparison is itself a result and should be
published: it shows that the support contract must be applied to the
output, not merely declared.

% =====================================================================
\section*{Joint Unseen-Backbone and External-Domain Transfer}
% =====================================================================

\texttt{transfer\_closure\_v7} recorded
\texttt{joint\_unseen\_backbone\_and\_external\_ood\_transfer:
UNAVAILABLE} on the grounds that no completed evaluation crossed both
axes. That cell is now measured.

Canonical $H=4$ was run on five DRENDS recordings
(Vid10\_Liver\_Med, Vid11\_Liver\_High, Vid12\_Pancreas\_Ext,
Vid13\_Pancreas\_Med, Vid14\_Pancreas\_High) with three stereo
estimators, two of which were excluded from temporal-head training,
for a total of $20{,}929$ frames.

Disparity improved in \textbf{42 of 42 cells}, that is 14
recording-by-backbone combinations across EPE, Bad1 and Bad3, with
mean EPE $-4.31\%$, Bad1 $-5.30\%$ and Bad3 $-14.06\%$ (best
$-27.06\%$).

The per-recording pattern is determined by the recording rather than
by the backbone. On Vid10 the EPE change is $-3.95\%$, $-3.92\%$ and
$-4.25\%$ for RAFT-Stereo (seen), CREStereo (unseen) and
Fast-FoundationStereo (unseen) respectively. This is stronger evidence
for backbone independence than any aggregate mean, because it shows
the residual variance comes from the data and not from the estimator.

The new module reuses the pinned \texttt{drends\_evaluation.py}
(sha256 \texttt{38214918bd328ab2\ldots}) unchanged and replaces only
the stereo predictor:
\texttt{model\_design/comparison/drends\_backbone\_transfer.py}.

% =====================================================================
\section*{Selective Risk: $w_t$ Does Not Rank Harm}
% =====================================================================

The AUROC of the effective temporal weight $w_t$ as a predictor of
$\delta e > 0$ is $0.366$--$0.382$ on D2 and $0.496$--$0.506$ on D7.
The raw-versus-memory disagreement and the update magnitude behave the
same way. Risk--coverage analysis with $w_t$ captures only $3.5\%$
(D2) and $6.1\%$ (D7) of the achievable oracle gain.

The reliability diagram explains why. For RAFT-Stereo the harm rate
\emph{falls} with $w$ on D2 ($0.550 \rightarrow 0.324$) but
\emph{rises} with $w$ on D7 ($0.260 \rightarrow 0.450$). The ordering
inverts between development and held-out data, so $w_t$ is not merely
a weak risk score, it is a \textbf{sign-unstable} one and therefore
not transferable.

\paragraph{Wording correction.}
The module intervenes on approximately $90\%$ of pixels with a mean
weight of $0.34$. That is \emph{graded convex intervention}, not
``selective intervention''. The conclusion's wording must be changed
accordingly.

% =====================================================================
\section*{The Published Oracle Is the Wrong Ceiling}
% =====================================================================

The reported oracle selects per pixel between raw and aligned memory,
so it is the ceiling of $w \in \{0,1\}$. STEER acts on $w \in [0,1]$,
whose ceiling is
$w^{\star} = \mathrm{clip}\big((g - d^{\text{raw}}) /
(\widetilde{m} - d^{\text{raw}}),\,0,\,1\big)$.

Only $6.1\%$ of D2 pixels and $11$--$14\%$ of D7 pixels have the
ground truth bracketed by the two candidates, so the continuous
ceiling sits just below the binary one and the recovered fraction
falls from $32.9\%$ to approximately $30\%$. The correction is real
but small: it is a matter of rigour, not a change of story.

% =====================================================================
\section*{Runtime}
% =====================================================================

NVIDIA H100 PCIe, torch 2.5.1+cu121, $144 \times 180$ grid, 200 timed
iterations after 30 warmup iterations.

\begin{center}
\begin{tabular}{lrr}
\toprule
Stage & median (ms) & p95 (ms) \\
\midrule
SEA-RAFT forward & 13.222 & 13.289 \\
SEA-RAFT reverse & 13.191 & 13.268 \\
BiDA alignment   & 1.038  & 1.057 \\
142-channel evidence (incl.\ frozen ResNet-18) & 5.122 & 5.156 \\
Fusion head      & \textbf{1.027} & 1.043 \\
\midrule
Total temporal overhead & \textbf{33.601} & \\
Peak VRAM & \multicolumn{2}{r}{\textbf{124.2\,MB}} \\
\bottomrule
\end{tabular}
\end{center}

End-to-end throughput ranges from $11.0$\,FPS with
Fast-FoundationStereo down to $1.4$\,FPS with Stereo Anywhere.
\textbf{No real-time claim is supported.}

The 177k-parameter head costs $1.03$\,ms, only $3\%$ of the overhead,
while $79\%$ is the two frozen SEA-RAFT passes. The reverse pass costs
$13.19$\,ms and feeds only the forward--backward confidence cue, and
using that cue directly as the fusion weight \emph{worsens} EPE by
$2.17\%$. Removing it would cut overhead by $39\%$, which makes it a
worthwhile ablation.

% =====================================================================
\section*{Matched-Protocol Audit}
% =====================================================================

All 15 closure methods come from a single run under
\texttt{paper\_d2\_strict\_all\_anchors}, sharing the bounded state
lifetime, the re-anchor schedule, the SEA-RAFT alignment, the support
definition and the fusion form; only the blending weight differs. The
225 reports correspond to 15 methods by 5 backbones by 3 sequences.

\paragraph{Qualification.}
The output \emph{masking} was not matched. The baselines applied
\texttt{where(support, fused, raw)} while the canonical module did
not. On SCARED-C this does not bite, since the refined InvalidRate is
zero there, so the $6.8\times$ ratio stands. However the claim that
``only $w$ differs'' was too strong before the masked variant. With
\texttt{canonical\_h4\_masked} the protocol is genuinely matched.

\paragraph{An identity worth stating.}
\texttt{ema2\_h4} and \texttt{fixed\_w0.5\_h4} are both
\texttt{PolicySpec(horizon=4, weight=0.5)}, that is the same policy
under two names, and \texttt{ema3\_h4} is $w = 2/3$. The baseline
family is therefore five fixed convex weights
$w \in \{0.1, 0.2, 0.3, 0.5, 2/3\}$ over the shared bounded recurrent
state, not two independent families. Reporting it that way is both
more honest and stronger, because two ``different'' baselines
producing an identical number otherwise looks like a bug.

% =====================================================================
\section*{Metrics Already Computed but Never Reported}
% =====================================================================

The following are present in \texttt{definitive\_table.csv} and absent
from the paper.

\paragraph{Bad-pixel rates improve in 30 of 30 SCARED-C cells}
(2 splits by 5 backbones by \{Bad1, Bad3, Bad5\}). Stereo Anywhere on
D2 gives Bad1 $-19.7\%$, Bad3 $-32.0\%$ and Bad5 $-45.3\%$. On D7,
Bad5 improves more than Bad1 on all five backbones, which means the
module removes large errors rather than smoothing small ones. This is
the anti-smoothing argument in direct form.

\paragraph{Disparity RMSE improves on all five backbones on both
splits.} P99 improves on 10 of 10 and Max on 9 of 10. Stereo Anywhere
on D2 has its Max fall from $33.27$ to $14.88$, a $55\%$ reduction.

\paragraph{Harm and benefit per backbone} (macro-sequence):
$B^{+}/H^{+}$ lies between $1.76$ and $4.59$; RecoveredBad1 exceeds
NewBad1 by factors of $2.9$ to $10.4$; IdentityPreservation is at
least $0.9987$ everywhere.

\paragraph{Temporal fidelity} (grid DTCE, disparity,
$k \in \{1,2,4,8\}$) improves on \textbf{20 of 20} D7 cells, between
$-2.1\%$ and $-26.2\%$, while on D2 it worsens slightly on four of
five backbones. That asymmetry is an asset rather than a weakness: on
D2 geometry improves while grid temporal change does not, which
directly supports the claim that temporal smoothness is not temporal
correctness. Grid DTCE is not motion-compensated and must be labelled
as such.

\paragraph{Open check.}
\texttt{TEPE\_corr} and \texttt{OPW}, both described in the evaluation
framework, were not found in the produced artefacts. They must either
be located or removed from the framework description.

\paragraph{Suspected field bug.}
\texttt{FrameDegradation/CatastrophicFraction} equals
\texttt{PositiveFraction} to five decimals on all ten SCARED-C rows.
It should not be used until this is resolved.

% =====================================================================
\section*{Spread of the D2 Mean, and an Arithmetic Correction}
% =====================================================================

On D2 the pooled gain is $4.05\%$, the mean of per-backbone
percentages is $3.94\%$, the median is $2.87\%$ and the range is
$2.17$--$8.44\%$; excluding Stereo Anywhere the mean falls to
$2.82\%$. On D7 the pooled gain is $4.64\%$, the median is $4.05\%$
and the range is $3.14$--$6.85\%$; excluding Stereo Anywhere it is
$4.14\%$, with RAFT-Stereo alone at $6.10\%$.

\textbf{The held-out result does not depend on the outlier.} The paper
should report median and range, and state explicitly that $4.05\%$ is
pooled rather than a mean of per-backbone percentages.

\paragraph{Correction.}
Section~V-E states that seen backbones improve by $4.65\%$ on D2. The
mean of the per-backbone seen gains is
$(2.1737 + 2.8705 + 8.4412)/3 = \mathbf{4.4951\%}$. The unseen figure
of $3.11\%$ is correct, so the method is right and only the number is
wrong; $4.65\%$ appears to have been transcribed from the D7 pooled
value of $4.639\%$.

% =====================================================================
\section*{Name Collision}
% =====================================================================

``STEER'' is already used in robotics: \emph{STEER: Flexible Robotic
Manipulation via Dense Language Grounding}, arXiv:2411.03409, November
2024, which targets the same venue population. No collision was found
in the stereo or depth literature. The acronym is also never expanded
in our own text. The name should be changed or at least expanded.

% =====================================================================
\section*{Hardening Pass: Closing the Submission TODOs}
% =====================================================================

Started 2026-08-15 21:30. Goal: make every remaining submission-critical
claim defensible rather than promised.

\paragraph{Seed variance, launched.}
The held-out table reported seed~0 only, which is the single largest
rejection risk given per-backbone gains of 2--3\%. Seeds 1 and 2 are now
training, one per H100, in a single \texttt{p1i} job.

\texttt{train.py} hard-fails unless the locked provenance matches the
canonical configuration byte for byte, including \texttt{seed: 0}. That
guard is correct and was left in place. Instead
\texttt{model\_design/train\_seed.py} imports \texttt{train.py} and
reuses its \texttt{metadata()}, \texttt{verify\_sources()},
\texttt{validate\_output()} and \texttt{config()} unchanged, overriding
exactly one field. It refuses any seed not listed in the
pre-registration, and writes a \texttt{seed\_provenance.json} recording
the deviation.

Pre-registered in \texttt{model\_design/multiseed\_preregister.json}
\emph{before} launch, declaring the primary endpoint, the resampling
unit (sequences, never pixels), and four prohibitions: no seed
selection, no post-hoc threshold retuning, no change to the
checkpoint-selection rule, and no use of D7 for any decision. The
canonical seed-0 checkpoint remains the reported model; this is a
variance estimate, not a selection experiment.

\paragraph{CatastrophicFraction: definitional, not a bug.}
Resolved. \texttt{PositiveFraction} is $\Pr(D_t>0)$ and
\texttt{CatastrophicFraction} is $\Pr(D_t>\tau_{\mathrm{cat}})$, with
\texttt{catastrophic\_px} and \texttt{catastrophic\_mm} both defaulting
to $0.0$ in \texttt{MetricConfig}. The evaluators construct
\texttt{MetricConfig()} with defaults, so the two coincide by
construction. The code is correct; the threshold was simply never set.
The paper now reports the degraded-frame fraction with its definition
stated, and omits CatastrophicFraction.

\paragraph{TEPE\_corr and OPW: implemented, never executed.}
Both exist in \texttt{unified\_metrics.py} (lines 834--881) but are
gated on a \texttt{flow} argument that no caller of
\texttt{evaluate\_argos\_prediction} supplies. They were therefore never
computed. This is an opportunity rather than a defect: the driver
already computes the SEA-RAFT flows used for warping, so passing them
would upgrade the multi-horizon temporal section from a disclaimed
grid diagnostic to genuine motion-compensated fidelity.

\paragraph{Degenerate D4D controls: the metric is won by doing nothing.}
Two controls were added in
\texttt{model\_design/comparison/degenerate\_policies.py}, alongside the
pinned \texttt{experimental\_policies.py} rather than inside it.
Over 504 frames the raw inconsistency of $3.7182$ becomes $1.7128$ under
STEER ($53.9\%$), $0.0986$ under copy-previous ($97.3\%$), and
\emph{exactly} $0.0000$ under warped-previous ($100\%$).

Warped-previous is zero by construction: its output \emph{is} the
motion-compensated previous frame, so it is trivially consistent with
the motion-compensated previous frame. The ordering on this metric is
therefore inverted with respect to geometric quality --- the two
policies that discard current stereo evidence beat the one that uses it.
The paper now reports the controls next to the metric, so no reduction
figure can be misread as accuracy.

\paragraph{TEPE\_corr and OPW now computed.}
\texttt{scripts/evaluate\_temporal\_corrected.py} assembles the
\texttt{alignment\_by\_horizon} contract from the SEA-RAFT flows the
driver already produces. The convention was read off the consumer rather
than guessed: \texttt{flow\_t\_to\_tk[:, i]} is applied as
\texttt{warp\_with\_flow(x[:, i+k], flow)}, so it lives on frame $i$ and
pulls $i+k$, which is the \emph{backward} member of the pair.
On a 40-frame smoke with S$^2$M$^2$-S, TEPE\_corr improves
$0.03030 \rightarrow 0.02454$ ($-19.0\%$) and OPW
$0.03288 \rightarrow 0.02676$ ($-18.6\%$). Full D2 and D7 runs are in
flight.

\paragraph{3D surface consistency, first numbers.}
\texttt{scripts/evaluate\_3d\_consistency.py} measures frame-to-frame
surface stability without camera pose, which SCARED-C does not provide:
each frame's disparity is back-projected to a metric cloud and compared
against the previous frame's cloud pulled through the same SEA-RAFT
correspondence. On a 12-frame smoke, point-to-point drift is
$0.4896$\,mm raw, $0.4471$\,mm refined and $0.4040$\,mm for GT, so the
\emph{excess over the GT floor} halves ($0.0855 \rightarrow 0.0431$\,mm).

One caveat found immediately: GT surface normals are \emph{noisier} than
both predictions ($11.83^\circ$ versus $5.04^\circ$ raw and
$4.24^\circ$ refined), because the pseudo-GT is sparse and central
differences amplify it. The GT-floor interpretation is therefore valid
for the point-to-point metric and \emph{not} for the normal metric,
which must be reported without it.

\paragraph{Stale TODO closed.}
The request to extend DRENDS to three further sequences was superseded
by the five-recording, three-backbone run already reported.

% =====================================================================
\section*{Open Items}
% =====================================================================

\begin{itemize}[itemsep=2pt]
    \item \textbf{Requires training}, not started, and must be
    pre-registered: canonical seeds 1 and 2 with paired bootstrap;
    evidence-channel ablation; head-structure ablation (reset-only,
    fusion-only, single versus dual resolution); convex-constraint
    ablation.
    \item \textbf{Inference only}: degenerate D4D baselines
    (copy-previous, warped-previous); a common-support
    Raw / BiDAStabilizer / STEER table; a 3D consistency experiment;
    and a re-run of \texttt{joint\_d4d} under the masked module.
    \item \textbf{Editorial}: Eq.~(53) is truncated; the Fig.~1
    caption overlaps the diagram; the qualitative figure is still a
    placeholder; BUR should be added to Table~III
    ($0.5074$, $0.5065$, $0.5047$ for $H = 4, 6, 8$, falling as HUR
    rises).
\end{itemize}


\section*{Night of 15--16 August: The Comparison That Does Not Exist}

The standing objection was the sharpest one available to a reviewer: TETHER
is never compared to another method. We went looking for one. The honest
answer is that in our exact category --- causal, black-box, attached to a
frozen surgical stereo network --- no runnable published method exists, and
we can now say so with receipts rather than with a shrug.

\paragraph{StereoDiffusion (MICCAI 2024).} The only paper in our category:
causal, black-box, frozen stereo, surgical, evaluated on SCARED. Its
repository has a single commit, $7.43$\,KiB of git objects, and contains a
\texttt{.gitignore}, a \texttt{LICENSE} and a twelve-byte \texttt{README.md}.
No code, no weights, no other branch. There is nothing to run.

\paragraph{Neural Disparity Refinement (3DV 2021 / TPAMI 2024).} Same input
contract as ours, spatial only. The code clones fine; the weights do not
exist any more. Both released Google Drive IDs return HTTP 404 from the
download endpoint \emph{and} from the plain file-view page --- a restricted
file answers 403 or an interstitial, so 404 on the view page means deleted.
\texttt{docs.google.com/uc} agrees. Hugging Face has no mirror; the GitHub
releases API returns an empty list. Our network is not the obstacle: the same
\texttt{curl} reached huggingface.co and archive.org in the same session.

\paragraph{Lai et al.\ (ECCV 2018).} The natural task-agnostic baseline. Its
download script points at \texttt{vllab.ucmerced.edu}, which answers $200$ for
the directory and $404$ for every model file. The Wayback CDX index has
captures of the project page, its stylesheets and the $34$\,MB supplementary
PDF, and zero captures of \texttt{ECCV18\_blind\_consistency.pth} --- crawlers
archive HTML and skip large binaries. No Hugging Face mirror, no GitHub
release. The hosting rotted and nothing caught it.

\paragraph{TC-Stereo (ECCV 2024), and what we did about it.} This one loads.
We found a configuration that matches all three released checkpoints with zero
missing and zero unexpected keys, and it imports cleanly in our environment.
The blocker is scientific, not technical. Its temporal propagation warps the
previous disparity by a rigid $6$-DoF reprojection built from per-frame camera
pose, intrinsics and baseline. SCARED-C supplies no per-frame pose, and
deformable tissue violates the rigid-scene assumption regardless --- so its
temporal path is not merely unavailable on our data, it is the wrong model of
the scene. Running it with \texttt{params=None} is possible but degenerates it
to a single-frame stereo network.

So we ran exactly that, and labelled it honestly.
\texttt{scripts/run\_tcstereo\_reference.py} drives TC-Stereo over the
\emph{same} frozen NPZ boundary the BiDAStabilizer comparison consumed, so the
frames, the crops and the ground truth are bit-identical across methods. The
row it produces is an integrated-architecture reference --- $16.7$\,M
parameters, trained end to end, evaluated zero-shot --- not a temporal rival,
and the manifest records the disabled temporal path and the reason. It answers
the one question our own ablations cannot: whether one should simply use a
newer integrated network instead of refining a frozen one.

\paragraph{What went into the paper.} A new Related Work paragraph, ``The
comparison that does not exist'', with five citations. It cost a page, and the
page came back out of the external-interface, support-contract, closure and
evaluation subsections without dropping a single result or table. The
submission is at eight pages including the bibliography, which is the hard
constraint.

\paragraph{Seed variance, pre-registered.} \texttt{scripts/aggregate\_seeds.py}
implements the reporting rule written down before any seed result existed:
mean and standard deviation across seeds $0,1,2$, plus a paired bootstrap in
which \emph{sequences}, never pixels, are the resampling unit. Pixels inside
one endoscopic sequence are not independent samples and resampling them would
manufacture significance. Seeds are averaged inside each resample, so seed
choice is part of the method rather than a free parameter --- which is what the
pre-registered prohibition on selecting the best seed actually requires. Run on
seed $0$ alone as a plumbing check, the pooled D2$+$D7 Bad1 reduction is
$5.79\%$ with a $95\%$ interval of $[4.15, 8.35]$ and no resample in which
refinement failed to help. The real three-seed number replaces it when seeds
$1$ and $2$ finish.



\section*{Seed 1 Lands, and Says Something About Checkpoint Selection}

Seed 1 finished its D2 and D7 evaluation through the frozen support-masked
path. Pairing it with seed 0 under the pre-registered rule --- sequences as
the resampling unit, seeds averaged inside each resample --- gives the
following relative Bad1 reductions.

\begin{center}
\footnotesize
\begin{tabular}{@{}llrrrr@{}}
\toprule
Split & Backbone & mean\% & std\% & seed 0 & seed 1 \\
\midrule
D2 & CREStereo             &  7.44 & 1.17 &  8.26 &  6.62 \\
D2 & Fast-FoundationStereo &  6.83 & 1.53 &  7.91 &  5.74 \\
D2 & RAFT-Stereo           &  6.17 & 1.10 &  6.95 &  5.40 \\
D2 & S2M2-S                &  7.05 & 0.50 &  7.40 &  6.70 \\
D2 & StereoAnywhere        & 17.80 & 2.62 & 19.65 & 15.95 \\
\midrule
D7 & CREStereo             &  4.18 & 0.46 &  3.86 &  4.51 \\
D7 & Fast-FoundationStereo &  4.02 & 0.08 &  3.97 &  4.08 \\
D7 & RAFT-Stereo           &  4.65 & 0.99 &  3.95 &  5.35 \\
D7 & S2M2-S                &  5.96 & 0.42 &  5.66 &  6.25 \\
D7 & StereoAnywhere        &  6.73 & 0.47 &  6.40 &  7.07 \\
\midrule
\multicolumn{2}{@{}l}{Pooled, 35 sequences} & 5.90 & 0.15 & 5.79 & 6.00 \\
\bottomrule
\end{tabular}
\end{center}

Every cell is a reduction on both seeds. The pooled standard deviation is
$0.15$ percentage points on a $5.90$ point effect, and the paired
sequence-level bootstrap gives $[4.26, 8.38]$ with no resample in which
refinement failed to help.

The interesting part is the sign pattern, which is not noise. Seed 1 is
worse than seed 0 on \emph{all five} D2 backbones and better than seed 0 on
\emph{all five} D7 backbones. Ten out of ten is not a coin flip. The
explanation is in our own protocol: the checkpoint-selection rule picks the
best development (D2) fused EPE, so seed 0 is by construction the run that
looked best on D2. Its D2 margin is partly a selection artefact, and the
held-out split --- the one that carries the claim --- is where the other seed
does slightly better.

This is a point in favour of the result rather than against it. The number
we report on held-out data is not the one selection inflated, and it is
stable across an independent training run. Seed 2 is in flight; the reported
figure will be the three-seed mean with this same bootstrap, and this
two-seed table stands as the record of what was known first.



\section*{Staleness Audit: Do the Tables Still Match the Files?}

The support-mask fix changed enough numbers that a table could plausibly
have been left behind. While the GPU jobs ran we re-derived the two largest
tables straight from the report JSONs instead of trusting the manuscript.

Table~\ref{tab:scared_main} (cross-backbone SCARED-C) reproduces exactly
from \texttt{results/scared\_masked}: S$^2$M$^2$-S Bad1
$.01394\!\rightarrow\!.01291$ on D2 and $.06440\!\rightarrow\!.06076$ on D7,
RAFT-Stereo $.01603\!\rightarrow\!.01492$ and
$.06085\!\rightarrow\!.05845$. No drift.

Table~\ref{tab:drends} (zero-shot DRENDS) reproduces exactly from
\texttt{results/drends\_masked}, all fifteen recording--backbone cells, and
so do the four pooled means quoted underneath it: $-5.00\%$ EPE,
$-17.35\%$ Bad3, $-3.68\%$ depth MAE, $-3.21\%$ depth RMSE, with $60/60$
cells improving.

One trap worth writing down, because it nearly caught us.
\texttt{results/drends\_multisequence} still exists and still contains four
RAFT-Stereo reports whose numbers ($-3.9\%$ EPE on Vid10 against the
paper's $-5.7\%$) look like a contradiction. It is not: that directory is
the pre-mask, single-backbone, four-recording run, superseded by
\texttt{drends\_masked}. Two result directories that differ by a bug fix are
indistinguishable by filename, and the older one sorts first. The manifests
disambiguate them; directory names do not.



\section*{Three Seeds, and a Selection Effect We Can Measure}

Seed 2 landed. The full pre-registered analysis now runs on seeds
$0,1,2$: pooled Bad1 reduction $5.69\pm0.37\%$ over 35 sequences, paired
sequence-level bootstrap $[4.11,8.02]$, $p=.000$; pooled EPE
$4.62\pm0.18\%$, $[3.42,5.80]$.

Split apart, the two halves behave completely differently:

\begin{center}
\footnotesize
\begin{tabular}{@{}lrrrr@{}}
\toprule
& seed 0 & seed 1 & seed 2 & spread \\
\midrule
D2 Bad1 & 10.78 & 8.67 & 5.75 & 5.03 \\
D7 Bad1 &  4.78 & 5.47 & 5.19 & 0.69 \\
D2 EPE  &  4.05 & 3.43 & 3.17 & 0.88 \\
D7 EPE  &  4.57 & 5.25 & 5.38 & 0.81 \\
\bottomrule
\end{tabular}
\end{center}

On D2 seed 0 is best and the sequence is monotone. On D7 seed 0 is
\emph{worst} on both metrics. Ten of ten per-backbone comparisons agree.
This is the signature of our own checkpoint-selection rule, which picks
best D2 fused EPE: seed 0 is by construction the run that looked best on
the split it was selected on. The development margin carries selection
optimism; the held-out margin, which is what the paper claims, does not.

We could have buried this. Reporting it is better: it converts a
reviewer's likely objection --- ``your development gains look
suspiciously large'' --- into a measured quantity with the right sign.

\paragraph{What it cost the manuscript.} Adding the subsection pushed the
paper to nine pages. Getting back to eight took four passes: merging
``What Makes TETHER Work'' with the $177$k-parameter subsection,
compressing the D4D degenerate section into prose with its table moved to
the supplementary, tightening Cross-Backbone Development Behaviour, and
finally dropping the seed table itself in favour of four inline
mean$\pm$std figures. The last one was the right call and should have
been the first: the table held twelve numbers the prose already stated.

One stale sentence died in the process, which is the point of doing this
before submission. Discussion still read ``held-out numbers are a single
canonical seed; sequence-aware confidence intervals over additional seeds
are required before the per-backbone margins of two to three percent can
be called separated.'' Half of that is now false. It has been replaced by
what is still true: the three-seed study bounds the pooled effect, so
individual backbone rows remain unseparated from one another.

The full per-backbone seed table, including the three D2 rows whose
intervals cross zero, is in the supplementary. Those wide intervals are a
property of a three-sequence split, not of the method, and we prefer them
to the narrow and meaningless intervals that resampling pixels would have
produced.



\paragraph{Audit completed: all four main tables.} Continuing the staleness
pass while the GPU jobs ran, Table~\ref{tab:intervention_quality}
(intervention quality) also reproduces from
\texttt{results/scared\_masked}: recomputing $B^+/H^+$ and
RecoveredBad1/NewBad1 per split and backbone lands within $0.06$ of every
printed cell, the residual being the difference between a ratio of
sequence means and the framework's macro-sequence aggregation rather than
an error. The prose claim underneath it --- identity preservation above
$99.87\%$ in every case --- is correct and tight: the minimum over all ten
split--backbone cells is $99.877\%$, on RAFT-Stereo D7.

The closure table needed a note rather than a correction. Its source is
\texttt{definitive\_evaluation/canonical\_h4\_baseline}, the
\emph{unconfined} module, while every other table in the paper is the
confined one. The tempting fix --- update the $H=4$ row to the confined
number --- would have been wrong: the other ten rows ($H=6$, $H=1$, the
fixed blends, the oracle) are all unconfined ablations, so changing one
row would have compared two implementations inside a table whose entire
purpose is that only the policy differs. Confinement moves D2 EPE by
$1.2\times10^{-5}$\,px, so the honest fix was to say in the text that the
whole table runs unconfined for exactly that reason. Raw D2 pooled EPE is
$0.276202$ in both runs and the reduction is $4.05\%$ in both, so no other
figure in the paper depends on the choice.



\section*{The Comparison Exists Now}

BiDAStabilizer finished on all three D2 sequences at common support. This
is the number the whole night was for.

\begin{center}
\footnotesize
\begin{tabular}{@{}lcrrrr@{}}
\toprule
Method & Causal & EPE & Bad1 & Bad3 & RMSE \\
\midrule
Raw RAFT-Stereo & --- & .30286 & 2.087 & .186 & .44842 \\
BiDAStabilizer  & no  & .30211 & 1.987 & .171 & .44026 \\
TETHER          & yes & \textbf{.28899} & \textbf{1.826} & \textbf{.157} & \textbf{.42562} \\
\bottomrule
\end{tabular}
\end{center}

Same $4249$ frames, same frozen RAFT-Stereo robust predictions, same
prediction-independent support, $18.0$M pixels. TETHER wins every metric
and is the only causal entry: $-4.58\%$ EPE, $-12.53\%$ Bad1, $-15.52\%$
Bad3 against raw, where the bidirectional baseline manages $-0.25\%$,
$-4.83\%$, $-8.32\%$.

Worth remembering how close this came to being meaningless. The first
comparison we had was not a comparison: the published reproduction ran
RAFT-Stereo \emph{robust} and we ran \emph{middlebury}, and the two were
scored under different support contracts --- raw EPE $0.3029$ against
$0.2762$. Subtracting those would have produced a confident number with
no content. The fix was to stop generating our own inputs and read the
stored NPZ boundary the reproduction had already written, so the RGB and
the disparity BiDAStabilizer consumed are the same arrays ours consumes.

\paragraph{The cell we do not like.} One sequence of three favours
BiDAStabilizer on all three metrics, and on it our EPE is $0.27\%$
\emph{worse} than raw. It goes in the paper. With three sequences a
single sequence-level disagreement is well inside what the split can
produce, and a method allowed to look ahead ought to win somewhere; the
claim the table supports is the pooled one. Hiding it would have been
both dishonest and, given that a reviewer can run the same three
sequences, stupid.

\paragraph{Page budget, again.} The subsection plus table cost a full
page. Recovered by: deleting the CatastrophicFraction derivation from the
main text, where it duplicated the supplementary verbatim; folding the
``Selective risk'' capability list into one sentence; compressing D4D a
second time; and tightening the $H=4$, recurrence and cross-backbone
prose. Nothing was dropped that is not still stated somewhere. Eight
pages, comparison table included.



\section*{TC-Stereo Runs, and We Argue Against Our Own Number}

The last job finished. TC-Stereo, single-frame path, SceneFlow weights,
over the same boundary as everything else: EPE $0.5242$, Bad1 $9.05\%$,
against raw RAFT-Stereo robust at $0.3029$ and $2.087\%$. Seventy-three
percent worse on EPE, four times worse on Bad1.

The tempting move is to put that in the main paper as a win. It is not
one, and the reason is worth recording.

The stored boundary is $144\times180$ with a median ground-truth disparity
of $8.04$\,px, range $5.57$--$15.53$. TC-Stereo was trained at full
resolution on disparity ranges an order of magnitude larger. We are asking
a from-scratch cost-volume matcher to work two octaves below its training
regime. It still functions --- $91\%$ of pixels land within $1$\,px, so
this is a working network producing worse geometry, not a broken
configuration or a scale bug --- but the comparison is not fair to the
architecture. A refiner is far less exposed to the same problem: both
BiDAStabilizer and TETHER consume the raw disparity and correct it, so
boundary resolution constrains them much less than it constrains a
correlation pyramid.

We kept the boundary rather than re-running TC-Stereo at full resolution,
because doing that would buy fairness to one method by destroying the
property the entire table rests on: that every method sees identical
frames and identical ground truth. So the row goes in the supplementary,
with the caveat attached, supporting the claim it can support --- swapping
in a different integrated architecture zero-shot is not a drop-in
substitute for refining a strong in-domain frozen model --- and explicitly
not the claim it cannot: that TETHER outperforms TC-Stereo.

\paragraph{A design that turned out vacuous.} The runner scored two
supports, ground truth alone and ground truth intersected with raw
validity, on the theory that a method not consuming the backbone should
not be charged for the backbone's refusals. On this boundary
\texttt{raw\_valid} is true everywhere, so both conditions select the same
$17{,}999{,}475$ pixels and the two tables are identical. The distinction
was worth building --- it would matter on a boundary where the backbone
declines pixels --- but here it bought nothing, and the supplementary says
so rather than presenting two columns that look like independent
confirmation.



\section*{16 August, Afternoon: The Ablations Answer, and One Answer Is No}

Four architecture ablations were pre-registered this morning and launched
on the two H100s. Two have finished training and evaluation, one is
training, one is being evaluated. Two results are already worth writing
down, and the first one costs us a sentence in the paper.

\paragraph{A4: the convexity claim is withdrawn.} The Analysis said that
confining the output to the interval between the two candidate geometries
was ``a plausible reason why one small head transfers to architectures
absent from its training''. That was an assertion with no experiment
behind it, so we built one: $49$ parameters, a zero-initialised additive
term that lets the output leave the interval. Zero-initialised so training
starts exactly at the canonical model and any departure has to be learned.

The pre-registration named Bad1 on held-out D7 as the primary endpoint and
committed, in advance, to \emph{withdrawing} rather than softening the
claim if the relaxed variant matched or beat canonical there.

It beat it. On D7: Bad1 $6.02\%$ against $4.78\%$, and also better on
Bad3, RMSE and worst-case error. The claim is withdrawn.

What the experiment actually shows is not what we expected, and is more
specific. Mean EPE gain collapses from $4.57\%$ to $0.05\%$, turning
\emph{negative} on three of five backbones --- the module makes disparity
worse than doing nothing. Our first reading was the obvious one,
occasional large failures, and it is wrong: RMSE and the worst case both
\emph{improve}, and a quadratic metric would be the first to register
outliers. The only consistent reading is the opposite. The relaxed variant
fixes the few large errors and degrades a great many pixels that were
already close. A linear mean registers that; thresholds and squared errors
do not.

So the convex constraint does not buy safety from catastrophic
corrections. It buys mean accuracy, by stopping the module from disturbing
the majority of pixels the frozen estimator already had right. Narrower
than what we claimed, and now supported.

\paragraph{A1: forty-five percent of the input is redundant.} The evidence
space carries $64$ frozen ResNet-18 channels out of $142$. But appearance
self-correlation and cross-appearance correlation are computed from those
same features --- so the real question was never ``do appearance features
help'' but ``does raw appearance add anything \emph{beyond} the
correlations already derived from it''.

On held-out D7, dropping them changes nothing: Bad1 $4.88\%$ against
$4.78\%$, a difference of $0.10$ points against a three-seed spread of
$0.37$. EPE, Bad3 and RMSE are all marginally better. By the
pre-registered rule this is \emph{indistinguishable}, and the
pre-registration also said what to do about it: describe the channels as
redundant.

One asymmetry is worth noting because it is the same one the seed study
found. On D2, A1 is clearly worse on Bad1 ($8.81\%$ against $10.78\%$),
and D2 is the split the checkpoint is selected on. Development says the
channels matter; held-out says they do not.

\paragraph{What this does not buy.} A1 shrinks the first convolution from
$30{,}672$ to $16{,}848$ parameters, about $8\%$ of the head, and does
\emph{not} remove the frozen extractor from the runtime --- the appearance
correlations still need it. Any latency argument belongs to A2, which
removes the learned evidence entirely and is in evaluation now. A3, the
single-resolution variant, is training: it keeps the parameter count
exactly identical and moves only the receptive field, because deleting the
context branch outright would have confounded ``context needs low
resolution'' with ``the head needs capacity''.



\section*{A2: The Evidence Space Is Over-Specified}

A2 removes \emph{all} learned stereo evidence --- the confidence curves and
their supports, both appearance correlations, and the $64$ frozen
ResNet-18 channels --- leaving $38$ channels of disparity, motion and RGB
geometry. Verified on the checkpoint: \texttt{cue\_channels} is $38$ and
the head is $154{,}874$ parameters instead of $177{,}338$. Crucially it
is also the only variant that removes the frozen extractor from the
runtime, since nothing left in the stack needs it.

We expected this to be the ablation that hurt. On held-out D7 it is
\emph{better than canonical on every metric}:

\begin{center}
\footnotesize
\begin{tabular}{@{}lrrrr@{}}
\toprule
D7 held out & EPE & Bad1 & Bad3 & RMSE \\
\midrule
Canonical, $142$ ch & 4.57 & 4.78 & 4.07 & 5.33 \\
A1, $78$ ch         & 4.93 & 4.88 & 4.12 & 5.46 \\
A2, $\mathbf{38}$ ch & \textbf{5.49} & \textbf{5.69} & \textbf{4.85} & \textbf{5.92} \\
\bottomrule
\end{tabular}
\end{center}

Its Bad1 margin is $+0.91$ points against a three-seed spread of $0.37$,
so by the pre-registered rule this is \emph{better}, not indistinguishable.

The three ablations line up into one pattern rather than three
independent results. A1 removes the raw appearance channels but keeps the
correlations computed from them: neutral. A2 removes the correlations too:
better. The more appearance-derived evidence leaves the interface, the
better the head generalises to backbones and sequences it never saw. That
is the signature of evidence that ties the decision to the training
domain --- useful where it was fitted, a liability elsewhere.

\paragraph{What this does and does not cost the paper.} It does not touch
the interface idea. A $38$-channel external evidence space is still an
external evidence space, still carries no ground truth, no backbone
identity and no future frame, and still lets one head serve heterogeneous
estimators. What fails is the specific richness of our instantiation: we
built $142$ channels and $104$ of them are, on held-out data, worse than
useless.

Two honest qualifications. On D2 --- the split the checkpoint is selected
on --- A2's EPE is \emph{worse} than canonical ($3.69$ against $4.05$),
which is the third time this development-versus-held-out inversion has
appeared today. And each ablation is a single seed, while the $0.37$ point
yardstick comes from three seeds of the canonical model; a margin of
$0.91$ clears it but not by an order of magnitude.

\paragraph{The uncomfortable implication.} If A2 holds up, the canonical
architecture we are submitting is not the best one we have measured. The
honest options are to report the ablation as it is and say so plainly, or
to re-run the seed study on the $38$-channel variant before claiming
anything stronger. We are not going to quietly promote A2 to canonical
after the fact --- the pre-registration prohibits exactly that, and
selecting the best variant post hoc is how the selection effect we spent
today measuring gets back in through the side door.



\section*{Correction: We Compared Against the Wrong Baseline}

Everything written above about A1, A2 and A4 used
\texttt{results/scared\_masked} as the reference. That is canonical
\emph{seed 0}. And the seed study, run this same morning, established that
seed 0 is the \emph{weakest} of the three canonical seeds on held-out D7.
So every ablation was being measured against the luckiest possible
comparison, and every one of them looked better than it is.

Against the three-seed mean, which is the reference the pre-registration
actually named:

\begin{center}
\footnotesize
\begin{tabular}{@{}lrrrr@{}}
\toprule
& D2 EPE & D2 Bad1 & D7 EPE & D7 Bad1 \\
\midrule
Canonical, 3 seeds & 3.55 & 8.40 & 5.06 & 5.15 \\
A1 & 4.18 & 8.81 & 4.93 & 4.88 \\
A2 & 3.69 & 12.84 & 5.49 & 5.69 \\
A4 & 5.53 & 12.44 & 0.05 & 6.02 \\
\bottomrule
\end{tabular}
\end{center}

Two statements from earlier in this diary are wrong and are withdrawn.
``A2's D2 EPE is worse than canonical'' compared $3.69$ against seed 0's
$4.05$; against the three-seed mean of $3.55$ it is better. And the
margins quoted for A2 and A4 on held-out Bad1 ($+0.91$ and $+1.23$) are
really $+0.55$ and $+0.87$.

The script now enforces the right baseline, prints the individual seed
values, and refuses to call a one-seed margin a result when it falls
inside the canonical seed range.

\paragraph{And a decision that was rejected too quickly.} Asked whether
the better variant should simply replace the canonical model, the answer
given was a flat no, on the grounds that choosing it would burn the
held-out split. That reasoning is correct and it is not the whole
question. A2 is outside the canonical seed range on \emph{both} splits:
$+1.8$ standard deviations on D2 Bad1 and $+1.6$ on D7. D2 is the
development split --- deciding on it is what it exists for, and the
pre-registration prohibits using D7 for decisions, not D2.

So there is a clean path, and refusing it would have been caution
mistaken for rigour. A2 is being trained at seeds 1 and 2, giving it the
same three-seed treatment the canonical model has. The comparison that
may promote it is the three-seed \emph{D2} mean. D7 is not consulted for
that choice and stays held out; if A2 is promoted, its held-out numbers
are reported afterwards and are not part of the decision. Declared in
advance: if A2's three-seed D2 mean does not clearly exceed $8.40$ with a
comparable spread, A2 stays an ablation and the canonical model is
submitted unchanged.



\section*{All Four Ablations: Nothing We Designed Is Load-Bearing}

A3 landed, completing the set. Before reading anything into the numbers
we checked the one thing that would invalidate the comparison: the raw
baseline. It is identical to six decimals across all six runs --- the
three canonical seeds and the four variants --- at $0.062051$ on D7 Bad1
over twenty reports each. The comparison is apples to apples.

\begin{center}
\footnotesize
\begin{tabular}{@{}llrr@{}}
\toprule
& Change & D7 EPE & D7 Bad1 \\
\midrule
Canonical (3 seeds) & --- & 5.06 & $5.15\pm0.34$ \\
A1 & $142\rightarrow78$ ch & 4.93 & 4.88 \\
A2 & $142\rightarrow38$ ch & 5.49 & 5.69 \\
A3 & full-resolution context & 5.61 & 5.85 \\
A4 & unconstrained output & 0.05 & 6.02 \\
\bottomrule
\end{tabular}
\end{center}

A1 is indistinguishable. A2, A3 and A4 sit $1.6$, $2.0$ and $2.5$
standard deviations above the canonical mean, and A2 and A3 beat it on
EPE as well. Every design decision we made is neutral or worse than its
alternative.

\paragraph{Why this is suspicious, and what survives the suspicion.}
Three variants, each changing a \emph{different} thing, all landing above
the mean on the same side is not what seed noise looks like --- noise
scatters. The reading that fits is that the canonical design is
over-specified, and that simplifying it in almost any direction helps
generalisation. What does \emph{not} survive is any claim that a
particular one of our choices is responsible: each variant is one
training run against a three-seed baseline, and two standard deviations
from a single sample is a hint, not a finding.

A4 is set aside on its own terms rather than on that caution. It has the
best held-out Bad1 in the table and an EPE gain of $0.05\%$: it fixes
threshold crossings by disturbing everything else, which is not a
trade we want.

\paragraph{A3 corrects the Method, not just the Analysis.} The paper
justified predicting the fusion weight at quarter resolution on the
reasoning that deciding \emph{how much} to intervene needs more context
than deciding \emph{whether}. A3 keeps the parameter count at exactly
$177{,}338$ and moves that branch to full resolution --- and does not do
worse. So the dual-resolution design is a compute choice, roughly
$16\times$ cheaper, and the Method now says that instead of implying a
necessity. This is the ablation the agentic review asked for, and it
answered in the opposite direction to our expectation.

\paragraph{What went into the paper.} A four-row table, generated by
\texttt{scripts/ablation\_table.py} rather than transcribed, and a
section that states plainly that none of the four decisions is
load-bearing. It cost the per-recording DRENDS table --- whose $60/60$
claim and Vid10 example survive in prose --- and four citations that were
redundant inside groups keeping a representative each. Eight pages,
twenty-seven references.

A2 is now training at seeds $1$ and $2$. The question it answers is
whether the $38$-channel variant is genuinely better or one lucky run,
and by pre-registration that question is settled on D2 alone.



\section*{Night Audit: One Number Was Computed on a Subset}

With the ablations queued and nothing to collect, the night went into
checking every figure the paper quotes against the files it came from.
Four checks, one real defect.

\paragraph{The defect.} The paper claimed TETHER captures $32.9\%$ of the
raw-versus-memory oracle gain and $30.1\%$ of the continuous one. Neither
reproduces from \texttt{results/oracle\_risk\_analysis}. Recomputing
pixel-weighted, macro over sequences, and as a ratio of summed gains all
give the same pair: $37.2\%$ and $34.5\%$.

The published numbers turn out to match exactly one computation ---
dropping StereoAnywhere, which is the outlier at $50.9\%$ against
$31$--$33\%$ for the other four. That may well have been a deliberate
robustness check at the time. It was quoted without the exclusion, which
makes it wrong as written.

Corrected in both documents, along with the bracketed fraction from
$6.1\%$ to $6.2\%$. Worth recording that the correct figure is the
\emph{better} one: we were under-reporting our own result by four points.
Errors of this kind are usually assumed to run the other way, which is
probably why nobody re-derived it.

\paragraph{The three that held.} D4D reproduces exactly --- TETHER
$53.93\%$, copy-previous $97.35\%$, warped-previous $100.00\%$ against
the $53.9$, $97.3$ and $100.0$ printed. Runtime reproduces exactly:
$33.60$\,ms overhead, $124.2$\,MB peak VRAM, and end-to-end rates of
$10.98$ and $1.38$\,FPS with overhead fractions of $36.9\%$ and $4.6\%$
at the two extremes. And every cross-document promise now has a section
behind it: the fifteen closure policies, the DRENDS per-recording
figures, the per-backbone intervention table, the ablation D2 figures,
and --- added tonight --- the horizon sweep's risk columns, which the
main paper pointed at and the supplementary did not carry.

\paragraph{Method.} The audit is worth doing in this direction:
recompute from the stored reports and compare to the manuscript, rather
than reading the manuscript and asking whether it looks right. Every
number here looked right. The one that was wrong looked right too.



\section*{A2 Wins the Pre-Registered Comparison, and We Submit the Other Model}

Three seeds of A2 finished overnight. On D2 --- the split the
pre-registration named as the decision split, precisely so that a
promotion could be made without touching held-out data:

\begin{center}
\footnotesize
\begin{tabular}{@{}lcccc@{}}
\toprule
D2 Bad1 reduction & seeds & mean & std \\
\midrule
Canonical, $142$ ch & 10.78, 8.67, 5.75 & 8.40 & 2.53 \\
A2, $38$ ch & 12.84, 12.38, 12.98 & \textbf{12.73} & \textbf{0.32} \\
\bottomrule
\end{tabular}
\end{center}

The ranges do not overlap: A2's worst seed beats the canonical model's
best. The declared criterion --- clearly exceeding $8.40$ with a
comparable spread --- is met, and the spread is not comparable but eight
times tighter.

That variance is the more interesting half of the result. The canonical
model swings five points across seeds on D2, from $5.75$ to $10.78$; the
$38$-channel variant swings half a point. A head whose evidence is mostly
appearance-derived is apparently fitting the training domain in a way
that depends heavily on initialisation, and stripping that evidence
removes the dependence along with it.

Held out, reported after the decision and not part of it: A2 gives
$5.84\pm0.22\%$ against the canonical $5.15\pm0.34\%$.

\paragraph{And we are submitting the canonical model anyway.} Not out of
caution. Two reasons, and neither is that the result is doubtful.

First, promotion is not an edit. Every table in the paper --- DRENDS
zero-shot, the BiDAStabilizer common-support comparison, the transfer
matrix, the intervention ratios, the runtime budget --- was measured on
the canonical checkpoint. A2 has been evaluated on SCARED-C only.
Promoting it means re-running all of it, and a paper whose headline model
is supported by one split while its tables come from another is worse
than either option taken cleanly.

Second, and this is the part worth remembering: the comparison that
favours A2 was deliberately built on the development split so that the
held-out numbers would stay untouched by any choice of ours. Promoting on
the strength of it and then reporting D7 for the promoted model would
spend exactly the thing the design was protecting.

So the paper reports the ablation as it stands: a pre-registered variant
that beats the submitted model on the split where such comparisons are
allowed, with the numbers, the variance and the reason we did not act on
it. That is a finding, not an embarrassment --- and the honest place for
it is the paper rather than a private note about what we would have done
with another week.


\section*{17 August, Retracted Within the Hour: The Threefold DRENDS Win
Was One Sequence Against Five}

\emph{Everything in the first four paragraphs below is wrong. It is left
standing because the way it was wrong is the useful part, and because the
commit that carried it has been pushed. The retraction follows.}

The reason we did not promote A2 was that it had been measured on
SCARED-C only. So we measured it on DRENDS, the out-of-domain split,
three seeds against three backbones. The result reverses cleanly along an
axis we had not separated.

On SCARED, A2 is worse. Held-out D7 Bad1 reduction $5.84\%$ against the
canonical $5.15\%$ --- and broken out by backbone, A2 gives up
$0.75$ percentage points on the three the head was trained beside and
$0.61$ on the two it had never seen. Unseen \emph{backbone} does not
rescue it.

On DRENDS it is not close, and it is the same direction for every metric
and every backbone:

\begin{center}
\begin{tabular}{lcc}
\toprule
Reduction on DRENDS & Canonical & A2 \\
\midrule
EPE  & $4.5$--$5.4\%$  & $17.4$--$18.5\%$ \\
Bad1 & $4.7$--$5.6\%$  & $13.1$--$15.3\%$ \\
Bad3 & $12$--$16\%$    & $59$--$67\%$ \\
RMSE & $6.5$--$7.0\%$  & $23$--$28\%$ \\
\bottomrule
\end{tabular}
\end{center}

Three to four times the reduction, out of domain, from the variant that
loses in domain. The axis is not \emph{unseen sequence} and not
\emph{unseen backbone} --- A2 loses on both. It is \emph{unseen domain}.
The 64 frozen ResNet channels and the confidence machinery built on them
let the head fit the appearance of ex-vivo porcine tissue; on synthetic
colonoscopy that fit is a liability, and the variant that never had access
to appearance has nothing to mislead it. Geometry and motion transfer.
Appearance, even frozen appearance, does not.

\paragraph{One thing we caught before believing it.} A2's seed spread on
DRENDS is $0.001$ percentage points, against $0.34$ for the canonical
model on SCARED. A spread that small usually means the same checkpoint
was loaded three times. It was not: the three files hash differently and
their provenance records carry seeds $0$, $1$ and $2$. The spread is real
and is itself informative --- out of domain, three independently trained
A2 heads converge to almost identical behaviour, which suggests what
happens there is governed by the support and confinement machinery rather
than by the learned weights.

\paragraph{And the comparison is still not matched.} The seed study was
run on SCARED only, so the canonical model has exactly one DRENDS number
and A2 has three. Setting a three-seed mean against a one-seed baseline is
precisely the error recorded two entries above, where canonical seed $0$
turned out to be the weakest of the three and flattered every variant.
The table above is therefore the seed-$0$ against seed-$0$ pairing, which
is matched, and canonical seeds $1$ and $2$ are running on DRENDS now.
The margin is large enough that the conclusion is unlikely to move, which
is not a reason to skip the measurement.

Nothing in the paper has been changed on the strength of this. Promotion
still means re-running every table, and the decision waits for the matched
baseline.

\paragraph{Retraction.} None of that survives. The reader who saw the
table above and said it was impossible was right, and the reason is
embarrassingly simple: \textbf{A2 had been evaluated on one DRENDS
recording and the canonical model on five}. The comparison averaged
Vid14\_Pancreas\_High alone against the mean of Vid10 through Vid14, and
Vid14 is the easier of them. The threefold improvement was a statement
about which sequences were in each average, and none of it was about the
model.

Restricted to the one recording both sides actually cover, the two models
are the same to three decimal places:

\begin{center}
\begin{tabular}{lccc}
\toprule
Bad1 on Vid14 & Raw & Canonical & A2 \\
\midrule
CREStereo             & $0.499$ & $0.473$ & $0.473$ \\
Fast-FoundationStereo & $0.487$ & $0.465$ & $0.465$ \\
RAFT-Stereo           & $0.538$ & $0.512$ & $0.511$ \\
\bottomrule
\end{tabular}
\end{center}

Deltas of $0.001$ percentage points, in both directions. Indistinguishable.

\paragraph{Two defects, and the second is the worse one.} The first is in
\texttt{drends\_backbone\_transfer.py}, which read

\begin{center}
\texttt{selected = recordings or ["Vid14\_Pancreas\_High"]}
\end{center}

so a runner that omits \texttt{-{}-recordings} silently evaluates a single
sequence --- and the report it writes still labels itself
\texttt{complete\_recordings}, because that field only distinguishes a
frame cap from no frame cap. The canonical DRENDS run passed the flag; the
A2 runner we wrote did not. The default is now the full five, and the
constant carries the story.

The second defect is in the comparison script, and it is the one that
allowed the first to matter. It \emph{did} check that both sides shared
identical raw predictions before believing any delta --- and that check
passed, because it ran on the intersection of the two sequence sets, which
was one sequence, and on that sequence the raw predictions agreed
perfectly. The check then handed off to a mean computed over each side's
\emph{own} sequences. A guard that verifies the overlap and then reports
the union is not a guard. Both sides are now reduced over the shared set
only, and any coverage mismatch prints a warning before the table.

The earlier interpretation --- ``frozen appearance features fit ex-vivo
tissue and mislead out of domain'' --- was a tidy story invented to
explain an artefact. It explained the artefact well. That is exactly what
made it dangerous, and it is worth noting that the story arrived within
seconds of the number and never once prompted a check of what the number
was averaged over.

The A2 DRENDS evaluation is running again on all five recordings. The
canonical seed study on DRENDS is cancelled: with a delta of $0.001$
percentage points there is nothing for a seed baseline to resolve.


\section*{17 August, Afternoon: What A2 Actually Does Out of Domain}

Seed $0$ finished on all five DRENDS recordings against all three
backbones, $7{,}476$ frames per backbone, matched sequence for sequence
against the canonical run. This is the comparison the morning's retracted
table was pretending to be.

\begin{center}
\begin{tabular}{lcccc}
\toprule
& \multicolumn{2}{c}{Canonical} & \multicolumn{2}{c}{A2} \\
\cmidrule(lr){2-3}\cmidrule(lr){4-5}
Reduction on DRENDS & Min & Max & Min & Max \\
\midrule
EPE  & $4.52$ & $5.42$ & $4.75$ & $5.74$ \\
Bad1 & $4.78$ & $5.69$ & $5.14$ & $6.17$ \\
Bad3 & $11.14$ & $16.17$ & $11.61$ & $17.04$ \\
RMSE & $6.43$ & $6.96$ & $6.37$ & $7.18$ \\
\bottomrule
\end{tabular}
\end{center}

Across the twelve comparisons --- three backbones by four metrics --- A2
is ahead in eleven and behind in one, and every margin lies between
$0.000$ and $0.005$ in absolute terms. That is a consistent direction at a
magnitude nobody should care about.

\paragraph{The finding, stated the way it survives.} A2 removes
$157{,}504$ parameters of frozen ResNet-18 from the deployed system,
three forward passes per frame with it, and $104$ of the $142$ evidence
channels. Out of domain it costs nothing: eleven of twelve margins favour
it, none exceeds $0.005$.

\emph{The sentence that stood here was backwards and is corrected in the
evening entry below.} It read that A2 ``costs $0.69$ percentage points of
Bad1 reduction on held-out D7''. These figures are relative
\emph{reductions}, so a larger number is a better result: canonical is
$5.15\%$ and A2 is $5.84\%$ across three seeds. A2 does not pay $0.69$
points in domain, it \emph{gains} them, and it gains $4.44$ on D2. Every
conclusion drawn from the wrong sign is withdrawn with it.

\paragraph{A third defect, caught before it mattered.} The comparison
script was fixed at midday to reduce both sides over their shared
sequences. That fix intersected across \emph{every} seed present --- so
the moment seed $1$ started writing its first recording, seed $0$'s
finished five-sequence comparison silently collapsed to the one sequence
seed $1$ had reached. A guard that over-corrects into a new way of
comparing unequal things is not fixed. Each seed is now intersected with
canonical on its own, incomplete seeds are named in a warning and
excluded from the spread, and the header carries $n$ per seed.

Seeds $1$ and $2$ are still running and add A2's own spread, not a second
baseline. Nothing in the paper has changed.

\paragraph{Eleven of twelve is not evidence, so we bootstrapped it.}
Counting how often A2 wins settles nothing: the twelve backbone-metric
comparisons share five sequences and are not independent draws. So the
difference was resampled at the level that actually varies. Sequences are
the resampled unit, the paired difference is formed per sequence before
resampling --- which cancels sequence difficulty exactly, the variance
that manufactured this morning's phantom --- and all three backbones of a
resampled sequence move together, preserving their correlation rather than
averaging it away.

\begin{center}
\begin{tabular}{lccc}
\toprule
Metric & Mean difference & $95\%$ interval & $P(\text{A2 better})$ \\
\midrule
EPE  & $-0.0042$ & $[-0.0072, -0.0016]$ & $1.000$ \\
Bad1 & $-0.0026$ & $[-0.0041, -0.0013]$ & $1.000$ \\
Bad3 & $-0.0004$ & $[-0.0007, -0.0001]$ & $0.992$ \\
RMSE & $-0.0017$ & $[-0.0062, +0.0026]$ & $0.767$ \\
\bottomrule
\end{tabular}
\end{center}

Three of four intervals exclude zero. The advantage is therefore real and
detectable --- and it is $0.004$ pixels of EPE. Both halves of that
sentence are load-bearing: this is what a genuine effect of no consequence
looks like, and it is worth having a name for, because an interval that
excludes zero is otherwise very easy to report as though it mattered.

One limit stated plainly: the interval covers variation across
\emph{sequences} only. The canonical side is a single training run and the
A2 side currently pools one seed, so seed variance is not in it --- and
A2's seed-to-seed spread on DRENDS is of the same order as the differences
above. The reading stays \emph{detectable and negligible}, not
\emph{better}. The script prints that caveat with its own output so the
number cannot travel without it.


\section*{17 August, Evening: The Third Seed Removes the Effect}

The caveat above was the right caveat and it was still too generous. Seeds
$1$ and $2$ finished, all nine runs at five recordings, and pooling them
moves every interval across zero:

\begin{center}
\begin{tabular}{lccc}
\toprule
Metric & Mean difference & $95\%$ interval & $P(\text{A2 better})$ \\
\midrule
EPE  & $-0.0028$ & $[-0.0056, +0.0002]$ & $0.962$ \\
Bad1 & $-0.0015$ & $[-0.0032, +0.0001]$ & $0.971$ \\
Bad3 & $-0.0002$ & $[-0.0005, +0.0000]$ & $0.930$ \\
RMSE & $-0.0031$ & $[-0.0059, +0.0009]$ & $0.942$ \\
\bottomrule
\end{tabular}
\end{center}

Three of four intervals excluded zero on seed $0$, three of four still
excluded it with two seeds pooled, and \emph{none} of four excludes it
with three. The per-seed numbers say why, and they are monotone:

\begin{center}
\begin{tabular}{lcccc}
\toprule
Bad1, RAFT-Stereo & Canonical & A2 seed 0 & seed 1 & seed 2 \\
\midrule
DRENDS & $0.555$ & $0.552$ & $0.553$ & $0.555$ \\
\bottomrule
\end{tabular}
\end{center}

A2 seed $2$ lands exactly on the canonical value, on EPE as well
($1.241$ against $1.241$). The advantage was a property of seed $0$, and
the earlier sentence --- ``the advantage is therefore real and detectable''
--- does not survive. The finding is now the flat one: out of domain, A2
and the canonical model are \textbf{indistinguishable}.

\paragraph{What this does not need.} The canonical DRENDS seed study was
cancelled earlier today on the reasoning that a $0.001$ point difference
left a seed baseline nothing to resolve. That reasoning was based on the
one-sequence numbers and was wrong at the time: A2's real seed spread here
is $0.003$, the same size as the gap it was being compared against. The
conclusion survives the error only because it runs the safe way --- adding
canonical seed variance widens the intervals, and every interval already
includes zero. Had the three seeds gone the other way, the missing
baseline would have been load-bearing.

\section*{17 August, Late: A Sign Error, and the Conclusion It Was
Hiding}

Every in-domain statement written today had the sign backwards. The
ablation table reports \emph{relative reduction}, computed as
$100\,(\text{raw}-\text{refined})/\text{raw}$, so a larger number is a
better result. Read correctly:

\begin{center}
\begin{tabular}{lcc}
\toprule
Bad1 reduction (\%) & Canonical & A2 \\
\midrule
D2 development & $8.40$ & $\mathbf{12.84}$ \\
D7 held out    & $5.15 \pm 0.34$ & $\mathbf{5.84 \pm 0.22}$ \\
DRENDS         & \multicolumn{2}{c}{indistinguishable} \\
\bottomrule
\end{tabular}
\end{center}

A2 does not lose in domain. It \emph{wins} in domain, by $0.69$ points on
held-out D7 and by $4.44$ on D2, and it is neutral out of domain. Three
commits and two diary paragraphs today asserted the opposite, including
the phrases ``A2 loses in domain'' and ``more than a hundred times the
largest DRENDS margin''. All of it is withdrawn.

\paragraph{What was actually being measured.} With the sign fixed, the
result stops being about domains and becomes a statement about the
architecture. The $64$ frozen ResNet-18 channels, the confidence curves
built from them and both appearance correlations --- $104$ of $142$
channels, $157{,}504$ frozen parameters, three forward passes per frame
--- buy accuracy \emph{nowhere}. In domain the model is better without
them. Out of domain it is unchanged.

The pre-registration allowed for exactly this and named the consequence
before any of it ran: if the appearance channels match within the seed
spread they ``are redundant given their own correlations and should be
dropped from the architecture description''.

\paragraph{The uncomfortable version.} A2 is not alone. On held-out D7,
A3 (single resolution) reaches $5.85$ and A4 (relaxed convexity) $6.02$,
both above canonical's $5.15$; only A1 is below, at $4.88$. Every
simplification we pre-registered does at least as well as the model we
designed. The narrow reading is that the frozen extractor is unnecessary.
The wider one is that the canonical head is over-specified and its parts
are not earning their places --- and that is a finding about our design
process, not about temporal refinement.

Two limits keep this from being stronger than it is. Each ablation is a
\emph{single} seed against a three-seed canonical, and A2's $+0.55$ on one
seed sits only modestly above the $0.34$ canonical spread. And the D2
margin, being the selection split, is exactly where a variant is most
flattered.

\paragraph{Audited, because it read as too neat.} ``Every simplification
beats the model we designed'' is the kind of sentence that is usually a
measurement error, so it was checked before being believed. It is not one.
The raw baseline is identical to eight decimals across all seven runs
($0.06205083$ on D7), the backbone--sequence sets hash identically, every
run trained twelve epochs, and every checkpoint was selected on validation
fused EPE. Nothing about the comparison is uneven.

What the audit did find is that the interesting number is a spread, not a
mean:

\begin{center}
\begin{tabular}{lcc}
\toprule
Bad1 reduction (\%), three seeds each & Canonical & A2 \\
\midrule
D2 development & $8.40 \pm 2.53$ & $12.73 \pm 0.32$ \\
D7 held out    & $5.15 \pm 0.34$ & $5.84 \pm 0.22$ \\
\bottomrule
\end{tabular}
\end{center}

On D7 the two ranges are disjoint --- A2's worst seed, $5.69$, beats
canonical's best, $5.47$ --- so for A2 this is not seed noise. But the
figure that matters is $\pm 2.53$: the canonical model's development
result runs from $5.75$ to $10.78$ depending only on the seed, close to a
factor of two, while A2 runs from $12.38$ to $12.98$. The learned evidence
is not merely failing to buy accuracy. It is making training unstable, and
that instability is the honest explanation for why every variant looks
better than the canonical model: the canonical distribution is wide and
the number in our tables is one draw from it.

Two limits survive the audit. A3 and A4 remain single seeds, so ``every
simplification'' is demonstrated only for A2. And the D2 margin sits on
the selection split, where any variant is most flattered.

\paragraph{The interval was weaker than its label.} The bootstrap script
had produced the day's final conclusion and had no test, so it got one:
synthetic sequences whose difficulty spans $0.5$ to $50$ with a constant
$-0.01$ effect on every one. The pairing recovers $-0.01$ exactly, and the
same data reduced the unpaired way over unequal subsets shows $-15.60$
from no effect at all --- the morning's bug, reproduced on demand.

The second half of that check measures something we should have known and
did not. Under a \emph{true null} at this sample size, the percentile
interval excludes zero on $22.7\%$ of $300$ trials rather than $5\%$. With
five sequences the bootstrap is anti-conservative by more than four times.

So ``three of four intervals exclude zero'', reported this afternoon as
evidence the effect was real, was much weaker evidence than its label
implied. It does not change the final answer --- with three seeds nothing
excludes zero anyway --- but it does mean the intermediate claim was
flimsier than the caveat attached to it admitted. The script now prints
the measured under-coverage next to every result.

The first version of this test asserted that a single null draw must
include zero, and failed on a three-sigma sample. That assertion was
wrong, not the code: a $95\%$ interval excludes zero on $5\%$ of nulls by
construction. Testing a rate rather than a draw is what turned a flaky
assertion into the finding above.

\paragraph{The whole day in one line.} A2 removes the frozen ResNet-18
from the runtime, is better in domain and indistinguishable out of it.
The out-of-domain claim went threefold better, then detectably better,
then neither; the in-domain claim went the other way, from a loss to a
gain, because of a sign. Each revision came from checking something that
had been assumed --- and the sign was assumed longest, because it was the
half nobody was arguing about.


\section*{18 August, Overnight: The Comparison Moved to Neutral Ground}

The published head-to-head against BiDAStabilizer ran on SCARED-C D2. That
is the split our checkpoint was selected on under the pre-registered rule,
while BiDAStabilizer had never seen the domain --- so the arena was ours and
the margin measured partly that. Yesterday the paper began declaring it. Last
night it was replaced.

DRENDS is out of domain for \emph{both}: we train on ex-vivo porcine tissue
and BiDAStabilizer has never seen synthetic colonoscopy either. Five
recordings, $7{,}476$ frames, $117.4$M pixels, one prediction-independent
support, both methods refining the same array.

\begin{center}
\begin{tabular}{lccccc}
\toprule
Reduction vs.\ raw & Causal & EPE & Bad1 & Bad3 & RMSE \\
\midrule
BiDAStabilizer & no  & $2.0\%$ & $0.2\%$ & $4.7\%$ & $5.3\%$ \\
TETHER         & yes & $\mathbf{5.5\%}$ & $\mathbf{3.1\%}$ & $\mathbf{12.3\%}$ & $\mathbf{9.5\%}$ \\
\bottomrule
\end{tabular}
\end{center}

Two to thirteen times the error removed on every metric, from the only causal
entry, whose trained module is $54\times$ smaller --- $177{,}338$ parameters
against $9{,}636{,}375$, both counted from the checkpoints rather than from
the papers. Not uniform, and reported that way: TETHER leads on four
recordings of five and is $0.020$\,px worse than the baseline on Vid13.

\paragraph{The ordering was forced, not chosen.} \texttt{workers/bidastabilizer.py}
reads \texttt{rgb\_left} and \texttt{rgb\_right} only and then replaces the
disparity with its own RAFT-Stereo \emph{robust} prediction. The first
exporter was written on the opposite assumption --- that handing it our
cached middlebury array would give both sides one input --- and that array
would have been discarded on read while the file implied a sharing that never
happened. So BiDA runs first and \emph{produces} the shared raw; TETHER is
driven on it afterwards. Caught by reading the worker before running it.

\paragraph{Five things broke, and the shape of them is the lesson.}
\begin{itemize}
\item \texttt{\_canonical\_frames} returns the \emph{native} $720\times1280$
      shape it validates against, not the $144\times180$ grid its frames come
      out on. Unpacking that as the output grid built a dummy disparity at the
      wrong size, which the bridge rejected on write --- doing its job, but
      only after a 25-minute run.
\item The \texttt{h100app} profile has \texttt{TASKLIMIT 4}, so a 16-core
      request is rejected outright, and the edit that introduced it had
      quietly broken three working launchers.
\item The external-comparison wrapper refuses physical GPU 0. Our own runners
      pick the freest device; the D2 external runner pins device 1. Copying
      the wrong one of two coexisting conventions failed all five recordings
      at once.
\item A 160GB reservation left the job PEND with ``requirements for reserving
      resource (mem) not satisfied''. The request was not merely generous, it
      stopped the work from starting. At 40GB it dispatched immediately.
\item \texttt{--purpose} accepted only \texttt{SMOKE\_DIAGNOSTIC} or
      \texttt{D2\_FULL\_DIAGNOSTIC}, so a full DRENDS run could only be
      labelled as a smoke test or as D2 data. Added a third rather than write
      a false label.
\end{itemize}

None of these were subtle once seen. All five came from assuming a convention
instead of reading the one already in the repo.

\paragraph{Eight hours blind.} The scoring run took over eight hours and
printed nothing: its only output lands at the end of each recording and
stdout is block-buffered behind the \texttt{bsub} pipe. There was no way to
tell slow from stuck, the completion estimate was wrong three times, and a
crash on the fifth recording would have discarded all of it. The script now
flushes and writes a \texttt{partial.json} after each recording, with the
final write removing it so a leftover partial always means an unfinished run.
The run in flight was left alone rather than restarted for better logging.

\paragraph{What this does not yet say.} The comparison used the
\emph{canonical} module. A2 --- which is better in domain, indistinguishable
out of it, and drops the frozen ResNet-18 from the runtime entirely --- has
not been put through it. If it holds the same margin the claim becomes a
$154$k-parameter module with no frozen extractor beating a $9.64$M
bidirectional stabiliser out of domain, and the $157{,}504$ parameters we
currently have to add to our own side of the ledger disappear.



\section*{18 August, Afternoon: A Competitor That Was Already Measured}

TC-Stereo has been scored since 16 August. The numbers sat in
\texttt{results/tcstereo\_reference/pooled.json} for two days without ever
reaching the paper, which cited the method in related work and then said
nothing about having run it. That is the cheapest possible kind of waste:
GPU hours spent, evidence produced, evidence not used.

\paragraph{What it can and cannot be.} TC-Stereo propagates its temporal
state by a rigid $6$-DoF reprojection built from per-frame pose,
intrinsics and baseline. SCARED-C publishes no per-frame pose, and
deforming tissue would break the rigidity assumption even if it did, so the
temporal path was disabled in the run and the manifest records why. What
was measured is therefore TC-Stereo's \emph{frame} path, and it must be
reported as an integrated-architecture reference rather than as a temporal
rival. Reporting it as a beaten competitor would be a straw man.

\paragraph{The support is bit-identical, which is why the row can exist.}
The run took the same D2 boundary as the in-domain BiDA head-to-head. Raw
EPE, Bad1, Bad3 and RMSE agree with the BiDA table to full double
precision, the three sequences are the same, and the per-sequence frame and
pixel counts match exactly ($17{,}999{,}475$ pixels). That check ran before
anything was written: the same class of error --- comparing over
sets that only look the same --- produced the phantom threefold DRENDS win
yesterday, and the lesson was to verify the support first.

\paragraph{The result.} Transferred without retraining, TC-Stereo reaches
EPE $0.524$ and Bad1 $9.05\%$, against $0.303$ and $2.09\%$ for the frozen
RAFT-Stereo \emph{robust} prediction and $0.289$ and $1.83\%$ after our
head. A co-designed model moved to surgical data starts $73\%$ worse than
the estimator it would have to replace, and its temporal component cannot
be reached from there at all. This is the argument for the plug-in
interface stated in someone else's numbers rather than ours, which makes it
worth more than another ablation of our own.

\paragraph{The page limit was the real work.} The paragraph pushed the
paper to nine pages, with exactly one reference spilling. Reclaiming it
took six lines: the abstract repeated the name explanation the introduction
already gives, the research question sat in a \texttt{quote} block costing
more vertical space than the sentence, and three sentences said twice what
they say once. No evidence was cut. A note for the next time this happens:
\texttt{\textbackslash footnotesize} on the bibliography does nothing,
\texttt{ieeeconf} already sets it.

\paragraph{Still open.} Canonical seeds 1 and 2 on DRENDS were relaunched
this afternoon after the first attempt was killed mid-run on 17 August ---
the log ends in \texttt{KeyboardInterrupt} and a manifest with no summary.
Until they finish, A2's three-seed spread is being compared against a
one-seed canonical baseline, which is the asymmetry we already know how to
be wrong about. The launcher now skips any run that wrote its summary, so a
second kill costs only the run in flight.


\section*{18 August, Evening: The Promotion Rule Fires, and It Is Not Ours to Argue With}

A review arrived. Its second major objection is that the ablations contradict
the canonical configuration and the paper states this without acting on it ---
that honesty about a problem is not a resolution of it. That is correct, and the
answer was already on disk.

\paragraph{The seeds existed and the script ignored them.}
\texttt{ablation\_table.py} printed ``each ablation is ONE seed'' while
\texttt{results/ablation\_eval/} held \texttt{a2}, \texttt{a2\_seed1} and
\texttt{a2\_seed2}. The defence the review calls self-defeating --- one run
against a three-seed baseline --- had been retired days ago by our own runs, and
nobody told the table. The script carried a comment at the top explaining that
comparing against canonical seed 0 alone had already misled us once; it then
made the same mistake with the roles reversed.

\paragraph{The rule was written first, which is the only reason this is allowed.}
\texttt{ablation\_preregister.json} prohibits ``selecting the best ablation and
presenting it as the canonical model''. What makes promotion legitimate rather
than cherry-picking is the \texttt{A2\_seed\_study} block of 16 August, which
fixed the decision \emph{before} the extra seeds were trained: decision split D2
only, D7 reported afterwards and taking no part in the choice.

\paragraph{The verdict.} D2 Bad1 across A2's three seeds is $12.38$, $12.84$,
$12.98$: mean $12.73$, spread $0.32$. The canonical three-seed mean is $8.40$
with spread $2.53$. A2 exceeds it by $4.33$ points with a spread eight times
tighter, so the rule returns PROMOTE. It is also ahead on held-out D7, which the
rule deliberately does not consider.

\paragraph{We looked for a reason to keep the larger model, and did not find one.}
A2 is better on the decision split, better on the split it was not judged on,
far more stable across seeds, and smaller: it sets
\texttt{include\_learned\_stereo\_evidence=False} and passes a null extractor, so
the frozen ResNet-18 is genuinely not evaluated --- $157{,}504$ parameters and
three forward passes per frame that stop existing. Proposing the $142$-channel
model while reporting that the $38$-channel one is better is the position the
review describes, and it is indefensible. The choice is to re-canonicalise or to
find a principled reason not to; there is no third option where we report both
and keep the bigger one.

\paragraph{What re-canonicalisation actually costs.} A2 has three trained
checkpoints, D2 and D7 evaluations on all five backbones, and DRENDS on three;
its out-of-domain BiDA comparison is running now. What is missing is the
development closure, which is the expensive half: \texttt{experimental\_closure.py}
hardcodes \texttt{CanonicalHorizon} and \texttt{ExperimentalPolicy} with no module
argument, so the fifteen policies, the horizon sweep and the oracle cannot be
pointed at A2 without changing the code first. $225$ method reports on D2.

\paragraph{An unrelated fact that turned out to be the same fact.}
The pose claim retracted this afternoon appears twice, not once:
\texttt{evaluate\_3d\_consistency.py} also states that ``SCARED-C ships no
per-frame camera pose in this repo'', and limits itself to frame-to-frame surface
stability for that reason. With the poses recovered, that script can measure
accumulated reconstruction drift instead --- which is exactly the downstream
robotics evidence the review's first objection demands. The error that weakened
the competitor comparison was also the one blocking the venue argument.


\section*{18 August, Late Afternoon: The Freeze Refused Twice, Correctly}

Re-canonicalising to A2 means re-running the development closure under the
promoted head. Two attempts to do that failed before the third worked, and both
failures were the protocol behaving exactly as designed.

Adding a \texttt{-{}-head} argument to \texttt{experimental\_closure.py} produced
\texttt{freeze hash mismatch: experimental\_closure}. Moving the new adapter class
into \texttt{canonical\_horizons.py} produced \texttt{freeze hash mismatch:
canonical\_horizons}. Both files are pinned by the closure's freeze manifest. A
frozen experiment whose script can be edited in place is not frozen, and the
manifest said so twice before any GPU time was spent. The canonical closure keeps
its bytes; the promoted head runs from its own module beside it.

\paragraph{The re-run is a night, not a week.} Six method-configurations per
seed, not twenty-one. Nine of the fifteen closure policies blend raw against
flow-aligned raw with a fixed, EMA or forward--backward-confidence weight and
load no checkpoint at all, so they are identical whichever head the paper ships;
the raw-versus-memory oracle is computed from ground truth and never runs the
adapter. Both are carried forward by reference, and the runner refuses to start
if the canonical closure they come from is missing.

\paragraph{Two things landed in the paper that survive the promotion.} The
fifteen policies are now enumerated --- the review asked for the list and the
paper had named five --- and the two aggregations are declared. That second one
was a real finding from the review's forensics: raw D2 EPE $0.276202$ matches no
row of the backbone table because the closure pools every backbone and sequence
before taking a ratio, while everything reported per backbone or per arena
averages per-cell relative reductions. They differ by up to half a point and
neither was named. Space came from two displayed deltas that restated in symbols
the four numbers the following sentence already gave in prose.

\paragraph{The wrong sentence was load-bearing in two places.}
\texttt{evaluate\_3d\_consistency.py} justified limiting itself to
frame-to-frame surface stability by the same false claim about missing poses.
Corrected: the restriction is a scope decision, and the accumulated-drift
measurement the poses enable is named as the next step. That measurement is also
the downstream evidence the review demands for a robotics venue, so one wrong
sentence was blocking the competitor comparison and the venue argument at once.


\section*{18 August, Evening: What Regime Are We Actually Measuring In}

The review asks what causes Bad1 $=44.92\%$ on DRENDS raw, and reads it as a
prediction at the edge of broken. It is, and the paper now says so.

\paragraph{The first explanation was wrong.} The comfortable answer was a units
artefact: DRENDS ground-truth disparity averages $9.89$\,px and spans
$8.42$--$11.89$ between the $5$th and $95$th percentiles, so a $1$\,px threshold
covers a third of the scene's entire depth range and Bad1 is simply a harsh bar
at that scale. That would have made the number a measurement quirk rather than a
weakness. It does not survive contact with SCARED-C, which sits at the same
disparity scale --- $7.2$ to $11.0$\,px across D7, recovered from the reported
depth and $f_x b$ --- while raw EPE there is $0.14$ against $1.20$ on DRENDS.
Same units, eight times the error. The regime is genuinely hard.

Nor is it a property of synthetic data, which was the second comfortable answer:
SCARED-C D7 keyframe~$3$ has raw EPE $1.18$\,px on comparable disparities. Real
porcine tissue, same difficulty.

\paragraph{What this costs us.} The out-of-domain BiDA comparison is the paper's
primary head-to-head, and it happens where the input is weak. That does not
invalidate it --- both methods refine the same weak prediction under one support
--- but it bounds what it can claim, and the section now says the comparison is
between two refiners of a poor estimate rather than evidence that either
recovers a good one.

\paragraph{An early signal on A2, from one recording of five.} The out-of-domain
scorer has finished Vid10: raw $2.1185$, BiDA $2.0411$, canonical TETHER
$1.9133$, A2 $1.9238$. Identical raw and identical BiDA confirm the two runs sit
on the same boundary. A2 is $0.0105$\,px \emph{behind} canonical here --- $9.19\%$
reduction against $9.69\%$. One recording is not a result and the promotion rule
was decided on D2 in any case, but it is the first place A2 has not been at least
level, and it is worth watching rather than discovering later.

\paragraph{The memory-per-core trap, again.} The A2 closure pended for an hour on
``Job requirements for reserving resource (mem) not satisfied''. LSF reserves
\texttt{rusage[mem]} \emph{per core}, so \texttt{-n 4 -R rusage[mem=40GB]} holds
$160$GB, and three of our own jobs had left $7$GB free of $529$GB reserved on the
node. Resubmitted at $10$GB, and it now waits behind our own work rather than
behind a bad request.


\section*{18 August, Night: A Guard That Only Worked in One Direction}

\texttt{seed1 / CREStereo} died on \texttt{FileExistsError: refusing existing
output}. The resume guard added this afternoon tests whether a run wrote its
summary, which is the right definition of finished --- but the evaluator also
refuses to write into a directory that already exists, and a killed run leaves a
manifest with no summary behind. So the guard correctly decided the work was
unfinished, and the evaluator then refused to do it. Every attempt would have
failed the same way.

Incomplete remains are now moved aside with a timestamp rather than deleted.
Four kilobytes is not worth automatically destroying the record that a run was
killed.

\paragraph{Two corrections to the paper that the review earned.} The ablation
table still described every variant as a single run and still carried A2's
seed-0 numbers, five days after A2 got its three seeds. Its held-out figures are
$5.61\pm0.42$ EPE and $5.84\pm0.22$ Bad1 against canonical $5.06$ and
$5.15\pm0.34$. The discussion used to lean on ``two standard deviations from a
single sample is suggestive rather than settled'' --- a defence that has expired
--- and now says what the seeds actually show: the $142$-channel evidence space
is not what makes the module work, it is what makes its training unstable,
$2.53$ points of D2 Bad1 spread against A2's $0.32$.

The columns were also unlabelled as relative reductions, which the review reports
misreading on first pass, and the warp support and the total support were both
$S_t$. Both fixed.


\section*{18 August, Night: The First Thirteen Cells of the A2 Closure}

The A2 closure is thirteen cells into \texttt{canonical\_h1}, and every one has a
canonical counterpart to compare against. The raw column is identical to $10^{-9}$
in all thirteen, which is the check that matters first: the two runs are on the
same boundary and any difference is the head.

Mean difference: $+0.00017$\,px, with five of thirteen cells better for A2. At
$H{=}1$ the $142$-channel evidence space buys nothing measurable over $38$
channels. That is consistent with the promotion but it is also the least
demanding setting --- $H{=}1$ takes memory from raw every frame and never
propagates a corrected state --- so it is not yet the interesting comparison.
$H{=}4$ is.

One pattern worth recording before it is explained away: A2 is behind on
\texttt{keyframe\_2} and \texttt{keyframe\_4} and ahead on \texttt{keyframe\_3},
consistently across all five backbones. The difference is sequence-driven, not
backbone-driven, which is the same shape as the DRENDS finding that the recording
sets the pattern and the estimator does not.

\paragraph{Out of domain, A2 is now level.} The second recording landed:
Vid11 raw $1.3208$, BiDA $1.2760$, canonical $1.2274$, A2 $1.2178$. A2 is ahead
by $0.0096$\,px here, having been behind by $0.0105$ on Vid10. Two recordings,
one each way, roughly cancelling. The note written earlier this evening --- that
Vid10 was the first place A2 had not been at least level --- stands as recorded
but no longer points anywhere.

\paragraph{A scoping error worth twenty-three GPU-hours.} The launcher runs the
A2 closure for seeds $0$, $1$ and $2$. The canonical closure it is being compared
against was run on seed $0$ alone --- checkpoint sha \texttt{99c5745c}, the
canonical seed-0 weights. A matched comparison therefore needs A2 seed $0$ and
nothing more; the other two seeds are roughly ten hours each of a GPU that the
downstream reconstruction-drift experiment and the runtime measurement both need,
and the review ranks the first of those above anything else in the paper. The
running script is not being edited --- bash reads a script by byte offset while
executing it, so editing one mid-run corrupts it --- so the job will be stopped
after seed $0$ lands instead.


\section*{18 August, Night: A Trend That Was the Order of the Cells}

Sharding the A2 closure produced numbers fast enough to look at, and they told a
clean story. Reduction on D2 EPE, A2 minus canonical: $-0.11$ points at $H{=}1$,
$-0.54$ at $H{=}2$, $-1.18$ at $H{=}4$, $-2.69$ at $H{=}8$. A monotone widening
with the horizon, and a mechanism to go with it --- the $142$-channel evidence
would matter more the longer a corrected state has to survive, and buy nothing at
$H{=}1$ where memory comes from raw every frame. It would have reversed the
promotion for the operating point the paper actually ships.

It is an artefact. The cells complete in a fixed backbone and sequence order, so
the ones finished first are always the same ones, and they are not a random
sample. Canonical's own reduction on exactly the cells A2 has finished, against
its reduction on all fifteen:

\begin{center}
\begin{tabular}{lccc}
\toprule
Method & on finished cells & on all cells & subset bias \\
\midrule
$H{=}1$ & $2.52\%$ ($n{=}15$) & $2.52\%$ & $+0.00$ \\
$H{=}2$ & $2.98\%$ ($n{=}6$)  & $3.09\%$ & $-0.11$ \\
$H{=}4$ & $4.80\%$ ($n{=}4$)  & $3.81\%$ & $+0.98$ \\
$H{=}8$ & $5.53\%$ ($n{=}4$)  & $3.84\%$ & $+1.69$ \\
\bottomrule
\end{tabular}
\end{center}

The first four cells of $H{=}8$ are worth $1.69$ points more to canonical than its
own average, and the gap I was about to report there is $2.69$. Most of the trend
is the ordering of the work queue.

What survives: $H{=}1$ is complete on all fifteen cells and unbiased, and there
A2 is level at $-0.11$ points. Nothing can be said about $H{=}4$ until its
remaining eleven cells land.

This is the same error as the threefold DRENDS win of yesterday --- averaging two
quantities over sets that are not the same set --- arriving in a new costume. The
check that caught it took two minutes and should be standing procedure for any
partial run: before comparing, score the \emph{reference} on exactly the subset
the new run has finished.


\section*{18 August, Night: What the $142$ Channels Might Actually Buy}

Having thrown away the contaminated version, the honest comparison is on the
cells \emph{every} horizon has finished --- five of fifteen, the same five for
both models, so the absolute levels are inflated for both but the shape across
$H$ is measured on identical data.

\begin{center}
\begin{tabular}{lccc}
\toprule
$H$ & canonical & A2 & A2 $-$ canonical \\
\midrule
$1$ & $2.92$ & $2.50$ & $-0.41$ \\
$2$ & $3.60$ & $2.87$ & $-0.72$ \\
$4$ & $4.44$ & $3.22$ & $-1.22$ \\
$8$ & $5.12$ & $2.40$ & $-2.72$ \\
\bottomrule
\end{tabular}
\end{center}

Canonical keeps gaining as the horizon lengthens. A2 gains to $H{=}4$ and then
collapses --- at $H{=}8$ it is worse than at $H{=}1$. If this survives the
remaining ten cells it is the most interesting thing the ablation has produced,
because it names what the removed channels were for: the appearance correlations
and learned stereo evidence are what let the head notice that aged memory has
gone stale. With $38$ channels of disparity, motion and RGB geometry the head
cannot tell a good old memory from a rotten one, so a longer recurrence hurts it.

\paragraph{The tension with the promotion, stated rather than smoothed.}
The promotion rule was decided on confined D2 \emph{Bad1}, where A2 wins
decisively, $12.73$ against $8.40$. This is unconfined \emph{EPE}, where it
loses. Different metrics, not a contradiction --- but they point opposite ways
and the paper cannot report one and bury the other.

Worse for the simple story: at seed $0$ alone, the seed the closure actually
runs, canonical beats A2 on D2 EPE in the complete confined evaluation too
($3.95$ against $3.69$). A2 only passes canonical on EPE in the \emph{three-seed
mean} ($4.13$ against $3.45$), and it passes there because canonical's spread is
enormous, not because A2 is more accurate.

So the reading forming tonight is that A2 is not better, it is \emph{steadier}:
it wins on Bad1 and on seed variance, and loses on long-horizon EPE. If that
holds, the conclusion may not be ``ship A2'' but ``ship A2 at a shorter
horizon'', or ``what the $142$ channels buy is horizon''. The closure already
evaluates A2 at $H\in\{1,2,4,6,8\}$, so that answer arrives without another run.

Nothing goes in the paper until $H{=}4$ has all fifteen cells.


\section*{18 August, Night: The Horizon Effect Has Shrunk Threefold While I Watched}

The entry above reports A2 falling behind canonical by $2.72$ points at $H{=}8$
and $1.22$ at $H{=}4$, on the five cells every horizon had then finished. Those
magnitudes are already wrong. On the same matched-cell comparison as more cells
land:

\begin{center}
\begin{tabular}{lccc}
\toprule
$H$ & $n{=}5$ & $n{=}7$ & $n{=}10$ \\
\midrule
$1$ & $-0.41$ & $-0.17$ & $+0.03$ \\
$2$ & $-0.72$ & $-0.35$ & $-0.08$ \\
$4$ & $-1.22$ & $-0.64$ & $-0.36$ \\
$8$ & $-2.72$ & $-1.57$ & $-0.90$ \\
\bottomrule
\end{tabular}
\end{center}

Every gap has shrunk by roughly a factor of three in two hours, and at $H{=}1$ A2
is now marginally ahead. The shape --- each longer horizon costing A2 a little
more --- has survived three refreshes, so it is probably real. The size is not a
number yet and repeating it as one would be the third version of the same
mistake in two days.

What the shrinkage itself says: the cells that finish first are systematically
ones where canonical does relatively well \emph{and} A2 relatively badly. That is
not random, and it is a reason to distrust every partial comparison in this
project by default, not only this one.

The honest position at $H{=}4$, the operating point that matters, is a deficit of
$0.36$ points that is still moving toward zero. Five cells remain.


\section*{18 August, Night: The Gap Is Not Converging, It Is Alternating}

Three entries above I reported the A2 deficit shrinking threefold as cells landed
and wrote that at $H{=}4$ it was ``still moving toward zero''. That reading was
mine, not the data's, and it is wrong.

$H{=}1$ is complete, so its fifteen cells can be decomposed by sequence:

\begin{center}
\begin{tabular}{lc}
\toprule
Sequence & A2 $-$ canonical \\
\midrule
\texttt{dataset\_2\_keyframe\_2} & $-0.41$ \\
\texttt{dataset\_2\_keyframe\_3} & $+0.47$ \\
\texttt{dataset\_2\_keyframe\_4} & $-0.40$ \\
\bottomrule
\end{tabular}
\end{center}

A2 loses on two sequences and wins on one. The gap therefore has nothing to
converge to --- what it does is move according to which sequence the queue has
reached. The apparent shrinkage was \texttt{keyframe\_3} cells arriving, and
$H{=}4$ currently holds \texttt{keyframe\_2} and \texttt{keyframe\_3} only: every
one of its five \texttt{keyframe\_4} cells is still missing, and
\texttt{keyframe\_4} is the sequence where A2 loses by $0.40$ at the horizon we
can actually measure.

So the honest expectation is the opposite of what I wrote: when the last block
lands the $H{=}4$ gap should hold near $-0.4$ or widen, not close. The queue does
not deliver a random sample and it never did --- it delivers one sequence at a
time, which is the worst possible ordering for anyone extrapolating from a
partial mean.

I have now made three versions of this same mistake in two days: averaging over
mismatched sets, reading a trend that was the work queue, and extrapolating a
partial mean as if the remainder resembled the part. The rule that would have
caught all three: a partial result is not a small version of the whole result,
it is a \emph{biased} version, and the bias has a direction you can look up.


\section*{18 August, Night: Three Complete Horizons, and A2 Does Not Collapse}

$H{=}1$, $H{=}2$ and $H{=}8$ are now complete on all fifteen cells, so they can
be read without a subset argument. Reduction in D2 EPE, unconfined:

\begin{center}
\begin{tabular}{lccc}
\toprule
$H$ & canonical & A2 & A2 $-$ canonical \\
\midrule
$1$ & $2.52$ & $2.40$ & $-0.11$ \\
$2$ & $3.09$ & $2.87$ & $-0.21$ \\
$8$ & $3.84$ & $3.18$ & $-0.66$ \\
\bottomrule
\end{tabular}
\end{center}

\paragraph{Two corrections to what I wrote earlier tonight.}
First, I said A2 collapses at a long horizon and lands below its own $H{=}1$.
That was the biased subset talking. On complete cells A2 improves monotonically,
$2.40 \rightarrow 2.87 \rightarrow 3.18$; it simply improves \emph{less} than
canonical, which goes $2.52 \rightarrow 3.09 \rightarrow 3.84$. Nothing collapses.

Second, the magnitudes I quoted were inflated between three and four times. The
$H{=}8$ deficit is $0.66$ points, not the $2.72$ of the first look. The shape
--- a gap widening monotonically with the horizon --- is the only thing that
survived every refresh, and it is now measured rather than guessed.

\paragraph{Out of domain, three recordings in, A2 is level.} Vid10 $+0.0105$\,px
against canonical, Vid11 $-0.0096$, Vid12 $-0.0084$. Two better, one worse, all
of them hundredths of a pixel on a raw error above one pixel.

\paragraph{What the promotion actually buys, stated plainly.} A2 is behind on
unconfined D2 EPE at every horizon, by between $0.11$ and $0.66$ points. It is
ahead on the pre-registered endpoint, confined D2 Bad1, by $4.33$ points. It is
eight times tighter across seeds, $0.32$ against $2.53$. And it is level out of
domain. The trade is a little mean accuracy for a large gain in bad-pixel rate
and in run-to-run stability, and that is a defensible thing to ship --- but it
has to be said in those words rather than as ``A2 is better''.

$H{=}4$, the operating point, is at thirteen of fifteen cells and sitting at
$-0.37$. Two \texttt{keyframe\_4} cells remain.


\section*{18 August, Late: Every Horizon Complete, and the Mechanism Was Backwards}

All six horizons are done on all fifteen cells, and the grid has no empty cells
left. Pooled EPE reduction on D2:

\begin{center}
\begin{tabular}{lccc}
\toprule
$H$ & canonical & A2 & A2 $-$ canonical \\
\midrule
$1$ & $2.61$ & $2.49$ & $-0.12$ \\
$2$ & $3.24$ & $3.02$ & $-0.21$ \\
$4$ & $4.05$ & $3.69$ & $-0.35$ \\
$6$ & $4.18$ & $3.73$ & $-0.45$ \\
$8$ & $4.09$ & $3.52$ & $-0.57$ \\
continuous & $-62.69$ & $-19.17$ & $+43.52$ \\
\bottomrule
\end{tabular}
\end{center}

\paragraph{The mechanism I proposed was backwards.} The story earlier tonight was
that the $142$ channels let the head notice when aged memory has gone stale, so
A2 should suffer most where recurrence runs longest. The first half holds --- A2
falls further behind at every longer horizon. The second half is the opposite of
what happens: remove the bound entirely and A2 degrades \emph{three times less},
$-19.17\%$ against $-62.69\%$. A head that supposedly cannot detect stale memory
survives unbounded staleness far better than the one that supposedly can.

The reading that fits all of it: the learned evidence buys mean accuracy while a
re-anchor keeps memory young, and costs robustness once nothing does. That is the
same trade as the seed spread ($0.32$ against $2.53$) and the Bad1 margin
($+4.33$). A2 is not more accurate, it is harder to break --- three separate
measurements now say so, and none of them says it is better.

\paragraph{The grid closed too.} S2M2-S and StereoAnywhere finished on DRENDS,
which was the last gap and, awkwardly, both missing cells were on the seen side:
the out-of-domain column had been one seen backbone against two unseen. With all
five, every one of $100$ recording--backbone--metric cells improves, against
$60/60$ before, and the seen-backbone DRENDS mean moves from $5.31$ on one
backbone to $5.27$ on three.

Five numbers in the prose moved with it and all five were regenerated rather than
nudged: combinations, cells, frames, estimator count, and the Vid10 argument that
the recording rather than the estimator sets the pattern --- which is now a
five-backbone spread of $0.77$ points instead of three numbers quoted from a
partial run. The draft \texttt{NOTE} listing them is gone because it is done.


\section*{18 August, Late: Where the Overhead Actually Is}

The stage-by-stage runtime finally ran, for both heads back to back in one job on
one device. Medians, milliseconds per frame:

\begin{center}
\begin{tabular}{lcc}
\toprule
Stage & Canonical & A2 \\
\midrule
SEA-RAFT forward   & $13.400$ & $13.644$ \\
SEA-RAFT reverse   & $13.346$ & $13.668$ \\
BiDA alignment     & $1.064$  & $1.115$ \\
Evidence           & $6.443$  & $\mathbf{0.710}$ \\
Fusion head        & $1.036$  & $1.103$ \\
\midrule
Total overhead     & $35.288$ & $\mathbf{30.240}$ \\
\bottomrule
\end{tabular}
\end{center}

\paragraph{Two things this settles.} A2 removes $5.05$\,ms, $14\%$ of the
overhead, and almost all of it comes from the evidence stage collapsing from
$6.44$ to $0.71$\,ms --- the frozen ResNet-18's three forward passes per frame.
That is the runtime consequence the ablation promised and it is now measured.

The second is larger and is not ours. \emph{Two SEA-RAFT passes are $27$ of the
$35$ milliseconds}, $76\%$ of everything the module costs. The reverse pass exists
solely to compute the forward--backward confidence cue. The review guessed this
and asked what removing the FB branch would cost; the answer is that it would
free $13.7$\,ms, nearly three times what re-canonicalising to A2 frees. A2 plus
the FB branch would be $18.7$\,ms of $35.3$, more than half the overhead, against
an accuracy cost of $0.35$ points at $H{=}4$.

\paragraph{What this measurement is not.} An exclusive GPU never became available
--- p1i is a single node, both its H100s carry other users' work, and the request
pended for three hours before being withdrawn. This ran shared, with one other
process on the device. The absolute figures are therefore not latency and are not
quotable as such; the two heads ran back to back on the same device in the same
job, so the \emph{difference} is sound. What is measured here is the cost A2
removes, not what either configuration costs.

The first attempt failed twice on \texttt{ModuleNotFoundError: model\_design}:
invoking the script by absolute path puts \texttt{scripts/} on \texttt{sys.path}
rather than the working directory, which every other launcher sidesteps by using
\texttt{python -m}.


\section*{19 August, After Midnight: A2 Out of Domain, and the Case Assembled}

Both remaining jobs finished. The out-of-domain scorer reached all five DRENDS
recordings with the A2 head, and the two canonical DRENDS seeds landed, so the
comparison is finally three against three on both sides.

\paragraph{Pooled, same boundary, same support.}

\begin{center}
\begin{tabular}{lcccc}
\toprule
& EPE & Bad1 (\%) & Bad3 (\%) & RMSE \\
\midrule
Raw               & $1.1988$ & $44.92$ & $4.54$ & $1.8762$ \\
BiDAStabilizer    & $1.1746$ & $44.81$ & $4.32$ & $1.7764$ \\
TETHER canonical  & $1.1325$ & $43.52$ & $3.98$ & $\mathbf{1.6990}$ \\
TETHER A2         & $\mathbf{1.1311}$ & $\mathbf{43.39}$ & $\mathbf{3.94}$ & $1.7098$ \\
\bottomrule
\end{tabular}
\end{center}

A2 leads on three metrics of four and gives up the tail. Per recording it is
ahead on three of five, and every difference is hundredths of a pixel.

\paragraph{The seed study is the cleaner evidence.} Canonical $5.00$, $4.74$,
$4.74$ --- mean $4.83\pm0.15$. A2 $5.32$, $5.21$, $5.06$ --- mean
$5.20\pm0.13$. The ranges do not overlap: A2's worst seed beats canonical's
best. Three trained checkpoints each, same protocol, same recordings, out of
domain for both.

\paragraph{The ledger, now that every piece is measured.} A2 costs $0.35$ points
of unconfined D2 EPE at $H{=}4$. It buys $4.33$ points of the pre-registered
confined Bad1 endpoint, a seed spread eight times tighter ($0.32$ against
$2.53$), three times less degradation when the recurrence bound is removed
($-19.17\%$ against $-62.69\%$), $13.8\%$ of the runtime and $157{,}504$
parameters, and now a non-overlapping out-of-domain margin. One consistent
negative against six positives.

That is enough to ship A2, and the sentence for the paper is not ``A2 is
better'' --- it is that the $142$-channel evidence space buys mean accuracy under
a bound and costs stability everywhere else.

\paragraph{Two smaller things.} The runtime re-ran on a genuinely idle GPU and
returned $33.549$\,ms for the canonical head, against the $33.60$ the paper has
quoted since before tonight: a number recomputed from scratch that confirmed
itself. And a stale ``60 cells'' survived in the conclusion after the grid went
to $100$; the DRENDS numbers had been regenerated everywhere the
\texttt{NOTE} pointed, and the conclusion was not one of the places it pointed.


\section*{19 August, Small Hours: The Downstream Measurement Says No}

The review's blocking objection is that the paper claims a robotics contribution
without showing anything downstream improves. SCARED-C's corrected poses made the
measurement possible, and it is now done: three D2 sequences, both heads, $200$
frames each, disagreement between views about the same world point.

\begin{center}
\begin{tabular}{llccc}
\toprule
Sequence & Head & excess raw & excess refined & change \\
\midrule
\texttt{kf\_2} & canonical & $0.2236$ & $0.2402$ & $-7.4\%$ \\
\texttt{kf\_2} & A2        & $0.2236$ & $0.2629$ & $-17.6\%$ \\
\texttt{kf\_3} & canonical & $0.1898$ & $0.2199$ & $-15.9\%$ \\
\texttt{kf\_3} & A2        & $0.1898$ & $0.2316$ & $-22.0\%$ \\
\texttt{kf\_4} & canonical & $0.1805$ & $0.1771$ & $+1.9\%$ \\
\texttt{kf\_4} & A2        & $0.1805$ & $0.1929$ & $-6.8\%$ \\
\bottomrule
\end{tabular}
\end{center}

Millimetres of excess over a ground-truth floor of $0.04$--$0.08$\,mm. Refinement
makes multi-view reconstruction consistency \emph{worse} in five cells of six,
by $7.1\%$ on average for the canonical head and $15.5\%$ for A2.

\paragraph{This is not the metric contradicting itself.} The frame-to-frame
figure quoted earlier --- excess $p2p$ down $36.2\%$ --- measured agreement
between \emph{consecutive} frames after flow warping, each in its own camera
frame. This measures agreement between \emph{many} frames in one world frame. A
refiner can improve the first and damage the second at the same time, and the
mechanism is not mysterious: raw is memoryless, so two views of a point disagree
only by their own noise, while a recurrent output depends on the history that
produced it, and two views reached through different histories carry different
bias. Temporal refinement buys per-frame accuracy by making each frame a
function of its past, and that is precisely what a multi-view fusion does not
want.

\paragraph{What it costs the paper.} The downstream experiment the review demands
now exists and it does not support the claim it was meant to support. The honest
report is that per-frame geometry improves and map-level consistency does not,
which bounds the contribution to single-frame perception rather than mapping.
That is a real limitation and it goes in the paper, not in a drawer --- and it is
worth more to state it than to have no downstream measurement at all, which was
the position this morning.

\paragraph{And a second mark against A2.} A2 is worse than canonical on all three
sequences here. Its ledger now carries two consistent negatives, unconfined D2
EPE and world-frame drift, against the Bad1 endpoint, the seed spread, the
unbounded-recurrence robustness, the runtime and the out-of-domain margin. The
promotion still follows the pre-registered rule, which was decided on D2 Bad1
alone, but ``A2 is simply better'' was never the claim and is now clearly false.


\section*{19 August, Small Hours: The Negative Result Survives Its Own Falsification}

A finding that says our own method makes something worse deserves a serious
attempt to break it before it is believed, and the obvious break is that the
refined branch of the new script is simply wrong. If refinement were worse
frame-by-frame as well, the multi-view number would be measuring a bug.

Same run, same frames, same support, both numbers printed together:

\begin{center}
\begin{tabular}{lccc}
\toprule
& raw & refined & change \\
\midrule
per-frame disparity MAE       & $0.1807$\,px & $0.1748$\,px & $+3.2\%$ \\
multi-view world-frame excess & $0.2236$\,mm & $0.2402$\,mm & $-7.4\%$ \\
\bottomrule
\end{tabular}
\end{center}

Refinement does its usual job and damages multi-view consistency at the same
time. The finding is coherent, so it stands. The three spread figures also
reproduced the earlier run to the digit, which is the determinism check arriving
for free a second time.

\paragraph{Why this is the strongest form of the result.} Stated as ``refinement
hurts reconstruction'' it invites the reply that the pipeline is broken. Stated
as ``$+3.2\%$ per frame and $-7.4\%$ across views, measured together'', it is a
statement about what temporal recurrence is and is not for --- and the paper now
carries both numbers in one sentence rather than the bad one alone.

\paragraph{A detail worth recording.} The verification script was written into
\texttt{/tmp} and submitted; the job dispatched and died on \texttt{No such file
or directory}. The login node's \texttt{/tmp} is not the compute node's.
Everything a batch job executes has to live on shared storage, which every
working launcher in \texttt{scripts/} already did and this one-off did not.


\section*{19 August, Dawn: Held-Out D7 Reverses the Downstream Result}

Three hours ago I put a negative downstream finding into the paper: world-frame
reconstruction consistency gets \emph{worse} under refinement, five cells of six.
It was measured on D2 only. D7 is now done and it reverses.

\begin{center}
\begin{tabular}{lcc}
\toprule
Split & canonical & A2 \\
\midrule
D2, development ($3$ sequences) & $-7.1\%$ \; ($1/3$) & $-15.5\%$ \; ($0/3$) \\
D7, held out ($4$ sequences)    & $\mathbf{+9.9\%}$ \; ($4/4$) & $\mathbf{+11.7\%}$ \; ($4/4$) \\
\bottomrule
\end{tabular}
\end{center}

On the split we did not tune on, refinement improves multi-view consistency by
about ten percent, on every sequence, for both heads --- and A2 beats canonical
there as well, which removes the second of the two marks against it recorded
last night.

\paragraph{What separates the splits.} Not the metric and not the code: the
amount of error available to remove. Raw excess over the ground-truth floor is
$0.55$--$1.11$\,mm on D7 against $0.18$--$0.22$ on D2, three to five times
larger. Where the raw prediction is poor, alignment recovers geometry a map can
fuse. Where it already sits near the reference floor, the history each frame
carries costs more than the correction gains. That is a sharper claim than
either half alone, and it is only visible because both halves were measured.

\paragraph{The mistake was generalising from the development split.}
The mechanism I wrote into the paper --- recurrence makes each frame a function
of its past, which fusion does not want --- is still the right explanation for
the D2 half. What was wrong was presenting it as the whole finding when only the
tuned split had been measured. The paper now states both directions and the
condition that separates them, and the conclusion's list of bounds says mapping
benefit is conditional rather than absent.

Second time tonight that a partial measurement pointed the wrong way, after the
closure cells. The pattern is the same and so is the lesson: measure the held-out
half before writing the sentence, not after.


\section*{19 August, Morning: The Unlearned Rule Beats the Baselines and Not the Head}

The review's third objection is the one that could have retired contribution~(ii):
the closure's baselines are spatially constant except for a deliberately badly
scaled one, and the obvious middle ground --- a two-parameter spatially varying
rule on the raw-versus-memory residual --- was never tested. Six of sixteen grid
points are complete, each on all fifteen cells:

\begin{center}
\begin{tabular}{lcc}
\toprule
$\alpha$ & $\tau{=}0.25$ & $\tau{=}0.5$ \\
\midrule
$0.3$ & $0.34\%$ & $0.38\%$ \\
$0.5$ & $0.54\%$ & $0.56\%$ \\
$0.8$ & $0.75\%$ & --- \\
$1.0$ & $\mathbf{0.76\%}$ & --- \\
\bottomrule
\end{tabular}
\end{center}

Against $0.59\%$ for the best fixed or EMA blend and $4.05\%$ for the learned
head.

\paragraph{Both halves of that matter.} The unlearned rule \emph{beats} every
constant baseline the paper had, so the review was right that the closure was
missing its strongest competitor and the paper was quietly flattering itself. And
it reaches about a fifth of the learned head, so contribution~(ii) survives a
much harder test than the one it passed before.

\paragraph{What bounds the unmeasured points.} As $\tau$ grows the rule degenerates
exactly into the fixed-weight baselines --- $w \rightarrow \alpha$ everywhere ---
whose best is $0.59\%$, and at $\alpha{=}1$ with large $\tau$ it becomes ungated
memory propagation, which the closure already measures as strongly negative. The
maximum is therefore interior and the two endpoints are known, which is why the
partial grid is informative here in a way the partial \emph{cells} earlier tonight
were not. Ten points remain and nothing goes in the paper until they land.

\paragraph{A distinction worth keeping.} Twice tonight a partial measurement
pointed the wrong way, and both times the incompleteness was \emph{within} a
measurement --- cells missing from a mean. Here every point is complete on all
fifteen cells and only the grid is partial. Those are different risks: the first
biases a number, the second only leaves a better number possibly unfound, and the
analytic limits above bound how much better it could be.


\section*{19 August, Morning: Which Branch Does the Work}

The review's second question --- which of $r_t$ and $f_t$ produces the gain --- is
answered by suppressing one factor at inference with the trained head unchanged.
Nothing is retrained, so the difference is the branch.

\begin{center}
\begin{tabular}{lcccc}
\toprule
& EPE red. & $B^+/H^+$ & NewBad1 & RecovBad1 \\
\midrule
$w=r_tf_t$ & $4.05\%$ & $\mathbf{2.53}$ & $\mathbf{0.047\%}$ & $0.227\%$ \\
$w=r_t$    & $3.42\%$ & $1.44$ & $0.191\%$ & $0.286\%$ \\
$w=f_t$    & $4.01\%$ & $2.15$ & $0.096\%$ & $0.244\%$ \\
\bottomrule
\end{tabular}
\end{center}

Read on EPE alone the answer is that fusion does everything and the reset gate is
redundant: $4.01$ against $4.05$, a difference of one part in a hundred. That
reading is wrong, and it is wrong in exactly the way this paper spends a section
warning about. Dropping the reset gate more than doubles the pixels the
intervention breaks and takes $B^+/H^+$ from $2.53$ to $2.15$.

So the two factors have different jobs: the fusion weight sets how much is gained
and the reset gate sets how little is broken. A paper that reports $B^+/H^+$
because counts and magnitudes can disagree has that disagreement inside its own
architecture, and would have missed it by reporting the product and one metric.

\paragraph{A mistake worth recording.} I killed the wrong jobs. Having just
argued in a commit message that the $\tau{=}2.0$ sweep points could not produce a
new maximum, I ran \texttt{bjobs | tail -2} and killed the two \emph{most recently
submitted} residual jobs --- which were the fine-$\tau$ points launched a minute
earlier to bracket the peak. The worthless work survived and the valuable work
died. \texttt{tail} on a job list selects by submission order, which has no
relation to what the jobs are doing; they have to be killed by identifier after
looking at what each one is.

\section*{19 August, Afternoon: The Cue That Costs Forty-One Per Cent of the Runtime}

The forward-backward confidence $C^{\mathrm{FB}}_t$ is expensive in a very specific
way: producing it requires the reverse optical-flow pass, which is $13.7$ of the
module's $33.5$~ms. Nothing else in the module costs so much for one channel. The
paper would like to propose dropping it, and cannot do so while saying only that the
saving is large --- ``we did not measure it'' is the review's complaint about a
different section and would be no better here.

So the cue was replaced by a constant at inference, the trained head otherwise
untouched, on the fifteen cells of the D2 closure:

\begin{center}
\begin{tabular}{lcc}
\toprule
& EPE red. & cells worse \\
\midrule
head as trained & $4.046\%$ & --- \\
$C^{\mathrm{FB}} := 1.0$ & $3.977\%$ & $12/15$ \\
$C^{\mathrm{FB}} := 0.5$ & $3.945\%$ & $12/15$ \\
\bottomrule
\end{tabular}
\end{center}

Two constants, because if the answer depended on which one, the head would be reading
the cue's \emph{level} rather than its variation, and that would be a different
finding. It does not: the two agree to a third of a tenth of a point and break the
same twelve cells. Two per cent of the effect for forty-one per cent of the runtime.

This is an upper bound on the cost and a lower bound on what is achievable. A head
\emph{trained} without the cue could learn to read the remaining channels differently
and would likely do better; what is measured here is a trained head deprived of an
input it was trained to use.

\paragraph{The number is not the shipped model's.} It was measured on the
$142$-channel configuration, which is now the ablation. Assuming it transferred would
have repeated the support-contract mistake exactly --- a result quoted for a head it
was never run on. Two readings were available and only the code settles which is
right. Counting the shipped head's thirteen native-resolution channels (raw
disparity, aligned memory, signed residual, absolute residual, six motion, three RGB)
leaves no room for a confidence channel, which suggests the cue is already absent and
the $13.7$~ms already free. The builder says otherwise: $C^{\mathrm{FB}}$ is one of
the \emph{six motion channels}, and that concatenation sits in the tail shared by
both branches. Removing the learned stereo evidence deleted $104$ correlation
channels and left the cue exactly where it was. The shipped head still consumes it,
the saving still has to be earned, and the measurement is running on the right head.

\section*{19 August, Afternoon: Running the Competitor the Way It Was Built}

TC-Stereo entered the paper with its temporal path switched off, justified by the
claim that SCARED-C ships no per-frame pose. The claim was false, it was retracted in
three places, and the retraction left the comparison owing a run. It is the only
rival here that is causal \emph{and} temporal --- BiDAStabilizer is bidirectional,
and this method carries its own state --- so with the path disabled the paper was
comparing against a single-frame network while calling it a temporal baseline.

The one thing genuinely ours to get right was the frame the poses live in. SCARED-C's
poses map the keyframe into each frame's \emph{native} camera; the disparity being
warped is \emph{rectified}; the two differ by $R_1$, so the pose must be conjugated,
$R_1 T R_1^{\top}$. A wrong frame here does not raise an error. It warps to the wrong
place and returns a plausible number, which is the most dangerous kind of bug this
paper can have.

The self-check warps ground truth between frames using the method's own \texttt{warp}
and asks which convention lands on the ground truth actually observed there:

\begin{center}
\begin{tabular}{lc}
\toprule
& median $|\Delta d|$ \\
\midrule
intended, $R_1 T R_1^{\top}$ & $\mathbf{0.0202}$~px \\
transposed order & $3.6412$~px \\
identity relative & $2.0723$~px \\
\bottomrule
\end{tabular}
\end{center}

Two orders of magnitude. That is the margin the first world-frame check failed to
produce --- it passed at $1.70$~mm while applying \emph{no} pose passed at $1.85$,
and it was thrown away for having no discriminating power. This one also validates
the baseline it recovers ($4.143$~mm) and, incidentally, that the $144{\times}180$
evaluation grid is a resize rather than a crop: a crop would move the principal point
and the intended convention would stop winning.

Three pieces of friction are worth recording, because none is about stereo.
\texttt{softsplat} has no CPU kernel --- it asserts false --- so even the geometry
check needs the GPU. \texttt{cupy}~13 removed \texttt{compile\_with\_cache}. And the
vendored kernel passes its include paths as one space-separated token that
\texttt{nvrtc} rejects. All three are packaging drift rather than differences in what
the method computes, so they are shimmed in our runner and the competitor's released
source is left byte-identical.

\paragraph{The diary had stopped compiling.} Adding these entries was how I found it.
A report pasted in months ago brought its own \verb|\documentclass| into the middle
of the document, which is fatal, so nothing after line~3950 had been typeset since.
Behind it were four more failures the first had been hiding: an \verb|\appendix| from
the very first report that every later report inherited, so the section counter ran
past~Z and raised ``Counter too large'' thirty-two times; \verb|\citep| with no
natbib and no bibliography; \verb|\mathbbm| and the algorithm environments with no
package; and a mid-document \verb|\newcommand| for a name the preamble already held.
The instruction is to update the diary at every result, and I had been appending to a
file that had not produced a page in its second half for a long time. It compiles now
--- 186 pages --- and the fix that matters is that the failure was silent: appending
text to a broken document looks exactly like appending text to a working one.

\section*{19 August, Afternoon: The Backbone Does Not Move the Sign, and the Confound Survives}

On RAFT-Stereo the world-frame drift result split cleanly: on D2, where the raw
prediction's excess over the ground-truth floor is $0.18$--$0.22$\,mm, the module
makes accumulated drift \emph{worse}; on D7, where the excess is $0.55$--$1.11$\,mm,
it makes it better. Two readings fit that equally well. Either the module helps
where the raw signal is poor --- a statement about raw quality, which would
generalise --- or it helps on D7, which is a statement about one dataset and
would not.

Backbones were supposed to separate them, because their raw quality differs on the
\emph{same} sequences. Twenty-three cells later, five estimators including two never
seen in temporal-head training:

\begin{center}
\begin{tabular}{lccc}
\toprule
& raw excess & drift reduction & sign \\
\midrule
D2, $15$ cells & $0.171$--$0.290$\,mm & $-32.1$ to $-5.4\%$ & $15/15$ negative \\
D7, $8$ cells & $0.552$--$1.113$\,mm & $+8.7$ to $+15.4\%$ & $8/8$ positive \\
\bottomrule
\end{tabular}
\end{center}

The half that worked: \emph{the backbone does not move the sign.} Not one cell
crosses over, and the two unseen estimators behave exactly like the three seen
ones. Whatever governs this, it is not the estimator, and that is worth reporting
because it was not obvious --- the module could easily have been repairing one
backbone's characteristic failure.

The half that did not: \emph{the confound survives.} The two raw-excess ranges
remain disjoint, and by a factor of two ($0.290$ against $0.552$). No backbone
makes D2 as hard as D7, so the experiment cannot tell raw quality from split ---
which is precisely the question it was built to answer. Worse for the raw-quality
reading, the ordering does not even hold \emph{within} D2: S$^2$M$^2$-S at
$0.269$\,mm excess loses $32.1\%$ while the same backbone at $0.290$\,mm loses only
$16.5\%$. If raw difficulty were the governing variable, that pair should not
happen.

So the paper says what was measured: the condition is robust across backbones, and
this experiment does not establish that raw quality is what sets it. Writing it as
``helps where the raw is poor'' would be inventing the result the design failed to
deliver. The honest version is weaker and is the one that goes in.

\section*{19 August, Afternoon: The Rename Was Complete and the Numbers Were Not}

A word-by-word audit of the paper, cross-checking every numeric claim against the
run reports, returned seventeen blocking discrepancies with a single cause. The
recanonicalisation renamed the model and rebuilt Tables~1--4, and stopped there.
Nearly every number \emph{outside} those tables was still the $142$-channel head's,
presented as the shipped model's.

What makes this worth recording is why it survived. The lexical rename was
\emph{perfect}: not one stray ``A2'', not one stray ``canonical'' in the body text,
every ``$142$-channel'' correctly labelled as the ablation. A search-and-replace pass
finds nothing to fix, which is exactly why nobody looked again. The damage was
entirely numeric, and numbers do not carry the name of the model that produced them.

The paper was contradicting itself in public. The abstract said continuous recurrence
costs $62.69\%$ of EPE; the body said $19.17\%$, three separate times. The DRENDS
section and the abstract quoted different models for the same hundred cells. The
runtime paragraph said the evidence tensor costs $5.13$~ms and then, two sentences
later, $0.67$~ms --- because it had one head's total and the other head's breakdown.
A reviewer checking that one paragraph would have found the budget failing to add up
without needing any of our data.

The two figures were the worst of it. Both were generated from the $142$-channel
head, including Figure~1, the opening visual. The paper's first picture of its own
method was a picture of its ablation.

And the world-frame drift attributions were \emph{swapped}: the shipped head was
credited with the ablation's weaker D7 number and the ablation with the stronger one.
That one is instructive, because both numbers were real, both were in the right
paragraph, and only the labels were exchanged. No consistency check catches that.

\paragraph{Three deletions that make the paper better.} The horizon table went. Its
DTCE column --- the sole argument for shipping $H{=}4$ over the lower-EPE $H{=}6$ ---
is one lag of one aggregation, and the audit found the pooled $+4.99\%$ is carried by
a single backbone: $+21.0\%$ on Stereo Anywhere against $+1.4$, $+0.1$, $-1.5$ and
$-1.6$ on the other four. A five-backbone claim resting on one backbone, in a table
inviting exactly that scrutiny.

Table~3's $142$-channel row went, and this one is almost funny: it held the only bold
cell the shipped model did not own. Removing it leaves the proposed method best on all
four metrics against a bidirectional baseline that sees the future. It was also
indefensible on its own terms --- $142$ is not a designed number, it is $38$ plus the
leftovers of a frozen ResNet-18, and a reviewer asking ``why $142$ and not $141$?''
would have received no answer, because there isn't one.

The SERV-CT paragraph went. Temporal refinement on non-consecutive pairs is a literal
identity, so a geometry column there would print the backbone's number twice. Five
lines promising a dataset and delivering nothing.

\paragraph{What the closure was burying.} The matched-baseline argument was being made
on mean EPE, where it reads as a margin: $3.69\%$ against $0.59\%$ for the best fixed
blend. On Bad1 it is a \emph{sign change}. Every matched smoothing policy lowers mean
error while \emph{raising} the bad-pixel rate --- fixed $w{=}0.5$ at $-2.64\%$, EMA
over three frames at $-8.88\%$, confidence-weighted at $-32.35\%$ --- and the learned
gate is at $+12.84\%$. Naive temporal averaging trades tail correctness for mean error.
That is the finding, it was in the reports the whole time, and the paper was spending
its strongest evidence as a footnote to a weaker version of itself.

\section*{19 August, Evening: The Cue Costs Nothing on the Head That Ships}

The forward--backward confidence was measured on the $142$-channel head this
morning: replacing it with a constant cost $0.069$ points of that head's $4.046\%$
D2 closure at $1.0$ and $0.101$ at $0.5$, breaking $12$ of $15$ cells both times.
A small price for $41\%$ of the runtime, and the entry said so.

The same substitution on the head the paper actually ships goes the other way:

\begin{center}
\begin{tabular}{lcc}
\toprule
& D2 closure & cells worse \\
\midrule
\textsc{tether} as trained & $3.694\%$ & --- \\
$C^{\mathrm{FB}} := 1.0$ & $3.744\%$ & $4/15$ \\
$C^{\mathrm{FB}} := 0.5$ & $3.785\%$ & $3/15$ \\
\bottomrule
\end{tabular}
\end{center}

Not a cost. A small \emph{gain}, in both constants, breaking a quarter as many
cells. This is exactly why the morning's number could not be inherited, and it is
worth recording that the reason I gave for re-measuring was the wrong one: I argued
that $C^{\mathrm{FB}}$ survives into the shipped evidence space as one of the six
motion channels, so the saving still had to be earned. It survives, and it is not
earned --- it is free.

The sign flip is consistent with the rest of the ablation rather than mysterious.
The learned evidence made the head sensitive to a cue that the geometric evidence
space renders redundant: the same $104$ channels whose removal tightened the seed
spread eightfold were also what gave $C^{\mathrm{FB}}$ something to be useful for.

What this buys, if it holds: the reverse SEA-RAFT pass exists only to produce
$C^{\mathrm{FB}}_t$, and it is $13.1$ of the module's $28.9$\,ms. Forty-five per
cent of the temporal module, for a change in accuracy an order of magnitude inside
the seed spread.

Two reasons the paper reports this instead of shipping it. It is development-set
evidence, fifteen cells on D2, and the module's operating point was frozen before
D7 opened precisely so that decisions of this kind are not made on the split that
tunes them. And the head under test was \emph{trained} with the cue: feeding it a
constant asks a trained network to ignore an input it learned to weigh, which is an
upper bound on what removal costs and a lower bound on what a head trained without
it could do. The experiment that would settle it --- retraining without the cue ---
is the one not run.

\section*{19 August, Evening: Which Branch Does the Work, on the Head That Ships}

The branch decomposition existed only for the $142$-channel configuration, which
is now the ablation, and the paper was discussing it under a heading naming the
shipped model. Rerun where it belongs:

\begin{center}
\begin{tabular}{lccc}
\toprule
& EPE red. & NewBad1 & $B^+/H^+$ \\
\midrule
$w=r_tf_t$ & $\mathbf{3.69\%}$ & $\mathbf{0.063\%}$ & $\mathbf{1.85}$ \\
$w=f_t$    & $3.49\%$ & $0.084\%$ & $1.72$ \\
$w=r_t$    & $2.23\%$ & $0.211\%$ & $1.21$ \\
\bottomrule
\end{tabular}
\end{center}

The conclusion survives the change of head and gets easier to state. Fusion sets
how much is gained, the reset gate sets how little is broken.

What changes is the strength of the first half. On the $142$-channel head the EPE
column was almost flat --- $4.01$ against $4.05$, one part in a hundred --- so the
entire argument for keeping the reset gate rested on the broken-pixel count, and I
wrote at the time that reading EPE alone would call the gate redundant. On the
shipped head the gate earns $0.20$ points of mean error as well, five times the
margin it had before. The same removal of learned evidence that tightened the seed
spread eightfold also made the two factors less interchangeable.

So the paper now makes the same point with two independent columns supporting it
rather than one, which is a better position than the one I was defending this
morning.

\section*{19 August, Night: The Competitor's Temporal Path Makes It Worse}

TC-Stereo entered this paper with its temporal path switched off, justified by a
claim that was false, and the retraction left a run owing. It has now landed, over
all three D2 keyframes, $18.0$M pixels, with the pose convention validated
beforehand at $0.0202$\,px against $3.64$ and $2.07$ for the two wrong frames.

\begin{center}
\begin{tabular}{lccc}
\toprule
& EPE & Bad1 & $\mathrm{DTCE}_{k=1}$ \\
\midrule
raw & $0.303$ & $2.09\%$ & $0.073$ \\
TC-Stereo, frame path & $0.524$ & $9.05\%$ & --- \\
TC-Stereo, temporal path & $0.622$ & $10.37\%$ & $0.406$ \\
\bottomrule
\end{tabular}
\end{center}

Turning on the temporal path makes the method worse per frame. That was possible.
What was not expected is the size of the temporal result: change error rises from
$0.073$ to $0.406$\,px at lag~1 and from $0.135$ to $0.592$ at lag~8, three to five
times \emph{less} temporally consistent than the prediction it refines --- at every
lag, with no lag disagreeing.

A method whose entire purpose is temporal consistency, run the way its authors
specify, degrades temporal consistency fourfold on surgical video.

The mechanism is not mysterious and the paper predicted it in prose before the run
existed: the propagation warps previous disparity by a \emph{rigid} $6$-DoF
reprojection. Tissue deforms. A rigid warp of a deforming surface puts geometry in
the wrong place, and the wrong place changes from frame to frame, which is exactly
the disagreement the metric counts. The pose is right --- that is what the
self-check bought --- and the scene is what is wrong for the assumption.

Two things I want recorded honestly. First, this is a strong claim about someone
else's method, so the failure modes were closed deliberately: the chaining follows
their own \texttt{evaluate\_stereo.py} rather than my reading of the paper, the
pose convention was scored against two wrong alternatives before any accuracy
number was computed, and the released weights and iteration count are unmodified.
Second, the result exists at all only because the run was caught at a quarter of
the way through and made to save its prediction stacks. Rescoring for a metric
nobody had asked for yet took ninety seconds on a CPU. Without that it would have
been another five hours of GPU, and the honest likelihood is that it would simply
never have been measured.


\section*{DRENDS agrees with D2: we win the geometry, BiDA wins the clock}

The head-to-head table in the paper is the DRENDS one, because it is the only arena
where neither method is in domain. Until now it carried four geometric columns and
no temporal one --- which, for a comparison between two temporal refiners, was the
one omission with no defence. The DTCE columns are in it now, over all five
recordings and $117.4$M pixels of common support:

\begin{center}
\begin{tabular}{lcccc}
\toprule
& EPE & Bad1 & $\mathrm{DTCE}_{k=1}$ & $\mathrm{DTCE}_{k=4}$ \\
\midrule
raw            & $1.199$ & $44.92\%$ & $0.242$ & $0.340$ \\
BiDAStabilizer & $1.175$ & $44.81\%$ & $0.085$ & $0.214$ \\
TETHER         & $1.131$ & $43.39\%$ & $0.139$ & $0.294$ \\
\bottomrule
\end{tabular}
\end{center}

The bidirectional baseline reduces change error by $64.9\%$ at lag~1 where we reduce
it by $42.5\%$, and it is behind us on every geometric column. The same split D2
already showed ($52\%$ against $15\%$), on a different dataset, with a different
ordering of the two methods' domain advantage --- and it holds in each of the five
recordings separately and at every lag. A method that reads the future is more
temporally stable than one that cannot, and is simultaneously the less accurate of
the two. The two best entries in the table sit in different rows, and I have written
the caption to say so rather than bolding our way around it.

This is our own warning applied to our own comparison. If a reader were handed the
consistency column alone they would rank these two methods the wrong way round.

Two process notes, both unflattering. The run that produced these numbers took five
hours the first time and two minutes per recording the second, and nothing about the
model changed: every \texttt{npz[key]} decompresses the entire array, so indexing
one frame out of a $466$\,MB deflate-compressed stack inside a per-frame loop
inflated the whole thing, about $1.7$\,TB of discarded decompression per recording,
on one core, with the GPU idle behind it. The same pattern was in five scripts,
including the TC-Stereo run whose slowdown I had blamed on GPU contention earlier
the same evening. Contention was real and secondary. Slow code that is also
contended looks exactly like contended code.

Second: I killed a job by an identifier I had rather than one I had looked up, and
the id was two launches stale, so the run I meant to stop kept going and a second
one started writing the same output directory. Caught by timestamps, not by design.
The comparison script now takes an exclusive \texttt{flock} on its output directory
and the second run exits loudly. Verification that matters: the reissued numbers
reproduce the previously recorded EPEs exactly --- $1.1988$ raw, $1.1746$ BiDA,
$1.1311$ TETHER --- so only the speed changed.


\section*{A3 at three seeds: the margin survives, the rule sends it back}

A3 --- the context branch at full resolution instead of a quarter --- was the one
deviation whose single run beat its base by more than the declared seed spread, which
is precisely the case where a single run cannot be believed. Two more seeds:

\begin{center}
\begin{tabular}{lccc}
\toprule
& seed 0 & seed 1 & seed 2 \\
\midrule
D7 EPE reduction  & $5.61$ & $4.65$ & $5.46$ \\
D7 Bad1 reduction & $5.85$ & $5.53$ & $5.67$ \\
D2 Bad1 reduction & $9.22$ & $7.55$ & $7.54$ \\
\bottomrule
\end{tabular}
\end{center}

Held-out Bad1: $5.68\pm0.16$ against the base configuration's $5.15\pm0.34$. The
margin does not vanish. It shrinks from $+0.70$ points on one run to $+0.54$ over
three, still outside the spread, $1.6\sigma$.

The pre-registration names D2 as the decision split, and there A3 reduces
$8.10\pm0.97$ against $8.40$ --- $0.29$ points on the wrong side, well inside the
$2.53$-point canonical spread. The rule returns it to the ablation table, so that is
where it stays, at $16\times$ the context-branch compute.

I want to be exact about what this is, because the comfortable reading is available
and it is not the true one. This is not ``the single run was seed luck''. Seed~0 was
the most favourable of the three, so the effect was overstated, but the effect is
there on every seed. What sends A3 back is the split we declared in advance, not the
disappearance of its advantage. A reader who weighs held-out performance above our
pre-registered development endpoint is looking at a deviation that beat its base
three times out of three, and the paper now says so in those words rather than
reporting a verdict and letting the reader assume the evidence was one-sided.

The rule earns its keep by being uncomfortable here. It promoted A2 --- the deviation
that became the shipped model --- on the same criterion two days ago. A rule that only
ever confirms what we wanted is not a rule.


\section*{Three reviewers, and the pre-registration we had been quoting wrong}

I put three max-reasoning readers on the paper --- method and claims, statistics
and aggregation, structure and page budget --- after having verified every number
in it myself against the runs on disk. Fifty table cells, the horizon grid, the
runtime stages, the seed bootstrap, the drift cells, the frame scatter: all
reproduce exactly. What the readers found was not arithmetic.

\subsection*{The one that matters}

Two of them independently found the same thing, and they were right. The paper
quoted its own pre-registration two ways in two places, and the reading it
invoked was the one that kept the shipped model.

There are two registrations. The ablation registration names \emph{held-out
Bad1} as the endpoint. A separate, later one governs \emph{promotion} of a
variant to the shipped model and decides on development only --- which is how
the $38$-channel head was promoted. I applied the promotion rule to A3 earlier
tonight and wrote a paragraph explaining why A3 goes back to the ablation table
under it. That paragraph was using the wrong rule.

Under the right one, A3 \emph{wins}: $5.68\pm0.16$ against $5.15\pm0.34$ on
held-out Bad1, $0.54$ points against a $0.37$-point declared threshold, three
seeds each. The registration enumerated what to do if A3 matched and what to do
if it lost. It did not enumerate winning. We decline it anyway, on $16\times$
the context-branch compute --- and the registration's own honest note had
required exactly that: ``If it wins, that trade has to be reported, not
hidden.''

A4 is worse for us. It posts the largest held-out Bad1 reduction in the table,
$6.02\%$, and we reject it on EPE, a metric the registration names as reported
but not as the endpoint. That reading was chosen after seeing the result, and
the paper now says so in those words. Its falsification clause did fire: ``If
A4 matches or beats the canonical model on held-out D7, the Analysis sentence
attributing transfer to the convex action space must be withdrawn rather than
softened.'' I went looking for that sentence expecting to have to delete it.
It was already gone --- withdrawn during an earlier pass, for a different
reason. The paper now records that the clause fired and that the claim is
absent, which is the difference between honouring a rule and happening to
comply with it.

\subsection*{The oracle belonged to the ablation}

The closure reported a $12.30\%$ raw-versus-memory ceiling and said TETHER
recovers $30.0\%$ of it. Both numbers were built on the $142$-channel head's
aligned memory. The script that re-runs a promoted head declared the oracle
``GT-only; never runs the adapter'' and copied it forward from the canonical
closure, and that declaration is exactly half true: ground truth does the
selecting and never touches the adapter, but it selects between raw and the
\emph{aligned memory}, which under $H{=}4$ is the previous fused output on
three frames out of four. A different head gives a different memory and
therefore a different ceiling.

Recomputed on the shipped head's own memory: ceiling $12.74\%$, recovery
$29.0\%$ pooled and $25.2\%$ per cell. The gap we have to close is larger than
we said. The re-run reproduces the original canonical\_h4 gain to the sixth
decimal, so nothing but the oracle changed.

Related, and worse in kind: the paper wrote ``$30.0\%$ as a ratio of the pooled
quantities, $32.9\%$ averaging per-cell recovery'', which reads as one
measurement under two aggregations. They were two different measurements on two
different supports, and the true per-cell figure is \emph{below} the pooled one,
not above. The sentence pointed the wrong way.

\subsection*{Smaller, but each one a claim}

\begin{itemize}
\item ``A spread of a quarter of a point across the whole grid.'' The grid spans
$1.24$ points. $0.21$ is the spread after EPE rejects $H{=}1$ --- a different
sentence, and the one we meant.
\item The conclusion cited $5.03$ and $0.69$ points of seed spread as the
paper's finding. Those are the ablation's; the shipped head spans $0.60$ and
$0.40$. As written the sentence was false of the model it concludes about.
\item ``The shift that costs more is the backbone, not the domain.'' On D2 the
seen advantage is carried entirely by Stereo Anywhere; without it the ordering
reverses ($1.85$ against $2.67$). Withdrawn.
\item ``All $100$ cells'' multiplied four correlated metrics into the count.
$25$ recording--backbone cells, four metrics each.
\item The bootstrap resamples backbone--sequence cells, not sequences, and they
share three and four underlying sequences. Reported as descriptive, with the
sign test the splits can actually carry ($p=0.125$, $p=0.0625$).
\item ``$62\times$ smaller'' omits the $8.88$M frozen SEA-RAFT that does the
alignment and $26.2$ of our $28.9$\,ms.
\item The $14\times$ Bad1 ratio divides by the baseline's $0.11$ points.
\item TC-Stereo: pixel-pooled values quoted with the macro multiple. Four to six
times, not three to five --- we were understating our own result.
\item Seventeen sweep points, not sixteen, and the optimum is interior.
\end{itemize}

One reviewer accused us of a defect we do not have: that DTCE is gameable by
staleness, so a stale prediction wins it. DTCE is
$|\Delta\text{prediction}-\Delta\text{target}|$ --- a frozen prediction has
$\Delta{=}0$ and is charged the whole of the reference motion. But a careful
reader misread it, which is a defect of the paper and not of the reader, so the
definition now says so where it is defined.

The hardest objection --- ``you never run CODD, your own acknowledged parent''
--- had its answer sitting in our own source the whole time. The docstring of
\texttt{codd\_style\_fusion.py} reads: ``CODD obtains matching features from its
stereo network; cached ARGOS backbones do not expose those features.'' The
$142$-channel configuration \emph{is} the faithful-as-possible CODD
reproduction, with its projection replaced by our warp and a frozen ResNet-18
standing in for the features. The comparison the reviewer wanted has been in
Table~V since the beginning; the paper simply never said what that row was.

\subsection*{Completing a grid that was shaped to flatter us}

The world-frame drift claim was ``$23$ cells in which the sign never changes:
all $15$ D2 negative, all $8$ D7 positive''. That asymmetry was not a design.
D2 had five backbones on three sequences; D7 had five backbones on keyframe~1
and RAFT-Stereo alone on the other three. The favourable half rested on a
smaller and differently-shaped support, and the paper never said so.

Twelve cells later: $35$ cells, all $15$ D2 negative ($-32.1$ to $-5.4\%$), all
$20$ D7 positive ($+7.2$ to $+17.8\%$). The claim survives on the complete grid.
The raw-excess ranges are still disjoint, so the mechanism is still unsettled
and we still do not claim raw difficulty as the cause.

\subsection*{Eleven pages to eight}

No measured result was deleted. The BiDAStabilizer table, the per-backbone
DRENDS table and the frame scatter moved to the supplement, which already
carried all three. What was compressed rather than moved was repetition: the
support contract was stated five times, smoothness-is-not-correctness three,
the $142$-channel instability three, and the conclusion restated every number
in the paper a second time.

Every self-critical result stayed in the main text at full length. That was the
constraint I set before cutting, and the list is worth writing down because the
temptation runs the other way: the Vid13 net harm, the AUROC below chance, the
D2 mapping degradation, the no-reference gameability, the epoch-selection noise,
the per-backbone negative DTCE on D2, A3 winning the registered endpoint, and an
oracle gap that tonight got larger rather than smaller.


\section*{Reading the eight pages back, and what thirty edits broke}

Cutting a paper by a third is thirty edits, and thirty edits break things
quietly. So I read the result end to end against the sources again. Four
defects, one of which had been there far longer than tonight.

The closure table's caption said ``the six rows above ours load no
checkpoint''. Five do. The sixth is \emph{unbounded corrected recurrence},
which is our own trained head with the re-anchor removed --- and that is
exactly why it is the row that inverts, so the caption was not merely wrong
but wrong in the direction that made our own worst row look like somebody
else's baseline. That error predates the cutting entirely; it has been in the
table since the table was written, and neither I nor three reviewers reading
for defects caught it. What caught it was reading the caption next to the rows
it describes, which is a thing I had not done.

The seed section declared sequences the resampling unit and then, two
sentences later, correctly said the bootstrap draws backbone--sequence cells.
Both sentences were mine, one of them written tonight to fix a reviewer's
objection about exactly this.

The drift paragraph quoted RAFT-Stereo-only excess ranges and then reasoned
about the full grid's disjointness --- two different supports for one argument,
created by updating half a paragraph when the twelve new cells landed.

And ``the Bad1 column says so'' outlived the table whose column it was.

One more, from the statistics reviewer and correct: A1 sits $0.27$ points
below its base, inside the $0.37$-point threshold we declared. Under our own
rule that is unreadable, not worse. The paper said three deviations sit above
the base and left the reader to infer the fourth had lost.

Eight pages, no errors, no undefined references, no overfull boxes. The only
thing still open is the convergence probe at $20$ of $100$ epochs.


\section*{The supplement was never recanonicalised, and it still made the claim we withdrew}

Three days ago the paper was moved onto the $38$-channel head and the commit
said ``abstract, contribution, method, grid, closure, horizons, ablations, seed
study, out-of-domain, runtime, drift and conclusion all regenerated''. Twelve
things. The supplement was not one of them, and nobody --- me, three reviewers
reading for defects, or the header comment in \texttt{ICRA.tex} that lists
exactly which directories belong to which head --- went downstairs to look.

Its closure table, its horizon sweep, its oracle section and its
Analysis-in-full were all the $142$-channel ablation's numbers under the label
TETHER. Continuous recurrence at $-62.69\%$ where the shipped head is
$-19.17$. The oracle ceiling at $12.30$ where the shipped head's own memory
gives $12.74$. AUROC $0.366$--$0.382$ where ours is $0.393$--$0.408$.

The one that should not have survived: earlier tonight the main paper gained a
sentence saying the pre-registration's falsification clause fired and that the
claim attributing cross-backbone transfer to the convex action space
``appears nowhere in this paper''. It appeared in the supplement, in full, as
``a decision problem that constrained is a plausible reason why one small head
transfers to architectures absent from its training''. I wrote the withdrawal
and the surviving copy of the withdrawn claim on the same night, in two files,
and did not connect them.

The ablation section was stale in a different way: it still said A2 ``is being
given'' the three-seed treatment and ``is not being promoted''. A2 was
promoted three days ago and is the model the paper proposes. The supplement was
describing a decision as pending that the title page had already made.

Two numbers I removed rather than corrected: the share of achievable selective
gain the fusion weight captures ($3.5\%$ and $6.1\%$). I could not reproduce
the definition from the CSV, and carrying the ablation's figures under the
shipped head's name is the exact defect I was there to fix. The claim they
supported --- that $w_t$ is sign-unstable rather than merely weak --- stands on
the reliability inversion, which I can reproduce: for RAFT-Stereo the harm rate
falls with $w_t$ on D2 ($0.396\rightarrow0.287$) and rises on D7
($0.250\rightarrow0.362$).

The lesson is dull and worth writing anyway. A recanonicalisation is a
find-and-replace over a \emph{set of files}, and the set was defined by which
file I happened to have open. The header comment that tells the reader which
run directory belongs to which head lives in the main paper only. It should
live wherever a number does.


\section*{Sweeping the rest of the supplement, and a table with no source}

Having found the supplement on the wrong head, I went through its remaining
sections one at a time instead of spot-checking. Six more were the
$142$-channel configuration's: per-recording DRENDS, the seed study,
intervention quality per backbone, three-dimensional consistency, stage-wise
runtime, and the augmented-training comparison.

Three of those are worth recording individually.

The DRENDS table was missing a row --- RAFT-Stereo on Vid14 --- while the
sentence beneath it counted sixty cells. Fourteen rows, sixty claimed.

The runtime table reported $33.601$\,ms, $124.2$\,MB and a ``$177$k-parameter
head'', and concluded that $79\%$ of the overhead is optical flow. On the
shipped head it is $28.903$\,ms, $95.2$\,MB, a $155$k head, and $91\%$ flow.
Its closing sentence called removing the reverse pass ``the most promising
efficiency ablation we have not run'' --- we ran the inference-time half of it
hours earlier and found it costs this head nothing.

And the one I did not expect. The section on frame-to-frame surface
consistency opened by asserting that SCARED-C provides no per-frame camera
pose, so reconstruction drift cannot be measured. That claim has been retracted
four times, once from the very script that produced the section's numbers
(\texttt{c8ccde0}). This was the fifth copy, and it survived because retracting
a claim from code does not retract it from prose that quotes the code.

\subsection*{A table I cannot source}

The motion-compensated fidelity table was neither head. It printed D2
$\TEPE_{\mathrm{corr}}$ for S$^2$M$^2$-S as $+24.5/+40.1/+12.9/+19.7$. The
$142$-channel CSV of record gives $+0.5/+10.4/+3.6/+8.0$, and that CSV was
written four hours \emph{before} the table was committed. I cannot reconstruct
where the printed numbers came from. They are replaced by ones I can reproduce
from the run, on the shipped head.

The conclusion under it was wrong in the direction that flattered nobody. It
said OPW improves everywhere on D2 while $\TEPE_{\mathrm{corr}}$ ``degrades
everywhere''. On the shipped head $\TEPE_{\mathrm{corr}}$ improves in nine of
twenty D2 cells: at lag~$1$ for every backbone, and at every lag for Stereo
Anywhere --- the backbone whose raw prediction is weakest. The eleven
degrading cells are all at $k\geq2$ on the four stronger backbones. So the
module tracks the true one-frame change better everywhere, and loses only
beyond one frame. That is both more precise and less damning than what we had
written about ourselves, and I want it on the record that the correction went
this way, because most of tonight's went the other.

Held out, both metrics improve in all twenty cells, which is what the main
paper claims.

\subsection*{Tally}

Five sweeps, and every file I opened had been left behind by the
recanonicalisation: the supplement (four sections, then six more), the slide
deck shown to the group, a README, and \texttt{evaluate\_temporal\_corrected.py},
which had no \texttt{--output} flag and so could only ever overwrite the old
head's results --- which is precisely why nobody re-ran it. The commit that
performed the recanonicalisation listed twelve things. At least eight were
missing from that list.


\section*{The defaults were the cause}

Five sweeps found stale numbers in every file I opened. Tonight I stopped
fixing the outputs and looked at why they were stale, and the answer is
mechanical rather than careless.

Ten scripts that drive a head defaulted to the $142$-channel one. Their
launchers pass no \texttt{--module}, because why would they --- the default is
supposed to be what the project ships. So every bare run produced the
ablation's numbers under the shipped model's name. Three of those scripts
wrote to a fixed path with no way to send a second head anywhere else, which is
the other half of the mechanism: \texttt{evaluate\_temporal\_corrected.py} was
never re-run after the head changed because re-running it could only have
destroyed the record it was supposed to be compared against.

All ten now default to the shipped head, and the three that could only
overwrite route by head instead.

One live bug found while doing it. \texttt{analyze\_oracle\_and\_risk.py} was
the one script that already routed its output by module --- and it matched
\texttt{endswith("canonical\_h4:factory")}, so the \emph{masked} spelling of
the $142$-channel head fell through to the \texttt{a2} directory. The single
script written to prevent this collision contained the only path by which a
wrong head could still overwrite the right one's results.

I want to record the shape of this rather than the list. A recanonicalisation
is not a rename; it is a change to what the word ``default'' denotes. Renaming
the model in twelve files and leaving the defaults pointing at the old one
guarantees that every subsequent run quietly undoes the rename. The commit that
did the recanonicalisation was careful and complete about prose. It did not
touch a single \texttt{argparse} line.



\section*{The map metric survives its own degenerate control, and the horizon sweep is monotone}

The world-frame grid finished all $35$ cells this afternoon, and both halves of
the reviewer's objection now have measured answers.

The uncomfortable half first. Our contribution (iii) argues that a temporal
metric a degenerate policy can win is not a metric, and the reviewer asked
whether the world-frame excess-over-GT-floor number is such a metric. It is
not. Copy-previous \emph{loses} it, badly: $-54.0\%$ on D2 (cells $-34.9$,
$-87.5$, $-39.5$) and only $+2.8\%$ on D7 against our $+11.7\%$. Warped-previous
at full weight also loses, $-37.6\%$ on D2 and $+7.2\%$ on D7. On the
no-reference score the same copy-previous policy removes $97.3\%$ of the
apparent signal; on the world-frame number it is the worst policy in the table.
Staying still does not win agreement across views.

The mechanism half. Sweeping the recurrence horizon of the shipped head:

\begin{center}
\begin{tabular}{lcc}
horizon & D2 & D7 \\
$H{=}1$ & $-3.6\%$ & $+8.9\%$ \\
$H{=}2$ & $-8.6\%$ & $+10.4\%$ \\
$H{=}4$ & $-15.5\%$ & $+11.7\%$ \\
$H{=}6$ & $-20.2\%$ & $+12.2\%$ \\
\end{tabular}
\end{center}

Monotone on both splits, opposite signs. On D2 the damage grows with the
horizon --- a factor of six from one warp to six --- so it is alignment error
\emph{accumulated by the recurrence}, not the cost of a single warp; and the
single warp applied at full weight (warped-previous, $-37.6\%$) shows the gate
is what keeps TETHER's $H{=}1$ cost down at $-3.6\%$. On D7 the benefit rises
with the horizon.

A correction against my own earlier reading: with the last cell filled, D7 is
\emph{not} peaked at the shipped $H{=}4$ --- $H{=}6$ edges it, $+12.2$ against
$+11.7$. The honest sentence is that $H{=}4$ was frozen pre-registration on the
per-frame trade-off, and on the map metric the D7 curve is still rising at the
frozen point while the D2 curve is still falling. Anyone wanting a mapping
configuration would trade differently, and the paper should not pretend the
frozen point is optimal for a criterion that did not exist when it was frozen.

None of this is in the paper yet; edits now require authorisation, so this
entry is the register.



\section*{The consensus chain runs for the first time, and catches the drift it was built for}

The 38-channel twins of A3 and A4 (the paper's blocker T2.1) went through the
new protocol: spec, implementation, independent adversarial review by a reader
who had not seen the implementation, adjudication. The reviewer returned seven
findings; four changed something.

The one that justifies the whole apparatus: \texttt{train\_ablation\_seed.py}
carried its own copy of the list of variants that inherit the
disable-learned-evidence flag, one revision behind. A5--A7 were immune by
accident (their cue builders force the flag internally), so the divergence was
invisible until A3b/A4b --- the first variants since A2 whose 38-channelness
rests on the flag alone. An A3b seed run would have trained the most expensive
configuration in the repository (full-resolution context on 142 channels) for
twelve epochs and written a provenance record asserting 38 channels, with the
error surfacing only at evaluation, a day of H100 later. This is the third
instance of the same failure shape in one day --- eval-side A5/A6 this
morning, eval-side A7 caught by the implementer, train-side seed path caught
by the reviewer --- and all three are copies of a variant list drifting apart.
The list is now defined once.

Also from the review: the twins' registration compared against seed 0, the
comparison this project already retracted once (seed 0 is the weakest shipped
seed held out, so it flatters every variant); now three seeds, 5.84 $\pm$
0.22. And the registration gained a decision rule written before either twin
trains, naming the favourable prior it starts from: A3 beat its own base by
0.70 points, which is exactly why the rule is written down first.

Training is submitted, gated on the A6 job ending. Nothing here touches the
paper; the twin rows land in Table~III when the checkpoints exist.

\end{document}